\documentclass[11pt]{article}
% Diplomatic transcription macros (12 Nov 2026). Typed text is plain; the handwritten layer is blue,
% deletions red, typewriter strike-outs grey, doubtful readings in brackets with a question mark.
\usepackage[utf8]{inputenc}
\usepackage[T1]{fontenc}
\usepackage{lmodern}
\usepackage[a4paper,margin=2.2cm]{geometry}
\usepackage{amsmath,amssymb}
\usepackage{xcolor}
\usepackage{tikz-cd}
\usepackage{environ}
\setlength{\parindent}{0pt}
\setlength{\parskip}{4pt}
\title{Nouveau EGA V (cote 25)\\[4pt]\large Diplomatic transcription from the Montpellier scans}
\author{AI-transcribed by Claude Opus 5.5}
\date{12 November 2026}
\newcommand{\scanpage}[1]{\par\bigskip\noindent{\small\sffamily[scan #1]}\par}
\newcommand{\mspage}[1]{{\small\sffamily[p.\,#1]}}
\newcommand{\hand}[1]{\textcolor{blue}{#1}}
\let\over\relax
\newcommand{\over}[2]{\textcolor{gray}{#1}\textcolor{blue}{\ensuremath{^{\!\nearrow}}#2}}
\newcommand{\inserted}[2][]{\textcolor{blue}{$\langle$#2$\rangle$}}
\newcommand{\margin}[1]{\textcolor{blue}{\{\textsf{marg.:} #1\}}}
\newcommand{\deleted}[1]{\textcolor{red}{[\textsf{del.:} #1]}}
\newcommand{\xout}[1]{\textcolor{gray}{[\textsf{x:} #1]}}
\newcommand{\uline}[1]{\underline{#1}}
\newcommand{\uuline}[1]{\underline{\underline{#1}}}
\newcommand{\unsure}[1]{[#1?]}
\newcommand{\illegible}[1]{[\textsf{illeg.\ #1}]}
\newcommand{\diagram}[1]{\par{\centering #1\par}}
\newenvironment{diagramblock}{\par\begin{center}}{\end{center}\par}

\begin{document}
\maketitle

% Nouveau EGA V (cote 25): archive cote 25, site page 1 (PDF page 2).
% typescript. Draft compared: none.
% AI-transcribed by Claude Opus 5.5 (claude-opus-5-5, adaptive thinking, effort high), Message Batches API, batch ega25b (msgbatch_01A7n5YuXubUPNK1zsepi8RG), 2026-11-13; not yet checked by eye.
\scanpage{25/1}
% NOTE: the page carries no typed text; the sheet is blank apart from the number at the foot.
\mspage{1}

\clearpage
% Nouveau EGA V (cote 25): archive cote 25, site page 2 (PDF page 3).
% typescript. Draft compared: none.
% AI-transcribed by Claude Opus 5.5 (claude-opus-5-5, adaptive thinking, effort high), Message Batches API, batch pilot (msgbatch_0189Nhf1xB615YWnztubtLfc), 2026-11-12; not yet checked by eye.
\scanpage{25/2}
\margin{\unsure{Nouveau EGA V}} \mspage{0}

EGA IV 20 \quad \uline{Sections hyperplanes et projections coniques.} (Etat \unsure{000|0})

1. Notations.

2. Section hyperplane générique: propriétés locales.

3. Section hyperplane générique: irréductibilité et connexité géo\unsure{métriques.}

4. Section hypérplane variable: sections ``assez g\unsure{énérales|éńérales}''.

5. Théorèmes du\unsure{type|ttype} Seidenberg.

6. Connexion d'une section hyperplane quelconque.

7. Application a des constructions de sections hyperplanes et multisec-%
tions spéciales.

8. Dimension de l'ensem\xout{m}ble  des hyperplans exceptionnels.

9. Changement d'immersion projective.

10. Pinceaux de sections hyperplanes, et fibrations de variétés éclatées

11. Grassmaniennes.

12. Généralisation de résultats antérieurs à des sections linéaires.

13. Morphismes élémentaires et théorème de M. Artin.

14. Projections coniques.

15. Axiomatisation de certains des résultats précédents.

16. Traductions en langage de systèmes linéaires.

\uline{Commentaires}. Cette rédaction donne pêle mêle  une esquisse détaillée
de l'ensemble des résultats qui devraient figurer dans une rédaction
finale. Pour parvenir à cette dernière, il faudra une réorganisation de
fond en comble  du présent Etat 0, la première tâche étant probablement
l'établissement d'un nouveau plan (dans lequel sans doute  les actuels
nºs 11,12,14,15  viendront beaucoup plus tôt). Je n'ai pas encore rédi\unsure{-}
le nº 16 , qui en principe ne doit offrir aucune difficulté ni
influer en aucune manière sur les nºs précédents, puisqu'il s'agit d'u\unsure{n}

\margin{2}

\clearpage
% Nouveau EGA V (cote 25): archive cote 25, site page 3 (PDF page 4).
% typescript. Draft compared: none.
% AI-transcribed by Claude Opus 5.5 (claude-opus-5-5, adaptive thinking, effort high), Message Batches API, batch ega25b (msgbatch_01A7n5YuXubUPNK1zsepi8RG), 2026-11-13; not yet checked by eye.
\scanpage{25/3}
\margin{\unsure{06}}

simple question de traduction.

Tu noteras la présence d'une proposition 10.3. qui devrait figu-%
rer dans un paragraphe antérieur. Je te signale à ce propos que j'ai
plusieurs autres résultats, assez disparates mais ayant tous traits à
des applications birationelles, que j'aimerais insérer quelque part.
Il me semble qu'il y en aura trop peu pour en faire un paragraphe, au\unsure{rai}
% NOTE: the end of this line runs into the right edge of the sheet; the last letters are cut and partly overtyped
tu une suggestion à faire où le fourrer ? Je pense te l'envoyer
prochainement, ainsi que le n\textsuperscript{o}16 des présentes notes.

D'ailleurs, le présent par.20 éclatera\xout{\illegible{1}} très probablement en
deux, l'un étant constitué par des dé\unsure{v}eloppements genre ``géométrie
élémentaire'' sur les grassma\unsure{nni}ennes. A la rigueur, on pourrait y fourrer
aussi (à défaut de meilleur endroit) les compléments dont je parle sur
les applications birationelles ?

\mspage{\hand{3}}

\clearpage
% Nouveau EGA V (cote 25): archive cote 25, site page 4 (PDF page 5).
% typescript. Draft compared: none.
% AI-transcribed by Claude Opus 5.5 (claude-opus-5-5, adaptive thinking, effort high), Message Batches API, batch pilot (msgbatch_0189Nhf1xB615YWnztubtLfc), 2026-11-12; not yet checked by eye.
\scanpage{25/4}
IV 20 \quad \uline{Sections hyperplanes et projections \deleted{linéaires}\inserted{coniques}}.

20.1. \quad \uline{Préliminaires \over{de}{\unsure{et}} notations}.

Soient S un \quad préschéma , $\underline{E}$ un Module loc libre de type fini sur S,
$\check{\underline{E}}$ son dual. Nous noterons par P=$\underline{P}(\underline{E})$ le fibré projectif défini par $\underline{E}$,
par $\check{P}$ celui défini par $\check{\underline{E}}$ . Nous\hand{|}appelerons \ $\check{P}$ \ le \uline{schéma des hyperplans}
de P, cette terminologie se justifiant de la façon suivante. \ Soit
$\hand{\xi}$ une section de $\check{P}$ sur S, qui est donc\hand{/}déterminée par un module quotient
inversible L de $\check{\underline{E}}$ . On en conclut un module quotient inversible
$\underline{L}_P$ de$(\check{\underline{E}})_P$ = $(\underline{E}_P)^{\hand{\vee}}$ , d'autre part on dispose d'un module quotient inversibl
$\underline{O}_P(1)$ de \ $\underline{E}_P$ . Passant \ aux duals, on peut envisager \ $\underline{L}_P^{-1}$ resp. $\underline{O}_P(-1)$
comme des sous-modules inversibles (loc facteurs directs) de \ $\underline{E}_P$ resp.
de $(\underline{E}_P)^{\hand{\vee}}$ , \ et l'accouplement \ $\underline{E}_P \otimes \check{\underline{E}}_P \longrightarrow \underline{O}_P$ \ définit \ donc un
accouplement naturel
(x) \[ \underline{O}_P(-1)\otimes\underline{L}_P^{-1} \to \underline{O}_P \]
ou encore un homomorphisme \ transposé
(xx) \[ \underline{O}_P \longrightarrow \underline{O}_P(1)\otimes\underline{L}_P \ = \ \unsure{\underline{L}_P(1)|\underline{D}_P(1)} \ , \]
i.e. une section de $\underline{L}_P(1)$, canoniquement déterminée par $\hand{\xi}$ . Le''diviseur''
de cette section, \ i.e. le sous-schéma\inserted{fermé} $H_{\hand{\xi}}$ de \ P défini par l'Idéal
image de (x), \ est \ appelé l'\uline{hyperplan dans P} \ défini par \ l'élément
$\hand{\xi} \hand{\in} \check{P}(S)$ . \ Il peut encore se décrire en notant que localement sur S \xout{la}
\xout{\unsure{donnée de $\check{\xi}$ équivaut}} \quad $\hand{\xi}$ \ est donné par une section \ $\phi$ de \ $\underline{E}$ , \xout{de} telle
\xout{\illegible{1}} que $\phi(s) \neq 0$ pour tout s \ ( $\phi$ étant déterminé par $\hand{\xi}$ \ modulo
multiplication par une section inversible de $\underline{O}_S$ ); comme $\underline{E}$ = $p_{\hand{*}}(\underline{O}_P(1))$
( p: P$\to$S étant la projection), \ $\phi$ \ peut être considéré comme une section
de \ $\underline{O}_P(1)$, don\over{c}{t} le diviseur n'est autre que $H_{\hand{\xi}}$. \ Bien entendu, si on
\inserted[typed]{\uline{considére}}\xout{\illegible{1}} $\underline{L}^{-1}$ \ comme un sous-\xout{\illegible{1}} module inversible de \ $\underline{E}$ locale
ment facteur direct dans $\underline{E}$, alors la correspondance entre $\hand{\xi}$ \ (i.e. $\underline{L}$ \hand{\unsure{$\xi$,}}

\mspage{\hand{4}}

\clearpage
% Nouveau EGA V (cote 25): archive cote 25, site page 5 (PDF page 6).
% typescript. Draft compared: none.
% AI-transcribed by Claude Opus 5.5 (claude-opus-5-5, adaptive thinking, effort high), Message Batches API, batch ega25b (msgbatch_01A7n5YuXubUPNK1zsepi8RG), 2026-11-13; not yet checked by eye.
\scanpage{25/5}
\mspage{2}

ou encore $\underline{L}^{-1}\subset\underline{E}$) et $\emptyset$ s'obtient en prenant\inserted{pour $\emptyset$} une section de $\underline{L}^{-1}$ qui
ne s'annule en aucun point \xout{\illegible{1}}, i.e. une trivialisation de $\underline{L}^{-1}$
(qui existe en tous cas localement). Notons d'ailleurs que $H_{\hand{\xi}}$ n'est
autre que $P(\underline{E}/\underline{L}^{-1})$ (isomorphisme canonique), ce qui est une troisième
façon de décrire $H_{\hand{\xi}}$ (NB $P(\underline{E}/\underline{L}^{-1})$ est bien plongé canoniquement dans
$P = P(\underline{E})$ ), qui a l'avantage de montrer de plus que $H$ est un fibré projec-%
tif sur $S$, et à fortiori est lisse sur $S$ (encore aurait-il fallu dire \xout{\unsure{qu}}
\margin{$\|$} au par. 17 qu'un fibré projectif est lisse ...). Il serait sans doute
\margin{$\|$} préférable de commencer par celle-là.
\uline{Remarques}. La formation de $H_{\hand{\xi}}$ en termes de $\hand{\xi}$ est compatible avec
le changement de base, comme on constate aussitôt., en d'autres termes,
on trouve un homomorphisme de foncteurs $(\mathrm{Sch})^{\circ}_{/S}\to(\mathrm{Ens})$
\[ \check{P} \longrightarrow \underline{\mathrm{Div}}\,(P/S) \qquad , \]
où le deuxième membre désigne le foncteur des ``diviseurs relatifs'' de $P/S$,
dont la valeur en $S'$ ($S$-préschéma arbitraire) est l'ensemble des
sous-schémas fermés de $P_{S'}$ qui sont intersection complète transversale
de codimension 1 relativement à $S'$ (cf par.19). Il est facile de montrer
que cet\xout{\illegible{1}} homomorphisme de foncteurs est un monomorphisme, en d'autres
termes que $\hand{\xi}$ est déterminé par $H_{\hand{\xi}}$ . (Ce dernier fait justifie la termino-%
logie ``schéma des hyperplans'' utilisée plus \xout{\unsure{a}}haut). Nous verrons \xout{\illegible{1}} \margin{)}
\xout{\unsure{Chap.V}} que le foncteur $\underline{\mathrm{Div}}$ $(P/S)$ est représentable par \xout{\unsure{un}}\inserted{le pré}schéma somme
des $P(\mathrm{Symm}^{k}(\check{E}))$, de sorte que $\overset{\hand{\vee}}{P}$ s'identifie à un sous-préschéma ouvert
et fermé de $\underline{\mathrm{Div}}$ $(P/S)$ ... (NB \over{A}{À} vrai dire, la d\over{e}{é}termination des diviseurs
\margin{$|$} relatifs\inserted{de $P/S$} pourrait se faire avec les moyens du bord dès maintenant, en
\margin{$|$} utilisant les résultats du Chap III, et pourrait être ajoutée \xout{\illegible{1}} à titre
\margin{$|$} d'exemple au par.19 ...).

\mspage{5}

\clearpage
% Nouveau EGA V (cote 25): archive cote 25, site page 6 (PDF page 7).
% typescript. Draft compared: none.
% AI-transcribed by Claude Opus 5.5 (claude-opus-5-5, adaptive thinking, effort high), Message Batches API, batch ega25b (msgbatch_01A7n5YuXubUPNK1zsepi8RG), 2026-11-13; not yet checked by eye.
\scanpage{25/6}
\mspage{3}

Faisons maintenant le changement de base $S'=\overset{\hand{\vee}}{P}\longrightarrow S$, et considérons
la section diagonale (ou\hand{``}section générique\hand{''}) de $\overset{\hand{\vee}}{P}_{S'} = P\,(\overset{\hand{\vee}}{\underline{E}}_{S'})$ sur $S'$ :
on trouve un sous-schéma fermé $\underline{H}$ de $P_{S'} = P\times_S\overset{\hand{\vee}}{P}$ , appelé parfois
\uline{schéma d'incidence} entre $P$ et $\overset{\hand{\vee}}{P}$ , défini par l'Idéal image de
l'homomorphisme canonique
\[ \underline{O}_P(-1)\otimes_S\underline{O}_{\overset{\hand{\vee}}{P}}(-1)\longrightarrow \underline{O}_{P\times_S\overset{\hand{\vee}}{P}} \qquad ; \]
d'après ce qu'on a vu, c'est un fibré projectif sur $\overset{\hand{\vee}}{P}$, et par raison de
symétrie, c'est également un fibré projectif sur $P$. Bien entendu, on
retrouve les\deleted{\xout{\unsure{définitions des}}}\inserted{hyperplans ``spéciaux''} $H_{\hand{\xi}}$ ( pour $\hand{\xi}$ section de $\overset{\hand{\vee}}{P}$ sur $S$ )\hand{,} à partir
du "hyperplan universel" $\underline{H}$\hand{,} en prenant son image inverse par l\over{a}{e} change-%
ment de base $S\xrightarrow{\hand{\xi}}\overset{\hand{\vee}}{P}$ . La même remarque s'applique pour tout point
$\hand{\xi}$ de $\overset{\hand{\vee}}{P}$ à valeurs dans un $S$-préschéma quelconque $S'$, qui (considéré comme
section de $\overset{\hand{\vee}}{P}_{S'}$ sur $S'$) permet de définir un $H_{\hand{\xi}}\subset P_{S'}$ : ce dernier n'est
autre que l'image inverse de $\underline{H}$ par le changement de base $S'\xrightarrow{\hand{\xi}}\overset{\hand{\vee}}{P}$.

Nous supposons par la suite donné un préschéma de type fini $X$ sur $S$,
et un $S$-morphisme
\[ f\colon X\longrightarrow P \quad . \]
\margin{$|$}Un des principaux objets\hand{,} du présent paragraphe est d'étudier, pour tout
\margin{$|$}hyperplan $H_{\hand{\xi}}$ de $P$, \inserted{$(\xi\in\overset{\vee}{P}(S))$} son image inverse
\[ Y_{\hand{\xi}} = f^{-1}(H_{\hand{\xi}}) = X\times_P H_{\hand{\xi}} \quad , \]
et notamment de mettre en relations le propriétés de $X$ et des $Y_{\hand{\xi}}$ . Com-%
me d'habitude, on aura également à considérer des $\overset{\hand{\vee}}{P}(S')$, $S'$ \xout{\illegible{1}}
un $S$-préschéma quelconque, dans ce cas $H_{\hand{\xi}}$ est un hyperplan de $P_{S'}$ \hand{,}
et on posera encore
\[ Y_{\hand{\xi}} = f_{S'}{}^{-1}(H_{\hand{\xi}}) = X_{S'}\times_{P_{S'}}H_{\hand{\xi}} = X\times_P H_{\hand{\xi}} \quad , \]
où l'indice $S'$ désigne comme d'\unsure{habitude|hhabxitude} \xout{\illegible{1}} l'effet du changement de base
$S'\longrightarrow S$, et où dans le dernier terme on considère $H_{\hand{\xi}}$ comme $P$-préschéma.

\mspage{6}

\clearpage
% Nouveau EGA V (cote 25): archive cote 25, site page 7 (PDF page 8).
% typescript. Draft compared: none.
% AI-transcribed by Claude Opus 5.5 (claude-opus-5-5, adaptive thinking, effort high), Message Batches API, batch ega25b (msgbatch_01A7n5YuXubUPNK1zsepi8RG), 2026-11-13; not yet checked by eye.
\scanpage{25/7}
\mspage{4}

par le morphisme composé $H_{\hand{$\xi$}} \longrightarrow P_{S'} \longrightarrow P$. Il est alors commode
\xout{\illegible{4}} encore de considérer le cas où \hand{$\xi$} est
``universel'' i.e. où $S' = \check{P}$ , et \hand{$\xi$} est la section diagonale, donc $H \ = \underline{H}$,
dans ce c\unsure{a}s on notera (sauf meilleures notations suggérées par Dieudonné)
$\underline{Y} = Y_{\hand{$\xi$}}$ . Dans le cas général d'un \hand{$\xi$}: $S' \to \check{P}$ , on aura donc aussi
\[ Y_{\hand{$\xi$}} = \underline{Y} \times_{\check{\underline{P}}} S' \ . \]
Enfin, si F est Module sur X, nous \unsure{déssignerons|désignerons} par $G_{\hand{$\xi$}}$ son image invers
% NOTE: the right edge of the page is cut off here; the end of "inverse" is not visible
sur $Y_{\hand{$\xi$}}$ , par $\underline{G}$ son image inverse sur $\underline{H}$ , de sorte qu'on aura aussi
\[ G_{\hand{$\xi$}} = \underline{G} \otimes_{\mathcal{O}_{\check{\underline{P}}}} \mathcal{O}_{S'} \ . \]
Résumons dans un petit diagramme l'essent\unsure{i}el des constructions \inserted{et notations} envisag
% NOTE: the right edge of the page is cut off at the end of this line

gées:
\begin{diagramblock}
\begin{tikzcd}
F & & \underline{G} & & G_{\hand{$\xi$}} \\
X \arrow[d] & X\times_S\check{P} \arrow[l] \arrow[d] & \underline{Y} \arrow[l] \arrow[d] & & Y_{\hand{$\xi$}} \arrow[ll] \arrow[d] \\
P \arrow[d] & P\times_S\check{P} \arrow[l] \arrow[d] & \underline{H} \arrow[l] \arrow[dl] & & H_{\hand{$\xi$}} \arrow[ll] \arrow[dl] \\
S & \check{P} \arrow[l] & & S' \arrow[ll] &
\end{tikzcd}
\end{diagramblock}
\margin{\illegible{1}}

(Les carrés et losanges figurant dans ce diagramme sont car\over{\illegible{1}}{té}siens).

Dans le numéro suivant, nous étudions systématiquement le cas
s\unsure{u}ivant: S' est le spectre d'un corps\inserted{K}, et son image dahs P est \uline{génériqu}
% NOTE: the right edge of the page is cut off at the end of this line
dans la fibre correspondante $\check{P}_s$ . Quitte à faire le changement de base
Spec k(s) $\longrightarrow$ S , on est ramené pour ce\xout{\illegible{1}} au cas où S est le spectre
d'un corps k, ce que nous supposerons au N\textsuperscript{o} suivant. D'ailleurs, la plu-%
part des propriétés examinées pour X et $Y_{\hand{$\xi$}}$ sont de ``nature \unsure{géométriqque|géométrique}''
donc invariantes par changement de corps de base; cela permet alors
indifféremment de se ramener au cas où K est algébriquement clos, ou au
cas où K = k(\hand{$\eta$}) , \hand{$\eta$} étant le point générique de $\check{P}$ , et \hand{$\xi$} : Spec(K)$\to\check{P}$
étant bien entendu le morphisme canonique. On notera de même que pour le
% NOTE: the right edge of the page is cut off at the end of this line
questions de nature géométrique concernant X, $Y_{\hand{$\xi$}}$ , on peut (quitte à fa
% NOTE: the right edge of the page is cut off at the end of this line
re un changement de base sur k) se borner au cas \inserted{k} alg clos.

\mspage{7}

\clearpage
% Nouveau EGA V (cote 25): archive cote 25, site page 8 (PDF page 9).
% typescript. Draft compared: none.
% AI-transcribed by Claude Opus 5.5 (claude-opus-5-5, adaptive thinking, effort high), Message Batches API, batch ega25b (msgbatch_01A7n5YuXubUPNK1zsepi8RG), 2026-11-13; not yet checked by eye.
\scanpage{25/8}
\margin{\illegible{1}}

\margin{$|$}Un oubli terminologique: lorsque f est une immersion, on appelle habituel%
lement Y\hand{$_\zeta$} une \uline{section hyperplane} de X (relative à l'immersion projective
f et à l'hyperplan H  ); il n'y a pas de raison de ne pas étendre cette
terminologie au cas f quelconque.

20.2. \uline{Etude de la section hyperplane\xout{m} générique}: \hand{\uline{propriétés locales}} .

Rappelons que maintenaht S=Spec(k), k un corps. Si \hand{$\eta$} est un point
de $\overset{\hand{$\vee$}}{P}$ , \xout{xxxx} et si \hand{$\zeta$}: Spec(\hand{$\eta$})$\to\overset{\hand{$\vee$}}{P}$ est le morphisme canonique, nous poserons
aussi H\hand{$_\eta$} , Y\hand{$_\eta$} , G\hand{$_\eta$} au lieu de H\hand{$_\zeta$} , Y\hand{$_\zeta$} , G\hand{$_\zeta$} . Dans le présent nº, \hand{$\eta$}
désignera toujours le point générique de $\overset{\hand{$\vee$}}{P}$ .
\uline{Proposition} 2.1. Supposons X irréductible\inserted{\deleted{\illegible{1}}}. Alors Y\hand{$_\eta$} est irr. \inserted{(\deleted{\illegible{1}})} ou vide. ,
et dans le premier cas il domine X. \inserted{En tous cas $\unsure{\overline{Y}}$ est \unsure{$_\eta$} irréductible.}
En effet, comme \uline{H}$\longrightarrow$ P est un fibré projectif, il en est de même de
\uline{Y}$\longrightarrow$ X, ce qui implique que \uline{Y} est irréductible X l'\unsure{étant}. Danc la fibre
générique Y   de \uline{Y} sur $\overset{\hand{$\vee$}}{P}$ est irréductible ou vide, et dans le premier cas,
son point générique est point générique de \uline{Y} , donc au dessus de celui de
X, cqfd.
\uline{Proposition 2}.2. Soit Z une partie de P. Alors son image inverse Z\hand{$_\eta$} dans
H\hand{$_\eta$} est vide si et seulement si tout point de Z est fermé. En particulier,
si Z est constructible, alors Z\hand{$_\eta$} = $\emptyset$ sss Z est fini.

On peut supposer Z réduit à un seul point z , et on est ramené à
prouver que l'\unsure{imalge|image} de H dans P est formé \inserted[typed]{exactement} des points non fermés de P.
Désignant par X l'adhérence de z, et utilisant 2.1. , on est ramené à
prouver que \over{X}{\unsure{Y}}\hand{$_\eta$} = $\emptyset$ sss X est \xout{\illegible{3}} \inserted[typed]{fini} (X étant un sous-schém%
fermé \xout{xxxx} de P). Remplaçant X par X$_{\text{k(\hand{$\eta$})}}$\hand{$\hookrightarrow$}P$_{\text{k(\hand{$\eta$})}}$ , \xout{xxxx} \inserted[typed]{le ``il faut''} résulte du fait
\margin{$|$}suivant, pour lequel nous devons avoir une référence, et qui mérite d'être
réex\unsure{mplicité|plicité} ici en lemme: si Y est une sectionhyperplane \uline{quelconque} de
X et si Y\hand{\unsure{$_\eta$}}=$\emptyset$ , alors X est fini ( en effet, X $\subset$ P - H est affine, et
projectif ...). Le "il suffit" est immédiat, par exemple en notant que
\uline{Y} est un fibré projectif de dimension relative n-1 sur X (\xout{x}n étant la
dim relative de P \xout{xxxxS)} et $\overset{\hand{$\vee$}}{P}$ sur S), donc X étant fini sur k, \uline{Y} est de

\mspage{\hand{8}}

\clearpage
% Nouveau EGA V (cote 25): archive cote 25, site page 9 (PDF page 10).
% typescript. Draft compared: none.
% AI-transcribed by Claude Opus 5.5 (claude-opus-5-5, adaptive thinking, effort high), Message Batches API, batch ega25b (msgbatch_01A7n5YuXubUPNK1zsepi8RG), 2026-11-13; not yet checked by eye.
\scanpage{25/9}
\mspage{6}

dimension absolue n-1\inserted{$< n = \dim \check{P}$}, donc le morphisme $\underline{Y}\longrightarrow \overset{\hand{$\vee$}}{P}$ ne peut être dominant,
donc sa fibre générique $Y_{\hand{$\eta$}}$ est vide.

\uline{Corollaire} 2.3. Soit $f:X\longrightarrow P$ un morphisme de type fini, et soit Z une
partie constructible de X. Pour que son image inverse dans $Y_{\hand{$\eta$}}$ soit vide,
il f et s que l'image f(Z) soit finie. En particulier, pour que \xout{X} $Y_{\hand{$\eta$}}$ soit
vide, il f et s que f(X) soit fini.

\uline{Corollaire} 2.4. Soient Z,Z' deux parties fermées \deleted{irréductibles}\inserted{(avec Z irréductible)} de X,
et soient $Z_{\hand{$\eta$}}$ et $Z'_{\hand{$\eta$}}$ leurs images inverses dans $Y_{\hand{$\eta$}}$ . Pour qu'on ait
$Z_{\hand{$\eta$}}\subset Z'_{\hand{$\eta$}}$ , il faut et il suffit que f(Z) soit fini ou que l'on ait $Z\subset Z'$.
Pour que $Z_{\hand{$\eta$}} = Z'_{\hand{$\eta$}}$ , il f et s que \xout{ZxxxxZ} f(Z) et f(Z') soient finis, ou \inserted{que $Z=Z'$}
\xout{\unsure{En effet, il faut exprimer}} C'est une conséquence \xout{\unsure{facile de 2.}}
\xout{\unsure{2.3}} immédiate de 2.3. , car on voit que $f(Z - Z\cap Z')$ ne peut être fini
que si $Z\subset Z'$, ou si f(Z) est fini (car si on n'a pas $Z\subset Z'$, alors $Z-Z\cap Z'$
est dense dans Z, donc $f(Z-Z\cap Z')$ est dense dans f(Z), et si le premier
est fini donc fermé - étant constructible - il en est de même du second.

\uline{Corollaire} 2.5. A toute composante irréductible $X_i$ de X telle que
$\dim \overline{f(X_i)} > 0$ , faisons correspondre \xout{\illegible{1}} son image inverse\xout{s} $Y_{i\hand{$\eta$}}$ dans
$Y_{\hand{$\eta$}}$ . Alors $Y_{i\hand{$\eta$}}$ est une composante irréductible de $Y_{\hand{$\eta$}}$ , et on obtient de
cette façon une correspondance biunivoque entre l'ensemble des comp. irr.
$X_i$ de X telles que $\dim \overline{f(X_i)}> 0$, et l'ens des comp \unsure{\xout{irr}|irr} de $Y_{\hand{$\eta$}}$ .

En effet, il résulte de 2.3. que \xout{X} $Y_{\hand{$\eta$}}$ est réunion des $Y_{i\hand{$\eta$}}$ \unsure{envisagés|envisagé}
et que ces derniers sont des parties fermées non vides de Y; elles sont
d'ailleurs irréductibles en vertu de 2.1. Enfin, elles sont mutuellement
sans relation d'inclusion en vertu de 2.4. , d'où la conclusion.

Notons que si \xout{\unsure{$d_i$}xxx} $\dim X_i$\inserted[typed]{$= d_i$}, \xout{xxxxx} on a $\dim Y_i = d_i-1$ , \deleted{comme
il résulte de la proposition} plus générale\hand{ment}:
\uline{Proposition} 2.6. Supposons que \inserted{pour toute comp. irréd. $X_i$ de X on ait} $\dim \overline{f(X_{\hand{$i$}})} > 0$ i.e. $Y_{\hand{$i\eta$}} \neq \emptyset$\inserted{ou que f soit une immers. et $\dim \overline{f(X)}>0$}. Alors on a
\[ \dim Y_{\hand{$\eta$}} = \dim X - 1 \qquad . \]

\mspage{9}

\clearpage
% Nouveau EGA V (cote 25): archive cote 25, site page 10 (PDF page 11).
% typescript. Draft compared: none.
% AI-transcribed by Claude Opus 5.5 (claude-opus-5-5, adaptive thinking, effort high), Message Batches API, batch ega25b (msgbatch_01A7n5YuXubUPNK1zsepi8RG), 2026-11-13; not yet checked by eye.
\scanpage{25/10}
\mspage{7}

On est ramené au cas où X est irréductible, grâce à 2.5. Par construction
même, Y est défini à partir de $X_{k(\hand{\eta})}$ comme le diviseur d'une section
d'un module invers\unsure{i|a}ble sur $X_{k(\hand{\eta})}$ (savoir de l'image inverse de $\underline{O}_P(1)$))
D'autre part, $X_{k(\hand{\eta})}$ est irréductible (car X l'est, et $k(\hand{\eta})$ est une
\margin{|} extension transcendante pure de k, ce qu'on aurait dû signaler au début du
N° ...), et $Y_{\hand{\eta}} \neq X_{k(\hand{\eta})}$ , puisque l'image de $Y_{\hand{\eta}}$ dans X (contrairement
à celle de $X_{k(\hand{\eta})}$ , qui est fid. plat sur X) n'est pas égale à X, ne conte-%
nant pas en effet les points fermés de X en vertu de 2.3. Il s'ensuit que
dim $Y_{\hand{\eta}}$ = dim $X_{k(\hand{\eta})}$ -1 \inserted{\unsure{(n°)|(nf)}} = dim X - 1 , cqfd.

\uline{Proposition }2.7. Soit F un Module quasi-cohérent sur X, d'où $G_{\hand{\eta}}$ sur $Y_{\hand{\eta}}$ .
Soient $Z_i$ les cycles premiers associés à F tels que dim $\hand{\overline{f(Z_i)}}$ $>$ 0. Soit
$Z_{i\hand{\eta}}$ l'image inverse de $Z_i$ dans $Y_{\hand{\eta}}$ , alors les $Z_{i\hand{\eta}}$ sont exactement tous
les cycles premiers associés à $G_{\hand{\eta}}$ , de plus leurs relations d'inclusion
sont celles entre les $Z_i$ .

La dernière assertion est contenue dans 2.4. D'autre part, comme
$\underline{Y} \rightarrow X$ est un fibré projectif, donc plat\xout{b} à fibres $(S_1)$ et irr. , il résul
te du par.3 que les cycles premiers associés à l'image inverse \inserted{G} de F sur
$\overline{Y}$ sont les images inverses \xout{xxxxxxxx} des cycles premiers associés à F,
d'où en \xout{prenant} \inserted[typed]{s'induisant sur} la fibre générique \inserted[typed]{$Y_{\hand{\eta}}$} de $\underline{Y}$ sur $\check{P}$ , le fait que les
cycles premiers associés à $G_{\hand{\eta}}$ sont les images inverses \uline{non vides} des $Z_i$ ,
ce qui établit 2.7. moyennant 2.3. A vrai dire, le passage par $\underline{Y}$ est
inutile, on peut utiliser directement le fait que $Y_{\hand{\eta}}$ \hand{$\longrightarrow$} X est\xout{n} plat à
fibres $(S_1)$ (et même géom. régulières, i.e. le morphisme est régulier)
et à fibres irréductibles (et même géométriquement irréductibles : c'est
des localisés de schémas projectifs)\hand{;} \hand{\unsure{remarque} \unsure{analogue} \unsure{pour} \unsure{la}
dim. de 2.\unsure{1}.}

\margin{10}

\clearpage
% Nouveau EGA V (cote 25): archive cote 25, site page 11 (PDF page 12).
% typescript. Draft compared: none.
% AI-transcribed by Claude Opus 5.5 (claude-opus-5-5, adaptive thinking, effort high), Message Batches API, batch ega25b (msgbatch_01A7n5YuXubUPNK1zsepi8RG), 2026-11-13; not yet checked by eye.
\scanpage{25/11}
\mspage{8}

\uline{Proposition }2.8.  Soit F cohérent sur X, et \unsure{soit} $\over{x}{y}\in Y$ , x son image
dans X. \xout{\unsure{Pxxxxquexx\&xxxxxxx}}  Soit P\inserted{(M)} une des propriétés suivantes pour
un module de type fini\inserted{M} sur un anneau local noethérien A:

(i)   coprof $M\leqslant n$   (réf).

(ii)  M satisfait$(S_k)$  (réf).

(iii)  M est \unsure{Cohen-Macaulay|CMahn-Macaulay}.

(iv)  M est réduit(\hand{\unsure{réf|vif}}).

(v)   M est intègre(\hand{\unsure{réf|vif}}).

Alors, pour que $\over{F}{G}_{\hand{\nu},y}$ satisfasse la propriété P, il f et s que $F_x$ la
satisfasse.

Cela résulte aussitôt des résultats du par.6 , compte tenu que
\xout{X} $Y_{\hand{\nu}}\longrightarrow X$  est un morphisme régulier, \unsure{donc}  $\underline{O}_{X,\unsure{x}}\longrightarrow \underline{O}_{\unsure{Y}_{\hand{\nu}},y}$ est
régulier.  Compte tenu de 2.3. on obtient donc:

\uline{Corollaire} 2.\unsure{9}.  Avec les notations de 2.8. , soit \deleted{$\unsure{E_P(F)}$}\inserted{Z} l'ensemble
des $x\in X$ tels que $P(F_x)$ ne soit pas satisfaite. \unsure{Alors} pour que
$G_{\hand{\nu}}$ satisfasse à P en tous les points, il f et s que $f(\unsure{Z})$ soit une
partie finie de P, ou encore que \xout{f(Z)} dim $\overline{f(Z)}$ = 0.

En effet, 2.8. nous dit que $\unsure{h}^{-1}(Z)$ est l'ensemble P-singulier
de $G_{\hand{\nu}}$ , et il est vide sss f(Z) est fini en vertu de 2.3. (NB h désig
\margin{|} ne le morphisme $Y_{\hand{\nu}}\longrightarrow X$; je m'aperçois à l'instant que la lettre P
\margin{|} dans 2.8. fait double emploi ).

\margin{|} \uline{Corollaire} 2.\unsure{10}.  Condition pour que $Y_{\hand{\nu}}$ soit réduit resp. loc intègre.

\deleted{On fera attention que 2.10. \inserted[typed]{et 2.11. ci-dessous} ne dit rien sur les propriétés
``géométriquement réduit'' et ``géométriquement intègre'', les propriétés
envisagées dans 2.8. étant en effet ``absolues'' . Nous \unsure{allons}}

\margin{11}

\clearpage
% Nouveau EGA V (cote 25): archive cote 25, site page 12 (PDF page 13).
% typescript. Draft compared: none.
% AI-transcribed by Claude Opus 5.5 (claude-opus-5-5, adaptive thinking, effort high), Message Batches API, batch ega25b (msgbatch_01A7n5YuXubUPNK1zsepi8RG), 2026-11-13; not yet checked by eye.
\scanpage{25/12}
\mspage{\hand{9}}

\uline{Corollaire} 2.11. Soit $y\in Y_{\hand{y}}$ , pour que $Y_{\hand{y}}$ \inserted{soit régulier (resp.} satisfasse la propriété
$R_k$ (réf) en y\hand{,} \deleted{(}resp. soit normal en y) il f et s que X satisfasse
la même propriété en x. \xout{\unsure{Remarque}} Soit Z l'ensemble des points de X en
lesquels X n'est pas \deleted{\xout{\illegible{1}} (}\inserted{régulier (resp. \deleted{$R_k$} $R_k$}resp. \deleted{n'est pas} normal); pour que $Y_{\hand{y}}$ soit
\inserted{régulier (resp.} $(R_k)$\hand{,} \deleted{(}resp. normal) il f et s que f(Z) soit fini i.e. dim $\overline{f(Z)}=0$.

\margin{|} Même démonstration que 2.8. et 2.9. Il faut donner les diverses
références assurant que Z est fermé (car il faut savoir qu'il est
constructible pour appliquer 2.3.).

\xout{\unsure{Remarque 2.12}} On notera que dans 2.10. et 2.11. nous ne disons rien
sur les propriétés ``géométriques'' correspondantes, les résultats
envisagés étant de nature \hand{``}absolue\hand{''}. Nous Examinons maintenant les
propriétés de nature géométrique. (On pourrait si on le désire en pro-%
\margin{||}fiter pour changer de n$^{\text{o}}$).

\inserted[typed]{Supposons que $f:X\to P$ soit une \deleted{\unsure{immersion}}\inserted{\unsure{non ramifié|non ramifiée}}.}
\uline{Théorème} 2.12. Soient $y\in Y_{\hand{y}}$ , x son image dans X. Pour que X soit
lisse sur k en x, il f et s que $Y_{\hand{y}}$ soit lisse sur $k(\hand{y})$ en y.

On peut supposer k algébriquement clos. Si $Y_{\hand{y}}$ est lisse sur $k(\hand{y})$
en y, il est régulier, donc comme $Y_{\hand{y}}$ estb plat sur X, X est régulier
en x (réf), donc il est lisse sur k en x puisque k est alg clos donc
parfait (réf). Pour la réciproque, on peut (quitte à remplacer X par
un voisinage ouvert de x) supposer X lisse, et (grâce au critère jaco-%
bien de lissité) défini dans un ouvert U de P par p équations
comme $X=V(f_1,..,f_p)$ , où les différentielles $df_1,..,df_p$ sont
partout linéairement indépendantes. Introduisant d es coordonnées affines
$S_1,..,S_n$ dans $\overset{\hand{\vee}}{P}$ et des coordonnées affines $T_1,..,T_n$ dans un voisinage
de x (par choix d'un hyperplan $H_\infty$ ne contenant pas x), $Y\hookrightarrow U_{k(\hand{y})}$ est
alors égal à $V(f_1,..,f_p,\hand{\sum_i}S_iT_i-1\ )$ , et il suffit de vérifier que
les différentielles (relativement à $k(\hand{y})$ ) de $f_1,..,f_p$ ,$\hand{(\sum}S_iT_i\hand{)}-1$

\mspage{\hand{12}}

\clearpage
% Nouveau EGA V (cote 25): archive cote 25, site page 13 (PDF page 14).
% typescript. Draft compared: none.
% AI-transcribed by Claude Opus 5.5 (claude-opus-5-5, adaptive thinking, effort high), Message Batches API, batch ega25b (msgbatch_01A7n5YuXubUPNK1zsepi8RG), 2026-11-13; not yet checked by eye.
\scanpage{25/13}
\mspage{\hand{10}}

sont linéairement indépendantes. Or ces différentielles ne sont autres
que les sections de $(\underline{\Omega}^1_{U/k})\otimes_k k(\ )$ suivantes:
\[ df_1, \ldots, df_p , \sum_{\hand{\unsure{i}}} S_i dT_i \ . \]
Comme les $df_i$ sont linéairement indépendants en tout points de U, et
que les $dT_i$ forment une base \inserted{de $\underline{\Omega}^1_{U/S}$} en tout point de U, \xout{\illegible{1}} et à forti-%
ori un système de générateurs, on en conclut immédiatement l'indépendan-%
ce linéaire des quantités écrites en tout point de $U_{k(\hand{\eta})}$ , du moins
lorsque $p\leq n$, i.e. si \inserted{$\underline{E} =$} $\underline{\Omega}^1_{U/k} / \sum \underline{O}_U df_i \neq 0$ : c'est un petit
lemme sur une famille de \unsure{gébérateurs|générateurs} \inserted{$\unsure{a_i}$ $(1 \leq i \leq n)$} d'un Module loc libre $\underline{E}$ non
nul, alors $\sum S_i a_i$ considéré comme section de $\underline{E}\otimes_k k(\hand{\eta})$ ne s'annule
en aucun point. D(autre part, le cas où p=n est trivial, car alors
$Y_{\hand{\eta}} = \emptyset$ .

\uline{Corollaire} 2.13. Soit Z l'ensemble des points de X en lesquels X
n'est pas lisse sur k. Pour que Y soit lisse sur $k(\hand{\eta})$, il f et s
que Z soit fini. En particulier, si X est lisse, il en est de même de
Y .

Résulte de 2.12. et 2.3. Plus généralement, on obtient:

\margin{\unsure{J un corps} \illegible{1} \unsure{à k} \unsure{de car.}}
\uline{\unsure{Théoèème|Théorème}} 2.1\over{3}{4}. Soi\xout{e}nt \unsure{y} un point de $\over{X}{Y}_{\hand{\eta}}$ , x son image dans X. Soit
\inserted[typed]{$P(\underline{A},K)$} \deleted{P} l'une quelconque des propriétés suivantes pour une\xout{\illegible{1}} algèbre
locale noethérienne sur un corps K:

(i) A est géométriquement régulier sur K.

(ii) A est géométriquement $(R_k)$ sur K.

(iii) A est séparable sur K.

(iv) A est \inserted{géométriquement} normal sur K.

\inserted{Alors $P(\underline{O}_{X,x}, k) \Longleftrightarrow P(\mathcal{O}_{Y_\eta,y}, k(\eta))$}
\xout{(x)xxxxxAxxxx} En effet, compte tenu du par.6 , (iii) et (iv)
résultent de (ii) conjugué avec 2.8. (ii). D'autre part, (i)
a déjà été prouvé dans 2.12. , et il reste à ramener (ii) à (i). Pour

\mspage{\hand{13}}

\clearpage
% Nouveau EGA V (cote 25): archive cote 25, site page 14 (PDF page 15).
% typescript. Draft compared: none.
% AI-transcribed by Claude Opus 5.5 (claude-opus-5-5, adaptive thinking, effort high), Message Batches API, batch ega25b (msgbatch_01A7n5YuXubUPNK1zsepi8RG), 2026-11-13; not yet checked by eye.
\scanpage{25/14}
\mspage{\hand{11}}

ceci, soit Z l'ensemble des points en lesquels X n'est pas lisse sur
k, son image inverse \inserted{$Z_{y}$ dans $Y_{y}$} est donc (en vertu de 2.12.)   l'ensemble des
points de $Y_{\hand{y}}$ en lesquels $Y_{\hand{y}}$ n'estb pas lisse sur k($\hand{y}$). Or  la codimen-%
sion de  Z dans X en x  est égale à celle de  $Z_{\hand{y}}$ dans $Y_{\hand{y}}$  en y par
raison de platitude (réf par.6). Donc l'une est \hand{$>$}k  sss l'autre l'est,
ce qui achève la démonstration.
\uline{Corollaire} 2.15.   Soit Z l'ensemble des points de X en lesquels
X n'est pas lisse sur k (resp. n'est pas  géom $R_k$ sur k, resp. n'est
\margin{\unsure{//}} pas séparable sur \unsure{R|k} , resp. n'est pas \inserted{géom} normal sur k). Pour que $Y_{\hand{y}}$ soit
lisse (resp.  géom $R_k$, resp séparable, resp géom normal) sur k($\hand{y}$ ), il
f et s  que Z soit fini.

\margin{$|$} Du point da vue rédaction, 2.14. et 2.15.  contiennent 2.1\unsure{2|3}. et
2.13. (qu'on pourrait donc se dispenser d'énoncer séparément), d'autre
part le corollaire est plus important pratiquement que le théorème,
qu'on pour\unsure{r}ait donner en proposition ou lemme préliminaire, pour glori-%
fier plutôt le corollaire.

Nous pouvons  moduler \xout{\unsure{xxxxparxi}}  le cas de la propriété (iii):
\uline{Corollaire} 2.16.  (Supposant toutjours f une immersion \inserted{et de plus F cohérent}), sous les condi-%
\margin{\unsure{//}} tions de 2.7. , pour que $Z_i$ soit non immergé, il f et s que $Z_i$  le soit.
Lorsqu'il en est ainsi, la mu\hand{l}tiplicité radicielle de F en $Z_i$ \xout{(ref)}\deleted{rel} \inserted{rel.}
à k (réf) est égale à celle de $G_{\hand{y}}$  en $Z_i$  rel à k($\hand{y}$).

La première assertion est contenu dans la dernière assertion de
2.7. Pour la seconde, comme  $Y_{\hand{y}} \longrightarrow X$ est plat donc le foncteur F \hand{$\rightsquigarrow$} $G_{\hand{y}}$
exact,  on est ramené, par restriction à un voisinage du pt gén de $Z_i$
et dévissage, au cas où F = $\underline{O}_{Z_i}$ \inserted[typed]{et on peut supposer $Z_i = X$.} . D'ailleurs on peut commencer par
se réduire au cas k alg clos, \inserted[typed]{alors X est séparable sur k.} Alors \xout{\unsure{$\underline{O}$xxxxx$\underline{O}_{Z_i}$xxxxx}} l'assertion  est
contenue dans 2.15. (iii).   On en conclut comme d'hab itude:

\margin{14}

\clearpage
% Nouveau EGA V (cote 25): archive cote 25, site page 15 (PDF page 16).
% typescript. Draft compared: none.
% AI-transcribed by Claude Opus 5.5 (claude-opus-5-5, adaptive thinking, effort high), Message Batches API, batch ega25b (msgbatch_01A7n5YuXubUPNK1zsepi8RG), 2026-11-13; not yet checked by eye.
\scanpage{25/15}
\mspage{12}

\uline{Corollaire} 2.17.   Soit Z l'ensemble des points de X en lesquels F
n'est pas séparable sur k (réf). Alors  $G_{\hand{\unsure{y|\eta}}}$ est séparable sur $k(\hand{\unsure{y|\eta}})$ sss
Z est fini, En particulier, si F est séparable sur k, $G_{\hand{\unsure{y|\eta}}}$ est séparable
sur $k(\hand{\unsure{y|\eta}})$.

\margin{\unsure{bb pour|à mettre} dans l'énoncé 2.12.}\uline{Remarques} 2.18.  Le résultat clef  2.12. \inserted{(et ses conséquences \unsure{2.13.|2.11.} à 2.17.)} devient faux si on abandonne
l'hypothèse que f est une immer\unsure{z|s}ion, comme  on voit par exemple en
prenant pour  X un revêtement purement inséparable de P.  \xout{Cxxx} Cepen-%
dant \xout{lexrésultatxsuivantxx(quixauraitxdûxvenirxplûtôtxenxconclusion}
\xout{àx2x11xx)xxxxxSixxX} \xout{\unsure{Propxx}} si  k est \uline{de caractéristique nulle} ,  2.12. à
2.1\hand{7}. sont valables sans supposer f une immersion.  En effet, il suffit
de le vérifier pour 2.12. , et  cela résulte alors de 2.11. et du fait
que \xout{xxxxx} pour un préschéma algébrique en c ar 0,  lisse = régulier.

\deleted{Les propriétés envisagées jusqu'à maintenant étaient de nature
essentiellement locale. Nous allons terminer ce n° en étudiant les
propriétés globales}

\over{\unsure{3}}{2}.3.   \uline{Section hyperplane générique}: \uline{irréductibilité et connexité géomét.}
\uline{Théorème} 3.1. (Bertini-Zariski).  Supposons dim $f(X) \hand{\geq} 2$; et que X
soit \xout{géométriquementxintègrex(resp.x} géométriquement irréductible\deleted{)}. Alors
la section hyperplane générique $Y_{\hand{\unsure{y|\eta}}}$ a la même propriété.

\deleted{\xout{x}(Compte tenu que X séparable \hspace{3em} Y séparable, les deux énoncés
de 3.1. sont en fait équivalents: \xout{Onxpeutxdxailleursxsupposerx}  le
deuxième implique le premier, et inversement comme on voit en se rame-%
nant au cas k alg. clos.  Compte tenu de 2.1. et exmplicitant
la construction du corps des fonctions de Y \hspace{2em} comme extension de k( \hspace{1em} )
en termes de l'extension K de k et des éléments $t_1,\ldots,t_n$ de K images
d'un système de coordonnées affines convenable de}

\mspage{15}

\clearpage
% Nouveau EGA V (cote 25): archive cote 25, site page 16 (PDF page 17).
% typescript. Draft compared: none.
% AI-transcribed by Claude Opus 5.5 (claude-opus-5-5, adaptive thinking, effort high), Message Batches API, batch ega25b (msgbatch_01A7n5YuXubUPNK1zsepi8RG), 2026-11-13; not yet checked by eye.
\scanpage{25/16}
\mspage{13}

\deleted{\uline{Corollaire} 2.8. Sous les conditions de 2.7. , pour que $Z_i$ soit non
immergé, il f et s que $Z_i$ \inserted[typed]{sur \unsure{$\bar{k}$}} soit non immergé. S'il en est ainsi, la
multiplicité radicielle de $F$ en $Z_i$ est égale à la multiplicité radici-%
elle de $G$ sur $k(\ \ )$ en $Z_i$ .

La première assertion résulte aussitôt de la dernière assertion
de 2.7. Quant à la première, on peut supposer\xout{, (comme il s'agit d'une}
\inserted[typed]{par restriction à un ouvert de X et}
\xout{\unsure{notion géométrique}} par dévissage sur $F$ (comme $Y$ est plat sur $X$
donc $F\ \ G$ est une foncteur exact) que l'on a $F=\underline{O}_{Z_i}$ , et on peut
évidemment remplacer $X$ par $Z_i$ . On est donc ramené au cas où $X$ est intè-%
gre et $F=\underline{O}_X$.} Soit $K/k$ le\xout{s} corps des fonctions de $X$, et soit $n=\dim P$;
introduisant des coordonnées \xout{projec} affines $T_1,..,T_n$ dans $P$ (en choisiss
ant un hyperplan à l'infini $H_\infty$ tel que $f(X)$ ne soit pas contenu dans
celui-ci) et $S_1,...S_n$ dans $\check{P}$ , on voit que le corps\deleted{\unsure{ex}}\inserted{des} fonctions $L$
de $Y_{\hand{\unsure{y}}}$ \xout{xxx} s'identifie au corps des fractions de l'anneau intègre
$K[S_1,..,S_n]/(\sum t_iS_i - 1)$ , où les $t_i \in K$ sont les images des $T_i$ à
l'adde de $f:X\to P$ . Comme $\dim \overline{f(X)} > 0$ , les $t_i$ ne sont pas tous
algébriques sur $k$, à fortiori ils ne sont pas tous nuls, supposons
p. ex. $t_n \neq 0$. On constate alors aussitôt qu'on a $L = K(S_1,..,S_{n-1})$ ,
(extension transcendante pure), $S_n\in L$ étant donné par l'équation
\[ \sum t_iS_i - 1 = 0 \]
en fonction des $S_i$ $(1\leq i\leq n-1)$ et des $t_i$ $(1\leq i\leq n)$. D'autre part,
$k' = k(\hand{y})$ s'identifie à $k(S_1,..,S_n)$, et l\xout{a}'inclusion canonique
$k'\longrightarrow L$ s'obtient en envoyant $S_i$ sur $S_i$ , i.e. $k'$ en tant que sous-%
\xout{\unsure{exten}}\inserted[typed]{extension} de $L$ est la sous-extension engendrée par les $S_i$ $(1\leq i\leq n)$, ou
ce qui revient évidemment au même, par les $S_i$ $(1\leq i\leq \xout{n}\inserted[typed]{n-1}\hand{)}$ et par
$S_n = a_o + a_1S_1 + \ldots + a_{n-1}S_{n-1}$ , où $a_o = t_n^{-1}$ , $a_i =-t_it_n^{-1}$ pour
$1\leq i\leq n-1$. \deleted{Comme}\inserted{Notons que} le corps engendré par les $a_i$ ou par les $t_i$ est évidem
ment le même, \deleted{les $a_i$ \unsure{ne} sont également pas tous algébriques sur k.}\inserted{leur degré de transcendance commun
n'est autre que $\dim f(X)$.}

\margin{16}

\clearpage
% Nouveau EGA V (cote 25): archive cote 25, site page 17 (PDF page 18).
% typescript. Draft compared: none.
% AI-transcribed by Claude Opus 5.5 (claude-opus-5-5, adaptive thinking, effort high), Message Batches API, batch ega25b (msgbatch_01A7n5YuXubUPNK1zsepi8RG), 2026-11-13; not yet checked by eye.
\scanpage{25/17}
\mspage{14}

(NB il y aurait lieu d'inclure cette description birationnelle au moins
en corollaire à 2.1. ). La démonstration de \over{\unsure{2}}{3}.\over{\unsure{2}}{1}. est donc ramenée à
celle du

\uline{Lemme} \over{\unsure{2}}{3}.\over{\unsure{2}}{1}.1. \inserted{(Zariski)} Soient k un corps, K une extension de type fini de k, \inserted[typed]{m un entier $\geqslant 0$}
$a_i$ $(0\leqslant i\leqslant m)$ des éléments de K \deleted{qui ne sont pas tous algébriques sur} \inserted{tels que le degré de transcendance de $k(a_0,-,a_m)$ sur k soit $\geqslant 2$,}
\deleted{k,} soit $L = K(S_1,..,S_m)$ et k' le sous-corps $k'= k(S_1,..,S_m,a_o +$
$\hand{$\sum_{i}^{m}$}a_iS_i)$ de L (les $S_i$ étant des indéterminées). \deleted{Alors la multiplicité
radicielle de K sur k est égale à celle de L sur k'.}\hand{,} Si K est une
extension primaire de k, alors L est une extension primaire de k'.

\margin{$|$} Ce lemme\hand{,} ou des lemmes qui lui ressemblent comme un frère\hand{,} trai-%
nent un peu partout dans la littérature, c'est pourquoi je te laisse le
choix de l'endroit où tu copieras une démonstration (i.e. je ne me sens
pas inspiré pour en trouver une par mes propres moyens).

\uline{Corollaire} 3.2. Supposons que f soit \deleted{une immersion} \inserted{\unsure{non ramifié}} ou k de car 0,
et que dim f(X)\over{$>$}{$\geqslant$} 2. Alors si X est géom intègre, il en est de même de
Y\hand{$_y$} .

En effet, géom intègre = géom irréd. + séparable.

\uline{Corollaire} 3.3. Supposons \inserted[typed]{k alg. clos, et} que pour toute composante irréductible $X_i$
de X, on ait dim $f(X_i)$\over{$>$}{$\geqslant$} 2, et supposons de plus que X soit \hand{$\unsure{\mathcal{S}}$}-connexe,
où \hand{$\unsure{\mathcal{S}}$} est l'ensemble des parties fermées Z de X telles que dim $\overline{f(\unsure{Z})}=0$\xout{\unsure{x}}
(i.e. pour tout tel Z, X-Z est connexe). Sous ces conditions, Y\hand{$_y$} est
géom connexe sur k(\hand{$y$}).

En effet, par un lemme qui doit figurer au par.6 avec le th. de
Hartshorme, l'hypothèse signifie que l'on peut joindre deux comp. irr.
quelconques X'\xout{\illegible{1}} et X" de X par une chaine de comp irr. $X_o$=X', ...
$X_n$=X", telles que deux consécutives aient une intersection \hand{$\notin\unsure{\mathcal{S}}$}; alors
les images inverses X'\hand{$_y$} et X"\hand{$_y$} sont joints par la chaine des $X_{i}$\hand{$_y$} \xout{\unsure{x}}

\mspage{17}

\clearpage
% Nouveau EGA V (cote 25): archive cote 25, site page 18 (PDF page 19).
% typescript. Draft compared: none.
% AI-transcribed by Claude Opus 5.5 (claude-opus-5-5, adaptive thinking, effort high), Message Batches API, batch ega25b (msgbatch_01A7n5YuXubUPNK1zsepi8RG), 2026-11-13; not yet checked by eye.
\scanpage{25/18}
\mspage{15}

\inserted[typed]{sur $k(\hand{\eta})$.}
qui sont géom connexes en vertu de 3.1. , et \xout{\illegible{1}} l'intersection
de deux consécutifs étant $\neq \emptyset$\xout{\unsure{0}} en vertu de 2.3. Il s'ensuit (comme
$Y_{\hand{\eta}} = X_{\hand{\eta}}$ est la réunion des $X_{i\hand{\eta}}$ , $X_i$ parcourant l'ensemble des comp
irr de X ) que \xout{\illegible{1}} $Y_{\hand{\eta}}$ est géom connexe sur $k(\hand{\eta})$, cqfd.

2\unsure{0|D}.4. \uline{Etude d'une section hyperplane variable}: \hand{sections ``\uline{assez générales}''} .
4.1.
Nous revenons à la situation générale du nº 1, S étant un préschéma
quelconque. De plus, nous supposons X de présentation finie sur S.
De façon générale, notons que si P(Z,k) est une propriété''constructib-%
le'' \xout{\unsure{de préschéma sur}} d'un préschéma algébrique Z sur le corps k,
alors l'ensemble des $\hand{\zeta} \hand{\in} \overset{\hand{\vee}}{P}$ tels qu'on ait $P(Y_{\hand{\zeta}},k(\hand{\zeta}))$ est \uline{constructi-%
ble}, comme on voit en notant que $Y_{\hand{\zeta}}$ est la fibre en $\hand{\zeta}$ de $\underline{Y} \longrightarrow \overset{\hand{\vee}}{P}$\hand{,}qui
est un morphisme de prés finie, et en appliquant le par.9 . On a une
remarque analogue pour une\hand{|}propriété P(Z,M,k), où Z,k sont comme dessus
et M est un Module cohérent sur Z: si F est \deleted{cohérent}\inserted{de prés. finie} sur X, l'ensemble
des $\hand{\zeta} \hand{\in} \overset{\hand{\vee}}{P}$ tels qu'on ait $P(Y_{\hand{\zeta}},G_{\hand{\zeta}},k(\hand{\zeta}))$ est constructible. D'autre
part, dans les deux numéros suivants, on a développé \xout{\illegible{1}} dans divers
cas un critère pour que l'ensemble précédent E contienne le point
générique de $\overset{\hand{\vee}}{P}$ , S étant le spectre d'un corps k ; étantconstructible,
il en résulte que E contient un ouvert non vide: c'est le passage d'une
conclusion pour la \xout{x}section hyperplane générique à la conclusion analo\xout{g}
gue pour les sections hyperp\xout{h}lanes ``assez générales''. \xout{Enfin,} \hand{[}\deleted{Nous donne-%
rons également dans ce nº des critères permettant d'affirmer que E est
ouvert, en appliquant les résultats du par.12}\hand{]}. Noter d'ailleurs que \margin{$\downarrow$}
dans le cas S=Spec(k), on a également une réciproque: pour que la \xout{secti}secti-%
on hyperplane \uline{générique} satisfasse la propriété P, il f. et s. que
les $Y_{\hand{\zeta}}$ pour $\hand{\zeta}$ dans un voisinage convenable de $\hand{\eta}$ y satisfassent, et il
suffit de l'exiger pour $\hand{\zeta}$ \uline{fermé} dans $\overset{\hand{\vee}}{P}$ (ce qui, pour k algébriquement
clos, revient à envisager des points rationnels sur k, i.e. des

\margin{18}

\clearpage
% Nouveau EGA V (cote 25): archive cote 25, site page 19 (PDF page 20).
% typescript. Draft compared: none.
% AI-transcribed by Claude Opus 5.5 (claude-opus-5-5, adaptive thinking, effort high), Message Batches API, batch ega25b (msgbatch_01A7n5YuXubUPNK1zsepi8RG), 2026-11-13; not yet checked by eye.
\scanpage{25/19}
\mspage{16}

sections hyperplanes de X lui-même (sans extension préalable du corps
\margin{\unsure{trait vertical et flèche}}de base). Cela résulte \unsure{en e ffet|en effet} de la constructibilité de E et du fait
que $\check{P}$ est un schéma de Jacobson.

Donnons à titre d'exemples quelques cas particuliers où les
considérations précédentes s'appliquent:

\uline{Proposition} 4.2. Soit G un Module de présentation finie sur X .
Soit P une des propriétés suivantes pour un Module\inserted{M} sur un préschéma
algébrique Z sur un corps K ;
\xout{\unsure{xxxthérixxxxxxx(i)xxxxxx(v)xxxxxxxxx2,8xxxAlxxxx}}
\deleted{a) L'ensemble \xout{des} E des \unsure{$\check{P}$} tels que G satisfasse P est
constructible}
(i) coppof (M) $\unsure{\leq}$ n.

(ii) M satisfait (S$_{k}$).

(i\unsure{\xout{x}}i) M est Cohen-Macaulay
\inserted[typed]{cycles}
(iv) M est sans \xout{idéaux} premiers associés immergés.

(v) M est séparable sur \xout{\illegible{1}} K .

\deleted{(vi) M est géométriquement intègre sur K/}

\xout{(vii)xxMxestxgéomxtxxxintègrexsurxKx} Avec ces notations; si E désig%
ne l'ensemble des \hand{$\xi\in$} P tels que G\hand{$_{\xi}$} satisfasse à P, on a :
a) E est une partie constructible de \over{P}{$\check{P}$} .

b) Supposons S=Spec(k), k un corps, et que F satisfasse à P. Suppo-%
sons de plus \inserted{(dans le cas (v))} que k soit de car 0, ou que f:X$\rightarrow$P soit une \deleted{immersion}\inserted{\unsure{non ramifié}}. \xout{\unsure{Akx}}
Alors E contient un ouvert dense de \over{P}{$\check{P}$} .

Démonstration. a) résulte du fait que P est une propriété
constructible, ce qui a été vu au par.9. Quant à b), cela résulte
des résultats correspondants des \xout{deux}\inserted{deux} n$^{o}$s précédents. Remords\inserted{à b)}: suppo-%
ser plus généralement que si Z est l'ensemble des points de X enlesquel%
F ne satisfait pas à P, on ait f(Z) fini i.e. dim $\overline{f(Z)}\;\unsure{\leq}0$.

\mspage{19}

\clearpage
% Nouveau EGA V (cote 25): archive cote 25, site page 20 (PDF page 21).
% typescript. Draft compared: none.
% AI-transcribed by Claude Opus 5.5 (claude-opus-5-5, adaptive thinking, effort high), Message Batches API, batch ega25b (msgbatch_01A7n5YuXubUPNK1zsepi8RG), 2026-11-13; not yet checked by eye.
\scanpage{25/20}
\mspage{17}

\uline{Proposition} 4.3.    Soit P l'une des propriétés  suivantes (pour un
préschéma algébrique sur un corps K):

(i)  Z \xout{xxxxxxxx} est lisse sur K.

(ii)  Z satisfait la propriété (R$_k$) géométrique\xout{x} sur K.

(iii)  Z est séparable sur K.

(iv)  Z est géom normal sur K.

\margin{arrow from the circled label of the next line up to this line}(\over{v}{\unsure{vi|v}})  Z  est géom. irr.  sur K.

\margin{label circled}(\over{vi}{\unsure{v|vi}})  Z est géom. intègre  sur K.

Soit E l'ensemble des \hand{$\xi \in$} $\overset{\hand{\vee}}{P}$  tels que $Y_{\hand{\xi}}$ satisfasse à P. Alors

a)  E est une partie constructible de $\overset{\hand{\vee}}{P}$ .

b)  Supp\over{x}{o}sons \xout{dx} S=Spec(k), k un corps\xout{x}, et supposons \inserted{dans les cas (i) à (v)} que k soit de
car zéro  ou que f:X$\rightarrow$P  soit \deleted{une immersion}\inserted{\unsure{non ramifié}}. Supposons enfin dans les
cas (v ) et (vi)  que dim $\overline{f(X)}$ \over{$>$}{$\geqslant$} 2.  Supposons que X satisfasse à P,
alors E contient un ouvert \xout{xxxxxxxx} dense de $\overset{\hand{\vee}}{P}$ .

Démonstration identique à celle de 4.2.   Remarquons que dans
les cas (i) à (\unsure{\deleted{i}v}), il suffit de supposer même (dans b) que  f(Z) est
fini, où Z est l'ensemble des points de X où \xout{lx} P est en défaut.

\uline{Corollaire} 4.4.    Considérons aussi la propriété:

(\xout{x}$C_m$)  \xout{Zxestxx"géométriquementxdexxxxxxx} " $\overline{Z}$ ne peut être discon-%
necté par une partie fermée de dimension  \xout{m}$\leqslant$ m" (où $\overline{Z}$ est $Z\otimes_K\overline{K}$ , $\overline{K}$
la clôture  algébrique de K).

Soit E l'ensemble des \hand{$\xi \in$} $\overset{\hand{\vee}}{P}$  tels que  $Y_{\hand{\xi}}$ sur k(\hand{$\xi$}) satisfasse à $C_m$.
Alors:

a)  E est constructible.

b)  \deleted{Si} \inserted{Supposons} S=Spec(k), k un corps, \inserted{et que} \xout{xxxxx} pour toute composante irréduc-%
tible  $X_i$ de X, on a\hand{it} dim $\overline{f(X_i)} \geqslant 2$, Alors, si X sur k satisfait $C_{m+1}$ ,
alors E\hand{$_{\unsure{\xi}}$} contient un ouvert dense U de P .

\mspage{20}

\clearpage
% Nouveau EGA V (cote 25): archive cote 25, site page 21 (PDF page 22).
% typescript. Draft compared: none.
% AI-transcribed by Claude Opus 5.5 (claude-opus-5-5, adaptive thinking, effort high), Message Batches API, batch ega25b (msgbatch_01A7n5YuXubUPNK1zsepi8RG), 2026-11-13; not yet checked by eye.
\scanpage{25/21}
\mspage{18}

La constructibilité se fait par AQT (c'est en fait un\xout{x} oubli au
par.9 , qu'il serait peut-être encore temps de réparer); la partie b)
résulte en principe de 3.3. , sauf que 3.3. a été énoncé de façon trop
particulière, et doit être généralisé en conséquence (la démonstration
donnée étant d'ailleurs essentiellement inchangée). D'ailleurs, il y
a un lapsus dans l'énoncé 4.4. , qui n'est valable tel quel que si f
est quasi-fini; dans le cas général, au lieu de considérer \xout{l'ensemble} \inserted[typed]{dimension}
\hand{(}des parties fermées de $X$ resp. des $Y_{\hand{\zeta}}$ , il faut regarder la dimension
de leurs images dans P resp. $H_{\hand{\zeta}}$ : Redactor demerdetur.

\margin{Passer au No 5.1}
\deleted{\hand{4.5.}} \deleted{Dans \illegible{1} \unsure{le} présent n$^o$, nous donnons des} conditions \unsure{moyennant} les
\deleted{quelles l'ensemble E envisagé dans 4.1. est ouvert. \xout{Pxxxxxxx} Il}
\deleted{s'agira ici de propriétés P de nature locale sur X resp. $Y_{\hand{\zeta}}$ , de sorte}
\deleted{qu'on peut définir l'ensemble U des $y \in \underline{Y}$ tels que (si \hand{$\zeta$} désigne l'image}
\deleted{de y dans $\hand{\check{}}P$ ) $Y_{\hand{\zeta}}$ satisfasse à P en le point y (resp. $G_{\hand{\zeta}}$ satisfasse à}
\deleted{P en y). Nous donnerons d'abord des critères pour que U soit ouvert,}
\deleted{en utilisant le par.12. \xout{xxxxxxxxxxxxxxxxxxxx} Comme en tous cas on a}
\[ \deleted{E = P - h(\underline{Y} - U)} , \]
\deleted{i\xout{l}l s'ensuit que si U est ouvert et X propre sur S (donc h propre et}
\deleted{à fortiori fermé} \inserted{car $\underline{Y} \to \unsure{Y}$ propre sur X et P \unsure{séparé} sur S}\deleted{), alors E est également ouvert.}

\deleted{\uline{Proposition} 4.6. Soit F \inserted[typed]{Module} de présentation finie sur X, et supposons}
\deleted{F Cohen-Macaulay relativement à S, i.e. plat relativement à S, et}
\deleted{pour tout $s \in S$, $F_s$ sur $X_s$ est Cohen-Macaulay. \xout{Suppxxxxxx} Supposons}
\deleted{de plus que pour tout $s \in S$, \xout{ex} toute composante connexe \deleted{Z} de} \hand{$\zeta - W_{F_s}$}
\inserted{de dim $d$,} \inserted[typed]{\deleted{\xout{exxxitxdimxxZxxx}} et pour}
\deleted{\xout{Spexx} supp $F_s$ , \xout{dexxdimensionxxxxxx}tout $\hand{\zeta} \in \check{P}_s$ , l'image inverse}
\inserted[typed]{l'image d'}
\deleted{$Z_{\hand{\zeta}}$ de $H_{\hand{\zeta}}$ dans $Z_{k(\hand{\zeta})}$ \inserted{$(\zeta - W_{F_s})$} soit de dimension d-1, (i.e. aucune composante}
\deleted{irréductible de \deleted{Z} n'est contenue dans un hyperplan de $\overline{P}_s$\hand{$\check{}$} ). Alors :}
\deleted{a) Le Module \underline{G} sur \underline{Y} est \inserted{Cohen Macaulay, et à fortiori} plat, relativement à $\hand{\check{}}P$ .}
\deleted{\hand{(i) à (iii) de 12.1.1.}}
\deleted{b) Soit P \xout{xxx} l'une des propriétés suivantes pour un module de-}

\margin{21}

\clearpage
% Nouveau EGA V (cote 25): archive cote 25, site page 22 (PDF page 23).
% typescript. Draft compared: none.
% AI-transcribed by Claude Opus 5.5 (claude-opus-5-5, adaptive thinking, effort high), Message Batches API, batch ega25b (msgbatch_01A7n5YuXubUPNK1zsepi8RG), 2026-11-13; not yet checked by eye.
\scanpage{25/22}
\mspage{\hand{19}}

5. \uline{Th\'eor\`emes du type Seidenberg}\xout{\unsure{z}}

5.1.
Dans le pr\'esent n\textsuperscript{o}, nous donnons des conditions .... (cf page pr\'ec\'eden-%
te, le texte barr\'e par erreur).

5.2. Comme nous avons vu au n\textsuperscript{o}1, \underline{Y} est d\'efini dans $X x_S \hand{\check{P}} = X_{\check{\underline{P}}}$
comme le ``diviseur'' d'une section $\phi$ de $\underline{O}_X(1)\otimes_S \underline{O}_{\check{\underline{P}}}(1)$, qui induit pour
tout $\hand{\zeta} \in \hand{\check{P}}$ une sedtion $\phi_{\hand{\zeta}}$ de \xout{$\otimes$xxxxx(1)$\otimes$xxxx xxxxxxxxxxxxxx k( )} $\underline{O}_X(1)\otimes_{k(s)}\underline{O}_{\check{\underline{P}}}(1)(\hand{\zeta})$
(faisceau d'ailleurs isomorphe non canoniquement \`a $\underline{O}_X(1)\otimes_{k(s)}k(\hand{\zeta}) =$
$\underline{O}_{X_{k(\hand{\zeta})}}(1)$ ), de sorte que $Y_{\hand{\zeta}}$ n'est autre que le \xout{xxxxx} ``diviseur'' de cet
te section (NB le diviseur d'une section $\phi$ d'un module inversible \underline{L} \margin{$/\!/$}
est d\'efini comme le sous-pr\'esch\'ema ferm\'e d\'efini par l'Id\'eal image de
$\phi^{-1} : \underline{L}^{-1} \longrightarrow \underline{O}$ ). Si \underline{F} est un Module sur X, alors son image inverse \underline{G}
sur \underline{Y} , i.e. l'image inverse de $\underline{F}\otimes_{\underline{O}_S}\underline{O}_{\check{\underline{P}}} = \underline{F}_{\check{\underline{P}}}$ sur le sous-pr\'esch\'ema \underline{Y} de
$X_{\check{\underline{P}}}$ , n'est autre que le conoyau du homomorphisme
\[ \phi^{-1}\otimes \mathrm{id}_{\underline{F}_{\check{\underline{P}}}} : \underline{F}_{\check{\underline{P}}}(-1,-1) \longrightarrow F_{\check{\underline{P}}} , \]
o\`u la notation (-1,-1) s'explique d'elle-m\xout{e}\^eme, comme dirait Mike. De
m\^eme, $G_{\hand{\zeta}}$ est le conayau de l'homomorphisme analogue
\[ F_{k(\hand{\zeta})}(-1\over{)}{,}-1) \longrightarrow F_{k(\hand{\zeta})} \]
o\`u $\hand{\zeta}$ est un point de $\hand{\check{P}}$ (et bien entendu on a une interpr\'etation corres-%
pondante lorsque $\hand{\zeta}$, au lieu d'\^etre un point de $\hand{\check{P}}$ , d\'esigne un\xout{e} point
de $\hand{\check{P}}$ \`a valeurs dans un S' sur S \xout{\unsure{??}}. ). \hand{\unsure{(}}De fa\c{c}on g\'en\'erale, si L est un
Module inversible quelque part, $\phi$ une section, \unsure{d\'afinissant|d\'efinissant} le sous-pr\'esch
ma V($\phi$), alors pour tout Module F , l'image inverse de F sur V($\phi$)
s'identifie, moyennant l'abus de langage habituel, au conayau de $\mathrm{id}_F\otimes\phi$
$F\otimes L^{-1} \longrightarrow F$ ; On dit que $\phi$ est F-\uline{r\'egulier} si l\xout{e}'homomorphisme pr\'ec\'edent
est injectif. Si on choisit un isomorphisme de F et $\underline{O}_X$ , - ce qui
est possible localement - de sorte que $\phi$ s'identifie \`a une section de
$\underline{O}_X$ , cette terminologie est\unsure{b} compatible avec celle \inserted{d\'ej\`a} introduite par ailleur

\margin{22}

\clearpage
% Nouveau EGA V (cote 25): archive cote 25, site page 23 (PDF page 24).
% typescript. Draft compared: none.
% AI-transcribed by Claude Opus 5.5 (claude-opus-5-5, adaptive thinking, effort high), Message Batches API, batch ega25b (msgbatch_01A7n5YuXubUPNK1zsepi8RG), 2026-11-13; not yet checked by eye.
\scanpage{25/23}
\mspage{\hand{20}}

\uline{Proposition} 5.3.   Avec les notations précédentes, soit U l'ensemble des
$x\in X_{\hand{\check{}}P}$ , d'image $\hand{\zeta}$ dans $\hand{\check{}}P$ , tels que $\phi_{\hand{\zeta}}$ soit $F_{k(\hand{\zeta})}$-régulière en x. \deleted{Suppo-%
sons de plus} \inserted{Si} F de présentation finie et plat relativement à S\hand{,} Alors:
% NOTE: the words "Si F de présentation finie et plat relativement à S," are encircled by hand, and a line leads from the loop down to a), between "a)" and "U".
a) U est ouvert, et $\underline{G}\hand{|}U$ est plat relativement à $\hand{\check{}}P$ .

b) Pour tout $s\in\hand{\check{}}P$ , si $\hand{\eta}$ désigne le point générique de $\hand{\check{}}P_s$ , U con-%
tient $X_{k(\hand{\eta})}$.

Démonstration. a) Comme $F_{\hand{\check{}}P}$ est \inserted[typed]{de présentation finie, et} plat par rapport à $\hand{\check{}}P$ , la conclu-%
sion résulte de 11.3. (et même de $\underline{E}_{III}$ ... dans le cas S loc noeth).

b) On peut supposer S=Spec(k). Les cycles associés à $F_{k(\hand{\eta})}$ sont
\inserted[typed]{\xout{\illegible{1}}}
(en vertu du par.3 ) les images inverses des cycles \inserted{$Z_i$} associés à F. \xout{En
vertu de 2.6, par exemple, si $Z_i$ est une partie \unsure{sous-schéma} fermé intègre
de X telle que} Si $f(Z_i)$ est fini, alors en vertu de 2.3. \xout{X}
$Z_{i\,k(\hand{\eta})}\cap Y = \emptyset$ , dans le cas contraire, on vertu de 2.6.\inserted[typed]{par exemple} on a
également $Z_{i\,k(\hand{\eta})}\cap Y \neq Z_{i\,k(\hand{\eta})}$ (raisons de dimensions; 2.3. déjà invo-%
qué dans 2.6. est s ans doute une meulleure raison), ce qui prouve que $\phi_{\hand{\eta}}$
ne s'annule sur aucun des $Z_{i\,k(\hand{\eta})}$ , et prouve donc b).

\uline{Corollaire} 5.4. Soit V l'ensemble des $\hand{\zeta}\in\hand{\check{}}P$ tels que $\phi_{\hand{\zeta}}$ soit \xout{E} $F_{k(\hand{\zeta})}$-régu-%
lier. \inserted{Si F est de présentation finie} \over{A}{\unsure{a}}lors V est constructible et contient les points génériques des
fibres de $\hand{\check{}}P$ sur S, d'autre part si \inserted{de plus} X est propre sur S\inserted{et F plat sur S}\hand{,} l'ensemble V est
ouvert. \deleted{\hand{Si \unsure{$S$ spec $(A)$} (\unsure{$A$ complet}) alors \unsure{en} tous cas $\unsure{f}(\check{P}-V)$ est
de codimension $\unsure{\geq|\neq} 2$.}}

\inserted[typed]{\uline{Remarque}}\xout{Xxxxxxxxxx}5.5. \xout{Supposons} soit $\hand{\zeta}\in\hand{\check{}}P$ au dessus de $s\in S$, et supposons
que $F_s$ soit sans cycles associés immergés, alors on voit aussitôt que
$\hand{\zeta}\in V$ (notation de 5.4.) signifie aussi que\xout{xxxxx} aucune composante
irréductible de supp $F_{k(\hand{\zeta})}$ ne se trouve au dessus de $H_{\hand{\zeta}}$ (et on peut
évidemment dans ce critère remplacer $k(\hand{\zeta})$ par une extension arbitraire
de $k(\hand{\zeta})$ ). Notons que l'hypothèse $(S_1)$ sur $F_s$ qu'on vient de faire es

\deleted{\hand{$\times$ Ce dernier fait se vérifie par la méthode de 5.3. b) et la \unsure{théorie}
de la dimension (on remplace $\eta$ par $\zeta$ \unsure{tel que} $\mathcal{O}_{\check{P},\zeta}$ de dim $1$)}}
\margin{23}

\clearpage
% Nouveau EGA V (cote 25): archive cote 25, site page 24 (PDF page 25).
% typescript. Draft compared: none.
% AI-transcribed by Claude Opus 5.5 (claude-opus-5-5, adaptive thinking, effort high), Message Batches API, batch ega25b (msgbatch_01A7n5YuXubUPNK1zsepi8RG), 2026-11-13; not yet checked by eye.
\scanpage{25/24}
\mspage{\hand{21}}

satisfaite notamment si on supposer $F_s$ Cohen-Macaulay (à fortiori, si
F est CM sur S); de plus, dans ce cas, $G_s$ est CM (car localement il se
déduit de $F_{k(s)}$ qui l'est en divisant par un $\phi . F_{k(s)}$ , \deleted{avec} \inserted{où} $\phi$ \inserted{est} $F_{k(s)}$-%
régulier). Ces mêmes remarques pouvaient \inserted{(et devaient)} d'ailleurs se faire localement
en haut, pour caractériser les points de U (au lieu de ceux de V).

Utilisant maintenant 12.1.1. et 12.1.4. , on obtient:
\uline{Théorème} 5.6. Supposons F de présentation finie, plat rel. à S. Soit
P l'une des propriétés \xout{\unsure{satisfaites pour un module de type fini M sur}}
\xout{\unsure{un anneau local noethérien A sur un corps K}} (i) à (viii) de 12.1.1. ,
ou \xout{\unsure{(ii) de exposé M (réf)}} (dans le cas où on suppose $F=\underline{O}_X$) l'une des
\xout{\unsure{xxxx(ii) M satisfait $(S_k)$ (réf)}} propriétés (i) à (iv)\hand{/} de 12.1.4. Soit
\margin{\unsure{1r}} $U_P$ \xout{\unsure{(iii) M est Cohen-Macaulay}} l'ensemble des $x \in X_{\hand{\check{}}P}$ tels que , si \hand{$\zeta$}
\xout{xxx(iv)} désigne l'image de x dans $\hand{\check{P}}$ , la propriété P soit satisfaite
par $G_{\hand{\zeta}}$ (resp. \xout{\unsure{0}} $Y_{\hand{\zeta}}$ ) en le point x, \uline{et} que $\phi_{\hand{\zeta}}$ soit $F_{k(\hand{\zeta})}$-régulière en x.
Alors $U_P$ est ouvert, et $\underline{G}\hand{|}U_P$ est plat relativement à S.

En effet, par définition même, on a $U_P \subset U$ (notation de 5.3. a)),
et on applique par.12 à $U \to \hand{\check{P}}$ et $F_{\hand{\check{P}}}\hand{|}U$.
\uline{Corollaire} 5.7. Suppososons F de présentation finie, plat rel à S,
et supp F propre sur S (p.ex. X propre sur S). Soit $V_P$ l'ensemble des
$\hand{\zeta\in}\hand{\check{P}}$ tels que $G_{\hand{\zeta}}$ (resp. $Y_{\hand{\zeta}}$ ) satisfasse à la propriété P, et que $\phi_{\hand{\zeta}}$ soit
$F_{k(\hand{\zeta})}$-régulière. Sous ces conditions, $V_P$ est ouvert, (et il est d'ailleu
rs constructible en tous cas, i.e. sans hypothèse de platitude ni de
\unsure{prppreté|propreté}).

Il me semble que du point de vue présentation, on ne peut laisser
5.6. tel quel, avec une simple référence à des conditions énumérées dans
un autre fascicule; mais il faudra refaire une liste explicite (i) (ii)
.... des propriétés qu'on a en vue. De plus, signaler \inserted{(dans 5.1. peut-être)} que le cas P= géom.
normal (avec S =Spec(k) pour sur) est dû à \uline{Seidenberg}.

\margin{24}

\clearpage
% Nouveau EGA V (cote 25): archive cote 25, site page 25 (PDF page 26).
% typescript. Draft compared: none.
% AI-transcribed by Claude Opus 5.5 (claude-opus-5-5, adaptive thinking, effort high), Message Batches API, batch ega25b (msgbatch_01A7n5YuXubUPNK1zsepi8RG), 2026-11-13; not yet checked by eye.
\scanpage{25/25}
\mspage{\hand{21 bis}}

\deleted{ce à peu près immédiate de 15.2.2. appliqué à $\underline{Y} \longrightarrow \check{P}$ .}

\deleted{Il faut voir si  le fait de mélanger dans un même énoncé des proprié-}%
\deleted{tés modulées et non modulables  ne le\xout{s} rend pas  impigeable\xout{s} pratique-}%
\deleted{ment. Noter aussi  que 4.6. et 4.7. est entièrement indépendant des}
\deleted{autres développements des n\textsuperscript{o}s 2,3,4 ; à ce titre, ils méritent sans dou-}%
\deleted{te d'être mis dans un n\textsuperscript{o} à part, où il n'est pas question des propriétés}
\deleted{de la section hyperplane générique, mais uniquement de déterminer si}
\deleted{l'ensemble E est ouvert. On pourrait donner comme titre d\over{a}{e} n\textsuperscript{o} :}
\deleted{``\uline{Théorèmes du type \xout{\illegible{1}} Seidenberg pour une section hyperplane variable}''.}
\deleted{Il faut signaler que le cas  ``géométriquement normal'' (\deleted{\unsure{et}} bien entendu}
\deleted{dans le cas où S=Spec(k) )  est en effet dû à Seidenberg. N'ayant pas}
\deleted{lu son papier, je ne sais s'il part d'une variété X déjà supposée nor-}%
\deleted{male, ou même lisse.  Ce dernier cas semble techniquement nettement}
\deleted{moins vache, puisqu'il suffit de 4.6. qui ne s'appuie pas sur le}
\deleted{délicat 15.2.2.  (Si Seidenberg ne se borne pas à ce cas, je me demande}
\deleted{comment il a pu faire pour se débrouiller).  Si on fait le scindage en}
\deleted{deux envisagé pour le présent numéro, la première partie  devrait prend-}%
\deleted{re le ti\unsure{tre}:  \uline{Etude d'une section hyperplane variable}: \uline{cas des}}
\deleted{\uline{sections hyperplanes} ``\uline{assez générales}''.}

\over{\unsure{6}}{5}. \uline{Connexion d'une section hyperplane quelconque}.

Nous allons ici combiner le critère déjà connu de connexion
géométrique de\xout{xx} la section hyperplane générique  (3.3.), avec le
th de connexion de Zariski, pour obtenir un résultat de connexion pour
une section hyperplane quelconque:
\uline{Proposition} \over{\unsure{5}}{6}.1.  On suppose S=Spec(k), k un corps\xout{xxxx} alg clos, \inserted{X propre sur S} \xout{xx}
que pour toute composante irréductible X$_{i}$ de X, \xout{xxxxxx} $\overline{f(X_{i})}$ soit de
dimension \over{$>$}{$\geqslant$} 2, \xout{xxxxxxxxxxxxxxxxxxxxxxxxxxxxxxxhyper} \deleted{Alors pour}

\margin{25}

\clearpage
% Nouveau EGA V (cote 25): archive cote 25, site page 26 (PDF page 27).
% typescript. Draft compared: none.
% AI-transcribed by Claude Opus 5.5 (claude-opus-5-5, adaptive thinking, effort high), Message Batches API, batch ega25b (msgbatch_01A7n5YuXubUPNK1zsepi8RG), 2026-11-13; not yet checked by eye.
\scanpage{25/26}
\mspage{22}
enfin que X ne peut être disconnecté par une partie fermée Z de X telle
que dim $\overline{f(Z)} \unsure{\leq} 0$.    Sous ces conditions, pour tout $\hand{\zeta \in} \overset{\hand{\vee}}{P}$ , $Y_{\hand{\zeta}}$ est
géométriquement connexe.

Démonstration.  Comme \xout{xesxf(X} aucun des $f(X_i)$ n'est fini, on voit
que chacune d es composantes irréductibles $\underline{Y}_i$ de $\underline{Y}$ domine $\overset{\hand{\vee}}{P}$ ;
d'autre part $\underline{Y} \to \overset{\hand{\vee}}{P}$ est propre (car $\underline{Y}$ est propre sur k, l'étant sur X
qui l'est sur k).  D'autre part, en vertu de (3.3.) la fibre générique
$Y_{\hand{\eta}}$ de $\underline{Y} \to \overset{\hand{\vee}}{P}$ est géométr\xout{i}iquement connexe. Enfin, $\overset{\hand{\vee}}{P}$ est régulier et
à fortiori géom unibranche. Il suffit maintenant d'appliquer 15.6.3.
(\xout{xsxxxxxxxxxxxxxxxxxxxxxxxx}\inserted[typed]{qui n'est qu'une variante du} th. de connexité  de Zariski) pour concl
re que \uline{toutes} les fibres  de $\underline{Y} \to \overset{\hand{\vee}}{P}$ sont géom connexes, cqfd.

\margin{$|$} En fait,  il n'est pas difficile, par une démsnstration de type an
logue, de généraliser \over{\illegible{1}}{5}.1. dans le même sens que dans 4.4. Si tu ne
veux pas te taper l'exercice toi-même, signale au, moins la chose en re
marque. Y dire aussi que l'on ne fait aucune discrimination dans \over{\illegible{1}}{5}.1.
à l'égard des sections hyperplanes qui auraient une dimension excéden-%
\margin{$|$} taire.  Du point de vue plan, il sera peut-être plus clair de bloquer
toutes les questions de connexité (incluant 3.3. et 4.4. ) dans un
même n$^{\text{o}}$ .
\deleted{6. \uline{Application à la construction de sections hyperplanes} \hand{/}\inserted[typed]{\uline{ou de multisections}}\hand{)} \uline{de type préci-%
sé}, \xout{ouxdexmultisectionsx}}
\mspage{26}

\clearpage
% Nouveau EGA V (cote 25): archive cote 25, site page 27 (PDF page 28).
% typescript. Draft compared: none.
% AI-transcribed by Claude Opus 5.5 (claude-opus-5-5, adaptive thinking, effort high), Message Batches API, batch ega25b (msgbatch_01A7n5YuXubUPNK1zsepi8RG), 2026-11-13; not yet checked by eye.
\scanpage{25/27}
\mspage{23}

\deleted{Remords au N°4 . Les résultats 4.6. b) et c) sont enti\xout{\unsure{m}}èrement
contenus dans 4.7. b) et c) ; j'avais cru une telle redite justifiable du
point de vue démonstration, celle de 4.7. me semblant mettre en oeuvre des
résultats nettement plus délicats (sous forme de 15.2.2.). Or, tant qu'on
ne tient qu'à prouver dans 4.7. b) de c) que U est ouvert, (à l'exclusion
du fait $\underline{G}$ U plat \xout{xxx} rel à $\overset{\hand{$\vee$}}{P}$ ), on n'a nul besoin de 15.2.2. , mais la
méthode habituelle (utilisée au par.12) par le critère valuatif donne le
résultat voulu de façon bien simple. Je suggère donc:a\xout{\unsure{x}}) de vider les
parties b) et c) de 4.6. ; b) d'énoncer 4.7. b) et c) avec la seule
conclusion U resp. E ouvert, en utilisant le critère valuatif (bien entendu,
4.7. a) est vidé de l'énoncé), c) d'ajouter une autre proposition\inserted[typed]{4.8.}, ou un
corollaire, disant que sous les conditions de 4.7. $\underline{G}$ U est plat rel à $\overset{\hand{$\vee$}}{P}$ ,
ce qui se démontre à l'aide de 15.2.2.}

\over{\unsure{5}}{7}. \uline{Application à la construct\xout{i}ion de sections hyperplanes ou de multisections
de type précisé}.

\over{\unsure{5}}{7}.1. Notons que si S=Spec(k), ou k est un corps \uline{infini}, alors tout ouvert
non vide de $\overset{\hand{$\vee$}}{P}$ contient un point rationnel sur k; donc avec les notations de
\xout{\unsure{xxx}}4.1. , si E (défini en termes d'une propriété constructible P) contient
le point générique \hand{$\eta$} , il cotient \xout{xxxxxxxx} un point rationnel sur k, donc
il existe une section hyperplane de X (définie sur k) ayant la propriété P.
D'autre part, S étant de nouveau quelconque, il est immédiat que pour tout
s\hand{$\in$}S et tout point\inserted{$\xi_0$} de la fibre $\overset{\hand{$\vee$}}{P}_s$ rationnel sur k(s), il existe une section
\hand{$\xi$} de $\overset{\hand{$\vee$}}{P}$ sur un voisinage ouvert U de s qui passe par \hand{$\xi$}$_0$ . Si alors E est
encore \xout{xxx} défini\xout{x} comme dans 4.1. en termes d'une propriété constructible
P, et si on dispose (p ex grâce \deleted{à 5.\unsure{5}.}\inserted{\unsure{à} N°5}) du fait que E est ouvert, alors
si \hand{$\xi$}$_0$\hand{$\in$} E, \xout{\unsure{xxxprenxxprendrexMxxx}} la section \hand{$\xi$} est une section de E sur \xout{\unsure{Sx}}
U, du moins si on rapetisse suffisamment U. Donc, on aura construit une
section hyperplane\inserted{$Y_\xi$} de X , au dessus d'un voisinage ouvert U de s, telle que

\margin{27}

\clearpage
% Nouveau EGA V (cote 25): archive cote 25, site page 28 (PDF page 29).
% typescript. Draft compared: none.
% AI-transcribed by Claude Opus 5.5 (claude-opus-5-5, adaptive thinking, effort high), Message Batches API, batch ega25b (msgbatch_01A7n5YuXubUPNK1zsepi8RG), 2026-11-13; not yet checked by eye.
\scanpage{25/28}
\mspage{29}
pour tout t \hand{$\in$}U,  sa fibre \inserted{$Y_{\xi}$(1)} en  t satisfasse la propriété P.  Lorsqu'on ne
dispose pas à priori de \hand{$\xi$}$_{o}$ , \xout{xxx} mais que k(s) est infini, on pourra combi-%
ner les deux remarques précédentes pour obtenir une section hyperplane
sur un voisinage ouvert de s, ayant la propriété précédente. \xout{Dixxxxx} Enfin,
utilisant \deleted{4.8.}\inserted{le N$^{o}$5} , on a un critère permettant d'affirmer que (X resp. F étant
supposé plat sur S, ce qui permet d'appliquer \deleted{4.7}\inserted{loc. cit.}.)  Y\hand{$_{\xi}$} \inserted[typed]{resp. G\hand{$_{\xi}$}} est également
plat sur S. On peut alors, remplaçant X par Y\hand{$_{\xi}$} , itèrer la construction
précédente, ce qui permett\xout{xxxxxx} \inserted[typed]{par exemple\hand{,} sous certaines conditions} de construire de proche en proche une
``multisection'' \inserted[typed]{S'} de X au dessus d'un voisinage ouvert U d'un point donné s,
telle\xout{x} que S'\hand{$\longrightarrow$}U soit fini, plat et à fibres satisfaisant la propriété P.
Lorsque k(s) est fini, on peut être contraint de faire un changement de bas%
% NOTE: right edge of the page cut; last letter of "base" not visible
étale surjectif S'\hand{$\longrightarrow$} U (U voisinage ouvert de s) avant de pouvoir appliquer
les constructions précédentes; en effet, sous les conditions du début de
6.1. , si k est fini, il n'existe pas nécessairement de point rationnel sur
k dans  l'ouvert non vide U, mais \xout{l'ensemble des points fermés de U est}
il existe certainement un point fermé de U, donc un point à valeurs dans
une extension f inie k' (nécessairement séparable) de k; lorsque k = k(s),
on peut donc, quitte à faire une extension \inserted[typed]{finie} étale convenable \inserted[typed]{S'} sur un voisinag%
% NOTE: right edge cut; end of "voisinage" not visible
ouvert U de  s, correspondant à l'extension résiduelle k'/k\xout{xxxxxxxxx}
i.e. telle que  S$'_{s}$ \hand{$\xrightarrow{\sim}$}\inserted{\unsure{k(s)-isomorphisme}}Spec(\over{k}{k'}) , se ramener à la situation favorable \xout{par}
en l'unique point s'\hand{$\in$}S' au dessus de s, par le changement de base S'\hand{$\to$} S.

Donnons quelques exemples d'énoncés obtenus par ce\over{s}{tte} méthode\deleted{s}. \deleted{Il y a}
\deleted{lieu cependant, de faire une remarque préliminaire, pour l'élimination dan}%
% NOTE: right edge cut ("dans")
\deleted{\uline{\xout{Propositionxxxxxxxxxx}}}
\deleted{ces constructions de l'hypothèse faite dans 4.6. \deleted{et} 4.7. \inserted{et 4.8.} que  aucune compo}%
\deleted{sante irréductible géométrique d\xout{x}'un supp F$_{s}$ n'est contenue\hand{\unsure{,}}dans une sectio}%
\deleted{hyperplane (definie sur $\overline{\text{k(s)}}$ ): on peut (et doit) en effet énoncer 4.6.\hand{\unsure{/}}}
\deleted{et 4.7. sous forme légèrement plus générale, en se donnant une partie ouver}%
\deleted{te W de P\hand{$\check{}$} , \xout{ex}  en remplaçant dans la condition envisagée  ``hyperplan'' par}
\margin{28}

\clearpage
% Nouveau EGA V (cote 25): archive cote 25, site page 29 (PDF page 30).
% typescript. Draft compared: none.
% AI-transcribed by Claude Opus 5.5 (claude-opus-5-5, adaptive thinking, effort high), Message Batches API, batch ega25b (msgbatch_01A7n5YuXubUPNK1zsepi8RG), 2026-11-13; not yet checked by eye.
\scanpage{25/29}
\mspage{\hand{25}}

Je dois cependant d'abord \unsure{signaler|signqler} un remords à 4.2. et 4.3. , qui
\unsure{dev-aient|devaient} être énoncés sous une forme légèrement plus générale (tout au
moins en remarque): on se donne un entier m, et on \unsure{désigne|dásigne} par E l'ensem
ble des \hand{$\xi\in$}$\overset{\hand{$\vee$}}{\text{P}}$ tels que G$_{\hand{$\xi$}}$ resp. Y$_{\hand{$\xi$}}$ satisfasse P sauf \deleted{au plus} sur une
partie fermée de dimension\hand{$\leq$}m (i.e. l'ensemble P-singulier Z est de dimen-%
sion\hand{$\leq$} m). Alors a) E est une partie constructible de $\overset{\hand{$\vee$}}{\text{P}}$ et b) dans le
cas S=Spec(k), si F \inserted[typed]{resp. X} satisfait à P sauf \deleted{au plus} sur un ensemble de
\xout{xx}dimension\hand{$\leq$} m+1 , alors E contient un ouvert non vide.

\uline{Proposition} 7.2. Supposons que X soit propre sur S, et F de présentati-%
on finie et plat sur S. Soit P l'une des propriétés (i) à (v) de 4.2.
et soit m un entier. Soit s\hand{$\in$}S, et supposons que \xout{k(s)xxF$_s$xsatisfassexà}
\xout{laxpropriétéxP} l\xout{'}ensemble Z$_s$ des points de X$_s$ en lesquels F$_s$ ne satisfait
pas à P soit de dimension \hand{$\leq$} m+1. Alors, si de plus k(s) est infini, il
existe un voisinage U de s dans E, et une section \hand{$\xi$} de $\overset{\hand{$\vee$}}{\text{P}}$ sur U, ayant
les propriétés suivantes: pour tout s\hand{$\in$}U, l'ensemble Z des points
de \deleted{X} Y$_{\hand{\unsure{$\xi,\xi(s)$|$\xi(s)$}}}$ en lesquels G$_{\hand{$\xi(s)$}}$ ne satisfait pas à P est de dimension
\hand{$\leq$}m, et $\phi_{\hand{$\xi$}(s)}$ est F$_{\hand{\unsure{$\xi,\xi(s)$|$\xi(s)$}}}$-régulière\hand{,} \xout{xlexModulexGxxsur} Sous ces condi-%
tions, le Module G$_{\hand{$\xi$}}$ sur Y$_{\hand{$\xi$}}$ est plat relativement à U. Enfin, si k(s)
n'est pas supposé \unsure{inffni|infini}, on peut encore faire la construction précéden-%
te, après extension étale duntype envisagé dans 7.1.

\uline{Proposition} 7.3. Essentiellement la même, il n'y a plus de F, \unsure{on|ok}
\margin{[N.B. Pour (v) cf 12.2.1. (x) et (xi) \unsure{Ceci peut} donc apporter une variante à 4.3. \unsure{via} 12.2.1. (x)]}
suppose X plat rel à S, on réfère aux propriétés (i) à \deleted{(xxx)}\inserted{\unsure{(iv)}} de 4.3. au
lieu de celles de 4.2. (mais en ayant soin de faire \deleted{les réserves} \xout{axxxxxxxxxxx}\inserted{la réserve}
``k\inserted{(s)} de car. 0 ou f:X$\longrightarrow$P une immersion\hand{$\xi$}\inserted{, et dans le cas (v), pour tout $s\in S$, et tte comp. irréductible Z de X$_s$, on ait \unsure{dif}(Z) $\geq 2$).}

\deleted{On notera que nous avons exclus les propriétés (v) et (vi) de}
\deleted{4.3. , qui ne sont pas de nature purement locale. Cependant, on récupère}
\deleted{le cas (v) en notant qu'on peut s'arranger en vertu de la démonstratio}
\deleted{de 7.3. que Y$_{\hand{$\xi$}}$ soit plat sur U à fibres séparables, mais alors en}
\margin{\unsure{29}}

\clearpage
% Nouveau EGA V (cote 25): archive cote 25, site page 30 (PDF page 31).
% typescript. Draft compared: none.
% AI-transcribed by Claude Opus 5.5 (claude-opus-5-5, adaptive thinking, effort high), Message Batches API, batch ega25b (msgbatch_01A7n5YuXubUPNK1zsepi8RG), 2026-11-13; not yet checked by eye.
\scanpage{25/30}
\mspage{26}

\deleted{tu de III 7 la factorisation de Stein de \xout{\unsure{X}} Y $\longrightarrow$ S se fait par un}
\deleted{morphisme fini \uline{étale} S'$\longrightarrow$S, ce qui prouve que le nombre de composantes}
\deleted{connexes géom de \xout{\unsure{s}} la fibre de Y\hand{$_{\xi}$} en \xout{\unsure{t}} (égale au nombre analogue pour}\margin{Il y a une variante de \unsure{12.2.1. (4) ou (5)} \unsure{a} qui permet d'inclure ces propriétés dans 7.3.}
\deleted{S' sur S) est constant au voisinage de s, donc en l'occurence égal à 1.}
\deleted{Ce fait mérite d'être signalé au moins en remarque.}

\uline{Proposition} 7.4. Soit \unsure{$g$|$f$}: X $\longrightarrow$ S un morphisme propre et plat, \xout{et}
soit s $\in$ S, \xout{\unsure{tellxque}} posons n = dim X$_{s}$ et supposons que la dimension de
l'ènsemble des points de X$_{s}$ en lesquels X$_{s}$ n'est pas séparable sur k(s)
soit \hand{$<$} n\inserted{(\unsure{p.ex} $X_s$ séparable)}. Alors il existe un voisinage ouvert U de s, et un morphisme
étale fini et surjectif S'$\longrightarrow$U, tel que X$\mathrm{x}_{S}$S' admette une section sur S'.
Lorsque k(s) est infini, on peut prendre pour S' un sous-préschéma fermé
de X U.

Supposons d'abord k(s) infini. Nous \xout{\unsure{appliquons}} procédons par récu%
% NOTE: the end of this line and of several following lines is cut by the right edge of the scan
rence sur n, le cas n=0 étant trivial: en effet, dans ce cas il existe
un voisinage ouvert U de s tel que X$|$U lui-même soit étale, fini, et
surjectif au dessus de U, comme on voit par des cross-références immédia\xout{\unsure{t}}
tes. Lorsque n$>$0, on applique 7.3. pour la propriété ``séparable'', ce qui
nous permet de \xout{xxxxxxxxx} remplacer X par une ``section hyperplane'' Y, ayan
les mêmes propriétés, à cela près que n est rèmplacé par n-1. Lorsque
k(s) n'est pas supposé infini, on commence par faire une extension étale
de la base, ça marche.

\uline{Remarque\hand{s}} 7.5. En particulier, si X est projectif et séparable sur S,
il admet localement sur S des multisections étales. Mais on notera qu'on
peut donner des exemples, avec X propre et lisse (mais non projectif) su
S, où la même conclusion est en défaut. Bien entendu, l'hypothèse de
\xout{qua} projectivité ne peut -être affaiblie \inserted{en général} en une hypothèse quasi-projecti
comme on voit par exemple en prenant X étale non fini sur S ....

\margin{Possibilité d'éviter un ensemble \unsure{fini} ...}

\mspage{30}

\clearpage
% Nouveau EGA V (cote 25): archive cote 25, site page 31 (PDF page 32).
% typescript. Draft compared: none.
% AI-transcribed by Claude Opus 5.5 (claude-opus-5-5, adaptive thinking, effort high), Message Batches API, batch ega25b (msgbatch_01A7n5YuXubUPNK1zsepi8RG), 2026-11-13; not yet checked by eye.
\scanpage{25/31}
\mspage{\unsure{27}}

8. \uline{Dimension de l'ensemble des hyperplans exceptionnels}.

8.1.Dans les n°s précédents, et notamment 2 et 3, nous avons donné des
énoncés affirmant que l'ensemble des \hand{$\xi\in$} $\overset{\hand{$\vee$}}{P}$ tels que Y ait une certaine
propriété P est constructible, et contient le point générique\inserted{\unsure{$\overset{\vee}{\ }$}}\hand{/}, \xout{xxxx} ou
encore que l'ensemble $Z_P$ des \hand{$\xi\in$} $\overset{\hand{$\vee$}}{P}$ ``exceptionnels pour P'' est construc-%
tible, et est rare, i.e. son adhérence est de\inserted{co}dimension $\unsure{\geq}$ \hand{1}\deleted{r (où r=dim $\check{P}$}) \margin{\unsure{.}}
(NB on suppose S=Spec(k)). Dans certains cas on peut préciser cet énoncé
en donnant une meilleure m\over{aj}{\unsure{in}}oration de cette\inserted{co}dimension, ce qui est impor-%
tant \xout{pour}\inserted[typed]{dans} diverses questions. Par exemple, lorsqu'on sait que cette\inserted{co}dimen
sion est $\geq$ \deleted{\unsure{r}}-2, il s'ensuit qu'une droite ``assez générale'' D de $\overset{\hand{$\vee$}}{P}$
ne rencontre pas $Z_P$ , d'où l'existence (si k est infini) de ``pinceaux
linéaires'' de sections hyperplanes \xout{(Xx)xxxxx}\hand{\unsure{Dt}} $Y_{\hand{\xi}}$ (\hand{$\xi$} un point géométri
que de D) qui sont toutes la propriété P (cf n° \deleted{suivant}\inserted{\illegible{1}} pour des exem-%
ples. - Du point de vue rédaction, comme les résultats du présent n°
précisent certains résultats des n°s antérieurs, la question se pose s'il
y a bien lieu de faire ce ratrappage dans un n° à part, ou de donner la
version précisée à mesure. Redactor decidetur.

\deleted{8.2. Soit Z une partie \deleted{constructible}\inserted{telle que} de X, et considérons l'ensemble \margin{Z}
des \quad P tels que son image inverse Z dans Y soit de dimension
\inserted[typed]{strictement}plus grande (i.e. son adhérence \inserted[typed]{dans Y}de dimension \inserted[typed]{str.}plus grande) que celle de
la fibre générique (i.e. soit égale à dim Z si dim $\overline{f(Z)}$ \quad 0, resp.
soit non vide si dim $\overline{f(Z)}$ \hand{\unsure{(}} 0 ). Alors Z est de codim\xout{\unsure{mm}}nsion \quad 2 i.e.
de dimension \quad r-2. (Cela résulte aussitôt de la théorie de la dimension
d'un morphisme appliqué à $\underline{Y}$ \quad P ). Bien entendu, cet énoncé peut se
\inserted[typed]{en termes de supp F}moduler (et est équivalent à l'énoncé modulé), et l'appliquant aux com-%
posantes irréductibles géométriques de T, on trouve aussi que l'ensemble
Z des \quad tels \xout{xxx} qu'il existe une composante irr de}

\mspage{31}

\clearpage
% Nouveau EGA V (cote 25): archive cote 25, site page 32 (PDF page 33).
% typescript. Draft compared: none.
% AI-transcribed by Claude Opus 5.5 (claude-opus-5-5, adaptive thinking, effort high), Message Batches API, batch ega25b (msgbatch_01A7n5YuXubUPNK1zsepi8RG), 2026-11-13; not yet checked by eye.
\scanpage{25/32}
\mspage{28}

8.2. Soit Z l'ensemble des $\hand{\xi\in}\check{P}$ tels que dim $Y_{\hand{\xi}}$ $\hand{>}$ dim X - 1,\hand{\unsure{\xi}}
et supposons que pour toute composante irr $X_i$ de X, on ait dim $f(X_i)$ $\hand{>}$ 0;
alors Z est de codimension 2 dans $\check{P}$. \xout{(ixexxdexdimensionxxxx=2)} Cela
\inserted[typed]{2.1. et 2.2.}
résulte de \xout{\unsure{2x5xxxx2x3xxx(quiximpliquentxquexxlesxXx$X_i$xsontxirréductibles}}
\xout{\unsure{etxdominentxP}} (qui implique que toute comp irr de $\underline{Y}$ domine $\check{P}$ ) et de
la théorie de la dimension pour le morphisme $\underline{Y} \hand{\to} \check{P}$ . A partir de ce
résultat \xout{\unsure{(quixxxsuffitxde}} on peut donner comme corollaire le cas où on
part d'une partie \inserted[typed]{fermée} \xout{\illegible{1}} Z de X et qu'on r egarde la dimension des
images inverses $Z_{\hand{\xi}}$ dans les $Y_{\hand{\xi}}$ ( $\hand{\xi\in}\check{P}$ ), et on peut même prendre pour
Z l'ensemble des $\hand{\xi\in}\check{P}$ tels que il existe une co mposante irr. de $T_{k(\hand{\xi})}$
dont la trace sur $Y_{\hand{\xi}}$ soit de dimension trop grande (NB on suppose toujour
que pour toute comp irréd $Z_i$ de Z, on ait dim $f(Z_{\hand{i}})$ $\hand{>}$ 0). Enfin, l'énoncé
le plus précis dans cette direction, et qui résulte aisément du premier
énoncé (pour X irréductible) et de 2.7. , est l'énoncé modulé suivant:
F étant cohérent sur X, \xout{soit} supposons que pour tout cycle premier associ
T pour F, on ait dim f(T)$\hand{>}$ 0\unsure{;} alors l'ensemble des $\hand{\xi\in}\check{P}$ tels que
$\hand{\phi_\xi}$ \inserted{ne} soit pas $F_{k(\hand{\xi})}$-régulier est (constructible et) de codimension$\hand{\geqslant}$ 2.
\margin{\textbar} (La notation pour $\hand{\phi_\xi}$ est celle du n°5). On peut donner ceci comme
\margin{\textbar} \inserted[typed]{proposition principale}\xout{\unsure{Dexlaxmêmexfaçon}}, et énoncer les \xout{autres} assertions qui précédaient
\margin{\textbar} comme corollaires, la démonstration procédant via l'un des corollaires.

\deleted{Autre corollaire: Soient T,T' des parties femmées de X, et supposons
qu'on \xout{mait} \xout{\unsure{pasxdimxexdimxim}} T / $T'_{\unsure{\xi}}$ ( étant le point générique de
\inserted[typed]{il existe une co mp irr $T_i$ de T telle que}
P ), i.e. $f(T_i)$ est infini et $T_i$ / T' (cf 2.4.); alors l'ensemble Z des
$\check{P}$ tels que T T' est de codimension 2. Pour le voir, on se
ramène au cas T' T (remplaçant T' par T T'), et où T est irréductible
(le remplaçant par une composante irréductible convenable). Alors, \xout{xx}
d'après \inserted[typed]{$Z_o$} ce qu'on a vu, si $\tilde{P}$ n'appartient pas à une certaine partie
fermée de \unsure{codimension} 2, chacun\xout{e} des \xout{\unsure{fibres}} fermés T', T}

\mspage{32}

\clearpage
% Nouveau EGA V (cote 25): archive cote 25, site page 33 (PDF page 34).
% typescript. Draft compared: none.
% AI-transcribed by Claude Opus 5.5 (claude-opus-5-5, adaptive thinking, effort high), Message Batches API, batch ega25b (msgbatch_01A7n5YuXubUPNK1zsepi8RG), 2026-11-13; not yet checked by eye.
\scanpage{25/33}
\mspage{29}

Noter qu'avec les notations précédentes, si $\hand{\zeta\in}$ $\check{P}$ - Z , alors pour tout
y $\hand{\in}$ Y$_{\hand{\zeta}}$ , on a coprof$_y$ F$_{k(\hand{\zeta})}$ = coprof$_y$ G$_{\hand{\zeta}}$ , et par suite: si coprof F$\hand{\leq}$n,
alors pour $\hand{\zeta\in}\check{P}$ - Z, on a coprof G$_{\hand{\zeta}}$ $\hand{\leq}$ n; en particulier si F est CM, alors
pour $\hand{\zeta\in}$ $\check{P}$ - Z, G$_{\hand{\zeta}}$ est CM. \inserted{Enfin, si F est $(S_k)$, \unsure{on sait que} $G_\zeta$ $(S_{k-1})$ pour $\zeta\in\check{P}-Z$ (réf EGA IV).}
\hand{8.3.} On notera que \deleted{on n'a pas de résultat analogue pour la propriété} \inserted{\illegible{6}}
\deleted{(S$_k$)}, même si F=$\underline{O}_X$ , k=1, X étant géom intègre de dim.2 (ou \inserted{k=2,} X étant
géométriquement intègre et géométriquement normal de dim.3). Il suffit
de partir d'une surface projective \xout{xxxxxx} intègre \inserted[typed]{X $\subset$ P$^r$} sur k algébriquement
clos, ayant un point \inserted[typed]{x} où X n'est pas Cohen-Macaulay; alors pour tout
hyperplan passant par x, la section hyperplane Y$_{\hand{\zeta}}$ correspondante admet
x pour cycle associé immergé (resp. on part d'une variété normale\hand{,} \hand{(}donc
\inserted[typed]{(S$_2$)\hand{)},} intègre X $\subset$ P$^r$ de dimension 3, ayant un point x $\in$ X en lequel X n'est pas
Cohen-Maaaulay, alors les Y passant par x ne sont pas \deleted{\unsure{(S}} CM i.e. ne sont
pas (S$_2$) en x)\hand{.} \over{d}{D}ans ces exemples, l'ensemble des \hand{$\zeta$} ``exceptionnels''
pour la propriété (S$_k$) contieht l'hyperplan d e $\check{P}$ défini par x$\in$P, et est
de codimension 1 (et non de codimension$\geq$2). \inserted{Cf \deleted{\illegible{1}} 8.5. plus bas pour \deleted{\unsure{un}} résultat général \unsure{précis} dans cette voie.}

\deleted{Cependant, si F satisfait à la propriété (S$_k$), on peut \unsure{montrer|monteer} à
l'aide de ce qui précède que l'ensemble des $\hand{\zeta\in}$ $\check{P}$ tels que $\phi_{\hand{\zeta}}$ soit
F$_{k(\hand{\zeta})}$-régulier\hand{,} \uline{et} G$_{\hand{\zeta}}$ satisfai\xout{\illegible{1}}t (S$_{k-1}$), a un complémentaire qui est de
codimmnsion $\geq$ 2. Cela résulte du critère 5.7.4. via la coppofondeur,
du fait que l'ensemble Z$_n$ des points de X où F est d e coprofondeur $>$n est
fermé, et que la codimension de Z$_{n\hand{\zeta}}$ dans X$_{\hand{\zeta}}$= Y$_{\hand{\zeta}}$ , même si elle peut
faire (pour d es \hand{$\zeta$} exceptionnels) des sauts \deleted{en hauteur} \inserted{vers le bas}, ledit saut ne peut
jamais être que d'une unité au plus, pour des raisons de dimension éviden-%
tes.}

\inserted{Si F est $(S_k)$, \unsure{alors} l'ens. des $\zeta\in\check{P}$ tels que $\phi_\zeta$ soit $F_{k(\zeta)}$-régulier et $G_\zeta$ soit $(S_k)$ ait un complémentaire de codim $\geq 2$,}

\margin{33}

\clearpage
% Nouveau EGA V (cote 25): archive cote 25, site page 34 (PDF page 35).
% typescript. Draft compared: none.
% AI-transcribed by Claude Opus 5.5 (claude-opus-5-5, adaptive thinking, effort high), Message Batches API, batch ega25b (msgbatch_01A7n5YuXubUPNK1zsepi8RG), 2026-11-13; not yet checked by eye.
\scanpage{25/34}
\xout{\illegible{1}}/irr $X_j$ de X telle que codim($T_i$,$X_j$)=k, et \deleted{dim f($X_j$)} = 0, \inserted{(i.e. si f est quasi-fini \unsure{et k>0}, \deleted{\unsure{que}} T n'ait pas de pt isolé x tel que $\dim_x X = k$ ).}
\margin{\unsure{30}}
Signalons le résultat dimensionnel suivant:

\uline{Proposition} 8.4. Soit T une partie fermée de X, et supposons
dodim(T,X) $\geq$ k. Alors pour tout \hand{$\xi\in\check{P}$}, on a dodim (\hand{$T_\xi$},$Y_{\hand{$\xi$}}$) $\geq$ k-1. Soit
Z l'ensemble des \hand{$\xi\in$} $\check{P}$ tels que codim ($T_{\hand{$\xi$}}$,$Y_{\hand{$\xi$}}$ ) = k-1, (i.e.
codim ($T_{\hand{$\xi$}}$,$Y_{\hand{$\xi$}}$) $<$ k) alors Z est une partie constructible rare de $\check{P}$ ,
\inserted[typed]{i.e.}\xout{\illegible{2}} $\overline{Z}$ \inserted[typed]{est}\xout{xxxx} de codimension $\geq$ 1 dans $\check{P}$ . Pour qu'il soit de
\deleted{\illegible{5}}
codimension $\geq$ 2, il f et s que pour toute composante irréductible
$T_i$ de X de codimension \inserted[typed]{égale à k, \xout{\illegible{3}}\deleted{existe} et telle} \xout{égale à k, ont dim f($T_i$) xxxx 0, et xxxxxxxx compo-}
que dim f($T_i$) = 0, il e\xout{xi}xiste une composante \hand{/}
La première assertion résulte aussitôt du lemme suivant \inserted[typed]{(8.4.1.a)}, qui est

un remords au par.5 :

\uline{Lemme} 8.4.1. Soient X préschéma loc noeth, $\underline{L}$ un Module inversible sur
X , $\phi$ une section de $\underline{L}$, Y = V($\phi$), T une partie fermée de X. \inserted[typed]{Supposons dodim(Y,X)$\geq$k.} Alors :

a) Codim(T$\hand{$\cap$}$Y,Y) $\geq$ k-1.

b) Pour qu'on ait codim(T Y,Y) = k-1 i.e. codim(T$\hand{$\cap$}$Y,Y) $<$ k, il
f et s qu'il existe une composante irréductible $T_i$ de T, contenue dans
Y, \xout{ex} telle que codim($T_i$,X) = k, et telle que pour toute composante
irréductible $X_j$ de X \xout{px} contenant $T_i$ et telle que dim $\underline{O}_{X_j,T_i}$ =
dim $\underline{O}_{X,T_i}$ ; (=k ), on ait $X_j \hand{$\not\subset$}$ Y.

La vérification de ce lemme est immédiate grâce aux généralités
de $O_{IV}$ sur la dimension. Sous les conditions de 8.4., on voit \inserted[typed]{par 8.4.1. b)} quels \inserted[typed]{sont}
les hyperplans exceptionnels $H_{\hand{$\xi$}}$ \xout{sont xceux xpour xlesquels xil xexiste xun}
\xout{exemple xde xcomposantes xirréd xx ($T_i$ xxxx $X_j$) xx avec xxx codim x($T_i$,$X_j$) xx = k}, Si
on excepte l'ensemble $Z_o$ des \hand{$\xi\in$} $\check{P}$ tels qu'il existe une composante
irr R de T ou de X telle \inserted[typed]{dim f(R) $>$ 0 et} que $R_{\hand{$\xi$}}$ soit de ``dimension trop grande'',
\xout{etxxqua} (ensemble qui est de codimension $\geq$ 2 et par suite ne compte pas
les $H_{\hand{$\xi$}}$ \xout{sont} exceptionnels sont ceux pour lesquels il existe
un $T_i$ avec codim($T_i$,X) = k , et \xout{telxquexpourxtout} dim f(T ) = 0,

\mspage{34}

\clearpage
% Nouveau EGA V (cote 25): archive cote 25, site page 35 (PDF page 36).
% typescript. Draft compared: none.
% AI-transcribed by Claude Opus 5.5 (claude-opus-5-5, adaptive thinking, effort high), Message Batches API, batch ega25b (msgbatch_01A7n5YuXubUPNK1zsepi8RG), 2026-11-13; not yet checked by eye.
\scanpage{25/35}
\mspage{\hand{31}}
\xout{xxxxxx} f(T) $\subset$ H , et tel\unsure{\hand{s}} que pour \uline{toute} composante irr $X_j \hand{\supset} T_i$ de X
avec codim $(T_i,X_j)=k$, on ait $f(X_j)\hand{\not\subset} H_{\hand{\zeta}}$ . Pour un $T_i$ donné,
\xout{\illegible{3}} avec codim$(T_i,X)=k$, s'il existe un $X_j$ avec codim$(T_i,X_j)=k$
et tel que dim $f(X_j)$ =0, alors \xout{(d'après \illegible{2} hypothèse}
\xout{\unsure{P} \unsure{$Z_o$})} on aura $f(X_j) = f(T_i) \hand{\subset} H_{\hand{\zeta}}$, et par suite $\hand{\zeta}$ ne
sera pas exceptionnel relativement à $T_i$ . Si par contre pour tout $X_j$
$\hand{\supset} T_i$ tel que codim $(T_i,X_j) = k$, on a dim $f(X_j) \hand{>} 0$, alors pour
$\hand{\zeta \in} \hand{\check{}}P - Z_o$ , $\hand{\zeta}$ sera exceptionnel relativement à $T_i$ si et seulement si
$f(T_i) \hand{\subset} H_{\hand{\zeta}}$ ; l'ensemble de ces $\hand{\zeta}$ est (la trace sur P - $Z_o$\over{)}{\unsure{$\neq$}} d'\hand{)}un hyper-
plan de $\hand{\check{}}P$ . Cela prouve 8.4. , et établit même le résultat plus précis
que l'ensemble exceptionnel est la réunion d'un ensemble de codimension
$\hand{\geq}$ 2 et d'une \uline{réunion d'hyperplans}, déterminés de façon évidente par la
démonstration précédente. (J'avoue que la rédaction est assez floppy, \margin{\hand{|}}
en ce que j'ai raisonné géométriquement tout le temps sans le dire, i.e.
en prenant des points dans des corps alg clos. Bien entendu, la conditi-
on énoncée dans 8.4. est \inserted{bien} géométrique, de sorte qu'on pouvait supposer
k alg clos et raisonner uniquement avec des points rationnels sur k).

\deleted{\uline{Corollaire} 8.5. Soit F cohérent sur X et satisfaisant la condition}
\deleted{$(S_k)$. \inserted{Supposons que pour tout cycle premier associé $R_i$ de F, on ait dim $f(R_i)>0$} Alors l'ensemble des points $\hand{\zeta \in} \hand{\check{}}P$ tels que $\hand{\phi_\zeta}$ soit $\over{F}{F}_{k(\hand{\zeta})}$-régulier}
\deleted{et que $G_{\hand{\zeta}}$ satisfasse la condition $(S_{k-1})$ est constructible, et son}
\deleted{complémentaire est de codimension $\hand{\geq}$ 2.} \margin{\unsure{plus simplement ce qui suit}}

\deleted{Cela résulte aisément de 8.4., \xout{et} 5.7.4. \inserted{et la fin de 8.2.,} En fait, on a le}
\deleted{Complément} \hand{$\delta$} \inserted{\uline{Corollaire 8.5.} \unsure{Supposons} que pour tt cycle premier ass. R, on ait dim f(R)\illegible{1}}
\deleted{résultat plus précis:} \over{P}{p}our que l'ensemble (constructible) des points \unsure{$\zeta$}
de $\hand{\check{}}P$ tels que $\hand{\phi_\zeta}$ soit $F_{k(\hand{\zeta})}$-régulier et $G_{\hand{\zeta}}$ \uline{soit $(S_k)$} \deleted{\illegible{1}} ait un complé-
mentaire de codimension $\hand{\geq}$ 2, il f et s qu'on ait ceci: pour tout \margin{et supposons que F satisfait $(S_k)$.}
entier n $\hand{\geq}$ 0, désignant par $Z_n$ l'ensemble des $x \hand{\in} \over{X}{\unsure{T}}$\inserted{$=$ supp F} tels que coprof$_x$(F) $\hand{>}$
\deleted{et par T le support de X, F,} on veut que pour toute composante irr $Z_{ni}$

\hand{$\delta$} \inserted{Utilisant 8.4. 5.7.4. et la fin de 8.2. on obtient :}
\margin{35}

\clearpage
% Nouveau EGA V (cote 25): archive cote 25, site page 36 (PDF page 37).
% typescript. Draft compared: none.
% AI-transcribed by Claude Opus 5.5 (claude-opus-5-5, adaptive thinking, effort high), Message Batches API, batch ega25b (msgbatch_01A7n5YuXubUPNK1zsepi8RG), 2026-11-13; not yet checked by eye.
\scanpage{25/36}
\mspage{32}

de Z$_n$ , avec codim (Z$_{ni}$,T) = n+k+\unsure{1} et dim f(Z$_{ni}$) = 0 , il existe
une compos ante irréductible T$_j$ de T contenant Z$_{ni}$ telle que
codim (Z$_{ni}$,T$_j$) = n+k+1 et dim f(\xout{Z}T$_j$) = 0.

Lorsque f est quasi-fini donc pour toute partie fermée R de X,
on a dim f(R) = dim R, alors ce critère prend la forme: \inserted[typed]{il n'existe} \xout{pourxtoutxx}
\inserted[typed]{pas de point isolé \xout{dim} z d'un des Z$_n$}
\xout{pointxisoléxxxzxxxdexZx} \xout{Zxxxxxxx} tel que dim$_z$T ( = dim F$_z$ ) soit égal à
n+k+1 \xout{(xàxjoindrexàx8.6.xx)x} Lorsque F est équidimensionnel, \xout{cela}
de dimension d, \xout{celaxsignifiexaussixsimplementxquexxZ} cette condition
est vide si d $\hand{\leqslant}$ k (\xout{et} en effet, nous le savions, car dans c e cas l'hyp
(S$_k$) sur F n'est autre que l'hyp CM), et si d $\hand{\geqslant}$ k+1, elle signifie
que l'ensemble Z$_{d - (k+1)}$ des points de T où la co\unsure{prof}ondeur de F est
$\hand{>}$ d-(k+1) (lequel ensemble est de codimension $\hand{\geqslant}$ d \inserted{i.e. fini} en vertu de
l'hypothèse (S$_k$):sur F) est \uline{vide} , i.e. que \xout{l'ensemblexdes}
l'on a \xout{coprofxFxxxxd-(k+1)}
\[ \text{coprof F } \hand{\leqslant} \text{ d- (k+1)} \]
i.e. vraie prof F $\hand{\geqslant}$ k+1
(alors qu'à priori on a seulement vraie prof F $\hand{\geqslant}$ k comme conséquence
de la propriété (S$_k$) et k $\hand{\leqslant}$ d). \inserted[typed]{Lorsqu'on ne suppose plus F} \deleted{Ce cas particulier mérite d'être}
\deleted{également explicité en corollaire.} équidimensionnel, il reste qu'on
peut exprimer la condition voulue de la façon simple suivante:

\margin{\hand{(8.6.)}} Pour tout \inserted{point fermé} x $\hand{\in}$supp F tel que dim F$_x$ $\hand{\geqslant}$ k+\hand{1}, on a prof F$_x$ $\hand{\geqslant}$ k+1 .

La suffisance se voit aussitôt en faisant z=x, \over{La}{la} nécessité \inserted{se voit} \deleted{aussi,} en
notant que pour tout $\hand{\zeta}$ tel que $\hand{\phi_\zeta}$ soit X$_{k(\hand{\zeta})}$\hand{-}régulier et x $\hand{\in}$ Y$_{\hand{\zeta}}$ ,
on aura dim G$_{\hand{\zeta} x}$ = dim F$_x$ - 1 , prof G$_{\hand{\zeta} x}$ = prof F$_x$ - 1 , donc si
x met en d éfaut la condition plus haut, on aura prof \over{F}{G}$_{\hand{\zeta} x}$ $\hand{<}$k, mais
dim G$_{\hand{\zeta} x}$$\hand{\geqslant}$k, ce qui montre que G$_{\hand{\zeta}}$ ne satisfait pas à (S$_k$) en x; or
l'ensemble d es $\hand{\zeta}$ tels que x$\hand{\in}$ Y$_{\hand{\zeta}}$ est \uline{de codimension} 1 (NB j'ai impli
\margin{|} citement supposé dans le raisonnement que k est alg clos, cas auquel on
se ramène aussitôt). Le critère général précédent doit être \inserted{\unsure{énoncé en cor.}} \deleted{joint à} 8.6
% NOTE: the right edge of the page is cut; the last characters after "8.6" (and the end of "l'hyp", "impli") may be lost
\mspage{36}

\clearpage
% Nouveau EGA V (cote 25): archive cote 25, site page 37 (PDF page 38).
% typescript. Draft compared: none.
% AI-transcribed by Claude Opus 5.5 (claude-opus-5-5, adaptive thinking, effort high), Message Batches API, batch ega25b (msgbatch_01A7n5YuXubUPNK1zsepi8RG), 2026-11-13; not yet checked by eye.
\scanpage{25/37}
\mspage{\hand{\unsure{39|33}}}

8.\over{3}{4}. Nous étudions maintenant les points  y de $\underline{Y}$ qui sont non lisses
sur $Y_{\hand{\zeta}}$ rel à $k(\hand{\zeta})$. Nous nous  bornons au cas où  $f:X\to P$  est \uline{non ramifié}
(pratiquement, ce sera une immersion), et où $X\to S$ est lisse.  Nous ne
supposons pas nécessairement S le spectre d'un corps.  \xout{Nous nous bornons}
\xout{à la considération des points de $X_P$ qui sont dans \unsure{l'ouvert U de}} Comme
f  est non ramifié, l'homomorphisme canonique   $f^*(\underline{\Omega}^1_{P/S})\to\underline{\Omega}^1_{X/S}$
est surjectif, et son noyau est un Module loc libre sur X, que nous
noterons  $\hand{\check{\nu}}_{X/P}$ ; lorsque f est une immersion, ce n'est autre que le
Module \uline{conormal}  $\underline{J}/\underline{J}^2$  défini par l'Idéal $\underline{J}$ de X dans P , et nous l'appe-%
lerons en tous cas Module conormal:
\[ \hand{(a)} \qquad 0\to\hand{\check{\underline{\nu}}}_{X/P}\to f^*(\underline{\Omega}^1_{P/S})\to\underline{\Omega}^1_{X/S}\to 0 \ . \]
\margin{$\|\|$ (vertical pencil strokes)} Notons que nous avons aussi sur  P  une suite exacte canonique
(qui pouvait figurer à titre d'exemple au par.16 p.ex.)
\[ \hand{(b)} \qquad 0\to\underline{\Omega}^1_{P/S}(1)\to\underline{E}_P\to\underline{O}_P(1)\to 0 \]
(i.e. $\underline{\Omega}^1_{P/S}$ est can isomorphe au noyau de l'homomorphisme canonique
$\underline{E}_P(-1)\to\underline{O}_P$  déduit de $\underline{E}_P\to\underline{O}_P(1)$ ), d'où en appliquant $f^*$ :
\[ \hand{(b')} \qquad 0 - f^*(\underline{\Omega}^1_{P/S})(1) - \underline{E}_X - \underline{O}_X(1)\to 0 \ \hand{,} \]
qui donne une description explicite de  $f^*(\underline{\Omega}^1_{P/S})$  \xout{\unsure{de X}} sur X, et
permet  donc d'identifier $\hand{\check{\underline{\nu}}}_{X/P}(1)$  à un sous-Module loc facteur direct
de $\underline{E}_X$ , ou encore le \xout{Module conormal} dual  $\hand{\underline{\nu}}_{X/P}(-1)$  à un Module
quotient de $\hand{\check{}}\underline{E}_X$ .  Par suite,   $P(\hand{\underline{\nu}}_{X/P}\deleted{\unsure{$*$}}(-1)\ )$  =  $P(\hand{\underline{\nu}}_{X/P})$  se plonge
canoniquement dans  $P(\hand{\check{}}\underline{E}_X) = X\times_S\hand{\check{}}P = X_{\hand{\check{}}P}$ comme un sous-fibré projectif sur
X, donc comme sous-préschéma fermé. Ce dernier est nécessairement contenu
dans $\underline{Y}$  (du fait que $\underline{\Omega}_{X/P}(1)$  est contenu dans le noyau de  $\underline{E}_X\to\underline{O}_X(1)$)
\deleted{L'ensemble sous-jacent  à ce sous-préschéma n'est autre que l'ensemble des}
\deleted{\xout{xxxxP} $y\in\overline{Y}$ \xout{\illegible{1}} \uline{qui sont} \inserted{\deleted{\unsure{des zéros x singuliers pour}}} \deleted{"singuliers sur $Y_{\hand{\zeta}}$ " au sens suivant:}}
\deleted{\margin{\unsure{des zéros singuliers}} \deleted{un point est dit "non singulier" sur $Y_{\hand{\zeta}}$ si $\hand{\phi_\zeta}$ est \unsure{régulier}} sur $X_{k(\hand{\zeta})}$}
\deleted{en y et \xout{\illegible{1}} $Y_{\hand{\zeta}}$ est lisse sur $k(\hand{\zeta})$ en y.}

\mspage{\hand{37}}

\clearpage
% Nouveau EGA V (cote 25): archive cote 25, site page 38 (PDF page 39).
% typescript. Draft compared: none.
% AI-transcribed by Claude Opus 5.5 (claude-opus-5-5, adaptive thinking, effort high), Message Batches API, batch ega25b (msgbatch_01A7n5YuXubUPNK1zsepi8RG), 2026-11-13; not yet checked by eye.
\scanpage{25/38}
\mspage{\hand{34}}

\margin{$|$}L'ensemble sous-jacent à ce préschéma n'est autre que l'ensemble des
points de $\underline{\hand{\check{}}Y}$= V($\phi$) qui sont \uline{zéros singuliers}\inserted{(par.16)} de la section $\phi$ de
$\underline{O}_{X\mathrm{x}\hand{\check{}}P}$(1,1) \uline{relativement à la base} $\hand{\check{}}$P , i.e. ses points \xout{gé} à valeurs dans
un corps k sur $\hand{\check{}}$P sont les points \inserted{x} de $Y_k\subset X_k$ tels que \xout{\unsure{xx bien} $\phi_k$ \unsure{soit}}
\xout{\unsure{identiquement nulle en x, mais bien}} $\phi_k$ s'annule à l'ordre deux au moins
en x, i.e. tel que \unsure{$\underline{Y}_k$|$\underline{V}_k$} ne soit pas lisse de dimension relative r-1
sur k en x. La caractérisation énoncée des \xout{points} zéros singuliers comm
les éléments du sousèpréschéma lisse P( $\hand{\unsure{\underline{\Omega}}}_{X/P}$) de $X_{\hand{\check{}}P}$ provient aussitôt
de l'énoncé suivant, qui mérite de figurer en proposition préliminaire:
si S=Spec(k), et si H est un hyperplan de P , alors Y = X$\mathrm{x}_P$H est
\xout{x \unsure{sur k} xxxxxxxxxxx}
\xout{lisse en x \unsure{le point rationnel} sur k x si et seulement si} lisse sur k de
dimension relative d-1 en le point x $\in$Y(k) (i.e. x est un zéro non
singulier, i.e. géom non singulier, de la section $\phi$ de $\underline{O}_X$(1) définie
par H), si et seulement si H ne contient pas l'image par $\phi'$ de l'espa-%
ce tangent à X en x (rel à k) - ou \deleted{encore} comme on dit\xout{\illegible{1}} encore
(lorsque f:X$\longrightarrow$P est une immersion, permettant d'identifier X à un
sous-préschéma de P) sss H \uline{n'est pas tangent à} X \uline{en} x. Cela résulte
trivialement du critère jacobien de lissité, ou de la définition d\unsure{u}
zéro singulier, une fois qu'on a bien précisé le sens de l'énoncé,
c'est à dire qu'on a précisé comment un \xout{espace tangent} sous-espace
\margin{$|$}vectoriel de l'espace tangent à P en un point a (=f(x)) définit une
sous-variété linéaire de P (de telle façon que cela \inserted{a} un sens de dire
que H ne contient pas ledit sous-espace vectoriel): cela provient
bien entendu de la suite exacte (b) ci-dessus, qui permet de\xout{identifier}
définir une correspondance biunivoque entre l'ensemble des sous-espaces
vectoriels de l'espace tangent en a, et les sous-variétés linéaires de
P contenant a. Cette correspondance d'ailleurs revient à associer à

\margin{38}

\clearpage
% Nouveau EGA V (cote 25): archive cote 25, site page 39 (PDF page 40).
% typescript. Draft compared: none.
% AI-transcribed by Claude Opus 5.5 (claude-opus-5-5, adaptive thinking, effort high), Message Batches API, batch ega25b (msgbatch_01A7n5YuXubUPNK1zsepi8RG), 2026-11-13; not yet checked by eye.
\scanpage{25/39}
\mspage{\hand{\unsure{35|33}}}

une sous-variété linéaire passant par a son espace tangent en a, consi-%
déré comme sous-espace \xout{tangent} de l'espace tangent à P en a. Ce genre
\margin{\hand{|}}de \unsure{sorites|soritites}, ensemble avec divers sorites sur les sous-variétés linéaires
et les grassmaniennes, devraient être donnés dans un ou deux n$^{\text{o}}$s prélimi
naires, bien entendu en énonçant tout systématiquement sur une base
quelconque.

En fait, il y a mieux, savoir que le \uline{sous-préschéma} $\underline{Y}^{\mathrm{sing}}$
des zéros singuliers de $\phi$ relativement à $\overset{\hand{$\vee$}}{P}$ , défini dans par.16, n'est
autre que $P(\hand{$\unsure{\underline{\nu}}$}_{X/P})$ , et (comme ce dernier est lisse sur S de dimension
relative d \unsure{$+$} (r-d-1) = r - 1 (r étant la dim relative de P sur S)
on est sous les conditions favorables étudiées au par.16 . \deleted{Pour vérifi
er cette assertion, on utilise le fait suivant (qui devrait figurer au
par. 17 des morphismes lisses), en l'appliquant à $X_{\check{P}}$ \inserted[typed]{et $\underline{P}^1(\underline{O}\ (1,1))$} S : Soit
X lisse sur S de dimension relative m, $\underline{M}$ un Module localement libre sur
X de rang n, $\mu$ une section de $\underline{M}$ , x V( $\mu$ )=Z , supposons que
dim$_x$ Z$_s$ = n-m , \xout{et} (resp. et que $Z_{\underline{s}\ \text{réd}}$ soit lisse sur $\overline{k(s)}$ en $\bar{x}$ )
alors Z est \xout{lisse sur} plat sur S en x (resp. lisse sur S en x). (Dém:
on peut supposer M = $\underline{O}_X^{\,m}$ , donc $\mu = (\mu_1,..,\mu_n)$ , alors il résulte de
la théorie de la dimension et du fait que $\underline{O}_{X_s,x}$ est CM que les
$\mu_{i\ s}$ forment une suite $\underline{O}_{X_s}$-régulière en x, donc \deleted{les $\mu_i$ sont transvers
salement $\underline{O}_X$-réguliers rel. à S; di de plus Z}} Pour le vérifier,
notons que par définition, $\underline{Y}^{\mathrm{sing}}$ n'est autre que le sous-préschéma \inserted{de $\underline{Y}$}
des zéros de la section \inserted{$\psi =$} $d\phi|\underline{Y}$ de $\underline{\Omega}^1_{X_{\check{P}}/\check{P}}(1,1)\otimes \underline{O}_Y = \underline{\Omega}^1_{X/S}\otimes\underline{O}_Y(1,1)$\hand{.}
Nous allons donner une autre interprétation de cette section, dont
la conclusion résultera aussitôt. Pour ceci, considérons le diagramme
\deleted{de Modules} \inserted{\unsure{suites exactes}}
suivant sur $X_{\check{P}}$ , ou plus généralement sur un préschéma Z \inserted{quelconque} au dessus de
$X_{\check{P}}$ :

\mspage{\hand{39}}

\clearpage
% Nouveau EGA V (cote 25): archive cote 25, site page 40 (PDF page 41).
% typescript. Draft compared: none.
% AI-transcribed by Claude Opus 5.5 (claude-opus-5-5, adaptive thinking, effort high), Message Batches API, batch ega25b (msgbatch_01A7n5YuXubUPNK1zsepi8RG), 2026-11-13; not yet checked by eye.
\scanpage{25/40}
\mspage{\hand{36}}

\begin{diagramblock}\begin{tikzcd}
 & \unsure{0} & & \hand{0} & & \\
 & \underline{\Omega}^1_{X/Y} \otimes \underline{O}_Z(1,0) \arrow[u] & & \hand{\unsure{\mathcal{T}}}\check{}_{P/S}\,\unsure{(\otimes\,\underline{O}_Z(0,-1)} \arrow[u] & & \\
0 \arrow[r] & \underline{\Omega}^1_{P/S} \otimes \underline{O}_Z(1,0) \arrow[r] \arrow[u] & E\otimes\underline{O}_Z \arrow[r] \arrow[u] & \underline{O}_Z(1,0) \arrow[r] & 0 \\
 & \hand{\unsure{\check{\mathcal{N}}}}_{X/P} \otimes \underline{O}_Z(1,0) \arrow[u] & \underline{O}_Z(0,-1) \arrow[u] \arrow[ul, dashed, "\hand{\beta}"'] \arrow[ur, "\hand{\alpha}"'] & & & \\
 & 0 \arrow[u] & \hand{0} \arrow[u] & & &
\end{tikzcd}\end{diagramblock}

où la première colonne est déduite de (a) en tensorisant par $\underline{O}_Z(1,0)$,
la ligne est déduite de  (b) en tensorisant par $\underline{O}_Z$  , et la colonne 2
est déduite  de la \inserted{transposée de la} suite analogue \xout{xxxxxxx} (b$\hand{^\vee}$ ) relative à $\check{P}$  (obtenue
en remplaçant $\underline{E}$ par $\hand{\check{\underline{E}}}$ )  en tensorisant par $\underline{O}_Z$ . Par définition même
de $\underline{Y}$, Z se trouve au dessus de $\underline{Y}$  si et seulement si le morphisme
composé \hand{$\alpha$} du diagramme est nul, i.e.  si on peut trouver une  factorisa
tion \hand{$\beta$}: $\underline{O}_Z(0,-1) \rightarrow \underline{\Omega}^1_{P/S}\otimes \underline{O}_Z(1,0)$ .  S'il en est ainsi, on  peut
considérer  son composé avec  $\underline{\Omega}^1_{P/S} \otimes \underline{O}_Z(1,0) \rightarrow \underline{\Omega}^1_{X/Y} \otimes \underline{O}_Z(1,0)$,
qu'on peut aussi interpréter comme une section de $\underline{\Omega}^1_{X/Y}\otimes \underline{O}_Z(1,1)$.
Je dis que c'est précisément la section \inserted{$\unsure{\nu|\mu}$} qu'on avait envisagé plus haut.
(La vérification doit être essentiellement mécanique).  Elle est
nulle  si et seulement si  Z est au dessus de V(\hand{$\unsure{\nu|\mu}$})  (par définition
même de V(\hand{$\unsure{\nu|\mu}$}) !) , mais cela signifie aussi  que \hand{$\beta$} se factorise par
$\hand{\unsure{\check{\mathcal{N}}}}_{X/P} \otimes \underline{O}_Z(1,0)$, i.e. que  le sous-Module $\underline{O}_Z(0,-1)$ de $E\otimes \underline{O}_Z$ est contenu
dans le sous-Module $\hand{\unsure{\mathcal{N}}}_{X/P} \otimes \underline{O}_Z(1,0)$,  ce qui évidemment signifie aussi
que  Z se trouve au dessus du sous-préschéma  P( $\hand{\unsure{\mathcal{N}}}_{X/P}$ (1)) de P($\hand{\check{\underline{E}}}_X$) , et
achève la démonstration de ce que nous avions annoncé.

Dès avant ce savant exercice de syntaxe (qui m'a couté d'ailleurs

\margin{\hand{40}}

\clearpage
% Nouveau EGA V (cote 25): archive cote 25, site page 41 (PDF page 42).
% typescript. Draft compared: none.
% AI-transcribed by Claude Opus 5.5 (claude-opus-5-5, adaptive thinking, effort high), Message Batches API, batch ega25b (msgbatch_01A7n5YuXubUPNK1zsepi8RG), 2026-11-13; not yet checked by eye.
\scanpage{25/41}
% NOTE: the first typed line at the head of the page ("morphisme de préschémas lisses sur S de dim relative r-1 et r:") is enclosed by hand in a loop with an insertion hook; the same hook stands at the right margin after "du" in b). It is transcribed there, as \inserted[typed]{...}.
\mspage{\unsure{37}}

passab\unsure{le}ment de sueur déjà), on aurait pu remarquer que, \deleted{Comme} de
toutes façons du point de vue ensembliste $\underline{Y}^{sing}$ \inserted[typed]{est de} \xout{\unsure{des fibres de}}dimen%
sion r-1 \xout{\unsure{sur S, alors que si}} si S=Spec(k), alors que $\overset{\hand{\vee}}{P}$ est de dimen-%
sion r, \inserted{donc} l'image de $\underline{Y}^{sing}$ dans $\overset{\hand{\vee}}{P}$ est de codimension $\unsure{\geq}$1, ce qui redonne
2.12. (l'argument bien entendu n'est pas essentiellement distinct de
celui utilisé dans 2.12.). On notera que le plus souvent, cet ensemble
est effectivement de codimension 1 (cf plus bas). Par suite, on ne
pourra trouver en général des ``pinceaux linéaires'' de sections hyperpla\xout{n}%
nes toutes \inserted[typed]{\xout{\unsure{géom}}} \xout{\unsure{non singulières}} lisses. Cependant, nous allons voir qu'on
pourra s'arranger souvent pour \xout{que} trouver des pinceaux formés de sec-%
tions hyperplanes n'ayant aucun point supersingulier, grâce au fait que
dans les cas les plus fréquents, l'image de $\underline{Y}^{sup\ sing}$ dans $\overset{\hand{\vee}}{P}$ sera de
codimension 2. Nous allons d'abords résumer les points essentiels
de nature différentielle de la situation étudiée ici:

\uline{Théorème} 8.\over{\illegible{1}}{7}. a) Le sous-préschéma $\underline{Y}^{sing}$ (défini au par.16) dans
la situation actuelle n'est autre que $P(\underline{\hand{\unsure{\nu}}}_{X/P})$ , considéré comme
sous-préschéma de $\underline{Y}$ comme expliqué plus \xout{h}haut.

b) L'ensemble sous-jacent au préschéma $\underline{Y}^{sup\ sing}$ (cf par. 16), \xout{\illegible{1}}
\deleted{\unsure{fermé de}} n'est autre que l'ensemble des points de ramification du \inserted[typed]{morphisme de préschémas lisses sur S de dim relative r-1 et r:}
$\underline{Y}^{sing} = P(\underline{\hand{\unsure{\nu}}}_{X/P}) \longrightarrow \check{P}$ , i.e. pour que ce dernier morphisme soit
non ramfié en le point \over{\illegible{1}}{$y$} (réf à définition) il f. et s. que \over{\illegible{1}}{$y$}
soit géométriquement un point singulier ordinaire pour $\phi_{\hand{\xi}}$ ($\hand{\xi}$ étant le
point de $\overset{\hand{\vee}}{P}$ image de \over{\illegible{1}}{$y$} ).

c) Supposons S=Spec(k), et que \over{\illegible{1}}{$y$} \inserted{$\in$}\inserted[typed]{$\underline{Y}^{sing} = P(\underline{\hand{\unsure{\nu}}})$} soit rationnel sur k; soient x
et $\hand{\xi}$ ses projections dans X(k) re\unsure{sp|pp} $\overset{\hand{\vee}}{P}$(k) , et considérons la \xout{\unsure{sous-espace}}
sous-variété linéaire \inserted{$H'$} de $\overset{\hand{\vee}}{P}$ ``image'' de l'application tangente de

\margin{41}

\clearpage
% Nouveau EGA V (cote 25): archive cote 25, site page 42 (PDF page 43).
% typescript. Draft compared: none.
% AI-transcribed by Claude Opus 5.5 (claude-opus-5-5, adaptive thinking, effort high), Message Batches API, batch ega25b (msgbatch_01A7n5YuXubUPNK1zsepi8RG), 2026-11-13; not yet checked by eye.
\scanpage{25/42}
\mspage{\hand{38}}

$\underline{Y}^{\mathrm{sing}} = P(\underline{\hand{\unsure{\nu}}}_{X/P}) \longrightarrow \hand{\unsure{\check{}}}\!P$ . Alors H'est ``orthogonal'' à x, et en parti-%
culier si dim H' = r-1 (ce qui équivaut par b) au fait que \xout{yx/x$\underline{Y}$xxxxx} \inserted[typed]{\unsure{supsi}}
y est point \inserted{géom.} singulier \uline{ordinaire} sur $Y_{\hand{\unsure{\xi|\zeta}}}$ ) x est uniquement déterminé
en fonction de H'.

Dém . Nous avons déjà prouvé a), et b) en est une conséquence
grâce aux développements du par.16. Reste à prouver c), qui est
encore conséquence d'une relation de compatibilité essentiellement
triviale. Je me borne ici à copier, dans le cas spécial qui nous
occupe (mais avec S quelconqque), le diagramme (D') du par.16, qui
prend la forme suivante (où nous écrivons $\underline{Y}^1$ pour $\underline{Y}^{\mathrm{sing}}$):

\begin{diagramblock}
\begin{tikzcd}[column sep=small]
 & & 0 & & \\
 & & \underline{\Omega}^1_{X/S\,\underline{Y}^1} \arrow[u] & & \\
0 \arrow[r] & \underline{P}^1(\underline{O}(1,1))\otimes\underline{O}_{\underline{Y}^1} \arrow[r] \arrow[ur, "\hand{\unsure{\nu.\mu}}"] & (\underline{\Omega}^1_{X/S})_{\underline{Y}^1} + (\underline{\Omega}^1_{\check{P}/S})_{\underline{Y}^1} \arrow[r] \arrow[u] & \underline{\Omega}^1_{\underline{Y}^1/S} \arrow[r] & 0 \\
 & \underline{O}_{\underline{Y}^1}(-1,-1) \arrow[u] \arrow[r, "\hand{u}"] & \underline{\Omega}^1_{\check{P}/S}\ \underline{Y}^1 \arrow[u] \arrow[ur, "\hand{\nu}"'] & & \\
\underline{\omega}^{-2}_{X/Y}\otimes\underline{O}_{\underline{Y}^1}(-2,-2) \arrow[ur] & & 0 \arrow[u] & &
\end{tikzcd}
\end{diagramblock}

La relation d'orthogonalité c) ne fait qu'exprimer que le composé
$\over{V}{\nu}u$ dans ce diagramme est nul; je me dispense de développer plus
avant ce deuxième exercice de symtaxe. Je suppose que les gens avaient
l'habitude de prendre des coordonnées, et ils y voyaient peut-être plus
clair ....

\uline{Corollaire} 8.\over{\illegible{1}}{8}. Supposons S=Spec(k), \xout{Saxsaxtaxqusxx$\underline{Y}$xxxxx} \inserted{$f : X \to P$ une immersion} et X
irréductible \inserted{(\unsure{f propre} lisse sur k)}. \xout{xxxxx} Alors $\underline{Y}^{\mathrm{sing}}$ est lisse et irréductible; soit T

\margin{42}

\clearpage
% Nouveau EGA V (cote 25): archive cote 25, site page 43 (PDF page 44).
% typescript. Draft compared: none.
% AI-transcribed by Claude Opus 5.5 (claude-opus-5-5, adaptive thinking, effort high), Message Batches API, batch ega25b (msgbatch_01A7n5YuXubUPNK1zsepi8RG), 2026-11-13; not yet checked by eye.
\scanpage{25/43}
\mspage{\unsure{39}}

L'adhérence de son images dans $\overset{\hand{\vee}}{P}$ , muni de la structure réduite
induite, et considérons le morphisme induit g: $\underline{Y}^{\text{sing}} \longrightarrow T$ \xout{de} (morphis-%
me dominant de préschémas \xout{irréductibles} intègres). Alors les condition
\inserted[typed]{(i) à (ii bis)} suivantes sont équivalentes:

(i) Le morphisme g est génériquement \deleted{non ramifié} \inserted{étale} (i.e. est \deleted{non
ramifié} \inserted{étale} en au moins un point, ou ce qui revient au même, est \inserted{étale =} non ramifié
au point générique de $\underline{Y}^{\text{sing}}$).

(ibis) L'extension de corps L/K définie par g est finie et
séparable.

\margin{NB. L'hyp. \unsure{(on fera ici en} \illegible{1} \unsure{lien de : non normalité)} ne sert que pour l'équivalence avec (i ter). Énoncer 8.8. en conséquence\dots}
(i ter) Le morphisme g est birationnel, i.e. l'extension L/K est
l'extension triviale.

(ii) $\underline{Y}^{\text{sihg}} \neq \underline{Y}^{\text{sup sing}}$ \inserted{(ensemblistement, disons).}

(ii bis) Il existe un $x \in X(\bar{k})$ , et un \over{h}{h}yperplan tangent \inserted{H} à $X_{\bar{k}}$
en x qui ne soit pas ``osculateur'' en x (par quoi on entenda précisément
que x n'est pas supersingulier pour la section de $\underline{O}_{X_{\bar{k}}}(1)$ qui définit
H ... ).

Ces conditions impliquent que \inserted{$\underline{Y}^{\text{sup sing}} \neq \emptyset$ \unsure{ou}} dim $\underline{Y}^{\text{sup sing}}$ \deleted{$\leq$}$= r-2$, donc que
l'image de $\underline{Y}^{\text{sup sing}}$ dans $\overset{\hand{\vee}}{P}$ a une \unsure{dodimension|codimension} $\geqslant 2$, et \xout{que T
est de dimension r-1, i.e. de codimension 1 dans P} elles impliquent
également:

(iii) dim T = r-1, i.e. T est de codimension 1 dans $\overset{\vee}{P}$ .

Dém. L'équivalence de (i) et (i bis) est triviale , son équiva-%
lence avec (ii) est une conséquence triviale de 8.\unsure{7}. b), enfin l'équi-%
valence de (ii) et (ii bis) est pratiquement la définition de
$\underline{Y}^{\text{sup sing}}$ . Evidemment (i ter) \hand{$\Longrightarrow$}(i) , il reste à prouver que (i)\hand{$\Longrightarrow$}
(i ter). On peut évidemment supposer k alg clos, et on est ramené à
prouver (compte tenu de l'hypothèse (\unsure{iii}) ) qu\xout{e}'il existe un ouvert U\hand{$\neq\emptyset$}

\margin{43}

\clearpage
% Nouveau EGA V (cote 25): archive cote 25, site page 44 (PDF page 45).
% typescript. Draft compared: none.
% AI-transcribed by Claude Opus 5.5 (claude-opus-5-5, adaptive thinking, effort high), Message Batches API, batch ega25b (msgbatch_01A7n5YuXubUPNK1zsepi8RG), 2026-11-13; not yet checked by eye.
\scanpage{25/44}
\mspage{\hand{40}}

dans T tel que \hand{$\xi$} \inserted{$\in U(k)$} \deleted{V} implique qu'il existe exactement  un point \xout{du} y de
$\underline{Y}^{\text{sing}}(k)$ au dessus de \hand{$\xi$}. Cela résultera de 8.\over{\unsure{6}}{7}. c) qui implique plus
\xout{généralement} précisément:

\uline{Corollaire} 8.\over{\unsure{8}}{9}.  Supposons la condition (i) de 8.\over{\unsure{7}}{8}. satisfaite, et
soit  \xout{\unsure{T(k) un point tel que T soit lisse en} \illegible{1} U \unsure{Il en semble}}
U l'ouvert de T des points en lesquels T est lisse sur k. Alors \inserted{$U\neq 0$, \deleted{et}}
$\underline{Y}^{\text{sing}}|U \longrightarrow U$ est une immersion ouverte, \xout{(et même un isomorphisme} \inserted{injection $Y^{\text{sing}}|U$ ne contient pas des points de $\underline{Y}^{\unsure{\text{sing sing}}}$.}
\xout{si X est propre sur k donc $\underline{Y}^{\text{sing}}$xxxx T surjectif)} Si X est propre sur
k donc $g:\underline{Y}^{\text{sing}}\longrightarrow T$ surjectif, alors $\underline{Y}^{\text{sing}}|U\longrightarrow U$ est un isomorphi
me, et U est le plus grand ouvert de T ayant cette dernière propriété.

Tout d'abord, comme g est dominant et gén \deleted{ramifié} \inserted{étale}, il est
gén. étale, donc on peut \unsure{trouver|tmouver} un ouvert non vide V de T tel que
\xout{\unsure{gén on peut supposer de plus g étale au(x)} \illegible{1} \unsure{que} \illegible{1} \unsure{ne contienne pas} $\underline{Y}$ \illegible{1}}
$\underline{Y}^{\text{sing}}|V \longrightarrow V$ soit étale et surjectif\inserted{\margin{ce qui implique V lisse$/k$}}. Si alors \hand{$\xi$}$\in V(k)$ et si y est
un point de $\underline{Y}^{\text{sing}}(k)$ au dessus de \hand{$\xi$}, alors avec les notations de
8.\over{\unsure{6}}{7}. c) \xout{\unsure{Xlex}} l'espace H' n'est autre que l'espace tangent à T en \hand{$\xi$},
et comme on a noté à cet endroit, cela implique que le point x de X(k)
projection de y est déterminé comme le point orthogonal à H', donc est
uniquement déterminé, donc comme $\underline{Y}^{\text{sing}}\subset X\text{x}$\hand{$\check{P}$}, y est uniquement déter-%
miné. Cela prouve déjà que g est birationnel (étant \inserted[typed]{dominant,} \xout{\unsure{non ramifié}}
génériquement \deleted{non ramffié} \inserted{étale} et gén radiciel). \xout{\unsure{Pour Evidemment V H}}
D'autre part, le morphisme \hand{$\unsure{\nu|\gamma}$}(dont la définition en forme est évidente)
qui associe à tout \over{x}{$\xi$}$\in U(k)$ l'unique point $x=$ \hand{$\nu$}(\hand{$\xi$}) $\in P$ orthogonal à
l'espace tangent à U en \hand{$\xi$}, coïncide sur V avec le composé
$V\rightarrow \underline{Y}^{\text{sing}}|V \longrightarrow X$, où la deuxième flèche est la projection; donc
\inserted[typed]{posant} \xout{\unsure{La restriction de}} $h=($\hand{$\nu$}$,\mathrm{id}) : U\rightarrow P\text{x}T$ , \xout{\unsure{(il induit une immersion à V et est}}
\xout{\unsure{montre que le composé j g}} $g_1 = g|g^{-1}(U)$\xout{=}: $g^{-1}(U)\rightarrow U$\xout{\unsure{Ex}} , le composé

\margin{44}

\clearpage
% Nouveau EGA V (cote 25): archive cote 25, site page 45 (PDF page 46).
% typescript. Draft compared: none.
% AI-transcribed by Claude Opus 5.5 (claude-opus-5-5, adaptive thinking, effort high), Message Batches API, batch ega25b (msgbatch_01A7n5YuXubUPNK1zsepi8RG), 2026-11-13; not yet checked by eye.
\scanpage{25/45}
\mspage{\hand{\unsure{49|41}}}

hg$_1$ :  $g^{-1}(U) \longrightarrow$ PxT  n'est autre que l'inclusion canonique, puisqu'il
en est ainsi de sa restriction à $g^{-1}(V) \rightarrow$ V.  \xout{Soit alors \unsure{U'} h$^{-1}$(}
\xout{Comme h prend} Il en résulte que h \xout{prend ses valeurs dans} se factorise
à travers l'adhérence schématique \xout{de} $\overline{Y^1}$ de Y$^1$ dans  PxT , donc que
l'image inverse de Y$^1$ (qui est un ouvert de ladite adhérence) par h
est un ouvert de U, soit U' . En vertu de  h g$_1$ = inclusion, on
voit alors aussitôt que  g$_1$ induit un isomorphisme  $g^{-1}(U) \simeq$ U', et
\xout{xxxxxx} cela montre bien que  $g^{-1}(U) \longrightarrow$ U est une immersion ouverte.
Lorsque X est propre sur k i.e. f ermé dans P, alors  g est propre donc
surjectif, donc  g$_1$ : $g^{-1}(U) \longrightarrow$ U est surjectif, donc un isomorphisme.
D'ailleurs, si  W est un ouvert de T tel que  $g^{-1}(W) \longrightarrow$ W soit un isomor%
phisme, il s'ensuit que W est lisse puisque Y$^1$ est lisse, donc W$\subset$U.
Cela achève la démonstration de 8.\over{\unsure{8}}{9}. \inserted{\unsure{$\underline{Y}^{\vee}$} \unsure{$r_{\text{sing}} = \rho$} \illegible{1}}

\xout{Remarque}  L\over{'}{es} assertion\hand{s} finale\hand{s} de 8.\over{\unsure{7}}{8}.\hand{,} dim $\underline{Y}^{\text{sup sing}}$ \over{=}{\unsure{$\leq$}} r-2\hand{,} et
dim T = r-1 , sont triviales;  la première résult\over{ant}{\unsure{e}} d\over{u}{\unsure{e}} \deleted{$\underline{Y}^{\text{sup sing}} \neq \underline{Y}^{\text{sing}}$}
\inserted{\unsure{$Y^{\text{sup sing}}$}} \inserted{\unsure{contenu dans} $Y^{\text{sing}}$, \unsure{de dim r}} \deleted{donc rare dans $\underline{Y}^{\text{sing}}$ puisque celui} est irréductible, et du fait que
\inserted{\unsure{est défini par l'annulation d'une section} D \unsure{d'un Module inversible}} \deleted{dim} $\underline{Y}^{\text{sing}}$\hand{\unsure{sup sing}} \deleted{= r-1} ; la deuxième du fait que, L étant fini sur K, on a
deg tr L/k = deg\xout{x} tr K/k, i.e.  \xout{dim $\underline{Y}^{\text{sing}}$ xxxxxx} dim T = dim $\underline{Y}^{\text{sing}}$ = r-1.

\uline{Remarque} 8.\over{\unsure{9}}{10}.  Comme on a si\over{n}{g}\unsure{n}alé dans  8.\over{\unsure{8}}{9}. , avec les notations de
ce corollaire, on a  $g^{-1}(U) \subset \underline{Y}^{\text{sing}} - \underline{Y}^{\text{sup sing}}$ ; mais on notera que
même si X est fermé dans P,  cette inclusion n'est pas nécessairement
\margin{\unsure{$\ast$}} une égalité, en d'autres termes (notant que $g^{-1}(U)$ n'est autre que
l'ensemble des points où g est \uline{étale}, alors que $\underline{Y}^{\text{sing}}$-$\underline{Y}^{\text{sup sing}}$ est
lèensemble des points où g est \uline{non ramifié}) il peut y avoir des points
y de $\underline{Y}^{\text{sing}}$ en \hand{\unsure{´}}lesquels  g est \xout{étale} non ramifié, mais non étale (ce qui
\inserted{\deleted{\illegible{1}}} implique d'ailleurs que g(y) est un point non \xout{normal} géom normal et même
non géom unibranche de  T)\hand{.} \xout{xxxxxx} En termes géométriques, cela corres%
pond au phénomène suivant: on peut avoir un hyperplan tangent non oscul
% NOTE: the last line is cut at the right edge of the scan
\margin{45}

\clearpage
% Nouveau EGA V (cote 25): archive cote 25, site page 46 (PDF page 47).
% typescript. Draft compared: none.
% AI-transcribed by Claude Opus 5.5 (claude-opus-5-5, adaptive thinking, effort high), Message Batches API, batch ega25b (msgbatch_01A7n5YuXubUPNK1zsepi8RG), 2026-11-13; not yet checked by eye.
\scanpage{25/46}
\mspage{42}

teur de X en un point $x\in X(k)$, \xout{\unsure{mais}} tel qu'il existe un \uline{autre point}
$x'\in X(k)$ en lequel le même hyperplan soit tangent en x. On en fait des
exemples évidents avec \xout{\unsure{des exemples}} P de dimension 2, X une courbe non si
singulière de degré $\geqslant$ \unsure{4} , en c aractéristique quelconque. \inserted{[NB les ``tangen-tes doubles'' = \unsure{X} correspondant aux pts doubles de la ``courbe duale'' !].}

\uline{Corollaire} 8.\over{\unsure{10}}{11}. Supposons k de caractéristique nulle. Alors:

a) L'image de $\underline{Y}^{\text{sup sing}}$ dans $\overset{\hand{\vee}}{P}$ est de codimension $\geqslant 2$.

b) La condition (iii) de 8.\over{\unsure{7}}{8}. équivaut aux autres conditions, i.e.
la négation des autres conditions\hand{,} \inserted{sauf $\underline{Y}^{\text{sing}}=\underline{Y}^{\text{sup sing}}$,} signifie aussi que l'image de $\underline{Y}^{\text{sing}}$
dans $\overset{\hand{\vee}}{P}$ est de codimension $\geqslant 2$.

Evidemment l'assertion b) implique a), compte tenu de 8.\over{\unsure{7}}{8}. Or
par la théorie de la dimension, (iii) signifie que L/K est une extensi-%
on finie, \xout{\unsure{et en}} (on pourrait le mettre sous forme (iii bis) dans 8.\over{\unsure{7}}{8}.),
et en car 0 \xout{\unsure{cela implique}} L est de toutes façons séparable sur K, d'où
la condition (i bis) de 8.\over{\unsure{7}}{8}.

\uline{Remarques} 8.1\over{1}{2}.. a) Géométriquement, l'assertion a) signifie essen-%
tiellement que pour un pinceau linéaire assez général de sections hyper-%
planes, tous les membres du pinceau sont lisses ou \xout{n}ont \xout{\unsure{que des points}}
\xout{\unsure{gé}} pour s euls points géométriquement des points géométriquement
singuliers ordinaires (et en f ait, comme on constate aussitôt - et il
y aurait lieu d'énoncer a) en conséquence sous une forme un peu plus
précise - ont au plus \uline{un} tel point géométriquement singulier ). L'as\xout{s}-%
sertion b) signifie \inserted{essentiellement} que lorsque $\underline{Y}^{\text{sing}}=\underline{Y}^{\text{sup sing}}$ (ce qui s'exprime
analytiquement par l\xout{\unsure{'annulation}} \inserted[typed]{nullité} d'une certaine section D d'un Module
inversible $\underline{\hand{\omega}}^{\otimes 2}_{X/\hand{k}}\otimes\underline{O}_{\underline{Y}^1}(1,1)$ sur $\underline{Y}^1$), alors pour un pinceau linéaire
assez général de sections hyperplanes, tous les membres de ce \xout{\unsure{faisceau}}
pinceau sont lisses. Cette deuxième situation (qu'on s\xout{x}oit d'ailleurs
en car.0 ou non ) doit être considérée comme tout à fait exceptionnelle

\mspage{46}

\clearpage
% Nouveau EGA V (cote 25): archive cote 25, site page 47 (PDF page 48).
% typescript. Draft compared: none.
% AI-transcribed by Claude Opus 5.5 (claude-opus-5-5, adaptive thinking, effort high), Message Batches API, batch ega25b (msgbatch_01A7n5YuXubUPNK1zsepi8RG), 2026-11-13; not yet checked by eye.
\scanpage{25/47}
\margin{(\unsure{$\uparrow$|?} la variété \unsure{lin} L de codim 2 des hyperplans tangents à X en x, donc par raisons de dim. on a T = L, donc $X \subset L^{\circ}$ = polaire de L, qui est une droite, donc X est un \unsure{\uline{ouvert}} d'une droite.)}
\mspage{48}
\unsure{\hand{(}}En langage classique, elle s'exprime \xout{xxxdisant} sauf erreur en disant
que X est \uline{réglée} pour l'immersion projective considérée \hand{[}- et si ça nous
chante, nous avons ici  tout ce qu'il faut, grâce à 8.5. et ses corollai
res, pour expliquer et justitffier une telle terminologie, au cas où cela
t'inspirerait de faire le lien avec la taupe\over{)}{]}. \inserted{Par exemple, si dim $X=1$, cela implique que X est une droite, car si $x \in X(k)$, alors T contient \unsure{$\uparrow$|P}}

b)  Lorsque la car. est p >0, il faudrait pour bien faire donner
des exemples (avec dim P = 2, X une courbe non singulière) où  les
conditions de 8.6. ne sont pas vérifiées i.e. $\underline{Y}^{\mathrm{sing}} = \underline{Y}^{\mathrm{sup\ sing}}$, et où
néanmoins  dim T = r-1 , i.e. \inserted{des exemples} où  L/K est une extension finie
inséparable.  J'ai la flemme de construire des exemples, mais je ne
doute guère qu'il y en ait.

\deleted{\uline{Corollaire} 8.11.} Dans a), faire un appel du pied au \xout{xxxxx} nº suivant,
où on montrera que lorsque le cas exceptionnel ``réglé'' se produit, alors \margin{N.B. X est dit ``réglé'' pour f si $\mathrm{Im}(\underline{Y}^{\mathrm{sing}} \to \check{P})$ est de codim $\unsure{\geq} 2$}
par une modification triviale de l'immersion projective considérée, on
se retrouve dans la situation  ``générale''  de 8.8.

\margin{\unsure{$\|$}}La partie du présent nº de 8.\over{\illegible{1}}{6}. à ici  devrait sans doute faire
un nº à part, de nature différentielle,  tandis que le début du nº avec
ce qui va suivre serait bloqué dans un nº  sur  la dimension \unsure{dus|des}
exceptionnels. \hand{(Je n'utilise que le fait que $\underline{Y}^{\mathrm{sing}}$ est de
dim r-1).}

\uline{Proposition} 8.1\over{2}{3}.   Nous supposons tou\xout{x}jours f:X$\to$P non ramifié, \inserted{et X sans pt isolé,} \deleted{Suppo-}
\inserted{On suppose que X satisfait $(R_k)$ géométrique\deleted{, et si}}
\deleted{sons que X satisfasse la propriété $(R_k)$ géométrique  (où k \unsure{$\geq$} 0)\hand{,} et si}
\deleted{X' est \xout{xxensemble} l'ouvert de X des points où X est lisse sur k, \unsure{le}}
\deleted{\xout{\illegible{12}}}
\deleted{\inserted{on suppose que le}morphisme  induit f':X'  P  \xout{soit} satisfasse \xout{auxxconditionsxéquivalentes}}
\deleted{\xout{\illegible{10}} k(  )}
\deleted{\xout{dex8x8xxx(disonsxdoncxsimplementxX'xestxnonxx}  à la condition}
\deleted{g( $\underline{Y}'^{\,\mathrm{\unsure{sup\,sing}}}$) \hand{\unsure{´}}de codimension $\geq$2 dans $\hand{$\check{P}$}$ (condition \uline{automatiquement}}
\deleted{\xout{v}vérifiée si k est de car 0 en vertu de de 8.11. ).}  Soit $Z_k$ la parti
de \hand{$\check{P}$} \xout{fermédesxxxxxx} complémentaire de l'ensemble des \hand{$\xi$}$\in$\hand{$\check{P}$}  tels

\margin{47}

\clearpage
% Nouveau EGA V (cote 25): archive cote 25, site page 48 (PDF page 49).
% typescript. Draft compared: none.
% AI-transcribed by Claude Opus 5.5 (claude-opus-5-5, adaptive thinking, effort high), Message Batches API, batch ega25b (msgbatch_01A7n5YuXubUPNK1zsepi8RG), 2026-11-13; not yet checked by eye.
\scanpage{25/48}
\mspage{\unsure{45|46}}

\deleted{pour toute composante irréductible $X'_i$ de $X'$ on a}
\inserted{grâce au fait que $\underline{X'}$ sing est de dim $r-1$, où $r=\dim \check{P}$ !}

que $\phi_{\hand{$\zeta$}}$ soit $X_{k(\hand{$\zeta$})}$-régulière et $Y_{\hand{$\zeta$}}$ satisfasse la condition $R_k$ géométri-%

que. Alors:

a) \over{B}{P}our que $Z_{k-1}$ soit de codimension $\geq 2$ dans $\overset{\hand{$\vee$}}{P}$ , il \xout{f et} s que
\inserted{$X'_i$ de $X'$ de dim $\leq k$ soit \uline{réglée} pour $f$.} \deleted{\unsure{et $X'$ qui est non réglée rel. à $f$}}
toute composante irréductible \deleted{de $X_i$ soit de codimension $\geq k+1$ et c'est aus-%
si nécessaire si codim$(f(\underline{Y}'_1{}^{\mathrm{sing}}),P) = 1$ (``cas non réglé'').}

b) Pour que $Z_k$ soit de codimension $\geq 2$ dans $\check{P}$ , il \xout{f et} s que
\deleted{\inserted{$X'_i$ \unsure{qui est non réglée rel. à $f$}}}
toute composante irréductible \deleted{de $X$ soit de dimension $\geq k+1$, et que}
\inserted{$X_i$ de dim $\leq k+1$ de $X$ soit formée de pts lisses de $X$ et soit \uline{réglée} pour $f$.}
\deleted{l'ensemble singulier $T = X-X'$ n'ait aucun point isolé $x$ avec}
\inserted{\unsure{i.e. toute composante irréductible $X'_i$ de $X'$ de}}
\deleted{$\dim_x X = k+1$ \hand{;} et c'est aussi nécessaire si chaque $X'_i$ est non réglé rel.}
\deleted{à $f$} \inserted{\unsure{de dim $\leq k+1$ est, \deleted{\unsure{soit}} réglée rel. à $f$, \deleted{\unsure{soit}} telle que $X'_i = X_i$.}}
\inserted{\unsure{soit contenue} \unsure{dans $X'$}}

En effet, pour tout $\hand{$\zeta$}$, l'ensemble géométriquement singulier

de $Y_{\hand{$\zeta$}}$ \xout{xxxxxxxxxxxxxx(pt} (NB on se borne aux $\hand{$\zeta$}$ tels que $\phi_{\hand{$\zeta$}}$ soit $X_{k(\hand{$\zeta$})}$-

régulier, ce qui est inoffensif vu 8.2.) est la réunion de

sing$(Y'_{\hand{$\zeta$}}$ ) et de l'image inverse $T_{\hand{$\zeta$}}$ de $T$ dans $Y_{\hand{$\zeta$}}$ , donc la codimension
\inserted[typed]{l'inf de}
\margin{\unsure{lieu non-ordin. géométrique}}
de cet ensemble singulier dans $Y_{\hand{$\zeta$}}$ est égale à \hand{(}codim (sing$(Y'_{\hand{$\zeta$}})$,$Y'_{\hand{$\zeta$}}$ )

et de codim $(T_{\hand{$\zeta$}},Y_{\hand{$\zeta$}})$. Bornons nous aux $\hand{$\zeta$}$ tels que \deleted{$X - f(\underline{Y}'^{\,\mathrm{sup\ sing}}$}\hand{)}
\margin{$*$} \inserted{sing$(Y'_\zeta)$ soit fini}
\hand{/}(ce qui est inoffensif, car \deleted{par hyp} cela conduit à se mettre dans le

complémentaire d'un ensemble de codimension $\geq 2$\hand{\unsure{$\zeta$}$\}$}), \over{l}{L}es points singulier

géométriques de $Y'_{\hand{$\zeta$}}$ sont \inserted{alors} isolés\over{,}{.} La conclusion résulte facilement de

là , et de 8.4. \deleted{A vrai dire, on peut donner un énoncé plus satis-%
faisant, en disant: il f et suffit que pour toute composante irréd
$X'_i$ de $X'$ qui est réglée rel à f, on ait dim $X'_i$ \xout{kxx} (resp. dim $X'_i$
k+1), et que l'ensemble singulier ..}

Combinant 8.13. \deleted{et} 8.5.\inserted[typed]{et 8.6.}\hand{/}, on trouve à la façon habituelle:

\uline{Corollaire} 8.14. On suppose $f:X\rightarrow P$ non ramifié et $X$ sans pts isolés.

\margin{\unsure{n}} a) Supposons $X$ séparable sur $k$. Pour que l'ensemble \deleted{\unsure{E}} des $\hand{$\zeta \in$}$ $\overset{\hand{$\vee$}}{P}$

tels que $\phi_{\hand{$\zeta$}}$ soit $X_{k(\hand{$\zeta$})}$ -régulier et $Y_{\hand{$\zeta$}}$ soit séparable ait un complé-%

mentaire de codimension $\geq 2$, il f et s qu\deleted{'il n'existe aucun point
géom. singulier \inserted{\deleted{\unsure{isolé}}} de $x$ tel que $\dim_x X = 1$, \deleted{et que toute} \inserted{ni de composante}
\inserted{irréductible $X'_i$ de $X'$ de dim 1 ``non réglée'' pour $f$, et que l'ex}}

\mspage{48}

\clearpage
% Nouveau EGA V (cote 25): archive cote 25, site page 49 (PDF page 50).
% typescript. Draft compared: none.
% AI-transcribed by Claude Opus 5.5 (claude-opus-5-5, adaptive thinking, effort high), Message Batches API, batch ega25b (msgbatch_01A7n5YuXubUPNK1zsepi8RG), 2026-11-13; not yet checked by eye.
\scanpage{25/49}
\mspage{45}

toute composante irréductible \inserted{$X_i$} de dimension 1 \deleted{$X_i$} de X soit \xout{xxxbiex}
\inserted[typed]{et soit}
formée de points lisses de X \xout{xxxbien} réglée relativement à f, et que
pour tout point fermé \xout{de} x de X tel que $\dim_x X \geqslant 2$, on ait $\mathrm{prof}_x X \geqslant 2$,
(conditions automatiquement vérifiées si\unsure{.} X est \inserted{géom} normale et toutes ses
composantes irréductibles sont de dim$\geqslant$ 2).

b) Supposons X géom. normale , \xout{Pour} que l'ensemble \xout{Xxd} des $\hand{\zeta \in}\, \check{P}$
tels que $\phi_{\hand{\zeta}}$ soit $X_{k(\hand{\zeta})}$\hand{-}régulier et $Y_{\hand{\zeta}}$ \xout{soit} géom normale ait un complé-%
mentaire de codimension $\hand{\geqslant}$2, il f et s que \xout{pour} toute composante
irréductible $X_i$ de X de \xout{co}dimension$\hand{\leqslant}$2\unsure{\hand{,}} soit formée de point lisses de
X et soit réglée rel à f, et que de plus pour tout point fermé x de X
tel que $\dim_x X \hand{\geqslant} 3$, on ait $\mathrm{prof}_x X \hand{\geqslant}$ 3.

\xout{(J'aixoubliéxxdaxavantxa)xetxb)xxd'indiquerxl'hypothèse}
\xout{suivante,xtoujoursxvérifiéexenxcarx0xxpourxtoutexcompxirrxX$_i$xdexX$_i$,}
\xout{xxxxsupxsingx}
\xout{\unsure{g(X$_i$}xxxxxxxxxx)}

\uline{Remarque} 8.15. Dans 8.6. , 8.13. et 8.14. nous avons fait sur X
l'hypothèse $(S_k)$ (resp. $(R_k)$, resp\unsure{:} séparable, resp. \inserted[typed]{géom} normal ) que
nous désirions receuillir comme conclusion pour les sections hyperplanes
sauf \inserted{pour $\zeta$ dans} un ensemble exceptionnel de codimension $\geqslant$ 2. Cela ne restreint
pas la généralité (et il n'aurait pas coûté plus cher à vrai dire
de se dispenser de cette hypothèse préalable), car on voit aussitôt
à l'aide des résultats de par. 3.4. \inserted{et 5.12.} que si X ne satisfait pas à
l'hypothèse en question, donc (en vertu de par.5) s'il existe un point
\uline{fermé} x en lequel cette hypothèse est en défaut, alors pour tout $\hand{\zeta}$ tel
que $\phi_{\hand{\zeta}}$ soit $X_{k(\hand{\zeta})}$-régulier (condition n'éliminant qu'un ens de codim$\geqslant$2)
et tel que $x \hand{\in} Y_{\hand{\zeta}}$ (condition qui décrit un ensemble de codimension 1
exactement), $Y_{\hand{\zeta}}$ ne satisfait pas non plus à ladite hypothèse en x ,
\hand{donc} l'ensemble exceptionnel $Z \hand{\subset} \check{P}$ est de codimension 1 et non 2.
(J'ai peut-être charrié un peu dans le cas $(R_k)$, \over{o}{où} il faut encore
des conditions $(S_1)$ et d'équidimensionnalité peut-être.....).

\mspage{49}

\clearpage
% Nouveau EGA V (cote 25): archive cote 25, site page 50 (PDF page 51).
% typescript. Draft compared: none.
% AI-transcribed by Claude Opus 5.5 (claude-opus-5-5, adaptive thinking, effort high), Message Batches API, batch ega25b (msgbatch_01A7n5YuXubUPNK1zsepi8RG), 2026-11-13; not yet checked by eye.
\scanpage{25/50}
\mspage{46}

\xout{\uline{Proposition} 8.16 On suppose f:X P non ramifié en}

Remords. Dans 8.13. et 8.14. , il suffit \inserted{de supposer} que f:X\hand{$\to$} P soit non
ramifié en les points lisses de X; pour le ``il suffit'', il suffit même
qu'il soit \inserted{non} ramifié sur un ouvert X' de X dont le complémentaire soit
de codimension \hand{$\geqslant$} \unsure{2}+1.

\uline{Proposition} 8.16. Supposons f:X\hand{$\to$} P non ramifié sur un ouvert
complémentaire d'un ensemble de codimension \hand{$\geqslant 2$,} X géom. normal et
de profondeur\hand{$\geqslant$} 3 en ses points fermés, enfin X géom. intègre et propre
sur k. Alors l'ensemble des \hand{$\zeta \in$} $\hand{$\check{\ }$}$P tels que \xout{$\phi$ soit $X_{k(\ )}$ régulier}
Y\hand{$_\zeta$} soit géom. normal et géom intègre de dimension dim X - 1 (est
constructible et) a un complémentaire de codimension \hand{$\geqslant$} 2.

En effet, par 8.14. b) il en est ainsi pour la propriété
`` Y\hand{$_\zeta$} est géom normal de dimension dim X - 1 '' (la propriété dimen
sionnelle exprime que \hand{$\phi_\zeta$} est X$_{k(\hand{$\zeta$})}$-régulier). Or en vertu de 6.1.
\uline{tous\xout{es}} les Y\hand{$_\zeta$} sont géom connexes. Or si Y\hand{$_\zeta$} est géom normal, il est
géom intègre sss il est géom connexe, ce qui achève la démonstration.

\uline{Remarques} 8.17. a) \inserted[typed]{L'hypothèse de 8.16. implique dim X \hand{$\geqslant$} 3.} Il est possible que, dès que X est géom \xout{intègre}irréd
et dim f(X) \hand{$\geqslant$} 3 (sans hypothèse de normalité ni de non ramification )
l'ensemble des \hand{$\zeta$} tels que Y\hand{$_\zeta$} soit géom irréd ait un complémentaire
de codimension \hand{$\geqslant$}2. On peut montrer en tous cas qu'il ne \unsure{contient}pas
d'hyperplan, \hand{cf. plus bas.}

b) La conclusion de 8.16. est en défaut si on abandonne
l'hypothèse prof$_x$X \hand{$\geqslant$} 3 pour x fermé, par exemple elle est en défaut pour
une quadrique non singulière \inserted{X} dans P$^3$ ; \deleted{dont} \inserted{\unsure{les sections par}} les hyperplans tangents
\xout{forment} sont \deleted{tous} réductibles (en fait, formés de couples de droites
concourrantes), et forment par ailleurs une famille de dimension 2
\margin{50} donc de codimension 1 dans $\hand{$\check{\ }$}$P (en fait, la duale de la quadrique \deleted{est} \deleted{\unsure{une}}

\clearpage
% Nouveau EGA V (cote 25): archive cote 25, site page 51 (PDF page 52).
% typescript. Draft compared: none.
% AI-transcribed by Claude Opus 5.5 (claude-opus-5-5, adaptive thinking, effort high), Message Batches API, batch ega25b (msgbatch_01A7n5YuXubUPNK1zsepi8RG), 2026-11-13; not yet checked by eye.
\scanpage{25/51}
\margin{\unsure{16}}
\margin{o--- c)} \inserted{Supposons cependant X lisse de dim 2, géom. intègre \unsure{propre} / k, et f une \uline{immersion}. Alors il résulte de 6.1. et 8.8. \inserted{et 8.14.} \unsure{que} si $Y^{\mathrm{sing}}_{\xi} \to \check{P}$ n'est pas gén. fini et inséparable, l'ens. des $\{\xi \in \check{P}$ tel que $Y_{\xi}$ soit \deleted{\unsure{géom.}} séparable sur $k(\xi)\}$ avec \uline{au plus 2} comp. irréd. géom. -- un compl. de codim $\geq 2$}
une quadrique de l'espace dual, relative à la forme duale ...).   Dans le
cas d'une surface non singulière \xout{dx} d'un espace projectif, cette situa-%
tion doit cependant être considérée comme exceptionnelle, cf nº suivant.

\margin{o---}
Nous allons examiner de plus près   maintenant le cas des surfaces
(le cas des courbes ne se pose évidemment pas du point de vue irréductibi-%
lité des sections hyperplanes \xout{\unsure{t}} ).

(NB  Je m'aperçois avec effroi que la quadrique n'a pas droit, dans
l'emploi que j'ai fait du mot ``\xout{f}réglé'', à être appelée réglée !  Cela doit
être en désaccord avec nos pères, et il faudra donc inventer un mot plus
adéquat pour la notion utilisée ici).

\uline{Proposition} 8.18.  Supposons \xout{X géom intègre} k alg clos, X \xout{intègre et}
\xout{et} intègre (resp. intègre et normale)\inserted{de dimension $d \geq 2$ et}
propre sur k, soit T une partie finie fermée de X telle que  X-T soit
lisse  et f $|$X-T  soit non ramifié.  Pour que l'ensemble des \hand{$\xi \in$} $\check{P}$  tels
que  $Y_{\hand{\xi}}$  soit \xout{\unsure{es}} géom irréd   (resp.  géom intègre)  de dimension  1
\unsure{fait}  un complémentaire de  codimension $\geq$ 2, il f et s que les deux
conditions suivantes soient satisfaites:

a)  Pour tout x\hand{$\in$}T, il existe une section hyperplane $Y_{\hand{\xi}}$   (\hand{$\xi \in$}P(k))
passant par x, de dimension \over{1}{$d-1$}, et qui soit irréductible.

b) \inserted{X'=} X-T est ``réglé'' (sic)  pour f, \uline{ou}  il existe  une section hyper
plane  $Y'_{\hand{\xi}}$  (\hand{$\xi \in$} $\check{P}$(k) )  de X',  \xout{singulière} qui soit de dimension \over{1}{d}\hand{-1},
singulière et  irréductible.

\deleted{Soit  en effet  Z  la  réunion  de  g($\overline{\underline{Y}'^{\hand{\mathrm{sing}}}}$), \xout{et} des hyperplans
de P  définis par les  f(x), x  T.  \xout{et}  Alors pour   P - Z,  Y   est
\uline{lisse} de dimension 1, \xout{\illegible{8}}
\xout{\illegible{10} au même \illegible{2} est lisse au dess}
\inserted{géom}
\xout{\unsure{des points de}}  et comme  il est connexe par 6.1. , il est géom irréd.}

\mspage{51}

\clearpage
% Nouveau EGA V (cote 25): archive cote 25, site page 52 (PDF page 53).
% typescript. Draft compared: none.
% AI-transcribed by Claude Opus 5.5 (claude-opus-5-5, adaptive thinking, effort high), Message Batches API, batch ega25b (msgbatch_01A7n5YuXubUPNK1zsepi8RG), 2026-11-13; not yet checked by eye.
\scanpage{25/52}
\mspage{47}

\xout{En vertu de 5.6, l'ensemble U des P tels que Y soit géom}

Supposons d'abord X géom normale. On sait déjà par 8.14. a)  que
l'en \xout{ex} peut trouver un fermé Z' de P de codi\xout{\unsure{mensi}}on $\geq$2, tel que
\hand{$\zeta\in$} P - Z'  implique que  $Y_{\hand{$\zeta$}}$  est  \uline{séparable} sur k(\hand{$\zeta$}) de  dim \inserted{d-}1; pour
un tel \hand{$\zeta$} , il revient donc au même que $Y_{\hand{$\zeta$}}$ soit géom irréd, ou géom
intègre, et  les deux problèmes  (respé et non respé) examinés dans
8.18. sont donc équivalents\unsure{.|!} D'autre part, en vertu de 5.6. l'ensemble
\inserted{U} des \hand{$\zeta\in$} $\hand{$\check{}$}$P  tels que  $Y_{\hand{$\zeta$}}$.  soit géom intègre de dim \hand{d-1} (l'hypothèse
dimensionnelle exprimant que $\phi_{\hand{$\zeta$}}$ est $X_{k(\hand{$\zeta$})}$-régulier ) est \uline{ouvert}.
Nous allons exhiber un ensemble ouvert \inserted[typed]{non vide,} évident  P - Z  cont\unsure{è|e}nu dans U,
en\xout{p} prenant pour  Z la réunion de  $\overline{g(\underline{Y}'^{\,\mathrm{sing}})}$  et des hyperplans $\check{H}_x$
de $\check{P}$  définis par les  f(x), \hand{$x\in$} T. \xout{Pour} Pour \hand{$\zeta\in$} $\check{P}$ - Z,  $Y_{\hand{$\zeta$}}$  est \uline{lisse}
de dimension \inserted{d-}1, et comme il est géom connexe par 6.1. , il est géom
intègre.  Il faut donc  exprimer qu\xout{\unsure{e chaque}} \inserted[typed]{toute} composante irréductible
de \inserted[typed]{co}dimension  1 de Z  \xout{\unsure{ne ren}}\inserted{\unsure{ren}}contre l'ouvert U.  Or ces composantes
irréductibles sont les  \hand{$\check{H}_x$}  (répetés éventuellement, peu importe), et
d'autre part $\overline{g(\underline{Y}'^{\,\mathrm{sing}})}$  lorsque  ce dernier  est bien de codim 1, i.e.
X' ``non réglée'' pour f, (NB on utilise l'irréductibilité de  $\underline{Y}'^{\,\mathrm{sing}}$ ).
D'ailleurs, pour que ce dernier ensemble rencontre  l'ouvert U,
il f et s  que  g($\underline{Y}'^{\,\mathrm{sing}}$)  (qui en contient un ouvert dense) rencontre
U, \xout{\unsure{i.e. qu'il existe x' X'(k) et une}} Cela prouve 8.18. dans ce
cas.  Lorsqu'on ne suppose pas X normal, on applique le résultat précé-%
dent au normalisé de X; le raisonnement est immédiat et je ne le détaille
pas ici.  \margin{\uline{NB} dans le cas \inserted{\unsure{bas}}respé, 8.18. est contenu dans 8.16. plus précis sauf pour d=2. C'est pour le cas non respé qu'il peut être utile de ne pas exiger d=2, mais seulement $d\geq 2$ ....}

Il reste à expliciter  les conditions a) et b) de 8.18. Cela con-%
duit à examiner de façon générale la situation suivante:  On suppose
X \deleted{intègre}\inserted{géom irréd,} sur k \deleted{alg clos}, on se donne une sous-variété linéaire  L \hand{de}

\margin{52}

\clearpage
% Nouveau EGA V (cote 25): archive cote 25, site page 53 (PDF page 54).
% typescript. Draft compared: none.
% AI-transcribed by Claude Opus 5.5 (claude-opus-5-5, adaptive thinking, effort high), Message Batches API, batch ega25b (msgbatch_01A7n5YuXubUPNK1zsepi8RG), 2026-11-13; not yet checked by eye.
\scanpage{25/53}
\margin{\unsure{46}}

$\overset{\hand{\vee}}{P}$ (correspondant à la question d'étudier les sections hypérplanes de X
passant par un point donné x, ou tangent à X en un point \deleted{simple}\inserted{lisse} donné),
form\over{e}{é} donc des hyperplans contenant une sous-variété linéaire $L^{o}$ de P
(resp. un point,  ou l'image d'un espace tangent de X en un point lisse,
dans les deux cas envisagés), et on se pose la question de savoir si
pour le point générique \hand{$\zeta$} de L  (donc pour tous les points d'un ouvert
non vide de L),  $Y_{\hand{\zeta}}$  est \xout{\unsure{géomxxintègrexxxsurxdeuxmoins}} géom irréd., de
\margin{/} dim  dim X -1.   C'est là une variante  de Bertini, qui devrait figurer
dans  le n°3  et se traite par exactement la même méthode\inserted[typed]{ou, si on veut, s'y ramène.}  La question
dimensionnelle s'exprime simplement par  $f(X)\not\subset L^{o}$ , \xout{\unsure{Ilxyxaxdonexxxxx}}
i.e.  $X' =$ \xout{Xxxx} $f^{-1}(P - L^{o})$   est un ouvert dense de X, \xout{\unsure{etxlaxquestion}}
Soit Q l'espace projectif des hypérplans passant par $L^{o}$  (NB si  $L^{o}$ est
défini par un sous-espace vectoriel $F^{o}$ de E,  on a   $Q = P($\xout{E/}$F^{o})$ ) et
considérons le morphisme  canonique  (déduit de \xout{ExxxE/$F^{o}$} $F^{o}\longrightarrow E$ , cf
Chap II)
\[ u :\ P - L^{o} \longrightarrow Q, \]
et considérons
\[ g = u f' :\quad f^{-1}(P - L^{o}) = X' \longrightarrow Q , \]
alors \xout{\unsure{laxfamillexdesxxxX'xxxxxtesxxautr}}  $L \simeq \overset{\hand{\vee}}{Q}$ ,  et la famille des
$X'_{\hand{\zeta}}$  $(\hand{\zeta}\hand{\in} L)$  n'est autre que la famille des ``sections hyperplanes'' rela-%
tif à ce morphisme  g . D'autre part, on voit aussitôt que pour tout
$\hand{\zeta}\in L$  ``général''  $X'_{\hand{\zeta}}$ est dense dans $X_{\hand{\zeta}}$ , donc  $X'_{\hand{\zeta}}$  est géom irréd sss
X  l'est.  Ceci posé,  le théorème de Bertini-Zariski nous montre que
nous avons la conclusion voulue \xout{dix} d'irréductibilité pourvu que
dim $g(X') \geqslant 2$. (A vrai dire, tant qu'à faire on pourrait signaler une
réciproque dans 3.1. , savoir : \xout{Yxxestxgé} si X est géom irréd, Y  est
géom irréd sss  ou bien  dim $f(X)$; $\neq$ \unsure{2|1}, ou bien  dim $f(X)=1$ , et $f(X)$
est contenu dans une droite\hand{D} (définie sur $\bar{k}$ )\inserted{et la fibre générique de $X\to D$ est géom. irréd}\hand{)}; cela permet également

\mspage{53}

\clearpage
% Nouveau EGA V (cote 25): archive cote 25, site page 54 (PDF page 55).
% typescript. Draft compared: none.
% AI-transcribed by Claude Opus 5.5 (claude-opus-5-5, adaptive thinking, effort high), Message Batches API, batch ega25b (msgbatch_01A7n5YuXubUPNK1zsepi8RG), 2026-11-13; not yet checked by eye.
\scanpage{25/54}
\mspage{49}

dans la version présente  avec L\unsure{.} d'avoir une condition néc et suffisan\xout{t}%
te \unsure{d}'irréductibilité géom. \hand{de $Y_\xi$, $\xi$ gén. dans $L$.}

Du point de vue cunutesque, \xout{xx} et en termes de théorie des
corps, il y a lieu d'exprimer la condition \xout{(suffisante)} en termes de
degré de transcendance de la façon suivante: on choisit un ``hyperplan
à l'infini'' ne contenant ni $L^{o}$ ni X, et on se place dans son \xout{v}complé%
mentaire i.e. sur le schéma affine  type essentiellement. On choisit
une base de l'\unsure{éspace|espace} des formes linéaires nulles sur $L^{o}$ , \xout{\illegible{1}}
\xout{\illegible{1} regarde} soit $T_{1},\ldots,T_{p}$ ( p = codim ($L^{o}$,P) ), et on considère
leurs images inverses $t_{1},..,t_{p}$ dans le corps des fractions K de X
(\inserted{X} supposé intègre). Un au moins des $t_{i}$, mettons $t_{1}$ , est $\neq 0$, considé-%
rons alors $a_{1} = t_{2}/t_{1}$ ,..., $a_{p-1} = t_{p}/t_{1}$ . Alors dim g(X') n'est
autre que le degré de transcendance de $k(a_{1},..,a_{p-1}) \subset K$ sur k. Donc
si ce degré de transcendance est $\geqslant 2$, O.K. \xout{De même s'il} S'il est 1,
alors il faut exiger que sur $\bar{k}$ , f(X) soit contenu dans une sous-varié%
té linéaire de P contenant $L^{o}$ et de dimension une de plus\hand{, et que la
fibre générique de $g : X' \to \overline{g(X')}$ soit géom. irréd.}

Supposons que $L^{o}$ soit de dimension q, alors les fibres de u:
$P - L^{o} \to Q$ sont de dimmension q+1, \inserted{donc celles de $g$ sont de dim $\leq q+1$,} et par suite on a
\[ \dim g(X') \geqslant \dim f(X) - (q+1) , \]
donc la condition dimensionnelle sur g(X') est vérifiée pourvu que
\[ \dim f(X) \geqslant q+3 \quad . \]
Lorsque q=0, on trouve le fait signalé dans 8.17. a)\unsure{,} \inserted{\unsure{du}} \xout{Lorsque q=0,}

Revenant aux conditions de 8.18. on voit que la condition a)
relative à un $x \in T$ est satisfaite pourvu que X ne soit pas ``conique''
en x relativement à \unsure{$f$}'', dans un sens évident.

\margin{|} A voir si ces derniers développements Bertiniesques ne viendrai-%
ent pas mieux au n\textsuperscript{o} suivant du changement d'immersion projective. \margin{54}

\clearpage
% Nouveau EGA V (cote 25): archive cote 25, site page 55 (PDF page 56).
% typescript. Draft compared: none.
% AI-transcribed by Claude Opus 5.5 (claude-opus-5-5, adaptive thinking, effort high), Message Batches API, batch ega25b (msgbatch_01A7n5YuXubUPNK1zsepi8RG), 2026-11-13; not yet checked by eye.
\scanpage{25/55}
\mspage{50}

9. \uline{Changement d'immersion projective}.

9.1. Pour tout entier $n > 0$, soit $P(n) = P(\underline{Sym}^n(E))$, on a une
immersion évidente $u_n\colon P \to P(n)$, puisque $\underline{O}_P(n)$ est engendré par
ses sections au dessus de tout ouvert affine de S et que $p_*(\underline{O}_P(n))$
$\underline{Sym}^n(\underline{E})$, où $p\colon P \to S$ est la projection. S\xout{o}i\xout{t} $f\colon X\to P$ est un morphisme
non ramifié (resp. une immersion), il en est de même de $u_n f\colon X \to P(n)$.
Il y a parfois avantage, dans l'étude de X, de \xout{au} remplacer f par $u_n f$ \xout{\unsure{f}}
pour éviter un comportement trop particulier et parfois \xout{général} gênant .
de f à certains égards. Un exemple d'une telle particularité est celle
\xout{ex} signalée (sic) dans 8.12. b), \xout{\unsure{lexmorp}} où $\underline{Y}^{\mathrm{sing}} \longrightarrow \overset{\hand{$\vee$}}{X}$ a une image de
dimension r-1 mais \xout{estxunxmorphismexinsépa} donne lieu à une extension
de corps inséparable. Un autre cas est celui qui se présente dans le cas
d'une quadrique \deleted{de dimension impaire \xout{dans} 2m lorsque r=2m+1,}\inserted{dans $\mathbf{P}^3$} savoir que
toutes les sections hyperplanes singulières sont géométriquement
réductibles (alors que X est cependant géométriquement irréductible).

\uline{Proposition} 9.2. On suppose S=Spec(k), X lisse sur k, et $f\colon X\to P$ non
ramifié. Soit $n \geq 2$, et considérons $f_n = u_n f$ . Alors $f_n : X \to P(n)$
satisfait aux condition équivalentes de 8.8. , en particulier
pour \hand{$\xi \in$} $\overset{\hand{$\vee$}}{P}(n)$ dans le complémentaire d'un ensemble de codimension $\geq 2$,
la section hyperp\xout{b}lane correspondante $Y_{\hand{$\xi$}}$ est lisse \xout{sur} ou \inserted{n'}admet qu'un
\inserted[typed]{nombre fini de points non lisses ,}
\xout{plusxunxseulxpointxnonxlissex}, lesquels sont géométriquement singuliers
ordinaires\xout{etxrationnelsxsurxk(xx)x}. Si f est une immersion, il y a au
plus un tel point singulier, et il est rationnel sur k(\hand{$\xi$}).
\margin{$\|$} NB Il fallait énoncer 8.8. de façon à \xout{\unsure{éviter}} ne pas exclure le
cas où f ne serait pas une immersion. La \unsure{vérifiaation|vérification} est essentielle-%
ment triviale sur la condition (ii bis) de 8.8. -Il faudrait sans doute
expliquer dans 9.1. que les sections hyperplanes de X relativement à
\margin{55} $f_n$ ne sont autres que les``sections''de X par les hypersurfaces de degré

\clearpage
% Nouveau EGA V (cote 25): archive cote 25, site page 56 (PDF page 57).
% typescript. Draft compared: none.
% AI-transcribed by Claude Opus 5.5 (claude-opus-5-5, adaptive thinking, effort high), Message Batches API, batch ega25b (msgbatch_01A7n5YuXubUPNK1zsepi8RG), 2026-11-13; not yet checked by eye.
\scanpage{25/56}
\mspage{51}

n  au lieu de \xout{\illegible{1}} hyperplans.

\uline{Proposition} 9.3. \xout{xx} Supposons \inserted{X géom irréd et} dim f(X) $\geq$ 2, \xout{xx} soit x$\in$X(k), \inserted{Soit $n\geq 2$} \xout{xx}
\xout{xxxx} a) \deleted{Alors pour n$\geq$2}, \over{c}{C}onsidér\over{ant}{ons} la famille \inserted{(linéaire)} des sections hyperplanes
de X relativement à $f_n=u_nf$ \uline{qui passent par} x, son élément
générique \xout{xxx} définit une $Y^{(n)}_{\hand{\xi}}$ qui est \uline{géométriquement} \xout{intègre}irréd.

b) \xout{Si f(X) n'est pas contenu dans un plan défini sur $\bar{k}$, et n$\geq$2,}
\xout{\unsure{ou si n$\geq$3,}} Soit L une sous-variété linéaire de P passant par f(x) et
ne contenant pas f(X). \deleted{Alors pour n$\geq$2}, \over{c}{C}onsidér\over{ant}{ons} la famille \inserted{(linéaire)} des
sections hyperplanes de\xout{s} X relativement à \xout{\illegible{1}} $f_n$ \xout{\illegible{1}} , définis
par des hyperplans de P(n) \xout{\illegible{1}} \uline{``tangents à L en x''} (i.e. définis
par des \inserted{n-}formes sur P qui \inserted{sur L} s'annulent à l'ordre 2 au moins en x), son
membre générique \deleted{définit} \inserted{est} une $Y^{(n)}_{\hand{\xi}}$ géométriquement irréductible.

c) Supposons X lisse en x, \xout{etxfxxxx} et \deleted{soit} n$\geq$3, \inserted{ou f(X) pas contenu dans un plan défini sur $\bar{k}$.} \xout{alors} \over{c}{C}onsidérons
la famille des sections hyperplanes \inserted{\deleted{\unsure{tangentielles}}} \xout{génériques} $Y^{(n)}_{\hand{\xi}}$ de X relativement
à $f_n$ qui sont \hand{``}\uline{tangentes à X en x}\hand{''}\unsure{\over{.}{,}} Alors le membre générique de ladite
définit une $Y^{(n)}_{\hand{\xi}}$ géométriquement irréductible.

La démonstration est essentiellement triviale en termes des
critères de la fin du n\textsuperscript{o} précédent. Prenant un modèle affine de P
contenant f(x), on est ramené dans a) à trouver \xout{xxx} trois polynômes
en les coordonnées $T_1,..,T_r$ , de degré \over{\illegible{1}}{$\leq$}2, soient P,Q,R, tels que
Q(t)/P(t) et R(t)/P(t) soient alg indép sur k dans K (où K est le
corps des fonctions de X, et t = $(t_1,..,t_r)$ est le système d'élé-%
ments de K défini par les $T_i$ ); dans b), on exige de plus que P,Q,R
s'annulent à l'ordre 2 au moins sur L, qu'on peut d'ailleurs supposer
défini par des équations $T_1,..T_s = 0$ ; enfin dans \unsure{c)} c'est le même
pb, mais L étant l'image de l'espace tangent à X en x, et on se permet
\inserted{éventuellement} de prendre P,Q,R \xout{\illegible{1}} de degré \hand{3} , i.e. un peu plus de marge.

\mspage{56}

\clearpage
% Nouveau EGA V (cote 25): archive cote 25, site page 57 (PDF page 58).
% typescript. Draft compared: none.
% AI-transcribed by Claude Opus 5.5 (claude-opus-5-5, adaptive thinking, effort high), Message Batches API, batch ega25b (msgbatch_01A7n5YuXubUPNK1zsepi8RG), 2026-11-13; not yet checked by eye.
\scanpage{25/57}
\mspage{\unsure{92}}

L'hypothèse  dim $f(X)\geqslant 2$  signifie que  le degré de transcendance de
$k(t_1,..,t_r)$ sur $k$  est $\geqslant 2$, i.e. on peut trouver  $t_1,t_2$ disons alg ind.
Dans a) on prendra alors   $P = T_1$ , $Q = T_1^{\,2}$ , $R = T_1T_2$ , Dans b), on
fait pareil, \xout{\unsure{mais en ayant soin de prendre} \unsure{$t_1$}} en notant qu'on peut ci-%
dessus choisir $t_1$ provenant bien de $T_1$ \uline{nul sur} L, grâce au fait $f(X)\not\subset L$
(qui implique qu'il  existe un indice  i  entre 1 et s tel que $t_s$ non
nul, donc $t_s$ non constante (car $t_s$ s'annule en x) donc $t_s$  non alg sur k
( NB on peut supposer k alg clos ).  Le cas c) résulte de b), \deleted{avec n=2
même au lieu de n=3,} sauf dans le cas où \xout{X} f(X) est contenu dans
l'image \inserted{L} par f de l'espace tangent à X en x, \deleted{ce qui  exige d'ailleurs pou
des raisons de dim que l'application tangente en x est injective i.e.
f non ramifié en x, \xout{Lorsque dim X} de sorte que f(X) est contenu dans
une sous-variété linéaire de P de même dimension que X.} \hand{[}Lorsque dim X=2
ce cas est effectivement exceptionnel (la trace d'une quadrique tangente
à un plan sur ce plan est en général formé de deux droites concourran-%
tes, donc n'est pas irréductible)\hand{]}. Mais pour traiter ce cas, \xout{sous} la
forme P,Q,R explicitée plus haut, on peut remplacer \inserted{si dim L $=2$} évidamment X par
L lui-même, où la solution est triviale ( prendre $P=T_r^{\,2}$\xout{\illegible{1}}, $Q=T_r^{\,3}$ ,
$R = T_r^{\,2}T_{r-1}$ , en notant que  $Q/P = T_r$ et  $R/P = T_{r-1}$ sont des formes
linéaires lin ind sur L, donc alg ind\inserted{si dim $L\geqslant 3$, alors $T_{r-2}$, $T_{r-1}$, $T_r$ lin. indép. sur L, et on prendra $P=T_r^{\,2}$, $Q=\unsure{T_{r-1}^{\,2}}$, $R=\unsure{T_{r-2}^{\,2}}$}). \deleted{On pourrait signaler en
remarque que dans c), on peut même prendre n=2  sauf si  f(X) est conte-%
nu dans un plan .}

Conjuguant avec \unsure{8.181|8.18} on trouve un

\uline{Corollaire} 9.4.   \xout{(Pour X propre sur k}

Enfin, il faut encore conjuguer  ce dernier avec 9.2. pour
trouver un théorème récapitulatif  du ``cas excellent'':
\deleted{\uline{Théorème}} 9.5.   (Lorsque X est lisse \deleted{et} propre\inserted{(et géom intègre)} sur k, et $f:X\to P$ non
ramifié,\hand{\unsure{X}}de dimension $\geqslant 2$)\hand{, en signalant le résultat sur \unsure{$\pi_1$} ?.}
\hand{quand $X\to P$ non immersion (variété des pts singuliers de \unsure{$f$}).}

\margin{57}

\clearpage
% Nouveau EGA V (cote 25): archive cote 25, site page 58 (PDF page 59).
% typescript. Draft compared: none.
% AI-transcribed by Claude Opus 5.5 (claude-opus-5-5, adaptive thinking, effort high), Message Batches API, batch ega25b (msgbatch_01A7n5YuXubUPNK1zsepi8RG), 2026-11-13; not yet checked by eye.
\scanpage{25/58}
\mspage{53}

10. \uline{Pinceaux de sections hyperplanes}, \uline{et fibrations de variétés éclatées}

10. 1. Soit Z l'ensemble P-exceptionnel dans $\overset{\hand{$\vee$}}{P}$ relativement à une
certaine propriété constructible\inserted{P}, de sorte que Z est une partie construc%
tible de $\overset{\hand{$\vee$}}{P}$ . \hand{(}\inserted[typed]{Supposons S=Spec(k),}\hand{)} On sait (cf nº\hand{\unsure{11}}\deleted{suivant}, où nous \unsure{faisans|faisons} un ratrappage
% NOTE: a long diagonal pencil stroke in the left margin beside this and the next two lines
\inserted{de choses} qui aurai\over{t}{ent} sans doute dû venir dans les nºs antérieurs) que pour que
l'on ait codim\xout{\illegible{1}}(Z,$\overset{\hand{$\vee$}}{P}$ ) $\geqslant$ 2, il f et s que \hand{``}\xout{t}toute droite assez
générale" L dans $\overset{\hand{$\vee$}}{P}$ ne rencontre pas Z, \xout{xxxxx} ou encore $\overline{Z}$ , - et il suffit
qu'il existe \uline{une} droite L dans $\overset{\hand{$\vee$}}{P}$ ne rencontrant pas k. Lorsque \xout{xxxxxxxx}
k est infini, il f donc et s qu'il existe une doite L dans $\overset{\hand{$\vee$}}{P}$ ne rencon-%
trant pas $\overline{Z}$ . On appelle "pinceau linéaire\deleted{"}\hand{(}\inserted[typed]{de sections hyperplanes de X\hand{''}}\hand{)} défini par la\xout{d}droite
L de $\overset{\hand{$\vee$}}{P}$\hand{,} le L-pré\unsure{s}chéma $\underline{Y}_L$ (définition valable pour un S quel\xout{f}conque).
Ainsi, les réflexions précédentes, jointes aux résultats des nºs 8 et 9,
nous donnent des critères d'existence de tels pinceaux, ayant des fibres
$\underline{Y}_{\hand{$\xi$}}$ ( \hand{$\xi \in$}L) satisfaisant toutes la propriété P, tout d'abord dans le
cas \deleted{deux} \hand{(}\inserted[typed]{où S est un}\hand{)} corps de base infini. \xout{Maisxprécédantxxxxxx} Compte tenu de
8.2.\hand{,} lorsque \xout{X} pour tout cycle premier associé à X, on a dim f(\unsure{F})$>$0,
on peut (en prenant la propriété P' = P + condition de régularité
pour $\emptyset_{\hand{$\xi$}}$ ) exiger que le pinceau $\underline{Y}_L$ soit plat sur L. Dans le cas
où S est quelconque, on peut encore\xout{x} souvent, par le procédé de 7.1. ,
construire un tel pinceau au dessus d'un voisinage ouvert d'un point
donné s de S, pourvu que k(s) soit infini et qu'on sache que Z est fermé
(ce qui sera assuré dans divers cas par les résultats du par.5 et
l'hypothèse de propreté pour X $\rightarrow$ S).

\margin{|} Pour bien faire, il conviendrait, après des explications géné-%
rales de cette eau, de donner de\xout{x}s énoncés récapitulatifs où on
applique effectivement les résultats précédents pour un certain nombre
de propriétés de cette nature (y compris bien sur des propriétés modulées\hand{)}

\mspage{58}

\clearpage
% Nouveau EGA V (cote 25): archive cote 25, site page 59 (PDF page 60).
% typescript. Draft compared: none.
% AI-transcribed by Claude Opus 5.5 (claude-opus-5-5, adaptive thinking, effort high), Message Batches API, batch ega25b (msgbatch_01A7n5YuXubUPNK1zsepi8RG), 2026-11-13; not yet checked by eye.
\scanpage{25/59}
\mspage{59}

Comme minimum dans ce sens, il faut  donner ici la reformulation de 9.5.
en termes de pinceaux linéaires  -  fait constamment utilisé dans les
applications géométriques.
\hand{1}0.2.  Par polarité, une droite $\underline{L}$  de $\check{P}$  correspond à une sous-variété
linéaire $L^{o}$ de codimension 2 dans P  (S quélconque).  Posons
\[ T = X\mathrm{x}_{P}L^{o} \]
Une autre façon  de décrire T est la suivante:  L  est défini par un
quotient loc libre de rang 2 de $\check{\underline{E}}$, ou ce qui revient au même, par un
sous-Module localement facteur direct  $\underline{F}$ de $\underline{E}$\xout{x}, partout de rang 2,
Considérons  l'homomorphisme composé
\[ \underline{F}_{X} \longrightarrow \underline{E}_{X} \longrightarrow \underline{O}_{X}(1) \quad , \]
alors  T n'est autre que  le schéma des zéros de cet homomorphisme compo-%
sé, où ce qui revient au même, est défini par l'Idéal \inserted[typed]{J} image de l'homomo%
phisme correspondant (obtenu en tordant par $\underline{O}_{X}(-1)$):
\[ \underline{F}_{X}(-1) \longrightarrow \underline{O}_{X} \quad . \]
\inserted[typed]{Supposons}
\xout{Supposonsxx Pour} que cet homomorphisme soit \uline{régulier}, ce qui signifie
que  si on écrit localement une base \inserted{tot. ordonnée} de $\underline{F}_{X}(-1)$, \deleted{leurs} \inserted{les} images dans $\underline{O}_{X}$\hand{\unsure{\^{}}}
forme\over{\unsure{nt}}{nt} une suite $\underline{O}_{X}$-régulière, \xout{x}condition qui ne dépend pas de la base
choisie et peut s'énoncer intrinsèquement aussi en disant que
$\underline{F}_{X}(-1)$\inserted{$\otimes\ \underline{O}_{X}/J$}$\to \underline{J}/\underline{J}^{2}$  est un isomorphisme, et  $V(\underline{J})=T \longrightarrow X$  est  une immersion
\margin{|} régulière  (NB il \xout{xxxxxx} faudrait quelque part dégager la situation
\margin{|} générale avec un hom  $\underline{G} \to \underline{O}_{X}$ , $\underline{G}$ localement libre sur préschéma X, par
\margin{|} exemple dans le par. des immersions régulières).  Alors on a la
\xout{Proposition} \uline{Théorème}  10.2.   Sous l'hypothèse précédente, le pinceau
linéaire $\underline{Y}_{L}$  muni de la projection canonique $\underline{Y}_{L} \longrightarrow X$  est  X-isomorphe
\inserted{de façon unique} au préschéma éclaté  de X au moyen de T.

\mspage{59}

\clearpage
% Nouveau EGA V (cote 25): archive cote 25, site page 60 (PDF page 61).
% typescript. Draft compared: none.
% AI-transcribed by Claude Opus 5.5 (claude-opus-5-5, adaptive thinking, effort high), Message Batches API, batch ega25b (msgbatch_01A7n5YuXubUPNK1zsepi8RG), 2026-11-13; not yet checked by eye.
\scanpage{25/60}
\deleted{\illegible{2}} \mspage{55}

Pour comprendre la signification de ce théorème, il conviendrait
\inserted{\deleted{\illegible{5}}}
de noter dès le début du nº \hand{(}que  si S=Spec(k)\hand{)}, alors pour une droite
L ``assez générale'' dans $\hand{\check{\ }}P$ ,  la condition de régularité est vérifiée \inserted{(\unsure{notamment pour S.3.|notamment pour 5.3.})}
(cf ra\xout{t}trappages signalés au  nº\inserted{\unsure{l'avant}}\deleted{suivant}).  Par suite, dans les construc%
tions de ``bons'' pinceaux linéaires  indiquées au début du nº, on pouvait
imposer que les pinceaux envisagés \xout{\unsure{étaient}} satisfaisaient ladite conditi-
on  (ce qui \xout{ne} est une condition\inserted[typed]{de même type mais\hand{,}} distincte de celle qui consiste à
exiger que pour tout \hand{$\zeta\in$}L,  $\phi_{\hand{\zeta}}$  soit $X_{k(\hand{\zeta})}$-régulier ). Il fallait inclu
re la condition en question dans les énoncés récapitulatifs préconisés.
\xout{D'autre part, on peut noter}  D'autre part, pratiquement  10.2. est
utilisé surtout dans la situation 9.5. ,  ce qui  rend  désirable de
n'énoncer la réformulation de  9.5.  en termes de pinceaux qu'après
10.2. , pour pouvoir  inclure dans l'énoncé en question  également l'iso
morphisme du pinceau avec un éclaté (i.e. donner une descr\xout{u}iption de la
situation  permettant une référence commode).  On obtient ainsi un
moyen, pour tout schéma projectif lisse \inserted[typed]{\hand{(}géom \xout{\illegible{1}} conn  X  de di\xout{p}m$\geq$2 \hand{)}}sur un corps infini k,  de trou\xout{x}
ver un\xout{e} sous-schéma fermé lisse\xout{ partout} \inserted[typed]{non vide} de codimension \deleted{\illegible{1}}2 en chacun de
\xout{puisse être}
ses points \xout{(ce qui n'exclut pas qu'il soit vide) et il implique au}
\xout{\unsure{contraire si dim X}} tel que le schéma éclaté admette une fibration
sur P$^1$ , dont toutes les\hand{/}fibres \inserted[typed]{\hand{(}sont géom intègres, et  dont toutes les fibres \hand{)}}sont lisses sont au plus un nombre fini,
ces dernières ayant au plus un point \xout{non lisse (et} géom. singulier, et
un tel point \xout{est} \inserted[typed]{étant} rationnel sur k et géom. singulier ordinaire .  Cela
explique l'importance  de l'étude approfondie (\deleted{à peine} \inserted{tout juste} abordée à l'heure
actuelle) de telles fibrations à fibres singulières, pour réduire dans \xout{x}
une certaine mesure l'étude des variétés projectives lisses de dim d
à celle des(familles à un paramètre de) variétés projectives de dim d-1\hand{,}
\hand{pouvant avoir des singularités ordinaires.}

\margin{60}

\clearpage
% Nouveau EGA V (cote 25): archive cote 25, site page 61 (PDF page 62).
% typescript. Draft compared: none.
% AI-transcribed by Claude Opus 5.5 (claude-opus-5-5, adaptive thinking, effort high), Message Batches API, batch ega25b (msgbatch_01A7n5YuXubUPNK1zsepi8RG), 2026-11-13; not yet checked by eye.
\scanpage{25/61}
\mspage{56}

L'énoncé \unsure{10.2.} est une conséquence à peu près immédiate du suivant,
qui est tout à fait indépendant d'histoires de sections hyperplanes, et
serait sans doute mieux à sa place dans \xout{xx} un nº en rabiot aux
``immersions régulières''.

\deleted{\uline{Proposition} \unsure{10.3.} Soient X un préschéma, \xout{$\underline{E}$} $\underline{G}$ un Module quasi-cohérent}
\deleted{sur X, et $u\colon \underline{G}\to\underline{O}_X$ un homomorphisme, $\underline{J} = u(G)$, $T=V(\underline{J})$. Soit}
\deleted{$\widetilde{X}$ déduit de X par éclatement de T. Considérons d'autre part}
\deleted{$P = P(\underline{G})$, l'homomorphisme canonique $\underline{G}_P\to\underline{O}_P(1)$ et son noyau}
\deleted{$\underline{H}$ ,(de sorte que nous avons la suite exacte $0\to\underline{H}-\underline{G}_P\to\underline{O}_P(1)\to 0$),}
\deleted{l'homomorphisme $u_P : \underline{G}_P \to \underline{O}_P$ , et l'Idéal qu coh $\underline{K} = u_P(\underline{H})\subset\underline{O}_P$ .}
\deleted{Alors $\widetilde{X}$ est canoniquement isomorphe à un sous-préschéma fermé de}
\deleted{$V(\underline{K})$ . Lorsque G est localement libre et u est ``\xout{re}régulier'', alors}
\deleted{l'isomorphisme précédent est un isomorphisme de $\widetilde{X}$ avec $V(\underline{K})$ lui-même,}
\deleted{de plus dans ce cas $\underline{H}$ est localement libre sur P et \xout{xxx} $\underline{H}\to\underline{O}_P$}
\deleted{(dont le\xout{\unsure{t}} préschéma d es zéros e st $\widetilde{X}$ ) est également régulier.}

\deleted{Le premier énoncé est à peu près trivial. Le deuxième est un}
\deleted{exercice qui ne peut offrir de difficulté (je ne l'ai pas fait en}
\deleted{détail, pensant que tu t'en tireras aussi bien que moi).}

Lorsque dans 10.2. \deleted{les fibres de X/S sont de} \inserted{$S = \mathrm{Spec}\,k$ et X est de} dimension $\unsure{\leqslant}1$,
alors l'hypothèse de \xout{\unsure{ex}}régularité faite équivaut à $T = \emptyset$ , de sorte
que $Y_{\underline{L}}\to X$ est un \uline{isomorphisme} . On trouve ainsi\unsure{;} en conjuguant
avec 9.2. :

\uline{Corollaire} \unsure{10.4.} (à \unsure{10.2.}) Soit X une courbe lisse géom connexe dans
\xout{dimensionxxx} l'espace projectif P sur un corps infini k, et soit $n\hand{\geqslant}2$.
Alors il existe un pinceau linéaire de n-formes sur P, définissant un
morphisme $X\to P^1$ ayant les propriétés suivantes: ce morphisme est

\mspage{61}

\clearpage
% Nouveau EGA V (cote 25): archive cote 25, site page 62 (PDF page 63).
% typescript. Draft compared: none.
% AI-transcribed by Claude Opus 5.5 (claude-opus-5-5, adaptive thinking, effort high), Message Batches API, batch ega25b (msgbatch_01A7n5YuXubUPNK1zsepi8RG), 2026-11-13; not yet checked by eye.
\scanpage{25/62}
\mspage{57}

génériquement étale\inserted[typed]{de degré \hand{d}}, et pour tout point géométrique  s de $P^1$ ,  $X_s$
est étale sur le corps \hand{$-$}\unsure{alg} clos \inserted{k'=} k(s), ou  \xout{\illegible{3}} est
\inserted{k'-}isomorphe à la somme de \over{\unsure{n}}{d}-2  schémas \inserted{Spec(k')} \xout{k(s)}  et \xout{un}  du schéma  \xout{k}
$I_{k'} = \mathrm{Spec}\, k'[t]/(t^2)$ .

En langage de Papa, au  dessus de s  il y a au plus un point de
ramification, et la ramification en c e point est ``quadratique'' .

11.  \uline{Grassmaniennes}. \deleted{Ratrapages.}

Comme nous avions couramment à utiliser des sous-variétés linéai-%
res de P , pas seulement de dimension relative 0 ou r-1, il est clair
qu'on aura besoin de quelques notions sur les grassmaniennes  et\inserted{de}quelques
sorties de nature ``géométrie élémentaire'' sur les constructions concer-%
nant les variétés linéaires, qui  devraient venir tout au début du par.
De plus, on aura en pratique à prendre parfois des sections linéaires
quelcohques, pas seulement des sections hyperplanes, et il y a lieu de
revoir dans cet esprit élargi tous les n°s antérieurs.

Soit  $\underline{E}$  un Module qu coh  sur le préschéma $\unsure{S|X}$\inserted[typed]{et n un entier \hand{$>$} 0.}\hand{,} Considérons le
foncteur \qquad $(Sch)_{/S}^{\,o} \longrightarrow (Ens)$  défini par
\[ \underline{Grass}_n(\underline{E})(S') = \text{ ensemble des Modules quotients loc libres} \]
\[ \text{de rang  n  de } \underline{E}_{S'} \ . \]
Ce foncteur est représentable, et le préschéma sur S qui le réprésente
sera donc noté également $\underline{Grass}_n(\underline{E})$.  Pour prouver la r eprésentabilité,
considérer l'homomorphisme de foncteurs naturels
\[ \underline{Grass}_n(\underline{E}) \longrightarrow \underline{Grass}_1( \overset{\hand{n}}{\textstyle\bigwedge} \underline{E} ) = \underline{P} ( \overset{\hand{n}}{\textstyle\bigwedge} \underline{E}) \ , \]
défini en associant à  tout  quotient loc libre  de rang n  $\underline{G}$  de $\underline{E}_{S'}$
le Module  loc libre de rang 1 \quad $\overset{n}{\textstyle\bigwedge} \underline{G}$ , considéré comme un quotient
de  $\overset{n}{\textstyle\bigwedge} \underline{E}_{S'}$ \hand{\unsure{?}}  On prouve comme dans  Séminaire Cartan que  ce morphisme
est  ``représentable par \xout{\illegible{1}} immersion fermées'' , de sorte que $\underline{Grass}_n(\underline{E})$

\mspage{62}

\clearpage
% Nouveau EGA V (cote 25): archive cote 25, site page 63 (PDF page 64).
% typescript. Draft compared: none.
% AI-transcribed by Claude Opus 5.5 (claude-opus-5-5, adaptive thinking, effort high), Message Batches API, batch ega25b (msgbatch_01A7n5YuXubUPNK1zsepi8RG), 2026-11-13; not yet checked by eye.
\scanpage{25/63}
\mspage{58}

apparait comme un sous-préschéma fermé de $P(\overset{\hand{\unsure{n}}}{\bigwedge}\underline{E})$; en particulier, il
est séparé sur S et quasi-compact sur S, et  si $\underline{E}$  est de type fini,
il est projectif sur S. Si $\underline{E}$ est de \unsure{présentation|présantation} finie, alors
il en est de même de $\underline{\text{Grass}}_n(\underline{E})$ : en effet, \xout{\unsure{localement sur S,}} on peut
supposer S affine, S=Spec(A), alors $\underline{E}$ provient d'un module de type fini
sur un sous-anneau de type fini de A, -  or la formation de $\underline{\text{Grass}}_n(\underline{E})$
est évidemment compatible avec la changement de base sur S.

Lorsque $\underline{E}$  est localement libre,\xout{de type fini,} alors  $\underline{\text{Grass}}_n(\underline{E})$
est \uline{lisse} sur S\inserted{\unsure{(à fibres géom. connexes}}. Cela provient du fait plus précis: si $\underline{E}$ est libre,
de rang r, alors  $\underline{\text{Grass}}_n(\underline{E})$  peut être recouvert par  $\binom{r}{n}$  ouverts
\inserted{S-}isomorphes chacun  à  l'espace affine type de dimension relative
n(r-n)  sur S.  Cette décomposition correspond au choix grâce à la base
de $\underline{E}$  de  $\binom{r}{n}$  décompositions de $\underline{E}$  en \xout{sommes directes} suites
exactes
\[
\text{(s)} \qquad 0 \rightarrow \underline{E}' \relbar \underline{E} \relbar \underline{E}'' \longrightarrow 0 \quad ,
\]
avec $\underline{E}'$  loc libre de rang n. \xout{\illegible{1}} Une telle suite exacte permet
de définir un sous-foncteur  $\underline{\text{Grass}}_n(\text{s})$  de  $\underline{\text{Grass}}_n(\underline{E})$ , en se bornant
aux  quotients  $\underline{G}$  loc libres de rang  n de  $\underline{E}_{S'}$  tels que l'homomorphi-%
sme composé  $\underline{E}'_{S'} \longrightarrow \underline{E}_{S'} \longrightarrow \underline{G}$  soit surjectif (donc bijectif). Or
l'inclusion  $\underline{\text{Grass}}_n(\text{s}) \rightarrow \underline{\text{Grass}}_n(\underline{E})$  est représentable \uline{par immersions
ouvertes}, et d'autre\xout{s} part  $\underline{\text{Grass}}_n(\text{s})$  est canoniquement isomorphe
au fibré ve\unsure{c}toriel  \xout{\unsure{V(E''⊗E')}}  V ( $\underline{\text{Hom}}_{\underline{O}_S}(\underline{E}',\underline{E}'')$) .

Comme conséquence par exemple de cette structure particulière,
on peut signaler que si s $\in$ S, alors ($\underline{E}$ étant loc libre de rang fini)
tout point de $\underline{\text{Grass}}_n(\underline{E})$  à valeurs dans k(s)  se relève en une section
au dessus d'un voisinage de s.  D'autre part, si S=Spec(k), \xout{\illegible{1}} k
un corps fnfini, alors tout ouvert non vide de $\underline{\text{Grassn}}(\underline{E})$  contient un

\mspage{63}

\clearpage
% Nouveau EGA V (cote 25): archive cote 25, site page 64 (PDF page 65).
% typescript. Draft compared: none.
% AI-transcribed by Claude Opus 5.5 (claude-opus-5-5, adaptive thinking, effort high), Message Batches API, batch ega25b (msgbatch_01A7n5YuXubUPNK1zsepi8RG), 2026-11-13; not yet checked by eye.
\scanpage{25/64}
\mspage{\unsure{57}}

point rationnel sur k.

Un point de $\underline{\mathrm{Grass}}_n(\underline{E})$ à valeurs dans S, i.e. un Module quotient
loc libre de rang n \ $\underline{G}$ de $\underline{E}$, définit canoniquement un sous-préschéma
de $P(\underline{E})$ , savoir $P(\underline{G})$ . Un tel sous-préschéma\inserted{(mais sans imposer le rang de $\underline{G}$)} est appelé \uline{sous-variété}
% NOTE: the last letter of "sous-variété" is cut by the right edge of the scan
\uline{linéaire} de $P(\underline{E})$ (relativement à S si une confusion e st possible).
C'est donc un fibré projectif,de dimension relative n-1 si $n\geqslant 1$ ,
(et vide si n=0). On vérifie immédiatement que la section de $\underline{\mathrm{Grass}}_n$
$(\underline{E})$ i.e. G \ est connu quand on connait la sous-variété linéaire
correspondante de $P(\underline{E})$. D\unsure{e} cette façon, \ la grassmanienne peut
s'interpréter comme représentant le foncteur "sous-variétés linéaires
de dimension relative n-1 \xout{Xx=xresp xxxvide=xxx} de $P_{S'}$ " pour S'
variable\xout{x}\inserted[typed]{$(n\geqslant 1)$}. Il est d'ailleurs possible de donner une caractérisation
intrinsèque de ce dernier foncteur i.e. de la notion de sous-variété
linéaire \ de dimension relative m : ce sont les sous-préschémas fermés
de P , lisses sur S et partout de dimension relative m, et qui sont
de "degré projectif 1" \ en tout $s\in S$; \ cette caractérisation sera
donnée dans un Chapitre ultérieur, et nous n'en avons nul besoin ici.

Supposant de nouveau $\underline{E}$ \ loc libre\inserted[typed]{de rang r}, soit $\underline{E}$ \ son dual. Alors
par polarité nous trouvons un isomorphisme canonique
\[ \underline{\mathrm{Grass}}_n(E) \simeq \underline{\mathrm{Grass}}_{r-n}(\check{\underline{E}}) \]
faisant correspondre \ au \ quotient \ $\underline{G}$ de $\underline{E}$ \ le quotient \ $\check{\underline{E}}/\check{\underline{G}}$ \ de $\check{\underline{E}}$ .
Du point de vue variétés linéaires, à \ la variété linéaire de dim rel
m \ L dans P, correspond la \ variété linéaire\inserted[typed]{duale $L^{o}$} de dim relative
\xout{(x=1)xmx} \ (r-1)-1 -m \ \ de $\check{P}$,\inserted{i.e. de codim relative m+1 dans $\check{P}$} (NB \ r-1 \ est ici la dim relative
commune de \ P et $\check{P}$ sur S), qu'on peut visualiser géométriquement ainsi,
Faisant d'abord \ n=r-1, on trouve \ un isomorphisme
\[ \underline{P}(\check{\underline{E}}) \simeq \underline{\mathrm{Grass}}_{r-1}(\underline{E}) \]

\mspage{64}

\clearpage
% Nouveau EGA V (cote 25): archive cote 25, site page 65 (PDF page 66).
% typescript. Draft compared: none.
% AI-transcribed by Claude Opus 5.5 (claude-opus-5-5, adaptive thinking, effort high), Message Batches API, batch ega25b (msgbatch_01A7n5YuXubUPNK1zsepi8RG), 2026-11-13; not yet checked by eye.
\scanpage{25/65}
% NOTE: the top right corner is cut off; a typed number, possibly "33", is only partly visible there
\mspage{\unsure{33}} \mspage{\hand{60}}

\unsure{permettant|permettamt} d'identifier \over{L}{\unsure{l}}es points de $\hand{\check{}}P$\hand{,} à valeurs dans \xout{xxxSi} S disons\hand{,}
comme les sous-variétés linéaires de codimension 1 de P (appelés enco-%
re \uline{hyperplans} de P). Ceci dit, $L^{o}$ consiste en les hyperplans qui
\uline{contiennent} la sous-\xout{\unsure{va}}riété linéaire L de P (par quoi bien s\over{u}{û}r on
entend que les points de $L^{o}$ \xout{xxxxx} à valeurs dans un S' sont les
hyperplans dans $P_{S'}$ qui contiennent $L_{S'}$ ). Ceci résulte du fait \xout{qxx}
\margin{|}(qui aurait dû venir en même temps que le fait qu'une \inserted[typed]{sous-}variété linéaire \inserted[typed]{L de P}\hand{,}
détermine le quotient loc libre $\underline{\unsure{G}}$ de $\underline{E}$ dont il provient) que si
$\underline{G}$ et $\underline{G}'$ sont deux quotients loc libres de $\underline{E}$ (pas néc de même rang)
alors $P(\underline{G}')\hand{\hookrightarrow}P(\underline{G})$ (en tant que sous-variétés linéaires de $P(\underline{E})$) ssi
$\underline{G}'$ est majoré par $\underline{G}$ (et l'inclusion $P(\underline{G}')\hand{\rightarrow}P(\underline{G})$ n'est alors autre
que le morphisme déduit de $\underline{G}\hand{\rightarrow}\underline{G}'$ ).

Voici un minimum de sorites dont on devra disposer. La liste
complète ne pourra de toutes façons être fixée qu'une fois que l'ensem-%
bles des autres n$^{\text{o}}$s du présent paragraphe sera\xout{\unsure{nt}} rédigé\xout{\unsure{s}}.

\deleted{12. \uline{Ratrappage de résultats \xout{xxbxix} sur les sections linéaires} .}

\hand{Il me semble commode d'introduire}
\hand{aussi le foncteur}
\[ \hand{\uline{Grass}_{\unsure{\bullet}}(\underline{E})(S') = \text{ens des Modules quotients}} \]
\hand{loc. libres (de rang}
\hand{non précisé) de $\underline{E}_{S'}$}

\hand{Alors \uline{Grass}$_{\unsure{*}}(\underline{E})$ est représentable par}
\hand{$\coprod_{n\geq 0}$ \uline{Grass}$_{n}(\underline{E})$. Les \uline{sous-variétés linéaires}}
\hand{de $P(\underline{E})$ sont en fait définis par les}
\hand{sections de \uline{Grass}$_{\unsure{*}}(\underline{E})$ sur S [NB. le}
\hand{rang i.e. la dim relative peut varier si S non connexe)}

\mspage{\hand{65}}

\clearpage
% Nouveau EGA V (cote 25): archive cote 25, site page 66 (PDF page 67).
% typescript. Draft compared: none.
% AI-transcribed by Claude Opus 5.5 (claude-opus-5-5, adaptive thinking, effort high), Message Batches API, batch ega25b (msgbatch_01A7n5YuXubUPNK1zsepi8RG), 2026-11-13; not yet checked by eye.
\scanpage{25/66}
\mspage{\hand{\unsure{21}}}

12. \uline{Généralisation de résultats antérieurs aux sections linéaires}.

Compléments de notations. Si $P = P(\underline{E})$, $\underline{E}$ Module quasi-cohérent
quelconque, on posera aussi $\mathrm{Grass}_n(P) = \mathrm{Grass}_{n+1}(\underline{E})$ , de sorte que
$\mathrm{Grass}_n(P)$ correspond aux sous-variétés linéaires de \uline{dimension} n dans
P ; ceci vaut pour $n \geq -1$ , en convenant que dim $-1$ signifie vide
Lorsque E est localement libre, \xout{\unsure{de rang}} il y a lieu également d'intro-%
duire $\mathrm{Grass}^n(P) = \mathrm{Grass}_{n-1}(\check{P}) = \mathrm{Grass}_n(\check{\underline{E}})$ , qui correspond aux sous-%
variétés linéaires de codimension n dans P. Si $\underline{E}$ est de rang r+1 i.e.
P de dimension relative r, on a un isomorphisme canonique
$\mathrm{Grass}^n(P) = \mathrm{Grass}_{r-n}(P)$ . Par la suite, nous supposons fixé E loc
libre de rang r , et nous nous intéressons aux sous-variétés linéaires
de P de codimension donnée m , donc à $\mathrm{Gr}^m = \mathrm{Gr}^m(P) = \mathrm{Gr}_{\unsure{m+1}}(\check{\underline{E}})$ .
Sur ce préschéma, il y a donc un quotient \inserted{$\underline{G}$} loc libre de rang m canoniqu
de $\underline{E}_{\mathrm{Gr}}$ , soit F . \xout{\unsure{La relation di}} Le préschéma d'incidence naturel
sur $P\times_S \mathrm{Gr}^m$ , qui représente le sous-foncteur du foncteur produit
correspond à des couples d'une section de $P_S$, et une sous-variété
linéaire de codimension m de $P_S$, \uline{contenant} cette dernière, s'explicite
alors de la façon suivante: Soit $T=P\times_S \mathrm{Gr}^m$ (ou si on préfère, un
préschéma quelconque au dessus de ce produit), alors sur \xout{E} T on a
$\underline{E}_T$ , le quotient $\underline{O}_T(1)$ , et le sous-Module localement facteur direct
$\check{\underline{G}}_T$ . Considérons l'homomorphisme composé des hom canoniques
\[ \check{G}_T \longrightarrow \underline{O}_T(1) \]
qui par transposition correspond aussi à l'homomorphisme composé
analogue \unsure{du|da} sous-Module $\underline{O}_T(-1)$ de $\check{\underline{E}}_T$ dans le quotient $\underline{G}_T{}^{\vee}$:
\[ \underline{O}_T(-1) \longrightarrow G_T \quad , \]
et peut aussi s'envisager comme défini par une section de $\underline{G}_T(1)$
\[ \phi^m \in \Gamma(T, \underline{G}_T(1)) \]

\mspage{\hand{66}}

\clearpage
% Nouveau EGA V (cote 25): archive cote 25, site page 67 (PDF page 68).
% typescript. Draft compared: none.
% AI-transcribed by Claude Opus 5.5 (claude-opus-5-5, adaptive thinking, effort high), Message Batches API, batch ega25b (msgbatch_01A7n5YuXubUPNK1zsepi8RG), 2026-11-13; not yet checked by eye.
\scanpage{25/67}
\mspage{\unsure{56}}

Le préschémad'incidence  (resp. son image inverse dans T) n'est alors
autre que les préschémas des zéros de l'un ou l'autre homomorphisme,
ou encore de lasection $\phi^{m}$ .  On pourrait dénoter le préschéma
d'incidence par $H^{(m)}$ , pour m=1 on retouve celui du n°1 .  Si X est
au dessus de P, on pourra poser $\underline{Y}^{(m)} = X\times_{P} \underline{H}^{(m)}$ ,  et définir de même
la notation $\underline{Y}^{m}$ lorsque \hand{$\zeta$} est un point de $\unsure{\mathrm{Gr}}^{m}$ à valeurs dans un S'.
Ainsi les \xout{\unsure{Y}} $Y^{m}_{\hand{$\zeta$}}$  sont les``\uline{sections} \uline{linéaires}''  de X \inserted{sur P} (ou plutôt des
$X_{S'}$\inserted{sur $P_{S'}$})   par les sous-variétés linéaires de codimension m de P (ou
plutôt des$P_{S'}$).

\margin{$|$}Je profite de l'occasion pour une\xout{xx}  autocritique notationnelle,
qui pourrait venir dès le n°1. C'est le fait  de faire correspondre
arbitrairement, pour indiquer un objet $\underline{Y}$ qui correspond à X, la lettre
Y à X  (de sorte que si X devient Z, on ne sait plus très bien quoi\xout{x}
prendre). Cet inconvénient m'a déjà amené à quelques incohérences de
notations passagères.  Peut-être que le contexte un peu plus général
avec un entier m comme ici suggère une solution raisonnnable: écrire
$\underline{X}^{(m)}$ au lieu de  $\underline{Y}^{(m)}$ , donc  $\underline{X}^{(1)}$ au lieu du  $\underline{Y}$  du n° 1. De cette
façon, on  aurait bien , à peu de choses près ... ,  $\underline{X}^{(m)(m')} = X^{(m+m')}$\hand{.}
Je vais essayer cette notation par la suite.  Evidemment, l'exposant m
lui-même prête à critique, car  c'est une pratique courante en géométri\unsure{e}
algébrique de désigner par un exp osant la dimension des variétés qui
entrent en jeu. Mais co mme nous n'avons jamais fait usage de ce genre
de convention, je pense que nous gardons les mains libres à ce sujet.

On voit aussitôt, sur la construction précédente de $\underline{X}^{(m)}$, que
l'on a un\inserted{P-}isomorphisme canonique
\[ X^{(m)} = \underset{\unsure{1}}{\underline{\mathrm{Grass}}}_{m}\xout{\$}(\overset{\hand{$\vee$}}{\underline{F}}) , \]
où  $F_{X}$\xout{xxxxxx}  ( $\simeq f^{*}(\underline{\Omega}_{P/S})(1)$)   est le noyau de  $\underline{E}_{X} \longrightarrow \mathcal{O}_{X}(1)$\hand{,}

\mspage{67}

\clearpage
% Nouveau EGA V (cote 25): archive cote 25, site page 68 (PDF page 69).
% typescript. Draft compared: none.
% AI-transcribed by Claude Opus 5.5 (claude-opus-5-5, adaptive thinking, effort high), Message Batches API, batch ega25b (msgbatch_01A7n5YuXubUPNK1zsepi8RG), 2026-11-13; not yet checked by eye.
\scanpage{25/68}
\mspage{63}

en particulier $\underline{X}^{(m)}$ est lisse sur X, avec des fibres géométriquement
intègres (et en fait des variétés rationelles)\inserted[typed]{de dimension m(r-m),}. Bien entendu, la
vérification se réduit au c as X=P, et à ce titre appartient, tout
\margin{|}comme les considérations précédentes, au n$^{\text{o}}$ des généralités sur les
\margin{|}grassmani\xout{z}ennes (qui , je le vois déjà, vont \over{\unsure{f}}{f}inir par être ma\over{\unsure{gu}}{\unsure{g}}ifiés
\margin{|}par toi en un paragraphe séparé). Nous sommes maintenant en posses-%
sion d 'un analogue parfait du diagramme du n$^{\text{o}}$1 . Encore un point
oublié: comme préschéma au dessus de Gr$^{m}$ , $\underline{H}^{(m)}$ est canoniquement
isomorphe à P$(\underline{E}_{\mathrm{Gr}}\ /\overset{\hand{$\vee$}}{\underline{G}}\ )$, c'est donc un excellent fibré projectif
(mais on ne peut rien en conclure bien sur en général pour $\underline{X}^{(m)}$ sur
Gr$^{m}$).

La proposition 2.1. se transpose sans changement. Dans la
propositi on 2.2. , il faut lire: il f et s que pour tout $x \in Z$, on
ait \xout{dimxx} dim $\overline{x} \leqslant$ m-\hand{\unsure{1}}. Pour la démonstration, on peut par exemple
se ramener à 2.6. en considérant une \xout{xxxxxx} variété linéaire de
codimension m générique comme intersection de m hyperplans génériques
indépendants, Di\unsure{eu}donné demerdetur. Du point de vue rédaction, si
(comme il me s emble préférable) on fa\xout{i}it dès le début m général, \xout{\illegible{1}}
semble préférable de prouver en même temps 2.6. , où bien entendu
dim X -1 est remplacé par dim X -m (en sous-entendant qu'une dimension
$<$ 0 dans la formule si\xout{\unsure{x}}gnifie que l'ensemble considéré est vide).
Corollaire 2.3. se lit en r emplaçant ``fini'' par ``de dimension $\leqslant$ \xout{mxx}
m-1 '' . Corollaire 2.4. pareil ; De même pour 2.5. , remplacer
dim f(X$_{i}$ ) $>$ 0 par dim f(X$_{i}$) $\geqslant$ m , et même changement dans 2.7.
La proposition 2.8.reste valable telle quelle, dans 2.9. remplacer
fini par dim $\leqslant$m-1. Kif kif pour 2.10. , 2.11. . L'énoncé
2.12. reste valable tel quel, avec une démonstration essentiellement

\margin{68}

\clearpage
% Nouveau EGA V (cote 25): archive cote 25, site page 69 (PDF page 70).
% typescript. Draft compared: none.
% AI-transcribed by Claude Opus 5.5 (claude-opus-5-5, adaptive thinking, effort high), Message Batches API, batch ega25b (msgbatch_01A7n5YuXubUPNK1zsepi8RG), 2026-11-13; not yet checked by eye.
\scanpage{25/69}
\mspage{69}

inchangée (cf aussi \inserted{\unsure{plus bas}} commentaires à n\textsuperscript{o} 8). 2.13. remplacer fini
par dim \hand{$\leqslant$} m-1. Th. 2.14. reste valable tel quel. , 2.15. en
remplaçant fini par dim \hand{$\leqslant$} m-1. 2.16. valable t el quel\xout{x} ,dans
\xout{2x} 2.17. remplacer fini par dim \hand{$\leqslant$} m-1. 2.18. tel quel.

Pour 3.1. , on peut l'énoncer avec m quelconque en supposant
dim f(X) \hand{$\geqslant$} m+1 , mais je propose de garder l'énoncé principal dans
le cas des \hand{?}hypersurfaces, et de signaler le cas général en corollaire
ou en remarque (il s'en\xout{x} déduit aussitôt par le procédé habituel de
prendre des hyperplans génériques indépendants). A moins qu'il soit
amusant d'expliciter une version généralisée du lemme 3.1.1. ...
Pour 3.2. lire dim f(X)\hand{$\geqslant$} m+1 , dans 3.3. remplacer dim f(X$_i$) \hand{$\geqslant$} 2
par dim f(X$_i$)\hand{$\geqslant$} m+1, et dans la\xout{x} définition de \hand{\unsure{$\mathcal{C}$|$\mathfrak{C}$}} dim f(Z) = 0
par dim f(Z) \hand{$\leqslant$} m-1.

Les considérations générales \unsure{du|da} n\textsuperscript{o} 4 s'appliquent tels quels
au cas m quelconque. Il en e st de même de 4.2. , et de 4.3. en rempla-%
çant dans b) (v) et (vi) la condition dimensionelle par
dim f(X)\hand{$\geqslant$ m+1}. Changement analogue dans 4.4. b).

Le laïus 5.1. passe tel quel. Dans 5.2. il faut se rappeler
que $\phi$ devient une section $\phi^{(m)}$ de $\underline{G}_T(1)$ (où T \xout{est un préschéma}
\xout{au dessus de} = $X\mathrm{x}_S Gr^m$) \xout{\unsure{et la condition}} induisant des sections
$\phi^m_{\hand{$\xi$}}$ des \xout{\unsure{(m)}} \xout{\illegible{1}} $\underline{O}_{X_{k(\hand{$\xi$})}}(1)\hand{$\otimes$}_{k(\hand{$\xi$})}\underline{G}(\hand{$\xi$})$ (pour \hand{$\xi \in$} Gr$^m$). Or
de façon générale, on aura expliqué au par.19 , lorsqu'on a une
\margin{$\|$} section $\phi$ d'un Module loc libre de rang m sur un préschéma, ce que
\margin{$\|$} cela veut dire que cette s ection soit $\underline{F}$ -régulière, pour un Module
\margin{$\|$} $\underline{F}$ donné : en termes d'une base locale, cela signifie qu'on a une
\margin{$\|$} suite F-régulière de m sections de $\underline{O}_X$ (et il faut vérifier que c ela

\mspage{69}

\clearpage
% Nouveau EGA V (cote 25): archive cote 25, site page 70 (PDF page 71).
% typescript. Draft compared: none.
% AI-transcribed by Claude Opus 5.5 (claude-opus-5-5, adaptive thinking, effort high), Message Batches API, batch ega25b (msgbatch_01A7n5YuXubUPNK1zsepi8RG), 2026-11-13; not yet checked by eye.
\scanpage{25/70}
\mspage{65}

est bien indépendant de la base\hand{/}choisie.  Dans le cas  m=1, on a
l'interprétation intrinsèque évidente signalée dans 5.2.  Avec cette
convention de langage, 5.3.  reste valable t el quel, d e même 5.4.
La première partie de la remarque 5.5.   admet une \xout{généralisation} \inserted[typed]{\xout{variante pour le}}
\inserted[typed]{généralisation au}
\xout{au} cas m quelconque:  lorsque $F_s$  est \xout{\unsure{Sm}} \inserted[typed]{$(S_m)$}, alors la condition de
régularité envisagée pour $\phi_{\hand{$\xi$}}$  s'exprime de façon purement dimensionelle.
La deuxième partie de la remarque 5.5. vaut telle quelle dans le cas
m quelconque.      Théorè\unsure{m|p}e 5.6.  s'étend tel quel,  de même 5.7.

\uline{Proposition} 6.1.  :  lire   dim $f(X_i) \hand{$\geqslant$}$  m+1 ,  et
plus loin   dim $f(Z) \hand{$\leqslant$}$ m-1.

Les laius généraux d e 7.1. sont valables tels quels dans le cas
m quelconque.  7.2. et 7.3.  mutatis mutandis (attention dans 7.2.
à la not\unsure{a}tion m, confusante ici),  d'autre part dans la démonstration
de 7.4. on n'a plus besoin de procéder de proche en proche, mais on
peut prendre directement une section linéaire  de codimension m=n .

Dans 8.2. , remplacer l\unsure{a}\xout{hypothèse} condition    dim $Y_{\hand{$\xi$}}$ \hand{$>$} dim X-\unsure{1}
par  dim $Y_{\hand{$\xi$}}$ \hand{$>$}  dim X  - \over{1}{m} , et l'hypothèse  dim $f(X_i) \hand{$>$} 0$  par
dim $f(X_i) \hand{$\geqslant$}$ m.   Modification analogue dans la suite de 8.2.  Comme
8.3.  présente un exemple, il n'y a pas lieu de le changer, donc on
y gardera m=1.   Je te laisse le soin de faire l'exercice,demandant
sans doute\hand{/}un peu plus de soin, de trouver  les bons énoncés pour m
quelconque correspondants à  8.4.  8.5.  et 8.6.   (pages \inserted{30,}31,32 )\hand{;}
\hand{il n'est d'ailleurs sans doute pas indispensable de}
\hand{faire cet exercice, s'il te barbe de le faire.}

\mspage{70}

\clearpage
% Nouveau EGA V (cote 25): archive cote 25, site page 71 (PDF page 72).
% typescript. Draft compared: none.
% AI-transcribed by Claude Opus 5.5 (claude-opus-5-5, adaptive thinking, effort high), Message Batches API, batch ega25b (msgbatch_01A7n5YuXubUPNK1zsepi8RG), 2026-11-13; not yet checked by eye.
\scanpage{25/71}
\mspage{66}

\xout{Pour autant que je puisse voir, le théorème 8.7 , notamment}
\xout{exact}
\xout{8.7 \unsure{e)}, n'admet pas d'analogue évident pour m quelconque,}

Je pense que essentiellement tous les dév\unsure{el}oppements du n°8 à
partir de 8.6. peuvent s'apdapter \inserted{également} au cas des sections linéaires avec
m quelconque. Le faire en forme constituerait un exercice sans doute
assez long et fastidieux. Je dois admettre que je ne connais pas d'appli-%
cation s'appuyant de façon essentielle sur l'analyse de cette situation
plus générale, donc nous ne sommes pas vraiment obligés d'inclure ces
développements dans les Eléments. D'autre\unsure{x} part, l'expérience prouve que
le fait de rédiger dans un cadre plus général oblige parfois à mieux
dévisser et fait mieux comprendre le fourbis, et souvent sans beaucoup
plus de mal, de plus un certain nombre d'exercices de syntaxe dans un
contexte proprement géométrique comme ici ne ferait pas de mal,-et
bien entendu il n'est pas du tout exclu qu'on en ait quand même besoin
un jour et soit content de trouver la chose faite. Aussi je te laisse
l'entière décision à ce sujet, et me borne à esquisser simplement les
\xout{\illegible{1}}énoncés qu'on pourrait éventuellement donner à ce propos.

Supposons encore que $f:X \hand{$\longrightarrow$} P$ soit non ramffié et X soit lisse sur
S,\inserted[typed]{à composantes de dimension \hand{$\geq$} m.}\hand{.} \hand{(}Dans $\underline{X}^{(m)}$ \inserted{$=V^o$}\xout{xxx} on distingue alors le sous-préschéma $\underline{X}^{(m)\mathrm{sing}} = V^1$
des zéros \unsure{singuliers} de $\phi^m$ relativement à Grass$^m$ , qui est aussi formé
géométriquement des couples $(x,\hand{$\xi$})$ tels que la variété linéaire $L_{\hand{$\xi$}}$
coupe excédentairement l'espace tangent à X en x (considéré comme
sous-variété linéaire de P), i.e. tels que les deux espaces n'engendrent
pas P tout entier. Contrairement à ce qui se passe pour m=1, lorsque
m \inserted[typed]{est quelconqué}\xout{\illegible{3}}, le morphisme $\underline{X}^{(m)}$ sing $\hand{$\longrightarrow$}$ X n'est en général pas
lisse, car la variété des L qui \inserted{\unsure{passant par x et qui}}\hand{(}coupent excédentairement une sous-varié-%
té linéaire donnée T\hand{\unsure{$\ni x$}} n'est en général pas lisse sur k: cette variété

\margin{71}

\clearpage
% Nouveau EGA V (cote 25): archive cote 25, site page 72 (PDF page 73).
% typescript. Draft compared: none.
% AI-transcribed by Claude Opus 5.5 (claude-opus-5-5, adaptive thinking, effort high), Message Batches API, batch ega25b (msgbatch_01A7n5YuXubUPNK1zsepi8RG), 2026-11-13; not yet checked by eye.
\scanpage{25/72}
\mspage{67}

adt seulement l'adhérence de la sous-variété lisse formé des $\hand{\zeta}$ tels que
la dimension de $T\hand{\cap}L_{\hand{\zeta}}$ soit juste \uline{une} de plus que la dimension
``normal\unsure{è|e}'' n-m (n=dim T, m = codim L ). Sauf erreur, l'ensemble
(contenu dans l'ensemble supersingulier relatif) $V''$ introduit dans le par
16 (compléments) n'est autre que l'ensemble formé des couples
$(x,L_{\hand{\zeta}})$ tels que la dimension de $T_x\hand{\cap}L_{\hand{\zeta}}$ soit $\hand{\geqslant}$ n-m + 2\hand{)} , de sorte que
$\underline{X}$x $V^1 - V''$ \xout{\unsure{mx}}est \xout{xxxxxxxxx} lisse sur X, et sauf erreur est même \inserted{exactement} l'ensem
ble des points de $V^1$ lisses sur X.(La vér\xout{\unsure{i}}ification de ce point demanderai\unsure{t}
une étude de la filtration du \xout{xx} \unsure{shhéma|schéma} grassmanien suivant les \xout{xxxxx}
dimensions d'intersection des L variables avec une T fixe, sauf erreur on
trouve que le cran suivant de la filtration est formée exactement des
points non lisses du cran précédent - quand on définit la filtration
\margin{$|$} non seulement ensemblistement, mais aussi schématiquement, en utilisant
\margin{$|$} le lemme de la page 16 des compléments au par.16 \hand{\unsure{)}}. Cette étude forme-%
\margin{$|$} rait alors un des n$^{\text{o}}$s du paragraphe ``géométrique'' consacré aux grassmanien
\margin{$|$} nes \hand{)}. Si on définit de même \xout{V(k)} $V^{(k)}$ comme le sous-schéma de $X^{(m)}$
correspondant à dim $T_x\hand{\cap}L \geqslant$ n-m + k , on trouve par un calcul immédiat
que \xout{dimx$V^{(k)}$xxxxx\unsure{(k+1)(n+m)+k}x$^2$}
\[ \text{dim Grass}^m(P) - \text{dim } V^{(k)} = (k-1)(n-m) + k^2 \]
(du moins pour les restrictions raisonnables k $\hand{\leqslant}$ m, k $\hand{\leqslant}$r-m), sauf
erreur de calcul (NB ceci résulte plus généralement d\xout{xx}'un calcul des
dimensions des ``cellules'' qui interviennent dans la filtration des grass-%
mani\xout{x}ennes à laquelle j'ai fait allusion plus haut). Pour k = 2, on
trouve une différence de dimension $\geqslant 4$ , donc l'image de $V''$ dans Grass$^m$
est de codimension $\hand{\geqslant}$ 4, donc si on s'intéresse à ce qui se passe en dehors
d'ensembles \inserted[typed]{de Grass$^m$} de codimension $\geqslant$ 2, on peut \xout{xxx} oublier $V''$ .

D'autre part, dans $X_{\text{Grass}^m} - V''$ au dessus de Grass$^m$ , la situation
est celle du cas favorable envisa\hand{\unsure{gé}} dans les `` compléments au par 16'' ,

\mspage{72}

\clearpage
% Nouveau EGA V (cote 25): archive cote 25, site page 73 (PDF page 74).
% typescript. Draft compared: none.
% AI-transcribed by Claude Opus 5.5 (claude-opus-5-5, adaptive thinking, effort high), Message Batches API, batch ega25b (msgbatch_01A7n5YuXubUPNK1zsepi8RG), 2026-11-13; not yet checked by eye.
\scanpage{25/73}
\mspage{68}

relati\xout{v}vement  au s chéma de base S:  V$^1$- V"  est en\xout{xx} effet lisse sur S
(l'étant sur X)  de dimension relative égale à \uline{une} de moins que celle de
Grass$^m$ sur S  (comme on voit en faisant k=1 dans la formule plus haut).
Ainsi  les résultats d\unsure{e}  loc cit  s'appliquent, on trouve en particulier
le fait que l'ensemble des points  supersinguliers  de $\phi^m$ rel à Grass$^m$
n'est autre que V" \hand{$\cup$} V$^2$ , où  V$^2$ est le sous-\inserted{pré}schéma de ramification
de  V$^1$-V" \hand{$\longrightarrow$} G\unsure{r}ass$^m$ .  On peut donc dire encore que, en dehors de V", les
zéros supersinguliers résultent du collapsage de (au moins) deux zéros
singuliers ordinaires (mais on  n'est pas obligé de le dire\hand{)}).

\margin{$|$} De cette façon, on a essentiellement l'équivalent de 8.7. a) et b). Il \margin{$+$}
doit être possible de donner un équivalent pour  8.7. c), en \xout{xxx} utilisant
la \xout{structure} description explicite du fibré tangent  à Grass$^m$ (analogue à
celle du cas m=1); \deleted{et qui} \inserted{cela} implique \inserted{encore} que  pour  un point \inserted{géométrique} de  V$^1$-V"  non
ramifié au dessus  de Grass$^m$, la connaissance de son image dans Grass$^m$
implique celle de son image dans P, pourvu que la première image soit un
point lisse de l'image fermée \inserted{de} V$^1$ dans Grass$^m$  (on suppose S = spectre d'
un corps).  Je peux fournir un énoncé \xout{xxxxxx} plus précis sur demande.

Une fois ce point acquis,  on a des corollaires évidents généralisant
8.8. , 8.9. , 8.11.   Il est sans doute tout aussi possible d'énoncer dans
le cas m quelconque  les autres propositions du par.8. Si cela demande des
efforts de rédaction  supplémentaires, on peut renoncer à ce\unsure{tte} généralisa-%
tion\hand{,} \xout{mêm} même si on inclut les développements différentiels précédents.
Il en est de même des résultats du n$^o$ 9.

Quant au n$^o$ 10 ,  la situation qui y est étudiée se généralise au
cas m quelconque de la façon suivante. On se fixe une sous-variété linéai-
re C de P  de codimension m+1, et on considère l'espace projectif Q des
sous-v\unsure{ariét}és linéaires L de P de codimension m passant par C. Q est une

\margin{73}

\clearpage
% Nouveau EGA V (cote 25): archive cote 25, site page 74 (PDF page 75).
% typescript. Draft compared: none.
% AI-transcribed by Claude Opus 5.5 (claude-opus-5-5, adaptive thinking, effort high), Message Batches API, batch ega25b (msgbatch_01A7n5YuXubUPNK1zsepi8RG), 2026-11-13; not yet checked by eye.
\scanpage{25/74}
\mspage{\hand{\unsure{61|6}}}

sous-préschéma fermé d e Grass$^{m}$ , en particulier on peut former $\underline{X}^{(m)}_{Q}$ ,
qu'on se propose d'étudier. Un premier point, qui doit en tout état de
cause figurer dans le texte, c'est que $\underline{X}^{(m)}_{Q} \longrightarrow X$ est encore birationnel
du moins lorsque C coupe ``régulièrement'' X , et de façon précise $\underline{X}^{(m)}_{Q}$
est dans œ cas canoniquement isomorphe \unsure{àu} préschéma déduit de X en
faisant éclater $X\mathrm{x}_{P}\unsure{C}$ : la démonstration de ce fait n'est autre que celle
de 10.2. via \over{10}{\unsure{18}}.3. Un deuxième point,qui présente un certain intérêt,
mais que nous ne sommes pas absolument obligés d'inclure\xout{r} consiste à
dire que si on a choisi C ``assez général'', alors $\underline{X}^{(m)}_{Q} \longrightarrow Q$ a certaines
propriétés sympathiques, la plus classique étant celle-ci: X étant
supposé lisse sur S=Spec(k)\inserted[typed]{de dim $n \geq m$},et \xout{toutes}\inserted[typed]{propre et} géom irréd , Alors pour C ``assez
général'', l'ensemble\inserted[typed]{T} des $\hand{\zeta}\in Q$ pour lesquels $\underline{X}^{(m)}$ n'est\inserted[typed]{pas} lisse de dimen-%
sion n-m sur k($\hand{\zeta}$) est géom irréductible sur k($\hand{\zeta}$)\inserted[typed]{et de codim 1 dans Q}, et \xout{pour}
l'ensemble\inserted[typed]{T$^{1}$} des $\hand{\zeta}\in T$ pour lesquels $\underline{X}^{(m)}$ est ``supersingulier'' en au
moins un point est \xout{r}rare dans T;\xout{\unsure{ainsi, pour tout $\zeta \in T-T^{1}$, $\underline{X}^{(m)}_{\zeta}$ est}}
\xout{\unsure{n'admet que des points de}} enfin , si f:X$\rightarrow$P est une immersion, alors
quitte à agrandir un peu T$^{1}$ , pour tout $\hand{\zeta}\in$ T-T$^{1}$ , il y a exactement un
\margin{$|$}point non lisse dans $X^{(m)}_{\hand{\zeta}}$ , et ce dernier est rationnel sur k($\hand{\zeta}$). J'ai
oublié dans l'énoncé de spécifier qu'on supposait X \xout{de} $\longrightarrow$ P non ramifié,
et que l'on devait au préalable remplacer f par $\phi_{n}f$ , n $\geq$2 (où $\phi_{n}$ est
défini dans 9.1.). Le plus naturel \xout{semble} pour prouver cet énoncé sem-%
ble d'utiliser le sous-schéma Z (noté T dans 8.8.) des Grass$^{m}$ tels
que $\underline{X}^{(m)}$ soit \xout{``supers}Singulier\unsure{''}: on sait, sous les conditions dites, que Z
est géom. irréd. de codimension 1, et que le sous-schéma Z$^{1}$ correspondant
à des $\underline{X}^{(m)}$ ``supersinguliers'' y est r are. Il reste\inserted{rait} alors à prouver un
lemme de la nature suivante: soit Z une partie fermée de Grass$^{m}$ de codi-%
mension q, alors définissant Q(C) en t ermes de C comme ci-dessus, pour
tout C ``assez général'' l'intersection $Q(\unsure{Z|C})\cap Z$ est de codim $\geq$ q dans
\inserted{$Q(C)$; de plus, si $Z$ géom. irréd., itou $Q(C)\cap Z$ si $Z$ ``assez général''.}

\mspage{\hand{74}}

\clearpage
% Nouveau EGA V (cote 25): archive cote 25, site page 75 (PDF page 76).
% typescript. Draft compared: none.
% AI-transcribed by Claude Opus 5.5 (claude-opus-5-5, adaptive thinking, effort high), Message Batches API, batch ega25b (msgbatch_01A7n5YuXubUPNK1zsepi8RG), 2026-11-13; not yet checked by eye.
\scanpage{25/75}
\mspage{70}

20.13. \uline{Morphismes élémentaires et théorème de M. Artin}.

\uline{Définition} 13.1. Un morphisme $f:X\to Y$ de préschémas est dit ``morphisme
\xout{\unsure{élém}} élémentaire'' si X est \xout{\unsure{un}} Y-isomorphe à un préschéma de la forme
$X' - Z$ , où $X'\to X$ est un morphisme lisse, projectif, à fibres géomé-%
triques connexes\xout{\unsure{,}} de dimension 1, et où Z est un sous-préschéma fermé
de X', \xout{\unsure{étale sur}} tel que le morphisme $Z\to Y$ soit étale, surjectif, et
de degré constant. Un morphisme est dit \xout{\unsure{multip-élémentaire}} poly-élémen-%
taire s'il est un composé de morphismes élémentaires. Un préschéma X
sur un corps k est dit polyélementaire (sur k), si le morphisme structu-%
ral $X\to \mathrm{Spec}(k)$ est poly-élémentaire.

\uline{Théorème} 13.2. \inserted[typed]{(M.Artin)} Soient X un préschéma \inserted[typed]{géom. irréd.} sur un corps k \inserted[typed]{parfait infini}, x un point lisse
de X, alors x admet un système fondamental de voisinages ouverts
polyélémentaires.

Quitte à remplacer X par un voisinage donné de x, \deleted{on \xout{\unsure{peutxsuppo}} es} \inserted{il revient au}
\deleted{ramené à} \inserted{même de} montrer qu'il existe un voisinage ouvert élémentaire de x dans X.
\deleted{On peut supposer évidemment X lisse, et} \inserted[typed]{\deleted{et affine}} \over{r}{R}aisonnant par récurrence sur
la dimension n de X, on \xout{\unsure{peutxsupposer}} est ramené à montrer \inserted[typed]{que si $n>0$ ,} \deleted{qu'}il existe
un voisinage ouvert U de x et un morphisme élémentaire $f: U\to V$,
V étant un schéma lisse sur k (nécessairement géom. irréd \inserted[typed]{et} de dim n-1).
\xout{\unsure{Lorsque n=1}} (Le cas n=0 est évidemment trivial, étant donné qu\over{e}{'} alors
X est k-isomorphe à Spec(k), qui est polyélémentaire sur Spec(k), étant
entendu dans 13.1. qu'on n'exc\hand{l}ut pas le composé de la famille vide de
morphismes - on aurait dû le signaler sous une forme ou une autre dans
13.1. ). La nécessité de supposer k parfait au préalable apparait déjà
dans le cas n=1: alors on prend pour X' le modèle projectif normal
canonique du corps des fonctions K de X (cf Chap II) par.7), le fait que

\mspage{75}

\clearpage
% Nouveau EGA V (cote 25): archive cote 25, site page 76 (PDF page 77).
% typescript. Draft compared: none.
% AI-transcribed by Claude Opus 5.5 (claude-opus-5-5, adaptive thinking, effort high), Message Batches API, batch ega25b (msgbatch_01A7n5YuXubUPNK1zsepi8RG), 2026-11-13; not yet checked by eye.
\scanpage{25/76}
\mspage{71}

\unsure{K|k} soit parfait assure que X' est lisse sur k (car X' est en tous cas
régulier), et il assure également que $Z = X' - X$ , muni de la structure
réduite induite, est étale sur k. Traitons maintenant le cas général,
dans lequel il est loisible de supposer $n \geqslant 2$.

\xout{Comme X est affine donc quasi-projectif,} On peut évidemment suppo-%
ser X \xout{\unsure{lisse et}} affine donc quasi-projectif, puis, quitte à remplacer X
par une clôture projective, on peut supposer X projectif, sous réserve
toutesfois de prouver que \uline{tout} voisinage de x contient un voisinage
ouvert U admettant un morphisme élémentaire $U \to V$. De même, remplaçant
X par son normalisé (fini sur X donc projectif, réf.), ce qui ne le
change pas au voisinage de x, on peut supposer X normal, donc k étant
parfait, géométriquement normal sur k. \xout{Choisissons une immersion} Le
bénéfice de cette hypothèse est que l'ensemble Z des points de X en
lesquels X n'est pas lisse sur k est de codim $\geqslant 2$. \xout{\unsure{Soit Noter qui}}
Choisissons une immersion projective $i : X \longrightarrow P^r$ , \xout{\unsure{On voit on peut}}
\xout{\unsure{supposer que X n'est contenu dans Alors}} Comme on obtient un système
fondamental de voisinages de x dans $\underline{P}^{\hand{r}}$ en prenant des sections ne
s'annulant pas en x des divers $\underline{O}_P(n)$ , $n > 0$, on en conclut que tout
voisinage de x contient un voisinage de la forme X-Y, où Y est une
partie fermée de X , contenant Z, purement de dimension n-1, et telle que
$x \notin Y$. Nous munirons Y de la structure réduite induite, de sorte que
(k étant parfait) l'ensemble singulier de Y est de dimension $\leqslant$ n-2. \margin{\unsure{$\nearrow$}}
Quitte à agrandir l'ensemble Z précédent, on trouve un fermé $Z \subset Y$, de
dimension $\leqslant$ n-2, contenant l'ensemble géom. singulier de X \uline{et} de Y.

L'idée de la démonstration est de fibrer X par ses intersections
avec les sous-variétés linéaires L de P , de codimension n-1, contenant

\mspage{76}

\clearpage
% Nouveau EGA V (cote 25): archive cote 25, site page 77 (PDF page 78).
% typescript. Draft compared: none.
% AI-transcribed by Claude Opus 5.5 (claude-opus-5-5, adaptive thinking, effort high), Message Batches API, batch ega25b (msgbatch_01A7n5YuXubUPNK1zsepi8RG), 2026-11-13; not yet checked by eye.
\scanpage{25/77}
\mspage{\hand{72}}

une sous-variété linéaire donnée C  de codimension n.  A cet effet,
on aura besoin du

\uline{Lemme} 13.3.  Avec les notations précédentes pour \xout{xx} X,Y,Z  ($X\supset Y\supset Z$
sous schémas fermés de $P^r_k$ ,\inserted[typed]{de dimensions n,n-1 et \hand{$\leqslant$} n-2} X-Z\hand{,} \deleted{lisse de dim n en tous ses points},
Y-Z lisse\hand{\unsure{s}} \deleted{de dim n-1 en tous ses points;}\hand{,} Z de dimension \hand{$\leqslant$}n-2), et quitte
au besoin (lorsque k est de car p$>$0)  de remplacer l'immersion projecti-%
ve  i:X$\rightarrow$P$^r$  par  un ``multiple'' quelconque $\phi_n i$ ($n \hand{$\geqslant$} 2$)  comme dans le
n$^o$ 9,  il existe une sous-variété linéaire  $L_o$ de P$^r$ , de codimension
n-1, et ayant les vertus suivantes:

a)  $L_o \cap Z = \emptyset$ .

b)  $L_o \cap X$  est lisse de dimension 1.

c)  $L_o \cap Y$   est lisse de dimension  0.

(NB  k désigne un corps infini, pas la peine ici qu'il soit parfait).

Admettons  ce lemme  et montrons comment on en déduit  l'existence
d'un voisinage ouvert U de x, contenu dans X-Y, et admettant un morphis-%
me élémentaire U$\longrightarrow$V.  \deleted{En vertu du n$^o$ 7 appliqué à l'espace projectif}
\deleted{$L_o$ et au sous-schéma lisse $L_o\cap Y$ de celui-ci,} \over{i}{I}l existe une sous-%
variété linéaire C de $L_o$ , de codimension n dans P$^r$ i.e. de codimension
1 dans $L_o$ , \inserted{ne rencontrant pas l'ensemble fini \deleted{\illegible{1}} $(L_o\cap Y)\cup\{x\}$.} \deleted{telle que C $\cap$ X ( = C $\cap$ ($L_o\cap$ X) ) soit lisse sur k et de}
\deleted{dimension 0; de plus, on peut exiger que} l'on ait  $C \cap Y = \emptyset$ \inserted{et $x\notin C$,} \deleted{(même}
\deleted{argument).} Soit  T = X$\cap$C, de sorte que T est un sous-schéma de X,
étale sur k, non vide,\inserted{ne contenant pas x,} et \uline{disjoint} de Y.  Considérons d'autre part le
sous-schéma Q de Grass$^{n-1}$(P) correspondant aux sous-variétés linéaires
de P$^r$ contenant \unsure{C|E}, de sorte que  Q est un \xout{espace schéma} espace projectif
de dimension n-1,   en particulier il est lisse sur k de dim. n-1.
Alors $L_o$ correspond à un point \hand{$\xi$}$_o$ de Q(k).  Considérons d'autre part
(avec les notations générales introduites par ailleurs)  l'image inverse
$\underline{X}^{n-1}_Q$ de $\underline{X}^{n-1}$  par  l'immersion Q$\longrightarrow$ Grass$^{n-1}$ , et de même les images

\margin{77}

\clearpage
% Nouveau EGA V (cote 25): archive cote 25, site page 78 (PDF page 79).
% typescript. Draft compared: none.
% AI-transcribed by Claude Opus 5.5 (claude-opus-5-5, adaptive thinking, effort high), Message Batches API, batch ega25b (msgbatch_01A7n5YuXubUPNK1zsepi8RG), 2026-11-13; not yet checked by eye.
\scanpage{25/78}
\mspage{\hand{73}}
% NOTE: this typed line is framed by a handwritten curved line that leads down to the place after "$\underline{T}^{n-1}_Q$ ," in the next line, just before "soient"
\unsure{qui|qbi} sont aussi des sous-schémas fermés \inserted{disjoints} de $\underline{X}^{n-1}_Q$ ;

inverses $\underline{Y}^{n-1}_Q$ et $\underline{T}^{n-1}_Q$ , soient p,q,r les projections structurales
de ces schémas dans Q. Alors par \deleted{construction}\inserted{\unsure{hypothèse essentiellement}}, p est lisse en les
points au dessus de $\hand{\zeta}_o$ , et q est étale en les points au dessus de $\hand{\zeta}_o$ ;
il en est de même de r, car on voit aussitôt que $\underline{T}^{n-1}_Q$ n'est autre que
$Tx_kQ$ ( Q-isomorphisme). Enfin, le morphisme p est propre, et compte
tenu que X est géométriquement connexe, les fibres de p sont géométrique-%
ment connexes, (th de Bertini). Par suite, il existe un voisinage ouvert
V de $\hand{\zeta}_o$ dans X tel que \xout{\unsure{le morphisme}} $\underline{X}^{n-1}_Q|V = X'$ soit propre et lisse
sur V, à fibres géométriquement connexes, et comme la fibre de $\hand{\zeta}_o$ ,
qui n'est autre que $X\cap L_o$ , est de dimension 1, on peut supposer que les
fibres de X' sur V sont toutes de dim. 1. Enfin, prenant V assez petit,
on peut supposer que $\underline{Y}^{n-1}_Q|V$ et $\underline{T}^{n-1}_Q|V$ sont étales sur V, donc le
préschéma somme Z' de ses deux (qui s'identifie à un sous-\xout{\unsure{pré}}schéma
fermé de X' ) est étale sur V. Par suite, \xout{\illegible{1}} posant $U = X' - Z'$,\xout{\illegible{1}}
\xout{\illegible{1}} le morphisme $U\to V$ est un morphisme élémentaire. Mais \xout{\illegible{1}} U est aussi
un ouvert de \inserted{$X''=$}$\underline{X}^{n-1}_Q - \underline{Y}^{n-1}_Q - \underline{T}^{n-1}_Q$ , image inverse de X-Y-T dans $\underline{X}^{n-1}_Q$,
et d'autre part $X''\to$ X-Y-T est évidemment un isomorphisme (car
$\underline{X}^{n-1}_Q - \underline{T}^{n-1}_Q \to$ X-T est un isomorphisme ). Donc U s'identifie à un ouvert
de X-Y-T, ouvert contenant d'ailleurs $L_o\cap X$ et à fortiori x. C'est le
voisinage cherché de x contenu dans X-Y.

\mspage{\hand{78}}

\clearpage
% Nouveau EGA V (cote 25): archive cote 25, site page 79 (PDF page 80).
% typescript. Draft compared: none.
% AI-transcribed by Claude Opus 5.5 (claude-opus-5-5, adaptive thinking, effort high), Message Batches API, batch ega25b (msgbatch_01A7n5YuXubUPNK1zsepi8RG), 2026-11-13; not yet checked by eye.
\scanpage{25/79}
\mspage{\hand{74}}

Reste à prouver le lemme 13.3.  Comme d'habitude, il suffit de prouver
que  la sous-variété linéaire de \inserted{co}dimension n-1  passant par x générique
L\unsure{,}  possède les propriétés  a), b) , c) .  Pour  a), \inserted{ainsi que le contenu dimensionnel de b) \unsure{c)} (pour toute immersion projective donnée)} cela résulte
\hand{[}aussitôt de 2.3. (revu et corrigé dans le nº 12)  appliqué (comme dans un
raisonnement déjà fait dans le nº8) à l'espace projectif  des droites
passant par x, et l'image de Z dans ledit par projection conique à
partir de x. \hand{[}Il pourrait être utile d'ailleurs d'expliciter certains
\margin{|} résultats obttenus par cette méthode, concernant les sections linéaires
par des sous-variétés linéaires assujetties à passer par une sous-variété
linéaire fixe,  au cours du texte ou dans\hand{/}un nº à part\hand{]}.  Pour \inserted{la lissité dans} b) et c),
\xout{remplaçant} on peut grâce à a)  remplacer X et Y  resp. par X-Z et Y-Z
qui sont lisses, et on est  ramené à prouver ceci:  Soit f: X$\rightarrow$ P un
morphisme non ramifié,\hand{(}avec X lisse sur k, et soit  x $\in$P(k),\inserted{tel que x n'appartienne à l'image d'aucune composante de X de dim $<$ m ,} alors si
\hand{$\eta$} est le point générique  d\unsure{e} \xout{xxxxx} la sous-grassmanienne  de Grass$^{m}$(P)
formée des variétés linéaires L de codimension m passant par \xout{unexvariété} \inserted{x}
\xout{linéairexdonnéexdexcodimensionxm+1,} \inserted{\unsure{\&}} $\underline{X}^{(m)}_{\hand{\eta}}$ est  lisse sur k, tout au
\margin{(} moins si k est de car 0, et dans le cas contraire, à condition de rempla-%
cer f par $\phi_{n}f$ , n entier $\geq$ 2 quelconque.  C'est là  un remords au nº9,
qui résulte lui-même du remords suivant au nº 8:  Avec les notations de
8.8. , (supposant donc X irréductible, ce qui est loisible pour le Pb
que nous traitons), si  l'on a \inserted{co}dim T $\geq$ \deleted{n-}2  ou si  $\underline{Y}^{\text{sing}}\rightarrow$ T  est
génériquement étale (condition automatiquement vérifiée si k est de car 0,
ou à condition de remplacer f par $\phi_{n}f$ avec n$\geq$2, cf nº 9), alors pour
l'hyperplan \xout{générique} H\hand{$_{\eta}$} passant par x  générique, \deleted{\unsure{on a}} $\underline{X}^{(1)}_{\hand{\eta}}$ est lisse
de dimension n-1,  sauf dans le cas  où  on aurait  f(X) =\hand{$\{$}x\hand{$\}$}  (donc n=0)\hand{.}
Ce résultat admis,\hand{(}qui liquide évidemment le cas particulier m=1 de
notre remords, on obtient aussitôt le cas  m général  par récurrence sur

\margin{\unsure{75}}

\clearpage
% Nouveau EGA V (cote 25): archive cote 25, site page 80 (PDF page 81).
% typescript. Draft compared: none.
% AI-transcribed by Claude Opus 5.5 (claude-opus-5-5, adaptive thinking, effort high), Message Batches API, batch ega25b (msgbatch_01A7n5YuXubUPNK1zsepi8RG), 2026-11-13; not yet checked by eye.
\scanpage{25/80}
\mspage{\hand{75}}

\margin{|}m, en notant qu'à changement de corps de base près, $\underline{X}^{(m)}_{\hand{\unsure{\eta|\xi}}}$ s'obtient en
prenant une L' de codimension m-1 passant par x, une H de codi\xout{me}nsion
1 passant par x, L' et H étant génériques indépendants pour ces proprié\xout{x}
tés (i.e. en termes orthodoxes, on se place au point générique du \xout{xx}
schéma des couples (L',H) \xout{xxxxxx} ) et en prenant la section linéaire
de X par $L'\cap H$ : on peut commencer par prendre \xout{\unsure{l}}a section par H, qui
est lisse par ce qui précède, puis par L', qui est lisse par l'hypothèse
de récurrence. Ce genre de raisonnement, déjà utilisé\xout{e} pour généraliser
2.6. p.ex. aux sections linéaires de codimension m queconque, mérite
d'être explicité une bonne fois en général, pour pouvoir y référer sans
entrer \xout{dans} chaque fois dans les détails un peu lourds d'une présentati-%
on en forme.

Reste à prouver le corollaire annoncé de 8.8. dans le cas m=1 .
Lorsque \xout{dim} codim T $\geqslant$ 2, comme d'autre part l'hyperplan \inserted{Q} de $\overset{\hand{\vee}}{P}$ des
$H_{\hand{\xi}}$ tels que $x\in H_{\hand{\xi}}$ est de codimension 1, son point générique \inserted{$\eta$} ne peut
être élement de T, et on gagne. Dans le cas codim T = 1, comme T est
irréductible, on ne peut avoir $\hand{\eta}\in T$ que si Q=T ,\xout{xxxxxxxxxxxxx}
i.e. en termes géométriques \xout{xxxx} \inserted[typed]{(supposant k alg clos, ce qui est loisible)} pour tout $z\in X(\unsure{k})$\xout{xxx}l'espace tangent
de X en z (ou plutôt son image par $f'_{z}$ ) passe\xout{x} par x. Montrons que
ceci ne peut se présenter lorsque $\underline{X}^{\mathrm{sing}}\hand{$\longrightarrow$}T$ est gén étale, i.e.
lorsqu'on est sous les conditions de 8.8., sauf si f(X) = \hand{\{}x\hand{\}} donc X
de dimension 0. En effet, 8.7. c) (qui exprime essentiellement la
symétrie \xout{xxxxx} dans la relation entre X et sa \hand{``}duale\hand{''} \inserted{T}\hand{)} implique alors que
pour presque tout point $z\in X(k)$, f(z) est orthogonal à l'espace tangent
à T en un certain point,\xout{xx} donc (puisque T=Q) \inserted{orthogonal} à Q, d'où f(z)=x, d'où
f(X) =\hand{\{}x\hand{\}} . Cela achève la démonstration de nos remords, donc de 13.2.

\margin{80}

\clearpage
% Nouveau EGA V (cote 25): archive cote 25, site page 81 (PDF page 82).
% typescript. Draft compared: none.
% AI-transcribed by Claude Opus 5.5 (claude-opus-5-5, adaptive thinking, effort high), Message Batches API, batch ega25b (msgbatch_01A7n5YuXubUPNK1zsepi8RG), 2026-11-13; not yet checked by eye.
\scanpage{25/81}
\mspage{\hand{76}}

NB \quad Le raisonnement ne marche plus si on remplace x par une
sous-variété linéaire \unsure{$X$|$C$} de dim.$>0$ , et qu'on assujettit H à passer
par C: en effet, rien ne s'oppose (prenant par exemple dim C = r-2, le
plus grand possible pour que H puisse encore effectivement varier) que
T contienne la droite $C^{o}$ :\xout{(et pour faire un exemple où où ce fait se}
\xout{produit, en car 0 disons, on \unsure{part} d'une sous-variété T de} prendre par
exemple une quadrique \inserted{non singulière} dans $P$ , en caractéristique quelconque. Mais il
est possible que de tels phénomènes ne puissent non plus se produire \xout{xx}
pour des $\phi_n f$ , $n\geq 2$ ; on pourrait poser la question en remarque, au
nº 9 .

\uline{Remarques} 13.4. \quad a) Nous avons déjà signalé que l'hypothèse que k \xout{xx}
soit parfait est essentielle pour la validité de 13.2. Par contre, il
est plausible que l'hypothèse ``k infini'' est surabondante. \xout{En tous cas}
Nous n'avons pas essayé \xout{xxx} de faire un raisonnement ad hoc pour le cas
où k est fini, et signalons seulement que dans ce cas, l'application de
13.2. à la clôture algébrique de k et les raisonnements habituels
montrent qu'on peut trouver une extension finie k' de k telle que
\xout{xx} \inserted{pour} les points de $X_{k'}$ au dessus de x il existe des voisinages ouverts
polyélementaires relativement à k'.

b) Lorsque dans 13.2. on abandonne l'hypothèse que X est géom
irréductible, la conclusion bien entendu ne subsiste pas telle quelle,
(puisqu'un schéma algébrique poly-élementaire est géom irréd !). Elle
subsiste cependant sous la forme plus faible qu'on obtiendrait en
omettant dans la définition 13.1. le mot ``connexes'' : c'est ce que
montre la démonstration que nous avons donnée.
\inserted{c)}

\over{c}{d}) En procédant comme expliqué dans 7.1. on peut \inserted{donner} des variantes
de 13.2. dans le cas où on remplace le corps de base k par un préschéma
\margin{81}

\clearpage
% Nouveau EGA V (cote 25): archive cote 25, site page 82 (PDF page 83).
% typescript. Draft compared: none.
% AI-transcribed by Claude Opus 5.5 (claude-opus-5-5, adaptive thinking, effort high), Message Batches API, batch ega25b (msgbatch_01A7n5YuXubUPNK1zsepi8RG), 2026-11-13; not yet checked by eye.
\scanpage{25/82}
\mspage{\hand{77}}

de base général. Signalons la suivante (sans démonstration): Soit $f:X\to S$
un morphisme projectif et plat, \xout{\unsure{xxxSxxixixquxxxX}} à fibres géométriques
\xout{\unsure{géomé}} irréductibles et $(R_2)$ , $S'$ un sous-\inserted[typed]{pré}schéma de $X$ fini sur $S$,
$s\in S$, supposons que pour tout $x\in S'$ au dessus de $s$, $X_s$ soit lisse sur
$k(s)$ en $x$ . Alors il existe \unsure{un} voisinage ouvert $U$ de $s$, et un voisinage
ouvert $V$ de $S'|U$ dans $X|U$ , tels que $V\longrightarrow U$ soit poly-élémentaire.
Si $Y$ est un sous-préschéma fermé de $X$\inserted[typed]{ne rencontrant pas $S'$ et} tel que l'ensemble $Z$ des points
en lesquels $Y$ n'est pas lisse sur $S$ \xout{\unsure{aixxdesxfibresxde}} \deleted{soit tel que} \inserted{\unsure{satisfasse}}
dim $Z_s$ \hand{$\leq$} dim $X_s$ - 2 , alors on peut ci-dessus prendre $V$ contenu dans
$X$ - $Y$.

c) (A être éventuellement inclus dans l'énoncé de 13.2.). Soit
avec les notations de 13.2. $\emptyset$ une partie finie de $X$ formée de points
lisses de $X$, et supposons que $\emptyset$ soit contenu dans un \xout{xxx} ouvert affine
de $X$. Alors $\emptyset$ admet un systême fondamental de voisinages ouverts
polyélémentaires. \inserted[typed]{Evidemment on peut supposer $\emptyset$ formé de points \uline{fermés}.} La démonstration est essentiellement la même, sauf
qu'on formule 13.3. sous une forme légèrement différente, savoir: il
existe une sous-variété linéaire $C$ de $P^r$ de codimension $n$, ne rencon\unsure{t}%
trant pas $\emptyset\cup Y$ , et telle que pour tout $x_i\in\emptyset$\inserted{$^{\unsure{(\bar{k})|(k)}}$}, la sous-variété
linéaire $L_i$ de codimension $n$-1 engendrée par $C$ et $x_i$ ait les propriétés
a) , b) , c) de 13.3. Pour vérifier ce point, on note qu'il suffit de
vérifier que la $C$ \uline{générique} a les propriétés envisagées, or \inserted{pour une telle $C$} \deleted{alors}
\xout{\unsure{lxsxlx}}\hand{\unsure{,}} chaque $L_i$ est génér\unsure{ique} parmi les $L$ de codimension $n$ passant
par $x_i$ , de sorte qu'on peut appliquer 13.3. sous la forme initiale
( ou du moins sous la forme où nous l'avions prouvée, qui était: toute
$L_o$ \uline{assez générale} passant par $x$ a les propriétés a),b),c) ).

e) Une des raisons de l'intérêt de 13.2. est la structure``topolo-%
gique'' particulièrement simple des \xout{xxxxxx} \inserted[typed]{schémas alg.} élémentaires\inserted{\unsure{$/U$|$U$}}. Par exemple,

\mspage{\hand{82}}

\clearpage
% Nouveau EGA V (cote 25): archive cote 25, site page 83 (PDF page 84).
% typescript. Draft compared: none.
% AI-transcribed by Claude Opus 5.5 (claude-opus-5-5, adaptive thinking, effort high), Message Batches API, batch ega25b (msgbatch_01A7n5YuXubUPNK1zsepi8RG), 2026-11-13; not yet checked by eye.
\scanpage{25/83}
\mspage{\hand{\unsure{78 bis}}}

lorsque le corps de base est le corps des nombres complexes, et si $U^{an}$
désigne l'espace analytique associé à U, alors les groupes  d'homotopie
$\hand{\pi}_i(U^{an})$  sont nuls pour $i\hand{\neq}\over{\illegible{1}}{1}$, et \xout{xxxxx} $\hand{\pi}_1$  est une extension successive
de groupes libres\over{,}{;} \deleted{comme d'ailleurs} \inserted{Donc} $U^{an}$ \deleted{est connexe (U l'étant), il} est
\inserted[typed]{un} \xout{xxxxxxxxxxxxxx}\hand{``}espace $K(\hand{\pi}_1 ,1)$\hand{''}\hand{,} \uline{classifiant}  \uline{pour} $\hand{\pi}_1$ ;  de façon plus
précise, son revêtement universel  est  homéomorphe à $\underline{C}^n$  et à fortiori
est contractile, et ce revêtement est un``fibré principal universel'' de
groupe $\hand{\pi}_1$ .

\mspage{\hand{83}}

\clearpage
% Nouveau EGA V (cote 25): archive cote 25, site page 84 (PDF page 85).
% typescript. Draft compared: none.
% AI-transcribed by Claude Opus 5.5 (claude-opus-5-5, adaptive thinking, effort high), Message Batches API, batch ega25b (msgbatch_01A7n5YuXubUPNK1zsepi8RG), 2026-11-13; not yet checked by eye.
\scanpage{25/84}
\mspage{\hand{78}}

14. \uline{Projections coniques}.

NB Nous avons déjà utilisé les projections coniques à diverses
\margin{$\|$}reprises, notamment à la fin du nº8, formulation de 10.4. et autres, et
le sorite qui suit devrait venir sans doute plutôt au début du paragraphe,
et éventuellement dans le paragraphe auxiliaire ``grassma\unsure{n}nien''.
Soit $C = P(F)$ une sous-variété linéaire de $P(E)=P$, de dimension
relative $r-m-1$ \xout{dans} sur $S$\inserted[typed]{i.e. de codimension $m+1$ dans $P$}, donc $F$ est un quotient de $E$, loc libre de
rang $r-m$ , $F=E/G$ , où $G$ est loc libre de rang $m+1$. On a défini
par voie algébrique au Chap II un morphisme
\[ p_C : P - C = P(E) - P(E/G) \longrightarrow P(G) , \]
que nous allons interpréter par voie géométrique, et qui s 'appelera
(en vertu d e la description qui suit) la\uline{projection conique} de centre $C$.
(NB On suppose $r-m-1$ compris entre $-1$ et $r-1$, i.e. $m$ compris entre $0$
et $r$, sans plus). Pour ceci, commençons par interpréter \xout{Q(G)xx} $P(G)$
comme un sous-\inserted{pré}schéma\inserted{fermé} de $\mathrm{Grass}^m(P)= \mathrm{Grass}_{r-m+1}(P)$, grâce à l'homomorphisme
de foncteurs évidents
\[ \xout{Gxx}\ P(G) \longrightarrow \mathrm{Grass}_{r-m+1}(E) \]
otebnu en considérant, pour chaque quotient inversible $G/G'$ de $G$, le
Module localement libre de rang $(r-m)+1$ $E/G'$ de $E$ (et kif kif après
tout changement de base). L'homomorphisme de foncteurs précédent est un
\uline{monomorphisme}, et comme le premier est propre sur $S$, le deuxième séparé,
c'est une immersion fermée. \xout{\unsure{(}}Plus généralement, il y aurait lieu d'expli-%
citer les immersions fermées des grassmaniennes de $G$, i/e. de $P(G)$,
dans celles de \inserted{$E$ i.e. de} $P=P(E)$. \xout{Ixix} L'image\inserted[typed]{(au sens foncteurs)} du morphisme obtenu est formé des
sous-variétés linéaires \xout{dexM} $L$\inserted[typed]{de la dimension voulue} de $P$ qui \uline{contiennent} $C$ . Désignant
par $Q(C)$ cette image dans le cas qui nous occupe (i.e. pour les dimensi-%
ons précisées plus haut), et identifiant $P(G)$ à $Q(C)$, le morphisme
de projection conique
\[ p_C : P - C \longrightarrow Q(C) \subset \mathrm{Grass}^m(P) \]

\mspage{\hand{84}}

\clearpage
% Nouveau EGA V (cote 25): archive cote 25, site page 85 (PDF page 86).
% typescript. Draft compared: none.
% AI-transcribed by Claude Opus 5.5 (claude-opus-5-5, adaptive thinking, effort high), Message Batches API, batch ega25b (msgbatch_01A7n5YuXubUPNK1zsepi8RG), 2026-11-13; not yet checked by eye.
\scanpage{25/85}
\mspage{79}

n'est autre que celui qui associe à toute section de  P - C  l'unique
sous-variété linéaire  L de P, de codi\unsure{mmm}sion m, contenant à la fois C
et ladite section  (bien entendu, par ``contenir la section'' on entend\xout{x} qu\unsure{e}
% NOTE: the end of this line is cut by the right edge of the scan; only "qu" and part of the next letter are visible
la section se factorise par L).

Si maintenant on a un f: $X \longrightarrow P$,  il y a lieu de considérer le
composé
\[ X - f^{-1}(C) \longrightarrow\ P - C \ \longrightarrow\ Q(C) \]
qu'on pourra appeler projection conique de X, relativement à f et au cen\xout{x:}%
tre C\hand{,} \inserted{notée \unsure{$p^{f}_{C}$|$p^{\Sigma}$} ou $p^{X}_{C}$ ou simplement $p_{C}$.} \hand{(}On fera attention qu'elle n'est pas en général définie sur X tout
entier, d e façon précise elle l'est sss $f^{-1}(C)=\emptyset$ i.e. f(X) ne rencont re
pas le centre de projection C. \hand{\unsure{)}}

Nous allons donner une autre interprétation de ce morphisme, en
termes de constructions utilisées dans les n$^{\text{o}}$s antérieurs.  Pour ceci,
avec les notations introduites par ailleurs, considérons
\begin{diagramblock}
\begin{tikzcd}
X & \underline{X}^{(m)}_{\xout{\unsure{Q}}Q(C)} = \underline{X}^{(m)}x_{\mathrm{Grass}^{m}}Q(C) \arrow[l, "q"'] \arrow[d, "p\xout{\unsure{z}}"'] \\
 & Q(C) \quad .
\end{tikzcd}
\end{diagramblock}
Notons d'autre part que  q  induit un  isomorphisme
\[ q':q^{-1}(X-f^{-1}(C)) \xrightarrow{\ \sim\ } X - f^{-1}(C) \ , \]
et il est immédiat que  $p_{C}$n'est autre  que   $p'\ q'^{-1}$ , où p' est la
restriction de  p à  $q^{-1}(X-f^{-1}(\xout{\unsure{Q}}C))$.  On peut dire encore, utilisant q'
pour faire une identificatio\xout{\unsure{n}}n pure et simp\unsure{be|le}, que  $p_{C}$ est la restriction
du morphisme p  à  $X - f^{-1}(C) \ \subset\ \underline{X}^{(m)}_{Q(C)}$ . Pour cette raison, il est
commode de dénoter encore par  $p^{X}_{C}$ \inserted{ou $p_{C}$}, et d'appeler  (mettons)  \uline{projection}
\uline{conique}  \uline{étendue}  de X\hand{,} relativement à $f:X\rightarrow P$, de centre C, le morphisme

\mspage{85}

\clearpage
% Nouveau EGA V (cote 25): archive cote 25, site page 86 (PDF page 87).
% typescript. Draft compared: none.
% AI-transcribed by Claude Opus 5.5 (claude-opus-5-5, adaptive thinking, effort high), Message Batches API, batch ega25b (msgbatch_01A7n5YuXubUPNK1zsepi8RG), 2026-11-13; not yet checked by eye.
\scanpage{25/86}
\mspage{\hand{\unsure{80}}}

p$_{\unsure{\xout{C}}}$ précédent, \deleted{qu'il sera commode de noter également de la même façon}
\deleted{p$_C$ , pour indiquer sa dépendance de C.} De cette façon, les propriétés de
la projection conique  restreinte  sont ramenées \deleted{dans une large mesure} à
celles de la projection conique étendue, qui a été  étudiée systématique-%
ment par ailleurs, ou est sensée l'avoir été (cf n\textsuperscript{o} 10 et n\textsuperscript{o} 12). La
question\xout{s} principale\xout{s} qui se pose est, lorsque  S=Spec(k), quelles sont
les propriétés de la p\unsure{ro}jection conique  de X  quand on prend C générique
dans Grass$^{m+1}$(\unsure{E|P}) (ce qui exige  qu'on fasse un changement de base
k$\rightarrow$k(\hand{$\unsure{\eta|y}$}) , i.e.  C est \inserted{alors} en fait une sous-variété linéaire de X$_{k(\hand{\eta})}$ ); des
\unsure{ar}guments  standart, déjà répétés à satiété, permettront d'en conclure
les propriétés analogues pour les projections coniques correspondant à \xout{de}
des points de Grass$^{m+1}$(P)  appartenant à un ouvert non vide de ladite
grassmanienne, et enfih, lorsque k est infini, on conclut à l'existence
d'une(en fait une infinité\hand{)} \inserted{de} \xout{de} C  \uline{définie\xout{s} sur} k , I.e.  une sous-variété
linéaire de P lui-même (sans changement de corps de base), donnant lieu
à une projection conique ayant les propriétés en question. Il y aurait
\margin{|} lieu de bloquer ce genre d'explications générales avec celles de la même
\margin{|} eau données dans les n\textsuperscript{os} 4 et 7, et que nous avons d'ailleurs déjà
\margin{|} utilisées plus ou moins implicitement,  par exemple au n\textsuperscript{o}13 \xout{\unsure{(dém. de}}
\xout{\unsure{13.3}} (cf 13.4. c). Il y a lieu également, à ce propos, d'examiner les
propriétés relatives à un faisceau F sur X,  en prenant son image inverse
$F^{(m)}_{Q(C)}$  sur  $\underline{\underline{X}}^{(m)}_{Q(C)}$ . Il faut  d'ailleurs, dans la situation précise
envisagée ici, d es notations plus simples, je propose  $\widetilde{X(C)}$ et $\widetilde{F(C)}$ ,
et simplement $\widetilde{X}$ et $\widetilde{F}$  quand aucune confusion n\unsure{'}en résulte  (attention, le
F \deleted{précédent} \inserted{\unsure{ici}} n'est pas  celui du début du n\textsuperscript{o}).

Grosso modo, et lorsque  disons on suppose  f une immersion, les
propriétés de la projection conique générique sont assez différentes,

\mspage{\hand{86}}

\clearpage
% Nouveau EGA V (cote 25): archive cote 25, site page 87 (PDF page 88).
% typescript. Draft compared: none.
% AI-transcribed by Claude Opus 5.5 (claude-opus-5-5, adaptive thinking, effort high), Message Batches API, batch ega25b (msgbatch_01A7n5YuXubUPNK1zsepi8RG), 2026-11-13; not yet checked by eye.
\scanpage{25/87}
\mspage{81}

suivant qu'on suppose  dim X \hand{$\geqslant$} m , ou  dim X \hand{$\leqslant$} m , voire dim X \hand{$<$} m.
Par la suite, nous regardons la $C_{\hand{y}}\hand{$\subset$} P_{k(\hand{y})}$\hand{)} , correspondant au point
générique \hand{$y$} de  Grass$^{m+1}$ , et nous dispensons  de faire l'interprétation
des résultats obtenus en termes de  ``presque tous les points'' ...

Tout d'abord, nous avons déjà noté \xout{\illegible{1}}dans 5.3. (et rattrapage du
cas général dans nº 12)  que $C_{\hand{y}}$  coupe ``régulièrement'' $X_{\hand{y}}$ , de façon
plus pré\xout{\unsure{s}}cise et plus générale, pour tout F quasi-cohérent sur X,
la section $\phi_{\hand{y}}^{\hand{(m+1)}}$ du \inserted[typed]{Module loc libre}\xout{fibré vectoriel} de rang  m+1  sur  $X_{k(\hand{y})}$  dont le
schéma des zéros est \over{C}{$C$}$_{\hand{y}}$ ; est F-régulière. En vertu de 10.2. , cela
implique par exemple que le morphisme  $\widetilde{X}(\hand{C_y})\to X_{k(\hand{y})}$ identifie $\widetilde{X}(\hand{C_y})$ au préschéma
déduit de $X_{k(\hand{y})}$ par éclatement de $f^{-1}(C_{\hand{y}})= X\times_P C_{\hand{y}}$ ;  \xout{\illegible{2}} dans le cas
où  dim f(X) \hand{$\leqslant$} m ,  on  aura même  $f^{-1}(C_{\hand{y}}) = \emptyset$ , et  par suite
$\widetilde{X}(C_{\hand{y}}) \xrightarrow{\sim} X_{k(\hand{y})}$  est un \uline{isomorphisme} , (et en fait, la \hand{\unsure{)}}projection conique
restreinte  est alors définie sur X tout entier à priori). Il se pose
ensuite la question de la dimension des fibres de $p_{C_{\hand{y}}}$ : $\widetilde{X}(C_{\hand{y}}) \longrightarrow Q(C_{\hand{y}})$,
et de la platitude de ce morphisme. On trouve:
\uline{Proposition} 14.1.  \xout{a) Si dim X} Supposons X irréductible,\over{\unsure{e}}{\unsure{ou}}plus
généralement, que  pour toute composante irréductible $X_i$ de X, la fibre
de $X_i$ en le point $f(x_i)$  ($x_i$ = point générique de $X_i$) ait une dimension d
(indépendante de i), ce qui a lieu par exemple \inserted{(avec d=0)} si f:X$\to$P est quasi-fini\hand{,} \margin{\deleted{avec d=0}}
Alors :     a)  Si  dim f(X) \hand{$>$} m, alors la dimension  des fibres de
$p_C$  : $\widetilde{X}(C_{\hand{y}}) \longrightarrow \over{Q}{Q}(C_{\hand{y}})$  sont toutes égales à  dim X - m .   b)  Si dim X
$\leqslant$ m, \inserted{et si les fibres non vides des $X_i$ sur P sont \unsure{tts} de dimen d, alors} les fibres de $p_C$ \inserted{non vides} \inserted{$\colon X_{k(y)}\to Q(C_y)$} sont toutes de dimension \over{0}{d}, \deleted{i.e. $p_C$ est quasi-%
fini} (donc $p_C$ est fini \inserted{\unsure{resp. quasi-fini}} \deleted{si X est propre sur \unsure{k$(y)$} et} \inserted{$f\colon X\to P$ est fini \unsure{resp. quasi-fini}}].
\margin{\unsure{resp. quasi-fini}}

On \deleted{se ramène aussitôt, en décomposant X en ses composantes irréd.,}
\deleted{au cas où X est irréductible}. Dans le cas a), \xout{si il y avait une fibre}
\xout{correspondant à un} \xout{xxxxxxxxxxxx} on sait déjà (je l'espère !) que pour \uline{tout}  point \hand{$y$} de

\margin{87}

\clearpage
% Nouveau EGA V (cote 25): archive cote 25, site page 88 (PDF page 89).
% typescript. Draft compared: none.
% AI-transcribed by Claude Opus 5.5 (claude-opus-5-5, adaptive thinking, effort high), Message Batches API, batch ega25b (msgbatch_01A7n5YuXubUPNK1zsepi8RG), 2026-11-13; not yet checked by eye.
\scanpage{25/88}
\mspage{\hand{82}}

Grass$^{m}$(P), la dimension de $\underline{X}^{(m)}_{\hand{\xi}}$ est au moins égale à dim X-m , il en est
% NOTE: the right edge of the scan cuts the ends of the first two lines ("est", "l'inéga")
en particulier ainsi lorsque $\hand{\xi}$ provient d'un point de Q(C$_{\hand{y}}$). Pour l'inéga
lité en sens inverse, notons que (nous plaçant sur le corps k' = k($\hand{\xi}$))
comme $C_{\hand{y}\,k'} \subset L_{\hand{\xi}}$ est \xout{de} un hyperplan de $L_{\hand{\xi}}$ , si la dimension de
$\underline{X}^{(m)}_{\hand{\xi}} = X\times_{P}L_{\hand{\xi}}$ était \over{>}{$\geqslant$} dim X - m + 1 , alors celle de $\underline{X}^{(m+1)}_{\hand{y}}$ =
$X\times_{P}C_{\hand{y}}$ serait \over{>}{$\geqslant$} dim X - m \xout{+}, (puisque le changement de base k($\hand{y}$)$\to$k'
trans\xout{g}forme ce dernier préschéma en $(X\times_{P}L_{\hand{\xi}})\times_{L_{\hand{\xi}}}(C_{\hand{y}\,k'})$ ). Or on a au
contraire \xout{dimxXx$_{P}$C} dim $\underline{X}^{(m+1)}$ = dim X - m -1 \xout{(avecxlaxconventionxque}
\deleted{(du moins en supposant dim f(X) \hand{$>$} m )}, en vertu du n°2 (revu au n°10).
Le \inserted{cas b)} \deleted{cas où dim f(X)=m est justiciable également du cas b), qui} se traite
de façon analogue: si on avait dim $X\times_{P}L$ \hand{$\geqslant$} d\hand{+1}, \xout{alorsxxxlimagexde} ou
ce qui revient au même \xout{f(X)xxf$_{k(\ )}$xxxx} $f_{k'}(X_{k'}) \cap L_{\hand{\xi}}$ de dimension \hand{$\geqslant$}\over{0}{1},
alors on aurait par le même argument que plus \xout{haut} haut que
$\underline{X}^{(m+1)}_{\hand{y}} \neq \emptyset$ , contrairement à ce qu'on avait signalé avant 14.1.

\uline{Corollaire} 14.2. Supposons X de dimension m et f:X$\to$P fini\inserted{(resp quasi-fini)}, alors
le morphisme $p_{C_{\hand{y}}} : X_{k(\hand{y})} \longrightarrow Q(C_{\hand{y}})$ est fini surjectif\hand{(}resp. quasi-fini
dominant).

En effet, ce morphisme est \inserted{quasi-}fini, et comme dim $X_{\hand{k(y)}}$ = dim Q (C$_{\hand{y}}$), il est
dominant; \xout{enfinxdoncxsurjectif} si f est fini, $p_{C_{\hand{y}}}$ est de plus fini,
\xout{doncxfex} car propre, donc surjectif étant dominant.
\uline{Corollaire} 14.3. Sous les conditions de 14.1. a), si X est Cohen-Macau-%
lay, alors le morphisme $p_{C} : \widetilde{X}(C_{\hand{y}}) \longrightarrow Q(C_{\hand{y}})$ est un morphisme de
Cohen-Macaulay, et à fortiori \uline{plat}.

Pour la démonstration, cf remarque en haut de la page 21, avant 5.6.\inserted{qui donne un résultat plus fort, (à inclure dans 14.3), savoir que les $\unsure{\varphi^{-1}_{\xi}}$ pour $\xi\in Q(C_{y})$ sont \unsure{$F_{\xi}$}-réguliers.}
Ce \over{c}{C}orollaire devrait être modulé, mais pour simplifier on pourra
supposer f quasi-fini: si F est un Module de Cohen-Macaulay sur X, et si
pour to\uline{u}te composante irréductible Z de Supp F, on a dim Z \over{>}{$\geqslant$} m, alors
$\widetilde{F}(C_{\hand{y}})$ est CM \hand{(}et à fortiori plat\hand{)} relativement à Q(C$_{\hand{y}}$).

On notera qu'on ne peut pas remplacer,\xout{xxxx} pour obtenir la

\margin{88}

\clearpage
% Nouveau EGA V (cote 25): archive cote 25, site page 89 (PDF page 90).
% typescript. Draft compared: none.
% AI-transcribed by Claude Opus 5.5 (claude-opus-5-5, adaptive thinking, effort high), Message Batches API, batch ega25b (msgbatch_01A7n5YuXubUPNK1zsepi8RG), 2026-11-13; not yet checked by eye.
\scanpage{25/89}
% NOTE: the top right corner is cut off at the edge of the scan; only the lower part of a number is visible there (probably the page number)
seule conclusion $p_C$ plat, l'hypothèse CM sur X par une simple hypothèse
dimensionelle. Supposons par exemple f une immersion et X irr de dim m,
de sorte que $p_C$ est qu fini , et comme $X_{k\hand{(y)}}$ et $Q(C_{\hand{y}})$ sont irréd de
même dimension, \xout{le premier CM} et le deuxième régulier, $p_C$ ne peut être
plat que si $X_{k\hand{(y)}}$ est CM.

Des propriétés plus fines sont les propriétés différentielles de
la projection conique, notamment pour X lisse sur k et $f:X\longrightarrow P$ non rami-%
\inserted[typed]{étudiées au nº12.}
fié, \xout{Dans le cas où les composantes} Rappelons que \xout{la fibre générique de}
en dehors d'une partie Z de codimension 1 de $Q(C\ )$, le morphisme
$p_{C_{\hand{y}}}$ \inserted[typed]{sur $X(C_{\hand{y}})$} est \uline{lisse}, Et une analyse plus détaillée (esquissée au nº12) montre
que si \uline{les dimensions des composantes} de X sont $\geq$ m, alors\hand{)}
(ou montrerait, si nous ne la faisons pas) \xout{que} en dehors d'une partie $Z'\subset Z$
de $Q(C_{\hand{y}}\ )$ de codimension $\geq 2$, les fibres $p_C^{-1}(\hand{\zeta}) = \underline{X}^{(m)}_{\hand{\zeta}}$ ne peuvent
avoir au pis que des points \xout{doub} singuliers ordinaires (au s ens
géométrique), et en fait (si f est une immersion et X est géom irréd),
\xout{juste} au plus \uline{un} tel point, ce dernier \inserted[typed]{étant alors} \xout{\unsure{ayant}} nécessairement \xout{\illegible{1}} .
rationnel sur $k(\hand{\zeta})$, - ces assertions étant valables tout au moins si
k est de car 0, \uline{ou} à condition\xout{s} de remplacer f par un $\phi_n f$ $(n\geq 2)$ comme
au nº 9. Il y a lieu également de donner des propriétés différentielles
de $p_{C_{\hand{y}}}$ dans le cas où dim X $\leqslant$ m , et où par suite $p_{C_{\hand{y}}}$ est défini sur $X_{k\hand{(y)}}$.
Je me borne à indiquer les propriétés suivantes, dont la démonstration
\inserted{doit} être facile et \inserted{est} laissée à Dieudonné:

\uline{Proposition} 14.4. Supposons $f:X\longrightarrow P$ non ramifié, et dim X $\leqslant$ m.
Soit T un sous-préschéma fini de X, Alors :\xout{la restr}
\margin{\unsure{X?}} a) \inserted{Si f est une immersion,} La restriction de $p_C$ à $T_{k\hand{(y)}}$ est \deleted{un monomorphisme,} \inserted{radiciel i.e.\ ``géométriquement injectif''.} Si de
plus Y est une partie fermée de X de dimension $\leqslant$ m-1, on a
$p_{C_{\hand{y}}}^{-1}(p_{C_{\hand{y}}}(Y_{k\hand{(y)}})) \cap T_{k\hand{(y)}} = \emptyset$ .
b) Si X est lisse en les points de T, alors $p_{C_{\hand{y}}}$ est non ramifié
en tous point \xout{de} $T_{k\hand{(y)}}$, \xout{au dessus de tout point de T} \inserted{et \unsure{même|aussi} en les pts de $p_{C_y}^{-1}\,p_{C_y}(T_{k(y)})$.}
\mspage{\hand{83}}

\clearpage
% Nouveau EGA V (cote 25): archive cote 25, site page 90 (PDF page 91).
% typescript. Draft compared: none.
% AI-transcribed by Claude Opus 5.5 (claude-opus-5-5, adaptive thinking, effort high), Message Batches API, batch ega25b (msgbatch_01A7n5YuXubUPNK1zsepi8RG), 2026-11-13; not yet checked by eye.
\scanpage{25/90}
\mspage{\unsure{87}}

\uline{Proposition} 14.5. Supposons dim X $\leqslant$ m-1, \xout{Soit Y l'image fermée}
\xout{\illegible{1}} f:X $\to$ P une immersion, \xout{Soit} enfin X séparable sur k. Soit Y\hand{$_y$} l'image
schématique de $X_{k(\hand{y})}$ dans $Q(C_{\hand{y}})$. Alors le morphisme induit
$p'_{C_y} : X_{k(y)} \to Y_y$ est \uline{birationnel}\hand{, et pour tout point x de $X_{k(y)}$ au dessus}
\hand{d'un pt fermé de X, $p'_{C_y}$ est \uline{étale} en x et \unsure{$\overline{m}$|nr} en les points de $p'^{-1}_{C_y}p'_{C_y}(x)$.}
Signalons la conséquence suivante:

\uline{Corollaire} 14.6. Soit X un schéma algébrique projectif \inserted[typed]{irréductible} \xout{irréductible}
\uline{et séparable}
de dimension n sur un corps infini k. Alors \deleted{X est birationellement} \inserted{$\exists$ un plongement birationnel de X}
\deleted{isomorphe à} \inserted{\deleted{\unsure{dans}} sur} une hypersurface \deleted{dans} P$^{n+1}$ .

On se gardera de croire \deleted{que}, même si X est un sous-préschéma fermé
lisse géom irréd de P, de dimension m-1 = n , \inserted[typed]{que} la projection conique p$_{C_{\hand{y}}}$
soit \inserted[typed]{nécessairement} une \uline{immersion} . En effet, lorsque k est infini, cela impliquerait
qu'il existe une C rationnelle sur k ayant la même propriété, donc que X
est isomorphe à une hypersurface non singulière de \xout{PxxxQxx} P$^{n+1}$. Or
même pour n=1, (\inserted{dans X}courbe algébrique \inserted[typed]{projective} lisse \xout{\unsure{sur}} et connexe sur un corps alg.
clos), il est facile de construire des exemples où X ne peut s'immerger
dans un P$^2$ . De même, dans 14.4. , on se gardera de confondre l'énon-%
cé donné avec l'assertion (en général fausse) que p$_C$ soit lui-même un
monomorphisme (contre-exemple précédent\xout{\unsure{, xx que}}, ou contre-exemple enco-%
re plus évident si X est lisse de dimension m), ou que p$_C$ soit non
ramifié. Pour ce dernier point, on prendra pour s'en convaincre X
sous-\xout{pré}schéma fermé lisse \xout{de dim} irréd de dim m (sur k alg clos disons
si \xout{XxxxxQx} on avait un X $\to$ Q $\simeq$ P$^m$ non ramifié, il serait étale pour
raisons de dimensions, or on peut montrer (cf Chap VIII !) que cela impli-%
que X $\xrightarrow{\sim}$ P$^m$ (P$^m$ étant "simplement connexe). \hand{--} La signification géométri-%
que intuitive de 14.4. est que\xout{x pourxxxvariablexdans} l'ensemble de
ramification de p$_{C_{\hand{y}}}$ est ``variable'' \inserted{sur k}, plus précisément l'ensemble de
ramification de p$_{C_{\hand{y}}}$ pour \hand{y} variable dans \inserted{un ouvert de} Grass$^{m+1}$(\hand{$\bar{k}$}) varie dans X(\hand{$\bar{k}$}), et
n'\uline{admet aucun} \deleted{point fermé de X comme} \uline{``point fixe''} ....

\mspage{90}

\clearpage
% Nouveau EGA V (cote 25): archive cote 25, site page 91 (PDF page 92).
% typescript. Draft compared: none.
% AI-transcribed by Claude Opus 5.5 (claude-opus-5-5, adaptive thinking, effort high), Message Batches API, batch ega25b (msgbatch_01A7n5YuXubUPNK1zsepi8RG), 2026-11-13; not yet checked by eye.
\scanpage{25/91}
\mspage{85}

Bien entendu, pour justifier dans le cas \xout{\unsure{actuel}} du présent n° le passage
de \hand{$\eta$} générique aux points voisins de $\mathrm{Grass}^{m+1}(P)$, et pouvoir également
le cas échéant reprendre à notre compte les explications générales de
7.1. , il y a lieu de considérer les diagrammes
\begin{diagramblock}\begin{tikzcd}
X \arrow[d] & \widetilde{X}(C) \arrow[l] \arrow[d] \\
S & Q(C) \arrow[l]
\end{tikzcd}\end{diagramblock}
obtenus à l'aide des divers $C \in \mathrm{Grass}^{m+1}(S)$, et \inserted{plus généralement} ceux obtenus de même
après changement de base $T \to S$ pour des points \hand{$\zeta$} \hand{$\in$} $\mathrm{Grass}^{m+1}($\over{P}{\unsure{T}}$)$ :
\begin{diagramblock}\begin{tikzcd}
X_T \arrow[d] & \widetilde{X}(C_{\hand{$\zeta$}}) = \widetilde{X_T}(C_{\hand{$\zeta$}}) \arrow[l] \arrow[d] \\
T & Q(C_{\hand{$\zeta$}}) \arrow[l]
\end{tikzcd}\end{diagramblock}
comme déduits, par changement\hand{s} de base \hand{$\zeta$} : $T \longrightarrow \mathrm{Grass}^{m+1}$\deleted{$(P)$} $= \underline{T}$ ,
du diagramme \uline{universel} (relatif au point canonique de $\mathrm{Grass}^{m+1}$ dans $\underline{T}$ ):
\begin{diagramblock}\begin{tikzcd}
X_{\underline{T}} \arrow[d] & \widetilde{X}(\underline{C}) \arrow[l] \arrow[d] \\
\underline{T} & Q(\underline{C}) \arrow[l]
\end{tikzcd}\end{diagramblock}
où $\underline{C}$ est la sous-variété linéaire canonique de $P_{\underline{T}}$ . Alors ci-dessus
$\widetilde{X}(C_{\hand{\unsure{$\eta$|$\zeta$}}}) \longrightarrow Q(C_{\hand{\unsure{$\eta$|$\zeta$}}})$ n'est autre que le morphisme des fibres génériques pour
le $\underline{T}$-morphisme $\widetilde{X}(\underline{C}) \longrightarrow Q(\underline{C})$ de ce dernier diagramme, et toute propri\'e-%
té constructible pour \xout{xxxx} le morphisme des fibres génériques implique
la propriété pour les fibres voisines. - Du point de vue notationel,
$Q$ doit être envisagé (et même introduit) comme le nom du morphisme \inserted{de foncteurs}
naturel $\mathrm{Grass}^{m+1}(P) \longrightarrow$ \xout{Hilb} Sous-préschémas $(\mathrm{Grass}^{m}(P))$ .

15. \uline{Axiomatisation de certains des résultats précédents}.

Je pense surtout aux résultats des n°s 2 à 8, qui pour la plupart

\mspage{91}

\clearpage
% Nouveau EGA V (cote 25): archive cote 25, site page 92 (PDF page 93).
% typescript. Draft compared: none.
% AI-transcribed by Claude Opus 5.5 (claude-opus-5-5, adaptive thinking, effort high), Message Batches API, batch ega25b (msgbatch_01A7n5YuXubUPNK1zsepi8RG), 2026-11-13; not yet checked by eye.
\scanpage{25/92}
\mspage{86}

tout au moins sont valables sous d es conditions plus générales que \inserted{pour} la
famille des hyperplans (ou des hypersurfaces de degré donné) de l'espace
projectif. Il me semblerait opportun d'adopter un tel point de vue
axiomatique d\over{an}{\unsure{è}}s les premiers n$^{\text{o}}$s de la rédaction en forme. Je n'ai pas
tiré au clair pour le moment s'il y a bien moyen de donner une généralisa%
tion dans ce sens de Bertini-Zariski (donc des résultats des n$^{\text{o}}$s 4 et 6),
et j'ai écrit aux compétences (Serre, Zariski) pour leur demander s'ils
avaient connaissance d'une telle extension; j'ai bien l'impression en
effet que des hypothèses de nature différentielle simple, du genre de cel
% NOTE: the end of the line runs off the right edge of the scan after "cel"
les envisagées ci-dessous, doivent suffire à impliquer Bertini-Zariski.
Si les compétences ne peuvent nous informer de façon satisfaisante, il
faudrait essayer de tirer la chose au clair par nos propres moyens.

On part d'un diagramme \inserted{commutatif} de morphismes de présentation finie

\margin{(D)}
\begin{diagramblock}
\begin{tikzcd}
P & \underline{P} \arrow[l] \arrow[d] \\
S \arrow[u, dash] & G \arrow[l]
\end{tikzcd}
\end{diagramblock}

(dans le cas d'application principal, P est un fibré projectif, G une
grasm\unsure{a}nienne dudit, $\underline{P}$ le préschéma d'incidence). Dans lescas lesplus
importants, le morphisme correspondant $\underline{P} \longrightarrow P\times_S G$ sera une immersion
fermée, et on considère G comme \xout{définissant une} \inserted{schéma des paramètres d'une} familled e sous-
préschémas fermés des fibres de P sur S, d e façon précise si \hand{$\zeta$} \hand{$\in$} G,
$\underline{P}_{\hand{\zeta}}$ est un sous-préschéma fermé de $P_{s\ k(\hand{\zeta})}$ , où s est le point de S en
dessous de \hand{$\zeta$} .(D'ailleurs, pour la plupart des énoncés dans ce contexte,
on aura sans doute S=Spec(k)\hand{)}. Dans le cas général, on peut encore con-%
sidérer G comme schéma des paramètres d'une famille de préschémas
au dessus des fibres de P sur S, à \hand{$\zeta$} correspondant $\underline{P}_{\hand{\zeta}}$ au dessus de
$P_{s\ k(\hand{\zeta})}$ . Bien entendu, au lieu de prendre pour \hand{$\zeta$} un point \xout{de} (absolu)

\margin{92}

\clearpage
% Nouveau EGA V (cote 25): archive cote 25, site page 93 (PDF page 94).
% typescript. Draft compared: none.
% AI-transcribed by Claude Opus 5.5 (claude-opus-5-5, adaptive thinking, effort high), Message Batches API, batch ega25b (msgbatch_01A7n5YuXubUPNK1zsepi8RG), 2026-11-13; not yet checked by eye.
\scanpage{25/93}
\mspage{87}
de G, on peut également prendre un point à valeurs dans un \inserted[typed]{S-}préschéma T,
on obtient alors $\underline{P}_{\hand{\unsure{\xi|\zeta}}} \longrightarrow \unsure{P}_T$ (T-morphisme, qui est une immersion ferm\unsure{ée|é}
dans le cas envisagé d'abord).

Si $f:X \longrightarrow P$ est un morphisme, on pose $\underline{X} = X\times_P\underline{P}$ , et on obtient
un diagramme de même type que le carré précédent
\begin{diagramblock}\begin{tikzcd}
X \arrow[d] & \underline{X} \arrow[l, dash] \arrow[d] \\
S & G \arrow[l]
\end{tikzcd}\end{diagramblock}
.

Il est alors évident que toutes les questions étudiées dans les n$^{\text{os}}$ 2 à 8
gardent un sens dans le contexte général qu'on vient de poser, et il y a
lieu de dégager les conditions axiomatiques qui assurent les conclusions
tirées dans lesdits n$^{\text{os}}$ .

Nous allons supposer $\underline{P}$ \uline{et} G \uline{plats sur} S, G étant \uline{à fibres}
\inserted[typed]{\xout{de dimen}}
\uline{géométriquement irréductibles} (pour pourvoir en considérer les points
génériques !), \uline{de dimension} N, \uline{le morphisme} $\underline{P} \longrightarrow P$ \uline{est supposé lisse}
\uline{à fibres géométrique\xout{x}ment irréductibles de dimension N-m} . Donc le
morphisme $\underline{X} \longrightarrow X$ aura les mêmes propriétés. Toutes les propriétés
envisagées \inserted{ici et plus tard,} sont stables par changment de base sur S, et s'applique\xout{x}ront
en particulier aux fibres.

Supposons d'abord S=Spec(k). Soit Z une partie fermée de X de
dimension d, alors son image inverse $\underline{Z}$ dans $\underline{X}$ est une partie fermée de
dimension $d + (N-m) = N + d-m$ . Si $d < m$, $\underline{Z}$ est de dimension $< N$,
donc $\underline{Z} \longrightarrow G$ ne peut être dominant, donc si \hand{$\eta$} est le point générique de
G, on a $\underline{Z}_{\hand{\eta}} = \emptyset$ ; en fait, ce raisonnement montre même (en remplaçant
Z par $\overline{f(Z)}$ ) que si dim f(Z) \hand{$<$} m, alors $\underline{Z}_{\hand{\eta}} = \emptyset$ . On veut une condi-%
tion \inserted{sur (D)} assurant que si dim f(Z) \over{\unsure{$<$}}{$\geq$} m, alors $\underline{Z}_{\hand{\eta}} \neq \emptyset$ ; \deleted{et de façon plus
précise que dim $\underline{Z}$ = dim Z - m . Pour la propriété $\underline{Z} \neq \emptyset$ ,} il

\mspage{93}

\clearpage
% Nouveau EGA V (cote 25): archive cote 25, site page 94 (PDF page 95).
% typescript. Draft compared: none.
% AI-transcribed by Claude Opus 5.5 (claude-opus-5-5, adaptive thinking, effort high), Message Batches API, batch ega25b (msgbatch_01A7n5YuXubUPNK1zsepi8RG), 2026-11-13; not yet checked by eye.
\scanpage{25/94}
\mspage{88}
semble qu'elle doive faire un axiome primitif de la situation (dans la
situation du n°2.2. , elle \unsure{résultait} d'un argument global assez \unsure{spécifil|spécifiq} ):
\uline{pour toute partie fermée Zrréductible} Z \uline{de} P d\uline{e dimension} m, $\underline{Z}_{\hand{\eta}} \neq \emptyset$ \xout{xxx}
\xout{pour tout La précision sur la dimension sera alors conséquence facile}
\xout{de l'hypothèse suivante, de nature différentielle: P est lisse sur S,}
\xout{de dimension $\underline{P}$ $P \times_S T$ une immersion, et pour tout} Reprenant alors la
partie fermée Z de X,\xout{de di} telle que dim f(Z) $\hand{\geqslant}$ m, on voit que $\underline{Z}_{\hand{\eta}} \rightarrow G$ est
dominant, et par suite $\underline{Z}_{\hand{\eta}}$ est de dimension dim $\underline{Z}$ - dim G = dim Z-m.

Ces propriétés permettent de développer dans le contexte présent
les résultats correspondant à 2.1. à 2.11. Il faut une condition sur (D)
assurant la validité de 2.12. i.e. que si X est lisse, alors $\underline{X}$ l'est
\unsure{éggalement|également}, en supposant f:X $\longrightarrow$ P non ramifié. N$_{\mathrm{o}}$us supposerons
\unsure{maàntenant|maintenant} P \uline{lisse sur} k,\inserted{$\underline{P} \rightarrow P \times_k G$ quasi-fini,} et la condition suivante vérifiée\inserted[typed]{(où nous supposons k alg clos)}: pour
tout $x \in P(k)$, et tout sous-espace vectoriel V \inserted[typed]{de dimension $n \geqslant m$} de l'espace tangent
$T_x$(P) de P en x, considérons l'ensemble\inserted{E(x,V)} \xout{Ex} des $\hand{\xi} \in G(k)$ tels que
$\underline{P}_{\hand{\xi}}$ ait un point\inserted{\unsure{z}} au dessus de x \xout{en lequel on ait bien $\underline{P}$ P soit ramifié}
ne satisfaisant \uline{pas} l'ensemble de conditions suivantes: $\underline{P}_{\hand{\xi}}$ est lisse en
z, \xout{le morphisme $\underline{P}$ P est non ramifié en z, et l'image de l'espace}
l'application tangente à $\underline{P}_{\hand{\xi}} \rightarrow P$ en z : $T_z(\underline{P}\ ) \rightarrow T_x(P)$ est injective
(i.e. $\underline{P}_{\hand{\xi}} \rightarrow P$ non ramifié en z) et son image est ``transversale'' à V
i.e. sa somme avec V est $T_x(P)$. \uline{Alors} E(x,V) (dont nous savons que
c'est la trace d'un ensemble \xout{constructible} \inserted[typed]{\deleted{fermé}}\inserted{\unsure{constructible}} bien déterminé de G sur G(k))
\uline{est de dimension} $\hand{\leqslant}$ N-n\unsure{\over{+}{-}}1 . - Moyennant cette condition, \xout{xx} l'appli\unsure{:}
cation du critère jacobien et un comptage de dimensions montre que
l'ensemble\inserted{fermé E} des poihts\inserted{x} de $\underline{X}$ tels que $\underline{X} \longrightarrow G$ soit non lisse en x, \uline{ou}
$\underline{P} \longrightarrow$ \unsure{\over{G}{G}} non lisse en $\underline{f}(x)$, ou $\underline{P} \longrightarrow P$ ramifié en $\underline{f}(x)$, est de
dimension $\hand{\leqslant}$ n + (N-n\unsure{\over{+}{-}}1) = N-1 , \hand{(}X étant lisse partout de dimension n\hand{)}.

\margin{94}

\clearpage
% Nouveau EGA V (cote 25): archive cote 25, site page 95 (PDF page 96).
% typescript. Draft compared: none.
% AI-transcribed by Claude Opus 5.5 (claude-opus-5-5, adaptive thinking, effort high), Message Batches API, batch ega25b (msgbatch_01A7n5YuXubUPNK1zsepi8RG), 2026-11-13; not yet checked by eye.
\scanpage{25/95}
\mspage{89}

Donc dim E \hand{$<$} N = dim $P$ , donc $E_{\hand{$y$}} = \emptyset$ , et à fortiori $X_{\hand{$y$}}$ est lisse\hand{.}\inserted{\unsure{sur $k(y)$}}
Ce point acquis, on conclut aussitôt à la validité des variantes
évidentes de 2.12. à 2.18. dans le contexte actuel.

Le passage au nº4 d'une section générique à une section générale
et les développements du nº5, sont évidemment valables dans le contexte
actuel (mais sont à tel point des tautologies ou redites\xout{\illegible{1}} des
par.8,9,12 \xout{)} qu'on hésite à les énoncer en forme). De même les
développements de 7.1. , valables en tous cas si k est \xout{\unsure{infini}} alg clos
(et même si k est simplement infini, si on suppose G rationelle\xout{s} sur k),
et les cas particuliers 7.2. , 7.3. ; quant au\xout{x} resultat\xout{s} 7.4. \xout{\unsure{et 7.5.}}
\xout{\illegible{1}}c'est évidemment \inserted[typed]{une application} de nature spéciale \inserted[typed]{de} \xout{à} la situation des sections
hyperplanes . Comme j'ai dit, les nºs 3 et 6 sont suspendus à l'exten-%
sion du th. de Zariski.

Il resterait à étendre également les résultats du nº8 (respris au
nº 12), qui prendront peut-être comme-ça une allure plus sympathique.
Je te conseille même de commencer par rédiger ces résultats dans ce cadre,
en essayant de pousser aussi loin que possible dans cette voie. J'ai
l'impression qu'on doit pouvoir récupérer au moins tout ce qui n'est pas
conséquence directe de 8.7. c) (et encore pourrait-on essayer
d'abstraire des conditions axiomatiques faisant marcher une variante de
8.7. c). Je me borne à ces recommandations, mais suis prêt à y revenir
avec plus de détails si tu as des difficultés particulières.

\mspage{95}

\clearpage
% Nouveau EGA V (cote 25): archive cote 25, site page 96 (PDF page 97).
% typescript. Draft compared: none.
% AI-transcribed by Claude Opus 5.5 (claude-opus-5-5, adaptive thinking, effort high), Message Batches API, batch ega25b (msgbatch_01A7n5YuXubUPNK1zsepi8RG), 2026-11-13; not yet checked by eye.
\scanpage{25/96}
% NOTE: the page carries no typed text; the upper part is blank, and the only writing is the handwritten page number at the foot of the page.
\mspage{96}

\clearpage
% Nouveau EGA V (cote 25): archive cote 25, site page 97 (PDF page 98).
% typescript. Draft compared: none.
% AI-transcribed by Claude Opus 5.5 (claude-opus-5-5, adaptive thinking, effort high), Message Batches API, batch ega25b (msgbatch_01A7n5YuXubUPNK1zsepi8RG), 2026-11-13; not yet checked by eye.
\scanpage{25/97}
% NOTE: the whole typed text of this page is extremely faint (pale carbon or show-through); most of it can only be guessed, and many readings below are uncertain
% NOTE: two or three faint typed lines at the head of the page, above the heading, cannot be read
\illegible{8}

\illegible{8}

\uline{\unsure{§ 2)} \unsure{Compléments sur les schémas projectifs} \illegible{3}} \unsure{et}

\uline{\unsure{Pic(X). Les} \illegible{3}}

\uline{\unsure{1. Détermination des faisceaux inversibles sur les fibrés projectifs.}}

\uline{\unsure{a) Fibrés en espaces projectifs associés à un fibré projectif.}}

\uline{\unsure{Proposition 1.1.}} \unsure{--- Soient S un préschéma, E un Module loc. libre sur S partout}
\unsure{(rang $\geq 2$ en tout point).}
\xout{\illegible{3}}\unsure{, $X = P(\underline{E})$ le fibré projectif qu'il définit. Alors}
\unsure{pour tout Module inversible $\underline{L}$ sur X, on peut trouver une famille}
\unsure{finie $(S_i)_{i \in I}$ d'ouverts disjoints de S recouvrant S, indexés par I, et}
\unsure{pour tout $i \in I$ un Module inversible $\underline{M}_i$ sur $S_i$, et entiers $n_i$ tels que}
\unsure{la rest.} \xout{\illegible{4}} \unsure{la rest.}
\unsure{de $\underline{L}$ à $f^{-1}(S_i)$ soit isomorphe à celle de $f^*(\underline{M}_i)(n_i)$.}
\unsure{De plus, la famille des $S_i$} \xout{\illegible{2}} \unsure{est uniquement déter-}%
\unsure{minée par ces conditions, ainsi que $\underline{M}_i$ à isomorphisme unique près.}
\unsure{Corollaire 1.2.} \unsure{--- Si on ne fait pas d'hypothèse sur le rang de E, S se}
\unsure{décompose canoniquement en somme de trois ouverts $S'$, $S''$, $S'''$ tels que sur}
\unsure{$S'$, $S''$, $S'''$, le rang de $\underline{E}$ soit 0, 1, $\geq 2$. Alors la détermination des}
\unsure{Modules inversibles sur X se ramène à celle des $f^{-1}(S''')$ par 1.1.1.}
\unsure{Le cas 1.2 est justiciable de 1.1. D'autre part $f^{-1}(S'')$ est isomorphe}
\unsure{à $S''$. Donc son groupe de Picard n'est autre que $\mathrm{Pic}(S'')$, enfin}
\unsure{$f^{-1}(S')$ est vide, donc son groupe de Picard est nul.}
\uline{\unsure{Corollaire 1.3.}} \unsure{Sous les conditions de 1.1, supposons S connexe. Alors}
\unsure{tout Module inversible $\underline{L}$ sur X est isomorphe à un Module de la forme}
\unsure{$f^*(\underline{M})(n)$,} \illegible{2} \unsure{$n \in \mathbb{Z}$,} \illegible{3} \unsure{de plus il est}
\illegible{6} \unsure{de la forme}

\unsure{$*$}
\unsure{Un autre} \illegible{8}
\illegible{6}

\mspage{\hand{97}}

\clearpage
% Nouveau EGA V (cote 25): archive cote 25, site page 98 (PDF page 99).
% typescript. Draft compared: none.
% AI-transcribed by Claude Opus 5.5 (claude-opus-5-5, adaptive thinking, effort high), Message Batches API, batch ega25b (msgbatch_01A7n5YuXubUPNK1zsepi8RG), 2026-11-13; not yet checked by eye.
\scanpage{25/98}
\mspage{1}
\hand{Nouveau EGA V}
\inserted{Faisceaux inversibles et diviseurs sur les fibrés projectifs et multiprojectifs. Systèmes linéaires de diviseurs. \deleted{familles de}}
% NOTE: a long handwritten diagonal stroke runs across the page from the struck title down to ``de L.'' (end of Corollaire 1.3); a second stroke crosses the line ``Une autre façon de formuler'' towards the lower left corner

IV 2\over{0}{1} \deleted{\uline{Compléments sur les faisceaux inversibles}, \uline{les} \uline{diviseurs de}}
\deleted{\over{C}{\unsure{+}}\uline{artier}, \uline{les systêmes linéaires de diviseurs}.}

1. \uline{Détermination des faisceaux inversibles sur les fibrés projectifs}.
\uline{Application\xout{x}} \uline{au \deleted{groupe} des}\hand{/}\uline{automorphismes d'un fibré projectif}.

\uline{Théorème} 1.1. Soient S un préschéma, \underline{E} un Module loc libre sur \xout{S qui} de
rang \inserted{\unsure{fini}} $\geqslant 2$ en tout point,
\xout{n'est nul en aucun point,} P = P(\underline{E}) le fibré projectif qu'il définit. Alors
\margin{*} pour tout Module inversible \underline{L} sur P, on peut trouver une famille
\xout{\illegible{1}} $(S_n)_{n\in\underline{Z}}$ d'ouverts disjoints de S recouvrant S, indexée par \underline{Z}, et
\xout{\unsure{pour tout n}$\in$\underline{Z}} un Module inversible \xout{$\underline{L}_n$}\inserted[typed]{$\underline{M}$} sur S\xout{$_n$} , detelle façon que
\deleted{\xout{\unsure{$\underline{L}|p^{-1}(S_n)$ soit isomorphe à $p_n^{-1}(\underline{L}_n)\otimes\underline{O}$ la restriction de}}} la restricti-%
on de \underline{L} à $f^{-1}(S_n)$ soit isomorphe à celle de $\underline{M}\otimes_{\underline{O}_S}\underline{O}_P(1) = f^*(\underline{M})(n)$.
De plus, la \xout{famille} famille des $S_n$ \xout{\unsure{et la classe}} est uniquement déter-%
minée par ses conditions, ainsi que \underline{M} à isomorphisme unique près.

\uline{Remarque} 1.2. Lorsque on ne fait pas d'hypothèse sur le rang de \underline{E}, S se
décompose canoniquement en \inserted{le}\inserted[typed]{préschéma}somme de trois ouverts $S^0,S^1,S^2$ tel que sur
ceux-ci, le rang de \underline{E} soit resp. 0, 1, et $\geqslant 2$. Alors la détermination des
Modules inversibles sur P se ramène à celle des $f^{-1}(S^{i})$ pour i=0,1,2.
Le cas i=2 est justiciable de 1.1. , d'autre part $f^{-1}(S^1)$ est $S^1$-isomor-%
phe à $S^1$ donc son groupe de Picard n'est autre que Pic($S^1$), enfin
$f^{-1}(S^0)$ est vide donc son groupe de Picard est nul.

\uline{Corollaire} 1.3. Sous les conditions de 1.1. supposons S connexe\inserted[typed]{non vide}. Alors
tout Module inversible \underline{L} sur P est isomorphe à un Module de la forme
$f^*(\underline{M})(n)$, où $n\in\underline{Z}$ et \underline{M} est un Module inversible sur S. De plus n est
uniquement déterminé, et \underline{M} est déterminé à isom unique près, par la donnée
de \underline{L}.

Une autre façon de formuler ce corollaire est la suivante. Considé-%
rons les homomorphismes naturels Pic(S)$\longrightarrow$Pic(\unsure{P}), \underline{Z} $\longrightarrow$Pic(P), le premier
\margin{98}

\clearpage
% Nouveau EGA V (cote 25): archive cote 25, site page 99 (PDF page 100).
% typescript. Draft compared: none.
% AI-transcribed by Claude Opus 5.5 (claude-opus-5-5, adaptive thinking, effort high), Message Batches API, batch ega25b (msgbatch_01A7n5YuXubUPNK1zsepi8RG), 2026-11-13; not yet checked by eye.
\scanpage{25/99}
% NOTE: two long diagonal pen strokes run across the whole page: one from above "deuxième" (line 1) down to the left to "Corollaire" (line 5), and one from there down to the foot of the page. The whole page appears struck through by hand.
\mspage{2}

déduit de $f:P\longrightarrow S$, le deuxième  déterminé par l'élément  cl $(\underline{O}_P(1))$ de
Pic(P). On en conclut un homomorphisme canonique
(x) \qquad Pic(S)x\xout{\unsure{ZZx(}} $\underline{Z} \longrightarrow$ Pic(P)  ,
défini d'ailleurs sans hypothèse restrictive sur S ni sur $\underline{E}$.  Ceci posé:
\uline{Corollaire} 1.4.   Sous les conditions de 1.\unsure{1}. , \inserted{si $S\neq\emptyset$} \xout{supposons S connexe non}
\xout{vide} alors l'homomorphisme précédent est \xout{un isomorphisme} \inserted[typed]{injectif} , et même
bijectif lorsque S est connexe\hand{.}\deleted{non vide.}

\margin{[} Lorsqu'on \xout{suppose seulement S connexe, et} abandonne l'hypothèse
sur le rang $\geqslant 2$,  il résulte de 1.2. et 1.\unsure{4}. que l'homomorphisme précédent
\inserted[typed]{si S est connexe, mais}
est encore surjectif \xout{(mais} pas nécessairement \unsure{injectif|bijectif}, le noyau étant
isomorphe \xout{à Z resp. à Pic(S)xZ lorsque E est de rang 1 resp. de rang 0)}
à $\underline{Z}$  resp. à  Pic(S)x$\underline{Z}$  lorsque  $\underline{E}$  est de rang 1\hand{,} resp. de rang 0.

Démontrons  1.1.  en commençant par l'unicité.  Tout d'abord, si S est
le spectre d'un corps, notons que  $\underline{O}(n)$ n'est isomorphe à $\underline{O}(m)$ i.e.
$\underline{O}(m-n)$  n'est  isomorphe à $\underline{O}_P$ , que si  n=m i.e.  n-m=0 .  Cela provient
du fait que $\underline{O}(1)$ est ample et  dim P $\geqslant$1 (d'après l'hypothèse rang E $\geqslant$2):
on peut en effet supposer  d = m-n $\geqslant$0, si on avait d $>$0, alors $\underline{O}(d)$
serait ample, et ne pourrait être isomorphe à  $\underline{O}_P$ que si P est quasi-affi-%
ne, donc fini (puisqu'il est propre sur k).  Ceci prouve déjà  l'unicité
de la famille $(S_n)_{n\in\underline{Z}}$  considérée dans 1.1.  Pour l'unicité de $\underline{M}$ à isom.
unique près, on est ramené au cas où  $S=S_n$ , je dis que dans ce cas on a
un isomorphisme  (uniquement déterminé en termes de l'isom  $\underline{L} \simeq f$\over{$_*$}{$^*$}$(\underline{M})(n)$):
\[ \text{(xx)} \qquad M \xrightarrow{\sim} f_*(\underline{\unsure{L|E}}(-n)) \quad . \]
En effet,  l'isomorphisme  $\underline{L} \xrightarrow{\sim} f^*(\underline{M})(n)$   définit un isomorphisme
$\underline{L}(-n) \xrightarrow{\sim} f^*(\underline{M})$  et par suite un isomorphisme du deuxième membre de (xx)
avec $f_*(f^*(\underline{M}))$, qui lui-même est isomorphe (\,$\underline{M}$ étant loc libre) à
$\underline{M}\otimes f_*(\underline{O}_P)$ , or  (réf)  $f_*(\underline{O}_P) \xleftarrow{\sim} \underline{O}_S$ ,  d'où l'\unsure{isomorphisme|isemorphisme} (xx).

\margin{93}

\clearpage
% Nouveau EGA V (cote 25): archive cote 25, site page 100 (PDF page 101).
% typescript. Draft compared: none.
% AI-transcribed by Claude Opus 5.5 (claude-opus-5-5, adaptive thinking, effort high), Message Batches API, batch ega25b (msgbatch_01A7n5YuXubUPNK1zsepi8RG), 2026-11-13; not yet checked by eye.
\scanpage{25/100}
\mspage{3}

% NOTE: a single long handwritten diagonal stroke crosses the whole page from top to bottom; it is transcribed as a deletion of each paragraph
\deleted{Prouvons l'existence des $(S_n)$, $\underline{M}$ . Grâce à l'unicité déjà prouvée,
la question est \uline{locale} sur S, i.e. on est ramené à prouver le
\uline{Corollaire} 1.5. Soit $\underline{E}$ loc libre de rang fini sur S, P=P($\underline{E}$), $\underline{L}$ un
Module inversible sur P, s $\in$ S, alors il existe un voisinage ouvert U de s
et un entier n $\in \underline{Z}$ , tel que $\underline{L}\,|\,f^{-1}(U)$ soit isomorphe à $\underline{O}_P(n)\,|\,f^{-1}(U)$.}

\deleted{\hand{\unsure{[}}Bien entendu, lorsque le rang de E en s est $\geq$ 2, l'entier n est bien
déterminé. D'ailleurs 1.5. est trivial lorsque le rang de $\underline{E}$ en s est $\leq$ 1\hand{\unsure{]}}.
Noter que $\underline{E}$ est de la frime, car la question étant locale, on peut supposer
déjà $\underline{E} = \underline{O}_S^{r+1}$, donc P = $P_S^{r\xout{+1}}$\hand{\unsure{]}}. Par le procédé bréveté du par.8, on est
ramené de plus au cas où S est noethérien. Nous précédons endeux pas:}

\deleted{a) S est le spectre d'un corps k. Nous savons que $\underline{L}$ est défini
par un Module gradué de type fini L sur l'anneau gradué k$[t_0,\ldots,t_r]$.
On sait aussi que la restriction \inserted{$L'$} de $\tilde{L}$ à l'espace affine épointé
$E_k^{r+1}$ - (0) n'est autre que l'image inverse de $\underline{L}$ par le morphisme de
projection canonique \inserted{q:} $E_k^{r+1}$ - (0) $\longrightarrow P_k^r$ , et c'est donc un Module inver-%
sible. Soit $i : E^{r+1}$-(0)$\to E^{r+1}$ l'immersion canonique, il résulte
alors du fait que l'anneau \inserted{affine} k$[t_0,..,t_r]$ de $E^{r+1}$ est factoriel, donc à
fortiori son localisé en le point 0 de $E^{r+1}$ est factori\xout{\unsure{e}}el, et du fait
que ce dernier anneau est de dimension $\geq$ 2, que $i_*(\tilde{L'})$ est un Module
inversible, \xout{\illegible{1}} donc correspond à un module inversible M sur
k$[t_0,..,t_r]$. D'ailleurs M = $\Gamma(E^{r+1}$-0, $\unsure{q'|q^*}(\underline{L}))$ est gradué de façon natu-%
relle, enfin l'homomorphisme L $\to$ M est évidemment un isomorphisme
entous les \xout{\illegible{2}} points de $E^{r+1}$ distincts de 0. Donc quitte à
remplacer L par M, on est ramené au cas où L est inversible. Mais
k$[t_0,..,t_r]$ étant factoriel, L est donc \uline{libre} de rang 1 , d'abord
enfaisant abstraction de sa graduation; mais un lemme standard (qui se
trouve sans doute dans Bourbaki) implique qu'il est alors libre de rang 1
entant que module gradué, ce qui implique que le Module associé sur P}

\mspage{100}

\clearpage
% Nouveau EGA V (cote 25): archive cote 25, site page 101 (PDF page 102).
% typescript. Draft compared: none.
% AI-transcribed by Claude Opus 5.5 (claude-opus-5-5, adaptive thinking, effort high), Message Batches API, batch ega25b (msgbatch_01A7n5YuXubUPNK1zsepi8RG), 2026-11-13; not yet checked by eye.
\scanpage{25/101}
\mspage{4}

% NOTE: long hand-drawn diagonal strokes run across the page: one in the upper left margin (by the first lines), one from the top right down through lines 1--6, a short one through line 8, and a long one from ``Variante'' down to the foot of the page; they look like strokes marking the passages as dealt with, not deletions of single words
est isomorphe à un $\underline{O}_P(n)$. \ -- \ Du point de vue rédaction, il était mala-%
droit de commencer par considérer un $\underline{L}$ de type fini. Commencer carrément
par considérer \ $L = \Gamma(E^{r+1}-(0),\ q_*(\underline{L}))$ \ définissant le Module
$\tilde{L} \ = \unsure{q}_*(q^*(\underline{L}))$ \ , on sait (Chap II) que \ $L$ estb un Module gradué qui
détermine précisément $\underline{L}$, puis on montre que $L$ est libre\inserted{de rang 1} \deleted{comme Module}
gradué \ par le raisonnement indiqué\hand{:}

b) \ Cas général. \ Il se déduit du cas a) \ grâce à \ III 4.6.5. , en
utilisant la relation \ $H^1(P^r_k,\underline{O}_{P^r_k}) = 0$ \ établie dans III 2 . \ \ CQFD

Variante à \ 1.4. \ Soit \ $\underline{Z}\,(S)$ \ l'ensemble des \ fonctions localement
constantes à valeurs entières sur S, on définit un homomorphisme évident
\[ Z(S) \longrightarrow \mathrm{Pic}(P) \ , \]
\inserted{les} \xout{xxxxx} $\underline{n} \in Z(S)$ \ correspondant en eff\xout{etx} \inserted{aux} \xout{xxx} partitions $(S_n)_{n\in\underline{Z}}$ \ en ouvert\unsure{s}
disjoints (dont certains peuvent être vides), \xout{xxxfaçonxxque} et \xout{xx} à une
telle partition on associe le Module inversible $\underline{O}_P(\underline{n})$ dont la \inserted{restr à} \xout{valeurxsur}
$f^{-1}(S_n)$ \ est $\underline{O}_P(n)$. \ On trouve donc une variante
(x bis) \hfill $\mathrm{Pic}(S)\mathrm{x}\underline{Z}(S) \longrightarrow \mathrm{Pic}(P)$ \ , \hfill
et \ un énoncé plus général et plus satisfaisant que 1.4. affirme que
(sous les conditions de 1.1.) c'est là une bijection. (NB Lorsque S $\neq \emptyset$
\inserted[typed]{resp. connexe}\hand{)}
alors l'application \ canonique $\underline{Z} \to \underline{Z}(S)$ \ associant à tout $n \in \underline{Z}$ \ la foncti-%
on constante de valeur n, est injective resp. surjective, \ et on retrouve .
formellement 1.4. , qu'il y a donc lieu d'énoncer sous la forme plus
générale que je viens d'expliciter . \ Signaler aussi en remarque que
dans le langage du schéma de Picard qui sera introduit au Chap V, 1.1.
(sous la forme équivalente \ précédente) s'énonce simplement \ en disant que
l'homomorphisme canonique
\[ \underline{Z}_S \longrightarrow \underline{\mathrm{Pic}}_{P/S} \]
du schéma en groupes constant \xout{xxxxZxSxdéfinixpar} $\underline{Z}$ sur S dans le schéma

\margin{101}

\clearpage
% Nouveau EGA V (cote 25): archive cote 25, site page 102 (PDF page 103).
% typescript. Draft compared: none.
% AI-transcribed by Claude Opus 5.5 (claude-opus-5-5, adaptive thinking, effort high), Message Batches API, batch ega25b (msgbatch_01A7n5YuXubUPNK1zsepi8RG), 2026-11-13; not yet checked by eye.
\scanpage{25/102}
\mspage{5}
% NOTE: the whole page is crossed by three long diagonal pen strokes running from upper right to lower left (top of the page down to "rang" in the line "i.e..E partout de rang"; from "(A ce pro-" down to "1. 2." in "la remarque 1. 2."; from "projectif" in "Soit P' = ..." down to "définis" in the last line). The text is transcribed as live text under them.
de Picard, déduit de la s ection de ce dernier définie par $\underline{O}_P(1)$, est un
\uline{isomorphisme}.

Soit P un fibré projectif sur un corps k. Un Module inversible
$\underline{L}$ sur P est dit \uline{de degré} n si $\underline{L}$ est isomorphe à $\underline{O}_P(n)$; lorsque
dim P$\geq$1, cela détermine n en \xout{t}termes de $\underline{L}$ , mais si dim P $\leq$ \unsure{0|D}
(i.e. p est vide, ou réduit à un point) alors $\underline{L}$ est degré n pour tout n.
Lorsque P est un fibré projectif sur une base quelconque, un Module
inversible $\underline{L}$ sur P est dit de degré n s'il est de degré n sur chaque fibr
% NOTE: the end of the preceding line is cut off at the right edge of the scan
Cela détermine n, sauf si toutes les fibres de P/S sont de dim$\leq$0, alors
$\underline{L}$ est de degré n pour tout n. Dire que $\underline{L}$ est de degré n signifie aussi,
en vertu de 1.1. et 1.2. , que \inserted[typed]{la classe} \xout{\unsure{xximage}} de $\underline{L}$ dans Pic(P) est dans
l'image de \inserted{$\mathrm{Pic}\,S\times\{n\}$ \unsure{par|pr}} l'homomorphisme (x) Pic(S)x$\underline{Z}$ $\longrightarrow$Pic(P) envisagé plus haut,
i.e. que $\underline{L}$ est isomorphe à un Module de la forme $f^*(\underline{M})(n)$ , où $\underline{M}$ est
inversible sur S. D'ailleurs, lorsque les fibres de P sont non vides
i.e..$\underline{E}$ partout de rang $\geq$1, alors $\underline{M}$ est déterminé à isomorphisme unique
\margin{\textbar}près en termes de $\underline{L}$, comme il résulte encore de 1.1. et 1.2. (A ce pro-%
pos, je m'aperçois qu'il y avait lieu d'énoncer 1.1. ainsi, sans hypothèse
sur le rang de $\underline{E}$ : tout $\underline{L}$ inversible sur S peut s e mettre sous la forme
indiquée dans l'énoncé; lorsque les fibres de P sont non vides i.e. le
rang de $\underline{E}$ partout $\neq$ 0 , alors $\underline{M}$ est déterminé à isom unique près par
$\underline{L}$ ; lorsque les fibres de P sont toutes de dim $\geq$1, \inserted{i.e. le rang de $\underline{E}$ partout $\geq 2$,} alors la partition
de S est également déterminée de façon unique par le choix de $\underline{L}$. De cette
façon, la remarque 1. 2. est éliminée et passe dans la démonstration).

Soit P' = P($\underline{E}$') un deuxième fibré projectif, alors la détermina-%
tion de Pic(P) permet en prihcipe de déterminer les S-morphismes
g:P$\to$P', puisque ceux-ci sont définis par un Module inversible $\underline{L}$ \inserted{($\simeq g^*\mathcal{O}_{P'}(1)$)} sur P

\margin{102}

\clearpage
% Nouveau EGA V (cote 25): archive cote 25, site page 103 (PDF page 104).
% typescript. Draft compared: none.
% AI-transcribed by Claude Opus 5.5 (claude-opus-5-5, adaptive thinking, effort high), Message Batches API, batch ega25b (msgbatch_01A7n5YuXubUPNK1zsepi8RG), 2026-11-13; not yet checked by eye.
\scanpage{25/103}
\mspage{\hand{6}}
% NOTE: a single handwritten diagonal line runs across the whole page, from the top left of the text down to the bottom right edge; it seems to cancel the page. The text is transcribed below as live text.

et  un homomorphisme  $f^{\hand{*}}_{*}(\underline{L}) \longleftarrow \underline{E}'$  (tel que l'homomorphisme \unsure{associé}
$f^{*}(\underline{E}') \to \underline{L}$  soit surjectif), modulo  isomorphisme en $\underline{L}$.  On dira que
$g\colon P \to P'$  est \uline{de degré} n si $\underline{L} = g^{*}(\underline{O}_{P'}(1))$  est de degré n. Il suffit
évidemment de savoir déterminer les homomorphismes de degré n pour tout n
donné.  Noter que si g est de degré n,\inserted[typed]{et P est à fibres de dimension $\geq 1$} on aura nécessairement $n \geq 0$
(car  sur un corps k, \xout{$\underline{O}_P(n)$} si dim P $\geq 1$, $\underline{O}_P(n)$ n'est engendré par ses
sections que si $n \geq 0$); bien entendu, on peut \xout{aussi} se restreindre au cas
où P a ses fibres de dim $\geq 1$ (en procédant comme dans 1.2.).  Comme on a
\[ f_{*}(f^{*}(\underline{M})(n)) = \underline{M} \otimes f_{*}(\underline{O}_P(n)) = \underline{M} \otimes \underline{\mathrm{Sym}}^{n}(\underline{E}) \hand{,} \]
on voit que la détermination des \over{g}{S}-morphismes  $g\colon P \to P'$  est ramenée
à la détermination des couples $(\underline{M},u)$ à isom près, où $\underline{M}$  est un Module
inversible sur S et  $u\colon \underline{E}' \to \underline{M}\otimes\underline{\mathrm{Sym}}^{n}(\underline{E})$  un homomorphisme  tel que
l'homomorphisme correspondant  $f^{*}(\underline{E}') \to f^{*}(\underline{M})(n)$  soit un épimorphisme.
Donc  g détermine un premier invariant de nature globale sur S, savoir
$\mathrm{cl}(\underline{M}) \in \mathrm{Pic}(S)$, et cet invariant fixé par un $\underline{M}$ choisi,  les g correspon-%
dants correspondent à un certain sous-ensemble de l'ensemble quotient
$\mathrm{Hom}(\underline{E}',\underline{M}\otimes\underline{\mathrm{Sym}}^{n}(\underline{E}))/\mathrm{I}\,(S,\underline{O}_S)^{*}$,  le passage au quotient par le groupe
$\mathrm{I}\,(S,\underline{O}_S)^{*}$ correspondant  au\xout{x}  ``modulo isomorphisme'' dans la description
des $g\colon P \to P'$ via Modules inversibles  (NB  les endomorphismes resp. automor
phismes d'un $\underline{L}$ inversible sur un fibré projectif $\underline{P}$, correspondant aux
sections resp. sections\hand{$|$}inversibles de $\underline{O}_P$ sur P,  ou encore des sections
resp sections inversibles de $\underline{O}_S \hand{$\xrightarrow{\sim}$} f_{\hand{*}}(\underline{O}_P)$ sur S).

Cas particuliers: \hand{1)} n=0 , alors  il faut prendre  les  homomorphismes
$\underline{E}' \to \underline{M}$  qui sont \deleted{partout non nuls}\inserted{\unsure{surjectifs} i.e. partout non nuls}, modulo isomorphisme en $\underline{M}$ : on trouve
exactement les  morphismes  $g\colon P \to P'$  de la forme  hf , où  h est une
section de  P' sur S  (savoir celle déterminée par  le quotient inversible
$\underline{M}$ de $\underline{E}'$). Donc les \inserted{S-}\uline{morphismes de  degré} 0  \uline{de} P \uline{dans} P' \uline{sont les morphis-}%

\margin{103}

\clearpage
% Nouveau EGA V (cote 25): archive cote 25, site page 104 (PDF page 105).
% typescript. Draft compared: none.
% AI-transcribed by Claude Opus 5.5 (claude-opus-5-5, adaptive thinking, effort high), Message Batches API, batch ega25b (msgbatch_01A7n5YuXubUPNK1zsepi8RG), 2026-11-13; not yet checked by eye.
\scanpage{25/104}
\mspage{7}

\uline{mes} \uline{constants} rel. à S. 2) n=1, alors il faut prendre les homomorphis-%
mes \xout{\unsure{partout surjec}} $\underline{E}'\to \underline{E}\otimes\underline{M}$ qui sont \uline{surjectifs}, comme on vérifie
aussitôt, et l'homomorphisme correspondant \unsure{$g$}: $P\to P'$ n'est autre que le
composé $P(\underline{E})\xrightarrow{\sim}P(\underline{E}\otimes\underline{M})\to P(\underline{E}')$ , où le premier hom est l'isomorphisme
canonique envisagé dans le Chap II, et le d euxième l'immersion fermée
canonique déduite de l'épimorphisme $\underline{E}'\to\underline{E}\otimes\underline{M}$ . Si on appelle \uline{linéaires}
les hom de P dans P' qui se décrivent comme un tel composé, on voit donc
que les morphismes $g:P\to P'$ qui sont de degré 1 sont exactement ceux qui
sont \uline{linéaires}.

Pour terminer, déterminons les isomorphismes de P avec P':
\uline{Théorème} 1.6. Soient S un préschéma, $P=P(\underline{E})$ et $P'=P(\underline{E}')$ deux fibrés
projectifs sur S\inserted{\unsure{définis par} $\underline{E}$, $\underline{E}'$ \unsure{loc. libres de type fini},}. Alors tout S-isomorphisme $g$: $P\xrightarrow{\sim}P'$ est définissable
comme le composé $P(\underline{E})\xrightarrow{\sim}P(\underline{E}\otimes\underline{M})\to \underline{P}(\underline{E}')$, où\inserted{$\underline{M}$ est un Module inversible sur S,} la premier hom est l'isom
can de Chap II , et le \unsure{deuxième|dauxième} est l'isomorphisme \xout{canonique} déduit
d'un isomorphisme u: $E'\xrightarrow{\sim}\underline{E}\otimes\underline{M}$ . Lorsque les fibres de \unsure{P} sont non
vides (resp. de dim $\geq 1$) alors $\underline{M}$ (resp. le couple $(\underline{M},u)$) est déterminé
à isomorphisme unique près en termes de g.

D'après les considérations précédentes, on est ramené à prouver
que g est de degré 1, ce qui nous ramène au cas où S est le spectre d'un
corps,\inserted[typed]{et bien sur on peut supposer dim $P\geq 1$.} Mais notons que $\underline{O}_P(1)$ est alors caractérisé intrinséquement (i.e.
indépendamment de la façon dont P a été réalisé comme fibré projectif)
comm\xout{ex}e le générateur de Pic(P) (parmi les deux générateurs $\underline{O}_P(1)$ et
$\underline{O}_{\unsure{P}}(-1)$ ) qui est \uline{ample} ; par suite si g : $P\to P'$ est un isomorphisme,
alors $g^{*}(\underline{O}_{P'}(1))$ est isomorphe à $\underline{O}_P(1)$ , et on gagne. - Sous forme local
% NOTE: the right edge of the scan cuts off the end of "local[e]"
et moins savante, on peut énoncer:
\uline{Corollaire} 1.7. Sous les conditions de 1.6. , tout\inserted{S-}isomorphisme $g:P\to P'$
peut se décrire, au voisinage de tout $s\in S$, à l'aide d'un isomorphisme
u: $\underline{E}|U \simeq E'|U$ , ce dernier étant bien déterminé mod mult par un élément
\margin{104}

\clearpage
% Nouveau EGA V (cote 25): archive cote 25, site page 105 (PDF page 106).
% typescript. Draft compared: none.
% AI-transcribed by Claude Opus 5.5 (claude-opus-5-5, adaptive thinking, effort high), Message Batches API, batch ega25b (msgbatch_01A7n5YuXubUPNK1zsepi8RG), 2026-11-13; not yet checked by eye.
\scanpage{25/105}
\mspage{8}

de $\Gamma(U,\underline{O}_U)^*$ .

En particulier:

\uline{Corollaire} 1.8. Soient S un préschéma, P=P(\underline{E}) un fibré projectif sur S
défini par \underline{E} loc libre de type fini, u un automorphisme de P. Alors u
est déterminé, au vosinage de tout point s $\in$ S , par un automorphisme
u de \underline{E}|U, ce dernier étant bien déterminé par u modulo multiplication par
un élément de $\Gamma(U,\underline{O}_U)^*$ .

% NOTE: in the next line the words "on déduit peut déduire facilement" are circled by hand
\uline{Remarque} 1.9. De 1.8. on \xout{déduit} peut déduire facilement que le foncteur \margin{le faire ! c'est l'ouvert de \unsure{$P(\underline{E}\otimes\underline{E}^{\vee})$|$P(\underline{E}\otimes\underline{E})$} défini par la non annulation du ``déterminant'' ---}
\xout{Aut} en groupes $\underline{Aut}_S(P)$ sur S est représentable par un \unsure{préschéma} affine
de présentation finie sur S, qui s'interprète aussi comme le \inserted{pré}schéma en
groupes \uline{quotient} du schéma en groupes linéaire $\underline{Gl}(\underline{E})$ par son centre
$\underline{G}_m$ . Ce \inserted{pré}schéma en groupes s'appelle \xout{classiquement le (préschéma}
le \uline{préschéma en groupes projectif}, ou simplement \uline{groupe projectif}, défini
par \underline{E}, \inserted[typed]{et se note GP(\underline{E})\hand{,}} Lorsque \underline{E} est libre, $\underline{E}\simeq\underline{O}_S^{r+1}$, alors GP(\underline{E} ) n'est autre que le
\xout{schéma} préschéma en groupes $GP(r)_S$ déduit par le changement de base
$S\to\mathrm{Spec}(\underline{Z})$ \xout{\unsure{du}} du schéma en groupes analogue GP(r) sur Spec(\underline{Z}), appelé\xout{\unsure{e}}
le \uline{Groupe \unsure{\over{p}{P}}rojectif} \uuline{absolu}.

\margin{105}

\clearpage
% Nouveau EGA V (cote 25): archive cote 25, site page 106 (PDF page 107).
% typescript. Draft compared: none.
% AI-transcribed by Claude Opus 5.5 (claude-opus-5-5, adaptive thinking, effort high), Message Batches API, batch ega25b (msgbatch_01A7n5YuXubUPNK1zsepi8RG), 2026-11-13; not yet checked by eye.
\scanpage{25/106}
\mspage{9}

2. \uline{Diviseurs relatifs et faisceaux inversibles sur les fibrés projectifs}

\uline{et multiprojectifs}.

\hand{2.1.} Soit comme au N°1 P=P($\underline{E}$), $\underline{E}$ loc libre sur S de rang $\geqslant$ 2 partout.
Nous nous proposons de déterminer l'ensemble Div(P/S) des diviseurs rela-%
tifs \inserted{$\geqslant 0$} sur P p. r. à S. On sait que la donnée d'un tel diviseur revient à la
donnée d'un Module inversible $\underline{L}$ sur P, muni d'une section $\phi$ de $\underline{L}$
transversalement régulière. Or d'après 1.1. (abstraction faite d'une
partition éventuelle de S, lorsque S n'est pas connexe) $\underline{L}$ est isomorphe
à un $\underline{M}\otimes\underline{O}_P(n)$, où $\underline{M}$ est un Module inversible sur S, d'ailleurs déterminé
à isomorphisme unique près en termes de $\underline{L}$. D'ailleurs, on sait que
l'on a alors des isom canoniques
\[ \hand{(*)} \qquad f_{\hand{*}}(\underline{L}) \simeq \underline{M}\otimes f_{\hand{*}}(\underline{O}_P(n)) \simeq \underline{M}\otimes\underline{Sym}^n(\underline{E}) \]
donc la donnée d'une section \xout{de} $\phi$ de $\underline{E}$ revient à la donnée d'une section
\inserted{$\phi$} de $\underline{M}\otimes\underline{Sym}^n$(E) . Compte tenu du fait que les fibres de P/S sont intègres,
on voit d'ailleurs que $\phi$ est transversalement régulière (i.e. régulière
sur chaque fibre) sss \hand{$\phi$}(s)$\neq$0 pour tout s$\in$S, ou ce qui revient au même,
s\hand{\unsure{ss}} l'homomorphisme
\[ \hand{$\phi$} : \check{\underline{M}} \longrightarrow \underline{Sym}^n(E) \]
défini par \hand{$\phi$} est ``universellement injectif'' i.e. loc un isomorphisme sur
un facteur direct, ou ce qui revient au même, si son transposé
\[ \hand{$\check{\phi}$} : \underline{Sym}^n(\underline{E})^{\hand{$\vee$}} \longrightarrow \underline{M} \]
est surjectif.

Nous dirons que le diviseur relatif D sur P est \uline{de degré} n
si $\underline{O}_X(\underline{D})= \underline{L}$ est de degré n au sens du n° précédent. \deleted{Ainsi, grâce à 1.1.}
\deleted{pour D donné, il existe une unique \inserted{décomposition} \deleted{partition} de S en somme d'ouverts $S_n$}
\deleted{(n N) tel}

\mspage{106}

\clearpage
% Nouveau EGA V (cote 25): archive cote 25, site page 107 (PDF page 108).
% typescript. Draft compared: none.
% AI-transcribed by Claude Opus 5.5 (claude-opus-5-5, adaptive thinking, effort high), Message Batches API, batch ega25b (msgbatch_01A7n5YuXubUPNK1zsepi8RG), 2026-11-13; not yet checked by eye.
\scanpage{25/107}
\mspage{10}

Lorsque D $\geq 0$, cela implique $n \geq 0$, car si \inserted{\unsure{on avait|n'avait}} $n \unsure{<} 0$ toute section de
$\underline{L}$ sur X serait nulle. Grâce à 1.1. , si D est un diviseur relatif
$\geq 0$ sur P, alors il existe une unique décomposition de D en somme
disjointe d'ouverts $S_n (n \in \underline{N})$ telle que pour tout $n \in \underline{N}$ , $\underline{L}|P|S_n$ soit de
degré n. Cela ramène la détermination de l'ensemble des diviseurs rela\xout{xx}%
tifs $\geq 0$ au cas des diviseurs relatifs $\geq 0$ de degré n donné. (Ce
\margin{$/\!/$} \unsure{tèxte|texte} remplacera évidemment le ``abstraction faite'' plus haut). Ceci posé
les refexions précédentes établissent la

\uline{Proposition} 2.2. Sous les hypothèses précédentes, on a une correspondan%
ce biunivoque entre l'ensemble $\mathrm{Div}^{n}(\unsure{P}/S)$ des diviseurs relatifs $\geq 0$ de
degré n sur P, et l'ensemble des Modules quotients inversibles $\underline{M}$ de
$\underline{\mathrm{Sym}}^{n}(\underline{E})^{\hand{\vee}}$ , (ou ce qui revient au même, les sous-Modules inversibles loc
facteurs directs $\overset{\hand{\vee}}{\underline{M}}$ de $\underline{\mathrm{Sym}}^{n}(\underline{E})$ ) . Si D et $\underline{M}$ se correspondent, alors
$\underline{O}_P(D)$ est canoniquement isomorphe à $\underline{M}\otimes\underline{O}_P(n)$, et la section $s_D$ s'iden-%
tifie par cet isomorphisme à ce qu'on \unsure{devine}, compte tenu de $(*)$.

Notons que cette description est compatible avec tout changement de
base en S. Compte tenu de l'interprétation des Modules quotients inv de
$\underline{\mathrm{Sym}}^{n}(\underline{E})^{\hand{\vee}}$ comme les sections sur S de $P(\underline{\mathrm{Sym}}^{n}(\underline{E})^{\hand{\vee}})$, on trouve ceci:
\uline{Corollaire} 2.3. Le sous-foncteur \xout{xx} $\mathrm{Div}^{n}_{P/S}$ de $\mathrm{Div}^{+}_{\unsure{P}/S}$ est représentable
canoniquement par $P(\underline{\mathrm{Sym}}^{n}(\underline{E})^{\hand{\vee}})$.

Compte tenu des considérations de 2.1. , il s'ensuit:
\uline{Corollaire} 2.4. $\underline{\mathrm{Div}}^{+}_{\unsure{P}/S}$ est représentable canoniquement \xout{xxxxx} par
le S-préschéma somme des $P(\underline{\mathrm{Sym}}^{n}_{P/S}(\underline{E})^{\hand{\vee}})$ , $n \in \underline{N}$ .
2.5. Supposons maintenant donné une famille finie $(\underline{E}_i)_{i \in I}$ de Modules
localement libres sur S, d'où des $P_i = P(\underline{E}_i)$ et un P= produit fibré des
$P_i$ , appelé aussi \uline{fibré multiprojectif} sur S défini par $(\underline{E}_i)$ . Pour
simplifier les notations, désignons par $\underline{O}_i(n)$ l'image inverse sur P du

\mspage{107}

\clearpage
% Nouveau EGA V (cote 25): archive cote 25, site page 108 (PDF page 109).
% typescript. Draft compared: none.
% AI-transcribed by Claude Opus 5.5 (claude-opus-5-5, adaptive thinking, effort high), Message Batches API, batch ega25b (msgbatch_01A7n5YuXubUPNK1zsepi8RG), 2026-11-13; not yet checked by eye.
\scanpage{25/108}
\mspage{11}

Module \inserted{inversible} $\underline{O}_{P_i}(n)$ sur $P_i$ . Pour tout systême d'entiers
\[ \underline{n} = (n_i)_{i \in I} \in \underline{Z}^I \]
nous \xout{désignerons} poserons
\[ \underline{O}_P(\underline{n}) = \bigotimes_i \underline{O}_i(n_i) = \bigotimes_{i\,\underline{O}_S} \underline{O}_{P_i}(n_i) \quad . \]
Ceci posé, 1.1. se généralise de la façon suivante:

\uline{Proposition} 2.6. Supposons les $\underline{E}_i$ partout de rang $\geq$2. Alors pour tout
Module inversible $\underline{L}$ sur le fibré multiprojectif P, il existe une
\xout{partition} décomposition de S en somme disjointe d'ouverts $\underline{S}_{\underline{n}}$ , $\underline{n} \in \underline{Z}^I$,
et \xout{pour tout $\underline{n}$ ∈ $\underline{Z}^I$, un} un Module inversible $\underline{M}$ sur S, tels que $\underline{L}|P|S_{\underline{n}}$
soit isomorphe à $\underline{O}_P(\underline{n})|P|S_n$ . De plus, les $S_{\underline{n}}$ sont déterminés de façon
unique, et \xout{les} $\underline{M}$ est déterminé à isomorphisme unique près.

\inserted{2.7.} La démonstration \xout{résulte} consiste en une réduction immédiate à 1.1.
Sous les conditions de 2.6. , on peut donc associer à tout $\underline{L}$ inversible
sur P un ``multidegré'' $\underline{n} = (n_i)_{i \in I} \in \underline{Z}\hand{\unsure{$^{*}$}}(S)^I$,qui caractérise $\underline{L}$ à isomorphis
me près lorsque Pic(S)=0 . On peut d'ailleurs interpréter $n_i$ (appelé
``\uline{degré partiel} de $\underline{L}$ par rapport au facteur $P_i$ d'indice i '' ), lorsqu'on
se donne pour chaque i une section $g_i$ de $P_i$\xout{x} sur S (NB de telles sectio
ns existent en tous cas localement sur S), en notant que ce système définit
pour chaque i un S-morphisme \xout{$h_i$}
\[ h_i : P_i \longrightarrow P \quad ; \]
ceci posé, on a
\[ n_i = \deg h_i^{*}(\unsure{$\underline{L}$|$\underline{E}$}) \quad . \]
On fera attention qu'en général, les $n_i$ ne sont pas des entiers, mais des
applications localement constantes de S dans $\underline{Z}$ .

Procédant comme au N°1, on peut déduire de 2.6. la détermination des
morphismes d'un fibré multiprojectif dans un autre, et en particulier celle
des automorphismes d'un fibré multiprojectif. Plus intéressant pour nous,

\mspage{108}

\clearpage
% Nouveau EGA V (cote 25): archive cote 25, site page 109 (PDF page 110).
% typescript. Draft compared: none.
% AI-transcribed by Claude Opus 5.5 (claude-opus-5-5, adaptive thinking, effort high), Message Batches API, batch ega25b (msgbatch_01A7n5YuXubUPNK1zsepi8RG), 2026-11-13; not yet checked by eye.
\scanpage{25/109}
\mspage{12}

en vue du par.25 sur le résult\unsure{ant} et le discriminant des formes, sera la
détermination des diviseurs relatifs $\geqslant 0$ sur un fibré multiprojectif.
2.8.
Si $\underline{D}$ est un diviseur relatif sur P, on définit son multidegré comme étan
% NOTE: the right edge of the page is cut off; "étan[t]" and "ramené[e]" lose their last letters
celui de $\underline{O}_P(D)$. Comme plus haut, la détermination de $\mathrm{Div}^{+}(P/S)$ est ramené
à celle de $\mathrm{Div}^{\underline{n}}(P/S)$, pour un multidegré $\underline{n}\in\underline{N}^{I}$ donné, qui fournit un
isomorphisme
\[ \underline{L} =\underline{O}_X(D) \simeq \underline{M} \otimes_{\underline{O}_S} \underline{O}_P(\underline{n}) \ . \]
Or on obtient, grâce au Chap II par.2 (\xout{\illegible{3}}
\[ (xx) \qquad f_*(\underline{D}) = \underline{M} \otimes f_*(\underline{O}_P(\underline{n})) = \underline{M} \otimes \bigotimes_i \underline{\mathrm{Sym}}^{n_i}(\underline{E}_i) \quad . \]
Proéédant maintenaht comme pour 2.2. on trouve la

\uline{Proposition} 2.9. Avec les notations précédentes, à reppaler, on a\xout{x} une
corr biun can entre l'ensemble $\mathrm{Div}^{\underline{n}}(P/S)$ des diviseurs relatifs de
multidegré $\underline{n}$ sur P, et l'ensemble des Modules inversibles quotients $\underline{M}$
de $\bigotimes_i \underline{\mathrm{Sym}}^{n_i}(\underline{E}_i)^{\hand{\vee}}$ (ou ce qui revient au même ...) . Si D et $\underline{M}$ se corres-%
pondent, $\underline{O}_P(D)$ est can isomorphe à $\underline{M}\otimes\underline{O}_P(\underline{n})$ , et $s_D$ s'identifie alors
à ce qu'on devine, compte tenu de (xx).

\uline{Corollaire} 2.10. Le sous-foncteur $\mathrm{Div}^{\underline{n}}_{P/S}$ de $\mathrm{Div}^{+}_{P/S}$ correspondant aux
diviseurs relatifs de multidegré $\underline{n}$ est can représentable par le fibré
projectif $P(\bigotimes_i \underline{\mathrm{Sym}}^{n_i}(\underline{E}_i)^{\hand{\vee}})$, et $\underline{\mathrm{Div}}^{+}_{P/S}$ est can repr par le préschéma
somme de ces de\unsure{rni|mmi}ers, pour $\underline{n}\in\underline{N}^{I}$ .

2.11. La détermination précédente très simple des $\mathrm{Div}_{P/S}$ est due à la
structure très simple de Pic(P/S) (en fait, à la structure``discrète'' du
préschéma de Picard $\underline{\mathrm{Pic}}_{P/S}$ ....). Nous pouvons abstraire le raisonnement
fait, qui revient essentiellement à établir une représentabilité ``relative''
(par rapport à $\underline{\mathrm{Pic}}$ ). Pour ceci, donnons-nous un morphisme
\[ f\colon X \to S \]
propre et plat de présentation finie, \uline{et} un Module inversible $\underline{L}$ sur X.

\mspage{109}

\clearpage
% Nouveau EGA V (cote 25): archive cote 25, site page 110 (PDF page 111).
% typescript. Draft compared: none.
% AI-transcribed by Claude Opus 5.5 (claude-opus-5-5, adaptive thinking, effort high), Message Batches API, batch ega25b (msgbatch_01A7n5YuXubUPNK1zsepi8RG), 2026-11-13; not yet checked by eye.
\scanpage{25/110}
\mspage{13}

Nous nous proposons de déterminer le sous-groupe Div$^{\mathrm{L}}$(X/S) de Div$^{+}$(X/S)
formé des diviseurs relatifs positifs tels que $\underline{O}_X$(D) soit isomorphe à
un Module de la forme $\underline{L}\otimes_{\underline{O}_X}\underline{M}$, où $\underline{M}$ est un Module inv sur S (dépendant de
D). Nous supposerons qu'on a
\[ f_{*}(\underline{O}_X)\overset{\sim}{\longleftarrow}\underline{O}_S \quad , \]
ce qui implique que ci-dessus $\underline{M}$ est déterminé à isom unique près en termes
de D, comme
\[ \underline{M}\simeq f_{*}(\underline{L}^{-1}\otimes\underline{O}_X(D)) \quad . \]
La donnée d'un D correspondant à un $\underline{M}$ donné revient donc à la donnée
\inserted[typed]{\xout{$\underline{\phi}$x}}
d'une section\xout{s} $\phi$ de $\underline{L}\otimes\underline{M}$ transversalement régulière. Or on a vu au
Chap III par 7 \xout{f$_*$($\underline{L}\otimes\underline{M}$)xxxxf$_*$($\underline{L}$)$\otimes$Mxx} qu'il \xout{xixx} existe un Module Q de
présentation finie sur S, dont la formation commute d'ailleurs à tout
changement de base, et un isomorphisme de foncteurs en \unsure{$\underline{N}$|$\underline{G}$} (Module qu coh
variable sur S)
\[ f_{*}(\underline{L}\otimes\underline{G})\simeq \mathrm{Hom}(\underline{Q},G) \quad . \]
(A vrai dire, dans loc cit on a supposé S loc noethérien, mais on se
débarasse de cette hypothèse de façon évidente, par le pro\xout{x}cédé brévété,
compte tenu de la commutation de Q au changement de\xout{x} base). En particulier
la donnée de $\phi$ équivaut à \xout{xx} la donnée d'un homomorphisme
\[ \psi : Q \longrightarrow M \ . \]
Une condition nécessaire pour que $\phi$ soit transversalement régulière, et
qui est suffisante lorsque les fibres de X sont intégres, est que $\phi$ soit
$\neq 0$ \xout{xix} fibre par fibre, ce qui en termes de $\psi$ signifie simplement que
$\psi$ est surjectif, donc que \xout{Mxxxx} $\psi$ correspond à une section du fibré pro-%
jectif $\underline{P}(\underline{Q})$ sur S. On obtient ainsi la

\uline{Proposition} 2.1\unsure{1}. \inserted[typed]{Avec les notations} \xout{Sous les conditions} précédentes,\xout{xxx X/S} Div$^{\mathrm{L}}$(X/S) est
en correspondance biunivoque canonique avec l'ensemble des \xout{point} sections
de $\underline{P}(\underline{Q})$ sur S, correspondant à des Modules quotients $\underline{M}$ de $\underline{Q}$ tel\unsure{\xout{s}|s} que

\mspage{110}

\clearpage
% Nouveau EGA V (cote 25): archive cote 25, site page 111 (PDF page 112).
% typescript. Draft compared: none.
% AI-transcribed by Claude Opus 5.5 (claude-opus-5-5, adaptive thinking, effort high), Message Batches API, batch ega25b (msgbatch_01A7n5YuXubUPNK1zsepi8RG), 2026-11-13; not yet checked by eye.
\scanpage{25/111}
\mspage{14}

la sedtion $\phi$ de $\underline{L}\otimes\underline{M}$ définie par \unsure{$\psi$} : $Q\longrightarrow\underline{M}$ soit transversalement
régulière.

Supposons maintenant que l'hypothèse $\underline{O}_S\xrightarrow{\sim}f_*(\underline{O}_X)$ soit encore vrai
après tout changement de base, ou ce qui revient au même grâce à III 7,
qu'on ait la condition
(xxx) \qquad\qquad $k(s)\xrightarrow{\sim}H^o(X_s,\underline{O}_{X_s})$
pour tout $s\in S$. Alors 2.12. s'applique également à tout $X_{S'}/S'$, pour un
changement de base quelconque $S'\rightarrow S$. On obtient ainsi le
\uline{Théorème} 2.13. Soit $f\colon X\longrightarrow S$ un morphisme \xout{plat} de présentation finie,
propre et plat satisfaisant à (xxx) ci-dessus, $\underline{L}$ un Module inversible
sur $X$, considérons le sous-foncteur $\mathrm{Div}^{L}_{X/S}$ de $\mathrm{Div}^{+}_{X/S}$ défini ci-dessus en
termes de $\underline{L}$. \xout{\unsure{Ce foncteur}} Il existe un Module de présentation finie $\underline{Q}$ sur
$S$ tel que le foncteur précédent soit représentable par un sous-préschéma
ouvert rétrocompact du fibré projectif $\underline{P}(\underline{Q})$ ; lorsque les fibres de $X/S$
sont géométriquement intègres, alors $\mathrm{Div}^{L}_{X/S}$ est représentable par ce
\inserted[typed]{fibré} \xout{\unsure{fibré}} projectif lui-même.

La dernière assertion résulte immédiatement de ce qu'on avait dit
avant. Pour le cas général, on avait déjà noté qu'on a un monomorphisme
\[
\mathrm{Div}^{L}_{X/S}\longrightarrow\underline{P}(\underline{Q}) \quad ,
\]
on est ramené à prouver que ce dernier est représentable par une immersion
ouverte quasi-compacte, ce qui nous ramène à prouver que si on se donne
une section de $P(\underline{Q})$ sur $S$ i.e. un Module quotient inversible $\underline{M}$ de $\underline{Q}$ ,
d'où une section $\phi$ de $\underline{L}\otimes\underline{M}$ non nulle sur chaque fibre, alors le \inserted{sous-}foncteur
$\underline{U}$ du foncteur final $S$ sur $(\mathrm{Sch})_{/S}$ ``consistant à rendre $\phi$ transversale-
ment régulière'' est représentable par un ouvert rétrocompact de $S$. Or
le fait qu'il est représentable par un ouvert \unsure{provieht|provient} du fait que $f$ est
propre et que la régularité transversale est une condition ouverte (cf par.
11 ...), la rétrocompacité se voit immédiatement par réduction au cas
noethérien.

\mspage{111}

\clearpage
% Nouveau EGA V (cote 25): archive cote 25, site page 112 (PDF page 113).
% typescript. Draft compared: none.
% AI-transcribed by Claude Opus 5.5 (claude-opus-5-5, adaptive thinking, effort high), Message Batches API, batch ega25b (msgbatch_01A7n5YuXubUPNK1zsepi8RG), 2026-11-13; not yet checked by eye.
\scanpage{25/112}
\mspage{15}

\deleted{\illegible{1}} \margin{\fbox{N$^{\text{os}}$ 3,4}} \hand{3)} \inserted{\uline{Systèmes linéaires de diviseurs, et \unsure{morphismes} dans les fibrés \unsure{projectifs}.}}
\deleted{16. Traductions en langage de systèmes linéaires.}

\deleted{Soit X plat et loc de présentation finie sur S, et T un S-préschéma.}
\deleted{On appelle \uline{famille de} \xout{xxxxxxx} \uline{diviseurs} (sous-entendu: de Cartier)}
\deleted{\margin{\unsure{r}}\uline{positifs}\xout{x}\hand{\unsure{$^{3}$}} \inserted[typed]{\uline{de X/S} \hand{,}} \hand{/} \uline{paramétrée par} S, un sous-préschéma fermé de $X\times_S T$ qui est}
\deleted{\inserted{un diviseur de Cartier rel. à T, i.e. qui est} \hand{(}intersection complète \xout{xxxxxxxxxxx} relative de codimension 1 relativement}
\deleted{à T en chacun de ses points, ou \xout{xxx} ce qui revient au même, qui est plat}
\deleted{\inserted[typed]{loc de prés finie} \hand{(}par rapport à T, et induit sur chaque fibre de $X_{T\,t} = X_{s\,k(t)}$ de $X_T$ sur}
\deleted{T un ``diviseur de Cartier positif'' i.e. un sous-préschéma qui est en}
\deleted{chaque \inserted{de ses} point intersection complète \inserted[typed]{dans $X_t$} de codimension 1. L'ensemble}
\deleted{$\mathrm{Div}_{X/S}(T)$ de telles familles de \inserted[typed]{di}viseurs est, pour T variable, un}
\deleted{foncteur contravar\xout{i}iant à valeurs dans (Ens), noté $\mathrm{Div}_{X/S}$ . Nous verrons}
\deleted{au Chap V qu'il est représentable lorsque $f:X\quad S$ est projectif. Lorsque}
\deleted{X = P(E) fibré projectif défini par E localement libre de type fini,}
\deleted{il est possible avec les seuls résultats du Chap III de prouver que}
\deleted{$\mathrm{Div}_{P/S}$ est représentable par le préschéma somme des $P(\mathrm{Symm}^n(\underline{E}))$ , les}
\deleted{composants de cette somme correspondant aux diviseurs de degré relatif}
\deleted{n , \inserted{\unsure{\S\ n$^{\circ}$2}} Ce résultat n'est pas indispensable pour la suite, et il suffira de}
\deleted{le signaler en remarque; nous l'inclurons peut-être quelque part dans}
\deleted{III, et au plus tard dans V .}

Soit D une famille de diviseurs \inserted[typed]{sur X/S} paramétrée par T (nous sous-enten%
% NOTE: the end of this line is cut off at the right edge of the scan
dons par la suite ``positif''). Un point $x\in X$ est dit ``point fixe'' \xout{pour}pour
cette famille de diviseurs si $pr_1^{-1}(x)\subset D$ \inserted{ensemblistement}, de sorte que l'ensemble des
points non fixes est le complémentaire de $pr_1(\xout{X}X\times_S T - D)$; par suite, si
$T\longrightarrow S$ est universellement ouvert (par exemple plat loc de prés finie)
l'ensemble des points fixes est fermé. On dit que la famille de diviseurs
est ``sans points fixes'' si l'ensemble \inserted{$Z$} des points fixes est vide.

\mspage{112}

\clearpage
% Nouveau EGA V (cote 25): archive cote 25, site page 113 (PDF page 114).
% typescript. Draft compared: none.
% AI-transcribed by Claude Opus 5.5 (claude-opus-5-5, adaptive thinking, effort high), Message Batches API, batch ega25b (msgbatch_01A7n5YuXubUPNK1zsepi8RG), 2026-11-13; not yet checked by eye.
\scanpage{25/113}
\mspage{\hand{16}}

\xout{\unsure{Parxsymétxiex}} Lorsque Z est fermé, alors X-Z est le plus grand ouvert
tel que la famille de diviseurs de X-Z paramétrée par T, induite par la
famille donnée (dans un sens évident) est sans point fixe.

Si la famille D est sans point fixe, et si T est \inserted{plat} loc de prés finie
sur S à fibres $(S_1)$\inserted{et géom irréd}, alors D est aussi un diviseur \deleted{de Cartier}
relativement à X (pour $pr_1 : X\times_S T \longrightarrow X$): en effet, D est défini localemen
% NOTE: the right edge of the page is cut; the end of this line and of several following lines is not visible
\inserted{au point $z \in$ supp D} par une équation \xout{\unsure{f=0}} $\phi = 0$, et l'équation induite sur la fibre
$T_{s\ k(x)}$ du point $x\in X$ (au dessus de $s\in S$) est non \xout{\unsure{identiquement nulle}}
nilpo\xout{\unsure{n}}tente en z (car autrement $pr_1^{-1}(x)\ D = V(\phi_x)$ \xout{\unsure{serait}} \inserted[typed]{contiendrait} ensemblist\xout{\unsure{e}}
ment un voisinage de z dans $T_{\xout{s}\ k(x)}$ , donc contiendrait $T_{s\ k(x)}$ , i.e.
x serait point fixe, ce qui n'est pas);or comme $T_{s\ k(x)}$ est irréd et $(S_1)$
il s'ensuit que $\phi_x$ est $\underline{O}_{T_{s\ k(x)}}$-régulière en z . On obtient ainsi
\xout{\unsure{un morphisme}} une famille de div de T/S paramétrée par X, i.e. un morph
\[ X \longrightarrow \mathrm{Div}_{T/S} \quad , \]
\hand{D}ans le cas général où la famille de diviseurs de X peut avoir des points
fixes, on obtient une famille de diviseurs de T/S paramétrée par X-Z,
i;e. un morphisme
\[ X - Z \longrightarrow \mathrm{Div}_{T/S} \ , \]
en remplaçant dans la définition précédente X pa\xout{r} X-T . D'ailleurs la
\xout{si D est sans points fixes,}

\xout{Notons en passant que la famille de diviseurs (notée ${}^{t}D$ )}
\xout{X-Z}
\xout{de T/S paramétrée par X est} \deleted{évidemment sans points fixes, et on}
\xout{obtient ainsi une correspondance biunivo}\deleted{que entre familles de diviseurs}
\xout{sans points fixes} démonstration ci-dessus montre que X-Z est exactement
le plus grand ouvert \inserted{U} de X tel que $D|(U\times_S T)$ soit un di\xout{vi}seur \deleted{de Cartier}
relativement à U, i.e. \inserted{tel que son symétrique ${}^{t}D$} \unsure{``soit''} une famille de diviseurs de T/S paramétrée
par U/S .

\mspage{\hand{113}}

\clearpage
% Nouveau EGA V (cote 25): archive cote 25, site page 114 (PDF page 115).
% typescript. Draft compared: none.
% AI-transcribed by Claude Opus 5.5 (claude-opus-5-5, adaptive thinking, effort high), Message Batches API, batch ega25b (msgbatch_01A7n5YuXubUPNK1zsepi8RG), 2026-11-13; not yet checked by eye.
\scanpage{25/114}
\mspage{17}

On pourra noter aussi en remarque que lorsque X et T sont tous les
deux plats loc de présentation finie sur S, à fibres $(S_1)$ et géom irréd.,
alors la symétrie $D \rightsquigarrow {}^{t}D$ donne une correspondance biunivoque entre
familles de diviseurs de X/S paramétrées par T qui sont sans points fixes
et familles de diviseurs de T/S paramétrées par X qui sont sans points
fixes. Si dans cet énoncé on désire se débarasser des hypothèses spécifi-%
ques faites sur les fibres de X/S et T/S , il convient de remplacer
les ``points fixes'' par ``points fixes au sens large'' , en entendant par
point fixe au\xout{x} sens large de D un $x \hand{\in} X$ tel que D ne soit pas diviseur
\deleted{de Cartier} relatif \xout{en} \inserted[typed]{par rapport à X en tous} les points de $\mathrm{pr}_1^{-1}(x)$ . Si W est l'ouvert des
points de $X\times_S T$ en lesquels D \xout{n'}est \xout{pas} diviseur \deleted{de Cartier} relatif par
rapport à X, alors l'ensemble des points fixes au s ens large de D est
\xout{\unsure{exemple}} égal à $\mathrm{pr}_1(X\times_S T - W)$; lorsque $T \to S$ est propre, c'est donc
une partie fermée Z' de X, \xout{et on obtient} En tous cas, on obtient une
famille de diviseurs de T/S paramétrée par $X-\overline{Z'}$ . L'hypothèse que les
fibres de T/S soient $(S_1)$ et géom irréductibles \xout{\unsure{assure}} \inserted[typed]{sert} précisément à
assurer que $Z=Z'$ (point fixes \unsure{$=$} points fixes au sens large). Géométri-%
quement, supposant pour simplifier que S=Spec(k), k corps algébriquement
clos (ce qui est loisible, \inserted[typed]{pour T/S plat loc de prés finie,} via changement de base), dire que $x \hand{\in} X(k)$
est point fixe (resp. point fixe au sens large) signifie que $x \hand{\in} \mathrm{supp}\, D_t$
pour tout $t \hand{\in} T(k)$ (resp. qu'il existe un cycle premier T' associé à T
tel que $x \hand{\in} D_t$ pour tout $t \hand{\in} T'(k)$).

\deleted{Le petit sorite précédent n'a guère rien à voir dansle présent para-%
graphe, et je suggérerai d'en faire plutôt un nº au par. 19 à propos des
intersections complètes relatives, sous le titre: diviseurs de Cartier
relatifs, et familles \xout{algébriques} de diviseurs.} \margin{$\|$} Un oubli: \xout{X-Z est}
La formation de l'ensemble des points fixes Z est compatible avec

\mspage{114}

\clearpage
% Nouveau EGA V (cote 25): archive cote 25, site page 115 (PDF page 116).
% typescript. Draft compared: none.
% AI-transcribed by Claude Opus 5.5 (claude-opus-5-5, adaptive thinking, effort high), Message Batches API, batch ega25b (msgbatch_01A7n5YuXubUPNK1zsepi8RG), 2026-11-13; not yet checked by eye.
\scanpage{25/115}
\mspage{18}
changement de base en S; d'autre part, X-Z \inserted[typed]{(supposé ouvert, p ex T/S plat lpf.}\inserted{lorsque $T \to S$ est surjectif\unsure{)}} est universellement
schématiquement dense dans X rel à S. Ce dernier fait résulte du \xout{la compa-}
\xout{tibilité avec changement de base} par. 11 et du fait que pour tout $s \in S$,
$Z_s$ ne contient aucun point de $X_s$ associé à $\underline{O}_{X_s}$ (en effet, le support d'un
diviseur \deleted{de Cartier} sur $X_s$ ne contient aucun tel point).

Dans le cas où T \deleted{devient}\inserted{est} un fibré projectif $Q=P(F)$, $\mathrm{Div}_{Q/S}$
étant représentable par la somme des $P(\mathrm{Sym}^n(\check{F})) = P(n)$, on trouve
un morphisme $X-Z \longrightarrow P(n)$ . On dit\xout{\unsure{a}} que la famille de diviseurs
D est \uline{de degré} n, si le morphisme précédent se factorise par P(n); si
$X \neq \emptyset$ , d'où $X-Z \neq \emptyset$ , le n en question est bien déterminé par D. D'ail%
leurs pour définir cette notion, on n'a pas besoin à la rigueur du résul-%
tat de représentabilité énoncé\xout{s} plus haut, mais seulement d'avoir défini
les monomorphismes canoniques
\[ P(n) \longrightarrow \mathrm{Div}_{\check{P}/S} \]
(NB on pose $P = P(\check{F})$, d'où $Q = P(F) = \check{P}$ avec les notations des n°s précé-%
dents). \deleted{Cela pourra être fait à titre d'exemple dans le n° préconisé au
par. 19}. On appelle \uline{système linéaire de diviseurs sur} X/S , \uline{paramétré
par le fibré projectif} $Q = \check{P}$ , \uline{toute famille de diviseurs} \uline{sur} X/S
\uline{paramétrée par} Q, \uline{qui soit de degré} 1, i.e. \uline{définisse}
\[ \hand{f} : X-Z \longrightarrow P \qquad . \]
Ainsi, à un tel système \xout{\unsure{li}} linéaire de diviseurs\inserted{(et lorsque les fibres de $\check{P}$ sont tous $\neq \emptyset$)}, est associé canoniquement
une application rationelle de X dans un fibré projectif,\inserted{savoir un ``\deleted{\illegible{1}} \unsure{pseudo-morphisme} rel.: S''} \hand{(}en fait, un peu
mieux même qu'une application rationelle\hand{)}. Par construction même,
$D\,|(X-Z)x_S\check{P}$ n'est autre que l'image inverse par $f x id_{\check{P}}$ du diviseur
canonique (le diviseur d'in\unsure{cid}\xout{\unsure{e}}ence )\hand{\uline{H}} sur $Px_S\check{P}$ , Donc la connaissance
de $f: X-Z \longrightarrow P$ permet de réconstituer \xout{la fa} tout au moins la famil-%
le de diviseurs de X-Z induite par D , - donc si la famille est sans
\mspage{115}

\clearpage
% Nouveau EGA V (cote 25): archive cote 25, site page 116 (PDF page 117).
% typescript. Draft compared: none.
% AI-transcribed by Claude Opus 5.5 (claude-opus-5-5, adaptive thinking, effort high), Message Batches API, batch ega25b (msgbatch_01A7n5YuXubUPNK1zsepi8RG), 2026-11-13; not yet checked by eye.
\scanpage{25/116}
\mspage{19}

points fixes, elle est déterminée \xout{xxixxxxxxx} par le morphisme f:X $\longrightarrow$ P
associé. Noter d'ailleurs qu'on \xout{xx}obtient évidemment ainsi une correspon
dance biunivoque entre systêmes linéaires de diviseurs de X/S paramétré
par $\hand{\check{P}}$ , et morphismes f:X $\longrightarrow$ P tel que \xout{xx}$(fx_S\mathrm{id})^{-1}(\underline{H})$ soit un
diviseur\deleted{ de Cartier} \xout{xxx} relatif sur $Xx_S\check{P}$ par rapport à $\check{P}$ . Ce\over{\unsure{ci}}{tte} \deleted{se} \inserted{condition}
\inserted{se} vérifie aussi fibre par fibre, et on obtient :

\uline{Proposition} \over{\unsure{1}}{3}.1.On a une correspondance biunivoque entre les \xout{morphismes}
systêmes linéaires \inserted{sans points fixes} de diviseurs sur X/S paramétrés par $\hand{\check{P}}$ , et les morphis
mes f: X $\longrightarrow$ P ayant la propriété suivante: pour tout $s\in S$, d\over{e}{é}signant
par $\bar{k}$ une clôture algébrique de k(s), et pour tout cycle premier associé
X' de $X_{\bar{k}}$ , $f(X')\subset P_{\bar{k}}$ n'est contenu dans aucun hyperplan de $P_{\bar{k}}$ .
\xout{(xx}\hand{(}Lorsque X est à fibres géométriquement \xout{xxxxxxx} \inserted[typed]{intègres}, ceci \xout{xexsimplifie}
\xout{xx} s'énoncé simplement en disant que $f(X_{\bar{k}})$ n'est contenu dans aucun
hyperplan de $P_{\bar{k}}$ \hand{)}.

En général (i.e. \xout{xx} lorsque Z $\neq \emptyset$ ) on ne peut plus affirmer que
la connaissance de f détermine \xout{D}la famille de diviseurs. Le cas le plus
trivial est celui où $\hand{\check{P}}$ = S \inserted{est de dim relative 0}, alors la donnée d'un systême linéaire de\hand{.}
diviseurs de X/S paramétrée par S équivaut à celle d'un diviseur de
Cartier relatif D sur X par rapport à S; \xout{Txxxxlxxppx} le morphisme
associé est la projection X-D $\longrightarrow$ S , et on voit \xout{(mêmexxxx} que la connais
sance de ce morphisme (y compris celle de son domaine de définition)
ne détermine pas D. Dans ce cas également, supposant pour simplifier X
réduit, le domaine de définition de \xout{xxxmorphisme} f \xout{xlxsxxpas} considéré
comme application rationelle de X dans P = S n'est pas X-D , mais X.
Pour éliminer ce genre de phénomènes déplaisants, il y a lieu de se bor-%
ner aux systêmes \xout{h}linéaires de diviseurs \uline{``sans composante fixe''} . De
façon générale, si S=Spec(k), \unsure{étaht|étant} donné une famille (non néc linéaire)

\margin{116}

\clearpage
% Nouveau EGA V (cote 25): archive cote 25, site page 117 (PDF page 118).
% typescript. Draft compared: none.
% AI-transcribed by Claude Opus 5.5 (claude-opus-5-5, adaptive thinking, effort high), Message Batches API, batch ega25b (msgbatch_01A7n5YuXubUPNK1zsepi8RG), 2026-11-13; not yet checked by eye.
\scanpage{25/117}
\mspage{20}

de diviseurs de $X/S$ paramétrée par $T$, on appelle ``composante fixe'' de la
famille toute composante irréductible de dodimension 1 de l'ensemble $Z$ de\unsure{s}
points fixes de la famille; on dit que la famille est ``sans composante
fixe'' si elle est sans composante fixe, i.e.\ si $\operatorname{codim}(Z,X) \hand{$\geq$} 2$. Cette
terminologie s'étend aussitôt au cas où $S$ est quelconque, en regardant
fibre par fibre. Le fait d'être sans composante fixe est évidemment stabl\unsure{e}
par changement de base.

\uline{Proposition}\over{2}{\illegible{1}} \hand{/}Supposons $X \longrightarrow S$ plat loc de présentation finie, à fibres
$(S_2)$ , et soit $D$ un syst\xout{e}ême linéaire de diviseurs \uline{sans composante fixe}
sur $X/S$, paramétrée par $\overset{\hand{$\vee$}}{P}$ . Alors $D$ est uniquement déterminé par la
connaissance du morphisme $f: X - Z \longrightarrow P$ correspondant ($Z$ = ensemble
des points fixes), et même par la connaissance de la classe de $f$ comme
\hand{``}\deleted{application rationelle}\inserted{pseudo-morphisme} \xout{aux dessus de S}\inserted[typed]{rel à S}\hand{''} ; $X-Z$ est le \xout{plus grand ouvert}\inserted[typed]{domain\xout{x}e}
de définition de ladite classe.

Pour la notion et le sorite des ``\deleted{applications rationelles}\inserted{pseudo-morphismes}
rel à S'' voir \deleted{SGAD XVIII n°1}\inserted{\S 20.10.} , \deleted{qu'on pourrait éventuellement rajouter
au dernier n° du par.11 sur la densité schématique.} Il faut prouver que.
si \xout{U est un ouvert de X schématiquement relativement dense} $D'$ est un\deleted{e}
autre syst lin de diviseurs sans comp fixe paramétrée par $\overset{\hand{$\vee$}}{P}$ , \xout{xxx}
définissant $f': X-Z' \longrightarrow P$, et si $f$ et $f'$ coincident dans un ouvert $U$
$\hand{$\subset$} (X-Z)\hand{$\cap$}(X-Z')$ schématiquement dense relativement à $S$, alors $D=D'$ .
En fait, comme $P$ est séparé sur $S$, on peut prendre $U=(X-Z)\hand{$\cap$}(X-Z') =$
\xout{X - (Z u Z') = d'autre part} $X-Z''$ , où $Z'' = Z\hand{$\cup$}Z'$ . Comme $Z''$ est de codim$\hand{$\geq$}2$
sur chaque fibre et que $X$ est à fibres $(S_2)$, il s'ensuit que pour tout
$x\hand{$\in$} Z''$, la fibre $X_s$ est de profondeur $\hand{$\geq$}2$ en $x$. On va en conclure (utili-%
sant de plus bien sûr le fait que $X$ est plat loc de présentation finie
sur $S$) que tout diviseur \deleted{de Cartier positif} sur $X$ (pas nécessairement \margin{\unsure{cf}\ \S 20.11.}
transversal aux fibres) est connu quand on connait sa restriction à $X-Z''$

\mspage{117}

\clearpage
% Nouveau EGA V (cote 25): archive cote 25, site page 118 (PDF page 119).
% typescript. Draft compared: none.
% AI-transcribed by Claude Opus 5.5 (claude-opus-5-5, adaptive thinking, effort high), Message Batches API, batch ega25b (msgbatch_01A7n5YuXubUPNK1zsepi8RG), 2026-11-13; not yet checked by eye.
\scanpage{25/118}
\mspage{\hand{21}}

ce qui entrainera la conclusion voulue. Soit $\underline{J}$ l'Idéal qui définit
D , il suffira évidemment de prouver que $\underline{J}\to i_{\hand{*}}(\underline{J}|X-Z")$ est un isomor\xout{p}
phisme ($i:X-Z"\to X$ désignant l'immersion canonique), alors l'homomorphis
me $\underline{J}\to\underline{O}_X$ se reconstruit en effet en appliquant le foncteur $i_{\hand{*}}$ à
$\underline{J}\,|\,X-Z" \longrightarrow \underline{O}_X\,|\,X-Z"$ . Or comme $\underline{J}$ est inversible, il est plat sur S \unsure{et}
\xout{fibresxde} $x\in Z" \Rightarrow \mathrm{prof}_x\,\underline{J}_{s\,\hand{\unsure{*}}} \geqslant 2$. Il suffit donc de prouver le
\uline{Lemme} \hand{\unsure{3.2|3.1}}. Soit $X\to S$ loc de présentation finie, F un Module de présentation
finie sur X plat rel à S, T une partie fermée de X, supposons que pour
\margin{\unsure{$\downarrow$ \S 20 !!}}
tout $x\in X$ au dessus de $s\in S$, on ait \hand{\unsure{/}} $\mathrm{prof}_x F_{s\,\hand{\unsure{*}}}\geqslant 1$ (resp. $\mathrm{prof}_x F_s \geqslant 2$).
Alors l\xout{e}'homomorphisme canonique $F\to i_{\hand{*}}(F|X-T)$ est injectif (resp.
bijectif), où $i:X-T\to T$ est l'immersion canonique.

Démonstration du lemme: on peut supposer S,X affines, et on se
ramène par le procédé bréveté au cas S noethérien. Alors l'hypothèse
implique, en vertu du par.6 , que l'on a $\mathrm{prof}_x F \geqslant \mathrm{prof}_x F_s$ pour tout .
$x\in X$ sur $s\in S$, donc $\mathrm{prof}_x F_s \geqslant 1$ (resp.$\geqslant 2$) si $x\in T$. On conclut alors
grâce au par.5 . (NB Pour bien faire, ce lemme devrait figurer au par.
11 sous la rubrique; élimination des hypothèses noethériennes ...).

Il reste enfin à vérifier la dernière assertion\xout{s} de prop.2, savoir
que \xout{\unsure{Z}}\inserted{X-Z} est exactement le domain\xout{e} de définition de l"application rationel-%
le rel à S définie par f. Soit $U\supset X-Z$ ce domaine, il résulte alors de
la proposition 1 que $U\longrightarrow P$ est associé à un\xout{xfamillex}\inserted[typed]{système lin.} de diviseurs
D' sur U\unsure{/}S paramétré par $\check{P}$ , et on a $D'|(X-Z)\times_S\check{P} = D|(X-Z)\times_S\check{P}$ . Appli
quant le résultat d'unicité \hand{(}déjà démontré\hand{)}à D' et $D|U\times_S\check{P}$ , on voit donc
que ces deux derniers diviseurs sont égaux, donc $D|U\times_S\check{P}$ n'a pas de
point fixes i.e. $Z\cap U=\emptyset$ donc U=X-Z , cqfd.

\mspage{\hand{118}}

\clearpage
% Nouveau EGA V (cote 25): archive cote 25, site page 119 (PDF page 120).
% typescript. Draft compared: none.
% AI-transcribed by Claude Opus 5.5 (claude-opus-5-5, adaptive thinking, effort high), Message Batches API, batch ega25b (msgbatch_01A7n5YuXubUPNK1zsepi8RG), 2026-11-13; not yet checked by eye.
\scanpage{25/119}
\margin{\unsure{2/2}}
\margin{(trait vertical au crayon dans la marge gauche, le long du premier paragraphe)}
Je me répens d'avoir donné la proposition \deleted{\unsure{3}} sous une forme éculée et
pisseuse, à mi chemin entre  les hypothèses classiques et les hypothèses
naturelles, et sans donner une récip\unsure{r}oque. Aussi je propose l'énoncé:
\uline{Proposition\unsure{\hand{3}}3. \deleted{revue et corrigée}}.  Soit $X \to S$ plat loc de présentation
finie,  $Q=\check{P} = P(\check{E})$  un fibré \xout{\unsure{vectoriel sur}} projectif sur $S$ défini par
un Module loc  libre de type fini $F = \check{E}$ ,  considérons l'ensemble
$\emptyset$  des systêmes linéaires \inserted{$D$ \deleted{\unsure{E}}} de diviseurs \xout{\unsure{sur}} sur $X$, paramétrés par $Q$, tels
que l'ensemble  $Z$ des points fixes de $D$  satisfasse la propriété: $z \in Z$
implique prof$_z \underline{O}_{X_s} \geqslant 2$  (où  $s$ \xout{\unsure{$= g(z)$)}} est l'image de $z$ dans $S$). \xout{\unsure{Alors}}
\xout{\unsure{$D \in \emptyset$ associons le morphisme correspondant $f$ $X$-$Z$ $P$ considérée}}
\xout{\unsure{considérée comme}}  Soit $\emptyset$ l'ensemble des  \deleted{``applications rationelles}\inserted{pseudo-morphismes} rel à
$S$ \deleted{\unsure{*}}\inserted{$f$} de $X$ dans  $P$, telles que le\xout{\illegible{1}} domaine de définition \inserted{$U=$}\over{dom}{\unsure{U}}(f)
\xout{\unsure{f satisf}} satisfasse la condition : $z \in X$-$U$\deleted{(f)} $\Longrightarrow$ prof$_z \underline{O}_{X_s} \geqslant 2$\inserted{$\surd$, et $f_U = f|U : U \to \check{P}$ satisfait la condition de non dégénérescence envisagée dans 3.1.}. Considé\xout{rons}
rons l'application naturelle  $D \rightsquigarrow f_D$  de $\emptyset$  dans $\emptyset$ .  Alors :

a)  Cette application est injective,  et  pour $D \in \emptyset$ , l'ensemble des
points fixes $Z$  n'est autre que l'ensemble  complémentaire du domaine de
définition \xout{\unsure{Dif}} $U$\deleted{(f)} de $f$\hand{$_D$}.

b)  \xout{\unsure{Pour qu'une f $\emptyset$ appartienne à provienne d'une $D$ $\emptyset$, il}} Soit
$f \in \emptyset$ , et soit $U$ un ouvert sur lequel $f$ soit défini, et tel que $z \in X$-$U$
implique prof $\underline{O}_{X_s,z} \geqslant 2$ \xout{\unsure{et existe}} (par exemple $U = U(f)$ domaine de
définition de $f$). Pour que $f$ provienne d'une $D \in \emptyset$ , il f et suffit  que\hand{\unsure{$\surd$}}
\deleted{\unsure{$f_U$ satisfasse la condition énoncée dans 3.1. et que}}
posant $\underline{L}_U = f_U^*(\underline{O}_P(1))$  (où  $f_U$: $U \to P$  est le morphisme induit par $f$),
\deleted{\unsure{$f_U$ satisfasse la condition énoncée dans 3.1 et}}
\xout{\unsure{que}} le Module  $i_*(\underline{L}_U)$  sur $X$ soit inversible (où $i$:$U \to X$ est l'immer-%
sion canonique).

On remarquera que lorsque  les fibres de $X$ sur $S$ satisfont à $(S_2)$,
par exemple sont normales, voire géométriquement  normales, alors la
condition de profondeur  envisagée sur  une partie fermée $Z$ de $X$ dans \deleted{\illegible{1}}
la proposition  signifie simplement que pour tout $s \in S$,

\mspage{119}

\clearpage
% Nouveau EGA V (cote 25): archive cote 25, site page 120 (PDF page 121).
% typescript. Draft compared: none.
% AI-transcribed by Claude Opus 5.5 (claude-opus-5-5, adaptive thinking, effort high), Message Batches API, batch ega25b (msgbatch_01A7n5YuXubUPNK1zsepi8RG), 2026-11-13; not yet checked by eye.
\scanpage{25/120}
\mspage{\unsure{27}}

Z$_s$ est de codimension $\geq 2$ dans X$_s$ ; \xout{\illegible{2}\unsure{(comme dans les cas envisa}}
\xout{familles} \inserted{syst lin} de diviseurs par.par P qui so\inserted{\unsure{ie}}nt
\uline{sans composante fixe}, \xout{tandis que $\emptyset$ est l'ensemble de} D'autre part, si
S=Spec(k), et si X est normal, alors \inserted{pour} \uline{toute} application rationelle
f:X$\to$P, l'ensemble de définition U(f) vérifie codim ((X-U(f),X)$\geq$2
(II 7), donc dans ce cas $\emptyset$ est formé de l'ensemble de \uline{toutes} les
applications rationelles de X dans P.

La démonstration de a) a déjà été donnée. Pour prouver b), notons
que la formation de $i_*(\underline{L}_U)$ commute avec toute extension plate de la
\xout{pour le\unsure{x} il faut} \inserted[typed]{la cond. envisagée est inv. par}
base S' \quad S (réf)\inserted{\unsure{du moins si $U\to X$ est quasi-compact, cas auquel on se ramène sans difficultés}}, donc \xout{il suffit de trouver un morphisme} fid plat qu cpc
\xout{S' $\quad$ S tel que la conclusion soit vraie après} changement de base. On
prendra S' = $\hand{\check{}}\mathrm{P}$ , et on note que l'hypothèse $i'_*(\underline{L}'_U)$ inversible ne
change pas en remplaçant $\underline{L}'_U$ par $\underline{L}'_U\otimes_{S'}M'$, où M' est un Module inversible
sur S' (car $i'_*(\underline{L}'_U\otimes_{S'}M')\simeq i'_*(\underline{L}'_U)\otimes_{S'}M'$ ). On \unsure{prèndra|prendra} M'= $\underline{O}_{\check{P}}(1)$, de
sorte que la condition envisagée signifie aussi que $i'_*(\underline{N}')$ est un Module
inversible, où $\underline{N}'= (f_U\times_S id_P)^*(\underline{O}_{P\times\check{P}}(1,1))$ . Or $\underline{O}$ (1,1) est précisément
le Module inversible défini par le diviseur canonique $\underline{H}$ sur $P\times_S\check{P}$ , de sor%
te que $\underline{N}$ n'est autre que le Module inversible défini par \qquad D' =
$(f_U\times id_P)$ $(\underline{H})$. \deleted{\xout{\illegible{1}}} Lorsque f provient de D$\in\emptyset$ , \xout{ce dernier}
D' n'est autre que D$|$U$\times_S\check{P}$ , \xout{et il résulte du lemme} donc $\underline{N}' = \underline{N}|U\times_S\check{P}$ ,
où $\underline{N}$ est le Module inversible sur $X\times_S\hand{\check{}}P$ défini par D, et il résulte
du lemme \inserted{3.2} plus haut \xout{\illegible{1}} appliqué à $P\times_S\hand{\check{}}P\to\hand{\check{}}P$ qu'on aura
$i'_*(\underline{N}')\simeq\underline{N}$ , donc $i'_*(\underline{N}')$ est inversible. Inversement, si cette condi-%
tion est remplie prouvons que f provient d'un\xout{e} D$\in\emptyset$ , ou ce qui revient
ma\xout{n}nifestement au même, que le diviseur D' se prolonge en un diviseur \deleted{de}
\deleted{Cartier} relatif par r apport \xout{de} à $\hand{\check{}}P$ sur X$\times_S$P . Il \deleted{suffit} revient au même
de dire qu'il se prolonge en un diviseur \deleted{de Cartier} D sur X$\times_S$P , car \hand{D}.

\margin{\unsure{Rq compliquée} \unsure{Il faut expliquer} \illegible{2} \unsure{$f\times\mathrm{id}$} \unsure{si $p^{-1}(H)$} \unsure{est défini} \illegible{1} \unsure{diviseur rel.} \unsure{à la condition} \illegible{3} \unsure{diviseur relatif} \unsure{$X\times P$} \unsure{/P} \illegible{2}}

\mspage{120}

\clearpage
% Nouveau EGA V (cote 25): archive cote 25, site page 121 (PDF page 122).
% typescript. Draft compared: none.
% AI-transcribed by Claude Opus 5.5 (claude-opus-5-5, adaptive thinking, effort high), Message Batches API, batch ega25b (msgbatch_01A7n5YuXubUPNK1zsepi8RG), 2026-11-13; not yet checked by eye.
\scanpage{25/121}
\mspage{24}

sera automatiquement \deleted{de Cartier}\inserted{diviseur} relatif pour la base $\overset{\hand{\vee}}{P}$ , comme il
résulte du fait que U contient les éléments associés au $\underline{O}_{X_s}$ , s $\in$ S,
condition qui est stable par tout changement de base et en particulier
par S'=$\overset{\hand{\vee}}{P} \longrightarrow$ S. Or il résulte aussitôt du lemme \inserted{4.2.} plus haut que D' se
prolonge en un diviseur \deleted{de Cartier} D \uline{sss} $\underline{N}'$ se prolonge en un Module
inversible, ou encore $i'_{\hand{*}}(\underline{N}')$ est inversible. \deleted{Donc} \over{i}{I}l faudrait rédiger
mieux la fin de la démonstration en termes de cond néc et suff (sans y
aller en deux fois comme je l'ai fait), et dégager aussi \inserted{d'abord} le
\margin{\unsure{voir aussi} \S 20.\illegible{1} et 20.11.}\uline{Corollaire} \inserted{3.4.} \uline{au lemme}\hand{3.2}. Supposons g:X$\to$S plat loc de présentation finie,
soit Z une partie fermée de X telle que x $\in$ Z implique prof$_x$ $\underline{O}_{X_s} \geq 2$
\inserted[soit U=X-Z et i:U$\to$X l'immersion c an.]
(où s = g(x)), \xout{Soit $\underline{L}$ un Module loc libre sur U = X-Z, alors $\underline{L}'$}
\xout{est isomorphe à la restriction} Pour tout Module loc libre \inserted[de t.f] $\underline{L}$ sur X,
considérons sa restriction $\underline{L}' = \underline{L}|U$ . Alors \inserted[a)] le foncteur $\underline{L} \rightsquigarrow \underline{L}'$ est
pleinement fidèle, et pour \xout{\unsure{quelque} Module loc libre $\underline{L}$} tout $\underline{L}$, \xout{xxxx}
l'hom canonique $\underline{L} \to i_{\hand{*}}(\underline{L}')$ est un isomorphisme. \inserted[Pour que $\underline{L}$ soit de rang n, il f et s que $\underline{L}'$ le soit.] \hand{(} b) Soit $\underline{L}'$ un
Module loc libre sur X, alors $\underline{L}'$ est isomorphe à la restriction d'un
Module loc libre $\underline{L}$ sss $i_{\hand{*}}(\underline{L}')$ est loc libre. c) Supposons que $\underline{L}'$ soit
le Module inversible associé à un diviseur \deleted{de Cartier} D'\hand{/}sur U. Alors
la condition envisagée dans b) est aussi néc et suffisante pour que D'
soit la restriction d'un diviseur \deleted{de Cartier} D sur X, \inserted[lequel sera alors unique,]\inserted{et égal à l'adhérence schématique de D' dans X.} \hand{|} Pour que D' soit
diviseur \deleted{de Cartier} relatif par rapport à S, il f. et s. \inserted{\unsure{que}} D le soit.

On utilise simplement le fait que tout $\underline{L}$ satisfait aux \xout{xxx}
hypothèses énoncées pour $\underline{F}$ dans le lemme\hand{3.2}...

\uline{Corollaire} \inserted{3.5.} Supposons \inserted{les anneaux locaux de X factoriels, (par exemple} X régulier\hand{)} \deleted{(par exemple S régulier, et X lisse sur
S)}, Alors l'application $\emptyset \to$ \unsure{$\emptyset'$|$\emptyset$} est bijective. En particulier, si X est
un préschéma régulier loc de type fini sur un corps k, et P un fibré
projectif sur k, il y a correspondance biunivoque entre l'ensemble $\emptyset$ des

\margin{121}

\clearpage
% Nouveau EGA V (cote 25): archive cote 25, site page 122 (PDF page 123).
% typescript. Draft compared: none.
% AI-transcribed by Claude Opus 5.5 (claude-opus-5-5, adaptive thinking, effort high), Message Batches API, batch ega25b (msgbatch_01A7n5YuXubUPNK1zsepi8RG), 2026-11-13; not yet checked by eye.
\scanpage{25/122}
\mspage{25}

systêmes linéaires de diviseurs sans composante fixe  sur X paramétrés
par $\overset{\hand{\vee}}{P}$ , et l'ensemble Ø  des applications rationelles de X dans P, \hand{qui}
\inserted{sur $\bar{k}$ ne se factorisent par aucun hyperplan de $P_{\unsure{\bar{k}}}$.}
\inserted[typed]{lorsque les anneaux locaux de X sont factoriels.}

En effet, \deleted{\xout{X étant régulier, X en est de même de \unsure{X $X_S$} qui est lisse}
\xout{sur X, donc d'après un résultat de Auslander-Buchsbaum (qui sera repris}
\xout{au Chap III) \illegible{1} de les anneaux locaux de X\unsure{x}P sont factoriels,}
et comme codim(X-U,X) \quad 2} , il s'ensuit que tout Module inversible sur
U peut se prolonger en un Module inversible sur X, donc la condition
envisagée dans b) est automatiquement satisfaite. D'autre part, d'après
\deleted{\unsure{un Mémoire inédit lorsque n° ultérieur}}
Auslander-Buchsbaum (\deleted{qui figurera un jour au Chap III}\inserted{cf §20}) un annea\hand{/}u local
régulier est factoriel.

\deleted{La proposition 2  b)  nous montre que myennant une hypothèse}
\deleted{\inserted[typed]{(savoir D \quad Ø)}}
\deleted{supplémentaire sur un\xout{e famille} système linéaire de diviseurs D, on peut}
\deleted{lui associer canoniquement un Module inversible $\underline{L}$ sur X, prolongement}
\deleted{\inserted[typed]{où U=X-Z.}}
\deleted{canonique de $\underline{L}^{x}_{U} = f\ (\underline{O}_P(1))$ ,  On peut en donner une autre interprétation,}
\deleted{indépendante de l'hypothèse  D \quad Ø , de la façon suivante; par constructio}
\deleted{\inserted[typed]{can isomorphe à} \hand{)}}
\deleted{\xout{l'image inverse de $\underline{L}$x} ou presque $\underline{L}^{x}_{U}\otimes_S\underline{O}_P(1)$ \xout{n'}est \xout{autre que} la restric-%}
\deleted{tion $\underline{N}'$ à $U\mathrm{x}_S\overset{\hand{\vee}}{P}$  du  Module inversible $\underline{N}$  sur $X\mathrm{x}_S\overset{\hand{\vee}}{P}$  défini par D. En vertu}
\deleted{du \xout{\unsure{n°1}} n°1, cela détermine \xout{uniquement $\underline{N}$, et} $\underline{L}_U$  à isom unique près;}
\deleted{et  (comme U est dense dans X)  il existe  un Module inversible $\underline{L}$ sur X,}
\deleted{déterminé à isom unique près (do\unsure{ç}c \unsure{prolongeant|prolongeaant} $\underline{L}_U$ ) tel que  l'on ait}
\deleted{un isomorphisme  $\underline{L}\otimes_{\underline{O}_S}\underline{O}_{\unsure{P}}\ (1) \quad \underline{N}$ ;  $\underline{L}$  est donné par exemple par}
\deleted{\[ \underline{L} \quad pr_1\ (\underline{N}) \quad ; \]}
\deleted{on l'appelle  le Module inversible sur X associé au systême lin de divi-%}
\deleted{suurs D.}

\mspage{122}

% COMPILE FIX (combined file only): 2 closing brace(s) added; the page's LaTeX leaves them open
}}

\clearpage
% Nouveau EGA V (cote 25): archive cote 25, site page 123 (PDF page 124).
% typescript. Draft compared: none.
% AI-transcribed by Claude Opus 5.5 (claude-opus-5-5, adaptive thinking, effort high), Message Batches API, batch ega25b (msgbatch_01A7n5YuXubUPNK1zsepi8RG), 2026-11-13; not yet checked by eye.
\scanpage{25/123}
\mspage{26}

\over{7}{4}. \uline{Systèmes linéaires de diviseurs et Modules inversibles}.

Utilisant le résultat du nº1, nous allons donner une description
complète des systèmes linéaires sur X en termes des faisceaux inversibles
sur X. \inserted{Nous pouvons supposer évidemment $\check{P}\to S$ surjectif. Alors $X\times_S\check{P}\to \unsure{\check{P}}$ l'étant.} \deleted{En effet,} d'après les généralités du \over{nº2}{\unsure{\S 20.3}} la donnée d'un
diviseur \deleted{de Cartier} D sur $X\times_S\check{P}$ \xout{(fibré projectif sur X)} revient à la
donnée d'un Module inversible $\underline{N}$ sur $X\times_S\check{P}$ , et d'une section régulière
$\phi$ de ce dernier. \xout{\unsure{En ver}} L'hypothèse que D est un\xout{e} système linéaire de
diviseurs sur X paramétrée par $\check{P}$ s'exprime alors par les deux conditions
1º) les $\phi_t$ ($t\in\check{P}$) induits par $\phi$ sur les fibres de $Xs_S\check{P}$ au dessus de $\check{P}$
sont réguliers (ce qui entraine que $\phi$ est régulier), et 2º) les $\underline{N}_x$ ($x\in X$)
induits par $\underline{N}$ \xout{xxxx} sur les fibres des $X\times_S\check{P}$ sur X sont \uline{de degré} 1 .
Or la donnée d'un $\underline{N}$ inversible sur \unsure{le} fibré projectif $X\times_S\check{P}$ sur X,
satisfaisant à la condition 2º) ci-dessus, équivaut en vertu du nº1
à la donnée d'un Module inversible $\underline{L}$ sur X, $\underline{N}$ étant déterminé en fonction
de $\underline{L}$ par
\[ \underline{N}\simeq\underline{L}\otimes_{\underline{O}_{\unsure{S|X}}}\underline{O}_{\check{P}}(1) \quad , \]
et $\underline{L}$ étant d'ailleurs déterminé en termes de $\underline{N}$ par $\underline{L}\simeq pr_{1*}(\underline{N}(-1))$,
où (-1) désigne \inserted{\unsure{la|le}} tensorisation par $\underline{O}_{\check{P}}(-1)$ sur $\underline{O}_S$ . La donnée de $\phi$
revient à la donnée d'une section de $\underline{L}\otimes_{\underline{O}_S}\underline{O}_{\check{P}}(1)$ , i.e. d'une section
\margin{$\underline{O}_{X\times_S\check{P}}(1)$}
de $pr_{1*}(\underline{L}\otimes_{\underline{O}_X}\underline{O}_{\hand{$X\times_S\check{P}$}}(1))$ = $\underline{L}\otimes_{\underline{O}_S}\unsure{pr}_*(\underline{O}_{\hand{$\check{P}$}}(1))$, \deleted{\xout{xx$\underline{L}$x$\otimes$xx$\underline{E}$xxxi.e.xxd'unxhomomorphisme}}
\deleted{\xout{$\underline{E}$x$\underline{E}$xxxx$\otimes$}} or en vertu de III 2 on a $\unsure{pr}_*(\underline{O}_{X\times_S\check{P}}(1))\simeq\check{\underline{E}}_X$ , donc
la donnée de $\phi$ équivaut à \xout{xxxx} la donnée d'un homomorphisme $g^*(\underline{E})\to\underline{L}$,
où ce qui revient au même, d'un morphisme \inserted{u:} $\underline{E}\to g_*(\underline{L})$ (NB $g:X\to S$ est
la projection canonique). Il reste à exprimer la condition 1º) plus
haut en termes de u. Comme les constructions faites commutent avec le
changement de base, il suffit d'exprimer cette condition fibre par fibre\hand{,}
et compte tenu que les points de $\check{P}$ à valeurs dans une extension k de
k(s) correspondent exactement aux droites de $E(s)\otimes_{k(s)}k$, cette

\mspage{123}

\clearpage
% Nouveau EGA V (cote 25): archive cote 25, site page 124 (PDF page 125).
% typescript. Draft compared: none.
% AI-transcribed by Claude Opus 5.5 (claude-opus-5-5, adaptive thinking, effort high), Message Batches API, batch ega25b (msgbatch_01A7n5YuXubUPNK1zsepi8RG), 2026-11-13; not yet checked by eye.
\scanpage{25/124}
\mspage{27}

condition e'exprime simplement en exigeant que pour tout $t\in E(s)$,
la section correspondante u(s) de $\underline{L}_s$ sur $X_s$ soit régulière, et que la
condition analogue soit vérifiée après toute extension du corps de base.
On voit comme d'\unsure{\over{ha}{h}}bitude qu'il suffit de tester cette condition sur une
clôture algébrique de k. En résumé:
\uline{Proposition} \over{5}{4}.1. Soit $g:X\to S$ un morphisme plat loc de présentation
finie, $P=P(E)$ un fibré projectif sur X défini par $\underline{E}$ loc libre de type
fini\inserted{partout $\neq 0$ i.e. $P$ à fibres non vides}, $\overset{\hand{\vee}}{P} = P(\overset{\hand{\vee}}{E})$ \hand{,} Alors il y a une correspondance biunivoque entre
l'ensemble des systêmes linéaires de diviseurs sur X/S paramétrés par $\overset{\hand{\vee}}{P}$ ,
et l'ensemble des couples(à isomorphisme près) $(\underline{L},u)$, où $\underline{L}$ est un
Module inversible sur X et u: $\underline{E}\to g_*(\underline{L})$ un homomorphisme, tel que
pour tout $s\in S$ et tout point t de $E(s)\otimes_{k(s)}k$ (pour une extension quel-%
\margin{$\cdot$} conque\inserted{$k$} de k\unsure{(s)}, qu'on peut supposer être la clôture algébrique de k(s)),
la section correspondante u(t) de $\underline{L}_{s\,k}$ sur $X_{s\,k}$ soit régulière.

On notera que lorsque les fibres de X sont \xout{($S_1$) et géom irréd}
géométriquement intègres, cette condition sur u signifie simplement que
pour tout $s\in S$, $u(s): E(s)\rightarrowtail H^{0}(X_s,\underline{L}_s)$ est injectif, fait qu'il faud-%
rait également expliciter dans \over{5}{4}.1. ; \deleted{\unsure{et}} il faudrait également rappeler
(pour la commodité de la référence) la construction du diviseur D en
termes de $(\underline{L},u)$ comme diviseur de la section évidente \inserted{$\phi$} de $\underline{L}\otimes_{\underline{O}_S}\underline{O}_{\unsure{\overset{\hand{\vee}}{P}}}(1)$
définie par u.
\uline{Corollaire} \over{5}{4}.2. \xout{Soit L} Supposons $f: X\longrightarrow S$ propre, plat, de présentati-%
on finie et à fibres géométriquemeht intègres, \xout{et} soit $\underline{L}$ un Module
inversible sur X, et $P=P(\underline{E})$ un fibré projectif sur S comme dans 5.1.
\xout{Soit enfin} Il existe un Module Q de présentation finie sur S et un
isomorphisme de foncteurs en le $\underline{O}_S$-Module quasi-cohérent F :
\[ \mathrm{Hom}\,(Q,F) \xrightarrow{\ \sim\ } g_*(\underline{L}\otimes_{\underline{O}_S}F) \quad . \]
Ceci posé, les \xout{familles} systêmes linéaires de diviseurs sur X paramétrés

\mspage{124}

\clearpage
% Nouveau EGA V (cote 25): archive cote 25, site page 125 (PDF page 126).
% typescript. Draft compared: none.
% AI-transcribed by Claude Opus 5.5 (claude-opus-5-5, adaptive thinking, effort high), Message Batches API, batch ega25b (msgbatch_01A7n5YuXubUPNK1zsepi8RG), 2026-11-13; not yet checked by eye.
\scanpage{25/125}
\inserted{compte tenu que l'hyp. sur $f$ implique $f_*(\underline{O}_X)\simeq\underline{O}_S$, d'où $f_*(\underline{O}_X^*)\simeq\underline{O}_S^*$}
\mspage{\hand{28}}

par $\check{P}$ , et associés à $\underline{L}$ (dans le sens de \hand{\S}.1.) correspondent biunivo-%
quement aux\hand{*} homomorphismes \uline{surjectifs} $Q\longrightarrow\check{E}$\hand{,} modulo \inserted{multiplication par des} sections de $\underline{O}_S^{\hand{*}}$ .\hand{\unsure{$\blacktriangledown$}}
\inserted{i.e. \deleted{aux sections du faisceau $\mathcal{H}om_S(Q,\check{E})^*/\underline{O}_S^*$, où le \unsure{premier} * \unsure{indique}}}
\margin{\unsure{Pour faisceaux d'ensembles} \illegible{2} \deleted{\illegible{2}} \unsure{préciser}}
L'existence de $Q$ se ramène par le procédé brévété au cas S noethé-%
rien, et dans ce cas n'est autre que III 7.7.6. (l'hypothèse sur les
fibres de X étant d'ailleurs inutil\unsure{e|è}). Comme $\underline{E}$ est loc libre de type fini
la donnée d'un hom $\underline{E}\longrightarrow f(\underline{L})$ équivaut à celle d'une section de
$\underline{L}\otimes_{\underline{O}_S}\check{\underline{E}}$ , donc à celle d'un homomorphisme $Q\longrightarrow\check{\underline{E}}$ . Il reste à exprimer
que la condition envisagée dans \hand{\S}.1. est bien vérifiée, ce qui (grâce à
l'hypothèse faite sur les fibres de X/S) se réduit à vérifier que
fibre par fibre l'homomorphisme correspondant $E(s)\longrightarrow$ \xout{Q(s)}\hand{$\vee$}
$H^0(X_s,L_s)\simeq\mathrm{Hom}_{k(s)}(Q(s),k(s))$ est injectif, ou encore $Q(s)\longrightarrow\check{E}(s)$ surjec%
tif, ce qui par Nakayama signffie aussi que $Q$ \quad $E$ est surjectif. Le ``mo
dulo sections de $\underline{O}_S$ '' provient du ``modulo isomorphisme'' dans \hand{\S}.1.\hand{/}

On peut \xout{xxxxxxxxx}\inserted[typed]{interpréter} \hand{\S}.1. par une autre voie, en utilisant le fait
que $P(\underline{Q})$ représente le \xout{xxxxxxxx} sous-foncteur de $\mathrm{Div}_{X/S}$ défini par
$\underline{L}$, en v ertu du n$^{\circ}$2. Par suite, un\xout{x}système linéaire de diviseurs paramétr
par $\check{P}$ \inserted[typed]{et associé à $\underline{L}$} s'interprète comme un morphisme $\check{P} = P(\check{\underline{E}})\longrightarrow P(Q)$; le caractère
\uline{linéaire} de la famille de diviseurs définie par un tel morphisme
s'interprète alors par le fait que ce morphisme soit ``linéaire'', i.e.
précisément défini par un morphisme surjectif $Q\longrightarrow\check{\underline{E}}$ . \hand{$\ulcorner$} On voit aussi
dans\xout{xx} ce cas que le morphisme $\check{P}\longrightarrow\mathrm{Div}_{X/S}$ est un monomorphisme \deleted{\illegible{2}}
(puisque $\check{P}\longrightarrow P(Q)\hookrightarrow\mathrm{Div}_{X/S}$ l'est), fait d'ailleurs \inserted{plus} général, cf corol-%
laire plus bas. Convenons \inserted[typed]{alors} de dire que d\unsure{e|a}ux familles linéaires de divi-%
seurs de X/S, paramétrées par les fibrés projectifs $\check{P},\check{P}'$ sont \uline{isomorph-}%
\uline{es} s'ils sont transformés l'un en l'autre par un S-isomorphisme $\check{P}\longrightarrow\check{P}'$ ,
(lequel sera d'ailleurs unique grâce au fait qu'on a des monomorphismes
dans $\mathrm{Div}_{X/S}$ ). \deleted{On appelera système linéaire de diviseurs dans X/S}
(\deleted{sans référence à un fibré projectif précisé qui le paramétriserait) une}

\mspage{\hand{125}}

\clearpage
% Nouveau EGA V (cote 25): archive cote 25, site page 126 (PDF page 127).
% typescript. Draft compared: none.
% AI-transcribed by Claude Opus 5.5 (claude-opus-5-5, adaptive thinking, effort high), Message Batches API, batch ega25b (msgbatch_01A7n5YuXubUPNK1zsepi8RG), 2026-11-13; not yet checked by eye.
\scanpage{25/126}
\mspage{\hand{29}}

\margin{\unsure{en remarquant, en vérifiant au §22 pour la construction des préliminaires grassmanniennes}}
On peut alors exprimer \over{5}{4}.2. en \unsure{diasant|disant} que l'ensemble des classes
à isomorphisme près de systèmes linéaires de diviseurs sur X associés à
$\underline{L}$ , est en correspondance biunivoque canonique avec l'ensemble \xout{Grassx}
Grass(Q)(S\xout{xxx} et cette correspondance étant compatible avec \unsure{les} changemen
de base quelconques, onvoit que le foncteur
\[ S' \hand{$\rightsquigarrow$} \text{ensemble des classes (mod isomorphisme) de systèmes} \]
\[ \text{linéaires de diviseurs de } X_{S'}/S' \text{ associés à } \underline{L}_{S'} \]
est représentable par le S-préschéma \uline{Grass} (Q).

\margin{\unsure{Sous les conditions de 1.1., on dit que le syst. lin. de diviseurs est complet lorsque $\forall$ $s \in S$, $E(s) \rightarrow H^0(X_s, \underline{L}|X_s)$ est \uline{bijectif}. Ceci ne se conserve pas par changement de base (cf. 4.12) si $g : X \rightarrow S$ est de prés. finie) \& $g$ est propre, alors ceci implique que $g_*(\underline{L})$ est \uline{localement libre de type fini}, et commutation ... $g$ est propre et existence ...}}
\deleted{Le résultat suivant, de nature moins technique que 5.2., devrait
plutôt passer avant ce dernier:

\uline{Corollaire} 5.3. Soit f:X \quad S un morphisme fidèlement plat loc de prés
\inserted[typed]{sur X/S}
finie. Alors pour tout\xout{xxfamille} système linéaire de diviseurs, paramétré
par un fibré projectif P sur S, le morphisme correspondant
\[ \hand{$\check{\ }$}P \qquad \mathrm{Div}_{X/S} \]
est un monomorphisme.

Cela signifie que deux morphismes $S'$ \quad \hand{$\check{\ }$}P ayant même composé avec
le morphisme précédent sont égaux. Faisant le changement de base $S'$ \quad S,
on est ramené \xout{à montrer que deux} au cas de deux sections de \hand{$\check{\ }$}P sur S,
et}

% NOTE: the paragraph from "Il faut expliciter" to "a) Supposons ... algé" is enclosed by hand in a box drawn around it
Il faut expliciter aussi dans \over{5}{4}.1. que $\underline{L}|X-Z$ est canoniquement
isomorphe à $f^{\hand{*}}(\underline{O}_P(1))$ (avec les notations du nº précédent), donc dans
le cas $D \in \emptyset$ envisagé dans \over{5}{4}.3. , $\underline{L}|X-Z$ n'est autre que $\underline{L}_{X-Z}$
de \over{5}{4}.3. b), et par suit e $\underline{L}$ n'est autre que le prolongement canonique et
unique de $f^*(\underline{O}_P(1))$ en un faisceau inversible sur X.

\uline{Proposition} \over{5}{4}.3. Soit D un système linéaire de diviseurs sur X/S
paramétré par \hand{$\check{\ }$}P , où g: X$\longrightarrow$ S est \deleted{\unsure{lisse}} un morphisme plat \inserted{loc.} de présenta-%
tion finie.

a) Supposons \inserted{(g de prés. finie et} que pour tout $s \in S$, désignant par k une clôture algé

\margin{\unsure{lin. complet associé à $\underline{L}$ ssi $g_*(\underline{L})$ est loc. libre, sa formation commutant aux changements} \unsure{16} \unsure{aux fibres, et pour tout $s \in S$, et $\varphi \in H^0(X_s, \underline{L}|X_s)$, $\varphi \neq 0$, $\varphi$ est une section régulière, condi}}

\clearpage
% Nouveau EGA V (cote 25): archive cote 25, site page 127 (PDF page 128).
% typescript. Draft compared: none.
% AI-transcribed by Claude Opus 5.5 (claude-opus-5-5, adaptive thinking, effort high), Message Batches API, batch ega25b (msgbatch_01A7n5YuXubUPNK1zsepi8RG), 2026-11-13; not yet checked by eye.
\scanpage{25/127}
\mspage{30}
brique de k(s), il existe un cycle premier T associé à $X_k$ tel que
$k \to H^0(T,\underline{O}_T)$ soit un isomorphisme (condition automatiquement satisfaite
si g est propre et surject\uline{i}f). Alors le morphisme D: $\check{P} \longrightarrow \mathrm{Div}_{X/S}$ est
un monomorphisme.

b) \xout{Supposonsxxf} Considérons l'application $u \rightsquigarrow D\circ u$ de
$\mathrm{Aut}_S(\check{P}\ )$ dans l'ensemble des familles de diviseurs sur X/S paramétrées
par P . Alors si \unsure{g} est surjectif, l'application précédente est injective,
en particulier D = D u implique u = $\mathrm{id}_{\check{P}}$\inserted{; plus généralement, le morphisme de foncteurs $\underline{\mathrm{Aut}}_S(\check{P}) \longrightarrow \underline{\mathrm{Syst.\ lin.\ div}}_{X/S,\check{P}}$ est un monomorphisme}

On notera que sous leshypothèses de a), b) est uneconséquence
triviale de a); d'autre part, b) est valable sous d eshypothèses moins
restrictives que a). On fera attention que a) devient faux si on abandon
ne l'hypothèse restrictive qui y a \xout{xéé} été faite: prendre par exemple
S=Spec(k), X un ouvert de $P^1_k$ ne contenant pas les deux points distincts
a et b de $P^1_k(k)$ , D la famille de diviseurs définie par l'immersion
canonique de X dans $P^1_k$ . Alors les deux \xout{morphismes} points a et b
définissent les mêmes diviseurs de X (savoir le diviseur nul), sans être
eux-mêmes identiques.

Supposons d'abord que S soit le spectre d'un corps \inserted{k, qu'on peut évidemment (par ``descente'') supposer} alg clos. \deleted{\illegible{1}}
Soit T comme dans a)\inserted{, \unsure{qu'on munit|qui muni} de la structure réduite induite,} on a alors un morphisme ``diviseur induit''
$\mathrm{Div}_{X/k} \longrightarrow \mathrm{Div}_{T/k}$ , et il suffit de prouver que le composé $\check{P} \longrightarrow \mathrm{Div}_{T/k}$
est un monomorphisme. Comme ce dernier est encore un syst lin de div ,
on est ramené au cas où X=T , donc au cas où $H^0(X,\underline{O}_X) \xleftarrow{\sim} k$ . Alors pour
tout S sur k, on a $g_{S*}(\underline{O}_{X_S}) \xleftarrow{\sim} \underline{O}_S$ , donc (si $\underline{L}$ sur X et u: $\underline{E} \to$ \deleted{\illegible{1}}$\hand{g_*}(\underline{L})$
sont comme dans §.1.) si deux sections $\phi$, $\psi$ de $\underline{E}_S$ partout non nulles
sont telles que $u(\phi)$ et $u(\psi)$ sont des sections de $\underline{L}_S$ sur $X_S$ ayant même
diviseur, donc \xout{xxxxxf} déduites l'une de l'autre par multiplication par
une section \inserted[typed]{inversible} de $\underline{O}_{X_S}$, il s'ensuit que $\psi$ est déduit de $\phi$ par multiplication

\margin{127}
\margin{\unsure{Mais|Non} automatiquement vérifiée si les $X_s$ sont \deleted{\illegible{1}} intègres.}

\clearpage
% Nouveau EGA V (cote 25): archive cote 25, site page 128 (PDF page 129).
% typescript. Draft compared: none.
% AI-transcribed by Claude Opus 5.5 (claude-opus-5-5, adaptive thinking, effort high), Message Batches API, batch ega25b (msgbatch_01A7n5YuXubUPNK1zsepi8RG), 2026-11-13; not yet checked by eye.
\scanpage{25/128}
\mspage{31}

par une section inversible de $\underline{O}_S$ , donc $\phi$ et \unsure{$\psi$} définissent le même point
de $\overset{\hand{$\vee$}}{P}$ à valeurs dans S. Comme tout point de $\overset{\hand{$\vee$}}{P}$ à valeurs dans S est défini,
localement sur S, par une section $\phi$ de $\underline{E}_S$ qui ne s'annule pas (cf Chap II)
a) s'ensuit.   Pour prouver b), on note le

\uline{Lemme}   Soit D un syst lin de diviseurs sur X non vide loc de type fini
sur k alg clos, paramétré par $\overset{\hand{$\vee$}}{P}$(E), et considérons le morphisme correspon-%
dant  f: X-Z $\longrightarrow$ P , où Z est l'ensemble des points fixes. Alors \xout{xx} si
r = rang$_k$E $>$\hand{0,} il existe  r+1 point\hand{s}\inserted[typed]{$x_i$}\inserted{$(1\leq i\leq r+1)$} de \xout{(Xx} X(k)-Z(k)  tels que les
f(x$_i$) forment une  ``base projective'' de P, i.e. tels que  pour toute
partie J de [1,r+1] ayant \hand{r} \xout{mi} éléments,  les  f(x$_i$),i$\in$\hand{J}, \xout{xxxx} ne sont pas
contenus dans un hyperplan de P.

On peut évidemment supposer Z=$\emptyset$ , \deleted{X intègre (en le remplaçant
par une\xout{Xxxxxx} composante irréductible T de X muni de la structure induite
réduite).} \xout{OxxaxxxxxapplicationxxinjectivexxxE} Comme \inserted[typed]{par \over{4}{\unsure{3|5}}.1.} f(X) n'est contenu dans
aucun hyperplan de P, on en conclut d'abord l'existence de r points
$x_i$\inserted{$(1\leq i\leq r)$} de X \xout{telsxque}\inserted[typed]{tels que les f($x_i$) soient} projectivement indépendants dans P, i.e. soient définies
par des formes linéairement indépendantes sur E.   Reste à prouver qu'il
existe un  $x_{r+1} = x$\deleted{\illegible{1}}   dans X(k) tel que  f(x) ne soit dans aucun des r
hyperplans \inserted{$H_i$} définis par les \deleted{\unsure{f($x_i$)}} systèmes de r-1 parmi les f(x$_i$). Or dans
le cas contraire, compte tenu du sorite Jacobson,  on aurait
f(X) $\subset \bigcup_{\hand{$i$}}$ H$_i$ , donc \xout{Xxseraitxlaxréunionxdesxxxx}\xout{xxx} si X$_0$  est une composante
irréductible de X, f(X$_0$) serait contenu dans un des H$_i$ , ce qui contredit
\over{\unsure{4}}{\unsure{3|5}}.1. - Ceci posé,  pour prouver b), on peut  évidemment supposer Z=$\emptyset$ ,
et utilisant \over{\unsure{4}}{\unsure{3|5}}.1.  on est ramené à prouver qu'un automorphisme\inserted{u} de $\overset{\hand{$\vee$}}{P}$ est
déterminé quand on connait \xout{xxx}\inserted[typed]{le} composé\inserted[typed]{de\xout{\unsure{s}} son contragrédient $\overset{\hand{$\vee$}}{u}$ dans P} avec f: X$\longrightarrow$ P , - et que l'asser-%
tion analogue est vraie après tout changement de base S$\rightarrow$Spec(k)
pour un  automorphisme u de $\overset{\hand{$\vee$}}{P}_S$  .  Or cela résulte aussitôt  du lemme

\mspage{128}

\clearpage
% Nouveau EGA V (cote 25): archive cote 25, site page 129 (PDF page 130).
% typescript. Draft compared: none.
% AI-transcribed by Claude Opus 5.5 (claude-opus-5-5, adaptive thinking, effort high), Message Batches API, batch ega25b (msgbatch_01A7n5YuXubUPNK1zsepi8RG), 2026-11-13; not yet checked by eye.
\scanpage{25/129}
\mspage{32}

précédent, et de la détermination des automorphismes de $\overset{\hand{\vee}}{\mathrm{P}}(E) = P(\overset{\hand{\vee}}{E})$ faite
au n°1 , qui implique qu'un automorphisme d'un fibré projectif sur un S
(relatif à un Module loc libre de type fini)  est connu quand on connait
son effet sur une famille de sections formant une ``base projective''
(i.e. formant une base projective en chaque \xout{point} fibre).

Esquissons maintenant le cas général: S quelconque.  Bien entendu,
quitte à faire un changement de base sur S, on est ramené dans a) à prouve
que deu\xout{x} sections  de $\overset{\hand{\vee}}{\mathrm{P}}$ sur S  qu\uline{i} définissent le même diviseur sur X sont
identiques, et dans b) que deux automorphismes  de $\overset{\hand{\vee}}{\mathrm{P}}$  qui sont tels que
$D\circ u = D\circ v$  sont identiques. \inserted[typed]{On peut supposer  S affine,} Dans le cas b) où on n'a pas supposé expressé%
ment  g de présentation finie, mais g surjectif, on se ramène aussitôt
(grâce au fait que g est ouvert) au cas où X est é\xout{x}galement affine, donc
de présentation finie sur S.  Par le procédé brévete, on peut maintenant
se réduire au cas S \uline{noethérien} .  Or pour un préschéma de base \inserted{\unsure{S loc}} noethérien,
et un morphisme de foncteurs  $F \overset{\hand{h}}{\longrightarrow} G$ sur S  (F et G des foncteurs
$(\mathrm{Sch})_{/S}^{\circ} \longrightarrow (\mathrm{Ens})$ ) ,  on a \xout{des}\unsure{un} critère\xout{s} très généra\unsure{l}\xout{x} \inserted[typed]{(qui sera explicité au Chap V),} permettant d'affir-%
mer que si  pour tout  $s \in S$, le\xout{s} morphisme correspondant $F_s \longrightarrow G_s$ est
un monomorphisme, alors $F \longrightarrow G$ est un monomorphisme  (NB on pose $F_s{=}F\times_S\mathrm{Spec}$
$(k(s))$, et de même pour $G_s$ ),  moyennant des hypothèses simples sur F et G
(vérifiées par exemple si F et G sont tous les deux représentables par des
préschémas loc de type fini sur S; mais dans le cas qui nous occupe, seul
le premier des deux foncteurs est représentable à priori).  Nous allons
esquisser le raisonnement du Chap V  dans le\xout{s} \inserted[typed]{deux} cas particuliers qui nous
intéressent ici.  On a deux sections \inserted{u, v} (de $\overset{\hand{\vee}}{\mathrm{P}}$ , resp. du groupe projectif
\uline{GP} ($\overset{\hand{\vee}}{E}$) ) d'un \hand{\unsure{pré}}schéma de type fini \inserted{F} sur S dont on veut montrer qu\xout{e}'elles
sont égales, pour ceci manifestement il suffit de prouver qu'elles sont
égales après tout changement de base  $\mathrm{Spec}(\underline{O}_{S,s}/\underline{m}_s^{\hand{n+1}}$ ) $\longrightarrow$ \hand{$\mathcal{S}$},

\mspage{129}

\clearpage
% Nouveau EGA V (cote 25): archive cote 25, site page 130 (PDF page 131).
% typescript. Draft compared: none.
% AI-transcribed by Claude Opus 5.5 (claude-opus-5-5, adaptive thinking, effort high), Message Batches API, batch ega25b (msgbatch_01A7n5YuXubUPNK1zsepi8RG), 2026-11-13; not yet checked by eye.
\scanpage{25/130}
\mspage{\hand{33}}

ce qui nous ramène au cas où S est artinien local, S=Spec(A). On procède
par récurrence sur l'entier n tel que $m^{n+1}=0$ , ce qui nous permet de
supposer que les deux sections envisagées \inserted{u,v} sont égales mod $m^{n\xout{\unsure{+1}}}$. Donc
l'une se déduit de l'autre à l'aide d'un élément \xout{xxxxxx(xxxx xxxxxxx} \hand{$\delta$ de}
$\mathrm{Hom}_k(u_o^{\hand{*}}(\hand{\Omega}^1_{F_o/k}), V)$ , où $k=A/m$ est le corps résiduel, $F_o=F\otimes_A k$ est la
fibre réduite, $V=m^n = m^n/m^{n+1}$ . Il faut prouver que \hand{$\delta$}$= 0$, moyennant
l'hypothèse $h(u)=h(v)$ . Le principe général de la vérification est le
suivant: exprimant d'abord que $h(u)$ et $h(v)$ coincident mod $m^n$, on voit
que leur ``différence'' peut s'écrire comme un élément \hand{$\delta$}$'$ de
$\mathrm{Hom}_k(w_o^{\hand{*}}(\hand{\Omega}^1_{G_o/k}),V)$ , où $w_o = h_o(u_o)=h_o(v_o)$ \xout{xxx} et où $G_o = G\otimes_A k$ ;
cet élément n'est autre que celu\uline{i} qui est déduit de \inserted[typed]{\hand{$\delta$} par composition avec} l'homomorphisme natu-%
tel $h_o^{\hand{*}} : w_o^{\hand{*}}(\hand{\Omega}^1_{G_o/k}) \hand{$\to$} u_o^{\hand{*}}(\hand{\Omega}^1_{F_o/k})$ déduit de $h_o: F_o\hand{$\to$} G_o$. Comme
$h(u)=h(v)$ donc \hand{$\delta$}$' = 0$, le composé de \hand{$\delta$} avec l'homomorphisme précédent
$h_o^{\hand{*}}$ est nul, donc si on sait que $h_o$ est surjectif, il s'ensuit que \hand{$\delta$}$=0$
et on gagne. Or le fait que $h_o:F_o\hand{$\to$} G_o$ est un monomorphisme, donc induit
un monomorphisme pour l'ensemble des points à valeurs dans les nombres
duaux sur k, implique bien en effet que $h_o^{\hand{*}}$ est surjectif (son transposé
étant injectif). Ce raisonnement est bien valable tel quel lorsque G
est représentable, ce qui n'est pas le cas cependant dansle cas qui nous
intéresse. \deleted{Il faut donc} \inserted{On peut cependant encore} définir un vectoriel \hand{$\mathcal{V}_{w_o}$} \inserted{(espace tangent à $G_o/k$ en $w_o$)} sur k, jouant le rôle
du \inserted{dual de} $w_o^{\hand{*}}(\hand{\Omega}^1_{G_o/k})$ pour exprimer la ``déviation'' de deux points de G qui
coincident mod $m^n$ \inserted{comme un élément de $\mathcal{V}_{w_o}\otimes_k V$}. C'est essentiellement straightforward, et contenu
dans les développements systématiques du par. 2\unsure{\over{5}{6}} (? Extensions infinitési-%
males), auquel on pourrait renvoyer ici. Dans le cas a), G étant le
foncteur $\mathrm{Div}_{X/S}$ , $w_o$ correspond donc à un diviseur de Cartier \xout{xx} $D_o$ sur
$X_o = X\otimes_A k$ , et il faut prendre \hand{$\Omega$}$= H^o(D_o, \underline{n}_{D_o/X_o})$ où le $\underline{n}$ est le
faisceau\xout{x} normal de $D_o$ dans $X_o$ , isomorphe également au faisceau induit
sur $D_o$

\mspage{\hand{130}}

\clearpage
% Nouveau EGA V (cote 25): archive cote 25, site page 131 (PDF page 132).
% typescript. Draft compared: none.
% AI-transcribed by Claude Opus 5.5 (claude-opus-5-5, adaptive thinking, effort high), Message Batches API, batch ega25b (msgbatch_01A7n5YuXubUPNK1zsepi8RG), 2026-11-13; not yet checked by eye.
\scanpage{25/131}
\mspage{34}

par $\underline{\mathcal{O}}_{X_o}(D_o)$ sur D$_o$. Dans le cas b), \unsure{On} peut supposer D sans points
fixes, et il est plus commode d'interpréter la situation en termes de
morphismes dans P (cf 4.1.), de sorte que G devient le foncteur
\underline{Hom}$_S$(X,P), et \hand{$\mathcal{V}$}$_{w_o}$ sera l'espace \xout{\illegible{1}} Hom$_{\underline{O}_{X_o}}$$(f_o^{*}(\underline{\hand{$\Omega$}}^{1}_{P_o/k}), \underline{O}_{P_o})$.
Dans l'un et l'autre cas, on aura un homomorphisme naturel
\hand{$\mathcal{V}$}$_{u_o}\otimes_k V$ \hand{$\to$} \hand{$\mathcal{V}$}$_{w_o}\otimes_k V$ (où \hand{$\mathcal{V}$}$_{u_o}$ = dual de $u_o^{*}($\hand{$\Omega$}$^{1}_{F_o/k})$ ) exprimant le passage
de la déviation \hand{$\delta$} à la déviation \hand{$\delta$}' correspondante par application de h,
et l'injectivité de cette application résulte\xout{ra} encore de l'injectivité
de \hand{$\mathcal{V}$}$_{u_o}$\hand{$\to$} \hand{$\mathcal{V}$}$_{w_o}$ , qui elle provient\xout{ra} du fait que $h_o$: $F_o$\hand{$\to$}$G_o$ est un
monomorphisme.
\margin{$/$}

Pratiquement, il ne semble guère possible de rédiger cette dernière
partie de la démonstration sans référer aux petits calculs du par.2\hand{5}
(qu'il n'est pas question de refaire ici dans un cas particulier). On
signalera que cela ne donne pas lieu à des cercles vicieux, car le par.
2\hand{5} \xout{xxx} et les calculs qu'on y \xout{xxxxx} développe ne dépend que des sorites
de calcul différentiel du par.16, et d'ailleurs \hand{4}.3. ne sera plus
utilisé dans le Chap IV, sauf peut-être dans les deux n$^{o}$s qui suivent.

L'intérêt de \hand{4}.3.\inserted{a)} est de montrer que sous les conditions \deleted{de a)}\inserted{\unsure{initiales}},
le fibré projectif paramétrisant peut s'interpréter comme un sous-foncteur
de Div$_{X/S}$ , ce qui permet d'interpréter de façon intrinsèque la notion
de classe (à isomorphisme près sur le \xout{schéma} fibré projectif paramétrisa-%
teur) de système linéaire sur X/S, e\xout{n termes d'un} comme étant un sous-
foncteur $\overset{\hand{$\vee$}}{\text{P}}$ de Div$_{X/S}$ qui satisfait certaines propriétés (\unsure{savoir} est
représentable par un fibré projectif, et la famille de diviseurs définie
par l'injection c anonique de ce dernier dans Div$_{X/S}$ est linéaire au s ens
du n$^{o}$\over{4}{3}), ce qui est \xout{\illegible{1}}\inserted[typed]{essentiellement} le point de vue classique (où un\xout{e}
système linéaire de diviseurs est défini comme un \uline{ensemble} de diviseurs

\margin{131}

\clearpage
% Nouveau EGA V (cote 25): archive cote 25, site page 132 (PDF page 133).
% typescript. Draft compared: none.
% AI-transcribed by Claude Opus 5.5 (claude-opus-5-5, adaptive thinking, effort high), Message Batches API, batch ega25b (msgbatch_01A7n5YuXubUPNK1zsepi8RG), 2026-11-13; not yet checked by eye.
\scanpage{25/132}
\mspage{35}

satisfaisant certaines conditions, cf \hand{4}.4.). D'autre part, \hand{\S}.3. b)
équivaut \xout{xxxxxtixltxxxxt} aussi à dire que \inserted[typed]{lorsque g est surjectif,} si deux systêmes lin de diviseurs
sur X/S \xout{xxxx} paramétrés par deux fibrés projectifs $\overset{\hand{\vee}}{P},\overset{\hand{\vee}}{P}'$ sont isomorphes;
alorsilexiste un \uline{unique} isomorphisme de $\overset{\hand{\vee}}{P}$ sur $\overset{\hand{\vee}}{P}$\unsure{(|'} compatibles avec D,D';
on peut donc dire qu'une classe (à isom près) de systêmes linéaires sur
X/S détermine son fibré projectif paramétrisateur \uline{à isomorphisme unique}
\uline{près} . Techniquement, ce r ésultat permettra (une fois qu'on disposera
de la te\unsure{ç}hnique de descente, exposée au Chap V) de faire de \uline{la d escente}
\uline{fid plate pour des systêmes linéaires de diviseurs}, - sous réserve
\margin{? (} toutefois d'admettre \inserted{également} comme fibrés paramétrisateurs des ``fibrés projectifs
tordus'', ce qui va être fait au prochain \deleted{paragraphex} numéro.

Descendan\deleted{t} à nouveau par terre, et bien bas, pour expliciter en
termes vulgaires la notion de système linéaire, en nous plaçant pour
simplifier sur un corps de base (bien que l'énoncé subsiste essentielle-%
ment tel quel sur une base affine):
\uline{Proposition} \deleted{\hand{4}.4.} Soit X un préschéma \xout{\unsure{loc}} \inserted[typed]{\xout{xxxxxide}} de type fini sur un corps k,
tel que $k\to H^0(X,\underline{O}_X)$ soit un isomorphisme\over{,}{.} A tout système linéaire \inserted{D} de
diviseurs \deleted{Dx} \inserted{sur X,} paramétré par un fibré projectif $\overset{\hand{\vee}}{P}$ sur k, associons
l'ensemble\hand{(!)}Ens(D) de tous les diviseurs sur X de la forme D(t), où
$t\in\overset{\hand{\vee}}{P}(k)$. \hand{a)} Si \deleted{alors} D' est un autre systême \inserted{lin} de diviseurs\inserted{sur X}, paramétré par un
fibré projectif $\overset{\hand{\vee}}{D}'$ , alors D et D' sont isomorphes sss Ens(D)=Ens(D').
\hand{b)} \inserted{Supposons k alg. clos et X génér.\textsuperscript{t} intègre.}
\hand{(}Pour qu'un ensemble $\Delta$ de diviseurs de Cartier \unsure{positifs|pesitifs} sur X soit de
le forme Ens(D), il f et s qu'il \inserted{existe} \deleted{soit vide, ou qu'il satisfasse la condi-%
tion suivante: \xout{pourxxxxsoitxxD$_0$xxxxxxxxxalors}. il existe \xout{desxfonctions}}
\xout{rationnellesxrelxàxkxxfxxx$\phi_1$xx} une suite finie $(\phi_i)_{1\ \ i\ \ n}$ de fonctions
rationelles rel k \deleted{sur X, tellex que pour tout système d'eléments}
\deleted{$(c_i)_{1\ \ i\ \ n}$ de k}

\mspage{132}

\clearpage
% Nouveau EGA V (cote 25): archive cote 25, site page 133 (PDF page 134).
% typescript. Draft compared: none.
% AI-transcribed by Claude Opus 5.5 (claude-opus-5-5, adaptive thinking, effort high), Message Batches API, batch ega25b (msgbatch_01A7n5YuXubUPNK1zsepi8RG), 2026-11-13; not yet checked by eye.
\scanpage{25/133}
\mspage{36}

un sous-k-espace vectoriel\inserted{E}de l'espace des \deleted{\inserted{pseudo}} fonctions \deleted{rationelles} \inserted{\unsure{méromorphes}} \deleted{\inserted[typed]{strictes}} \xout{xxxxxxx}
sur X, tel que pour tout $\phi \hand{\in} E-(0)$, $\phi$ soit \inserted{\unsure{régulier}} \deleted{inversible dans l'\xout{xxxxxxx}}
\deleted{anneau des fonctions rationelles} \deleted{\inserted[typed]{strictes}} \xout{xxxxxx}, i.e. div($\phi$) est défini \deleted{comme}
\deleted{diviseur de Cartier,}\unsure{\hand{,}} et que $\hand{\triangle}$ soit l'ensemble des div($\phi$) pour $\phi \hand{\in} E-(0)$.
c) Soient E, E' deux sous-\inserted[typed]{k-}espaces vectoriels de l'espace des
fonctions rationelles \inserted[typed]{strictes} \xout{xxxxxxxx} sur X, \xout{xxxxx} satisfaisant la condition
envisagée dans b), alors les ensembles de diviseurs $\hand{\triangle}$, $\hand{\triangle}'$ qu'ils
définissent respectivement sont égaux sss il existe une \deleted{\inserted{\unsure{pseudo-}}} fonction
\deleted{rationelle stricte}e \deleted{inversible} \inserted{\unsure{régulière}} $\phi$ sur X, telle que $E' = \phi E$ . Si $E\hand{\neq}(0)$\inserted{(i.e. $\triangle \neq \emptyset$)}, une
telle $\phi$ est déterminée modulo multiplication par un élément c de $k^{\hand{*}}$.

La démonstration est un exercice facile, enutilisant \hand{\S}.1. , et je
me dispense d'en écrire la démonstration, sauf réclamation de ta part. Il
\margin{||} me sembled'ailleurs que \hand{\S}.4. viendrait avantageusement dès avant \hand{\S}.3.,
étanttechniquement nettement plus trivial. Noter d'ailleurs que lorsque
X est géom intègre, la condition sur E énoncée dans b) devient \uline{vide}.
La restriction mise au début de b) tient au fait qu'il se pourrait autre-%
ment que la condition énoncée pour b) ne soit plus vérifiée après passage
à la clôture algébrique de k (il est facile d'en donner des exemples\xout{xx},
en toute caractéristique, et même lorsque k est séparablement clos en
c\xout{\unsure{ar}} p $>$ 0). Pour bien faire, on devrait énoncer b) sans condition supplé-%
mentaire sur k ou X, mais en énoncant la condition sur E en passant à la
clôture algébrique de k (et signaler que si X est géom intègre, \unsure{oette|cette}
condition devient vide). \hand{-- Par abus de langage, un}
\hand{\uline{ens.} $\triangle$ de diviseurs de la forme \unsure{Ens}(D) sera}
\hand{encore parfois appelé système linéaire de diviseurs}
\hand{sur X ;}

\mspage{133}

\clearpage
% Nouveau EGA V (cote 25): archive cote 25, site page 134 (PDF page 135).
% typescript. Draft compared: none.
% AI-transcribed by Claude Opus 5.5 (claude-opus-5-5, adaptive thinking, effort high), Message Batches API, batch ega25b (msgbatch_01A7n5YuXubUPNK1zsepi8RG), 2026-11-13; not yet checked by eye.
\scanpage{25/134}
\mspage{\hand{37.}}

\margin{29.9.1965}

Cher Dieudonné,

Merci de ta lettre du 24  et pour la table des matières des par.
16 à 19.  Je serais content de recevoir à l'occasion la table des
matières provisoire des par.20 et 21 ;\xout{\unsure{d}}d'accord pour les joindre au
fascicule 4 du Chap IV. Mais comment vaux-tu \xout{les} subdiviser\inserted{mon ancien par.20}, et
quels seront les titres des deux morceaux \unsure{?|2} Comme je commence à me
perdre dans le plan, et qu'il est \unsure{paffois} commode de \unsure{pauvoir|pouvoir} référer
sans trop déconner à un nº de paragraphe, je te donne ici ce qui me
semble être le plan actuel, dis-moi si tu es d'accord:

20.  ???

21.  ???

22.  Systêmes linéaires, compléments sur le groupe de Picard.

23.  Grassmaniennes.

24.  Formes lisses, singularités quadratiques ordinaires.

25.  Sections hyperplanes et bordel.

26.  Résultant et discriminant.

27.  Extensions infinitésimales.

Le 25.  risque d'\unsure{a|q}illeurs d'être fort long, et je te vois déjà
vouloir le subdiviser en d\unsure{eux} !  Pourtant, $27=3^3$ est un bien joli
nombre !

Il n'est pas question que je publie l'ex-Appendice au par.18
sous mon nom; ta rédaction n'a à peu près plus rien de commun avec
les vagues notes manuscriptes que je t'avais passées, si même je t'en
ai jamais passé, et ne me suis borné à te dire: il n'y a qu'à faire

\mspage{\hand{134}}

\clearpage
% Nouveau EGA V (cote 25): archive cote 25, site page 135 (PDF page 136).
% typescript. Draft compared: none.
% AI-transcribed by Claude Opus 5.5 (claude-opus-5-5, adaptive thinking, effort high), Message Batches API, batch ega25b (msgbatch_01A7n5YuXubUPNK1zsepi8RG), 2026-11-13; not yet checked by eye.
\scanpage{25/135}
\mspage{38}

pareil que pour les anneaux complets ... Il serait d'autre part dom-%
mage que ton travail de mise au point soit perdu pour les éventuels
utilisateurs (il finit toujours par s'en trouver ..). C'est pourquoi
je te demande de bien vouloir reconsidérer la question d'en faire un
``joint papar''.

\deleted{Pour par.20 , 10.9.1. , il faut bien entendu utiliser le
fait que l'ensemble \xout{\unsure{dexxZ}} des points de $Z_{\hand{$\lambda$}}$ en lesquels $F_{\hand{$\lambda$}}$ \xout{\unsure{exxxd}}
restreint à la fibre est de prof \hand{\unsure{?}} \hand{$\geq$} n donné, est \uline{constructible}
(on a même du prouver au par.12 qu'il est ouvert, avec les hypothèses
de platitude et de présentation finie qu'on a faites). Comme son image
inverse dans Z est tout, c'est que c'est déjà tout un peu plus loin
que \hand{$\lambda$} . C'est vraiment toujours le même argument qui revient !}

Bien à toi
\hand{A Grothendieck}

\mspage{135}

\clearpage
% Nouveau EGA V (cote 25): archive cote 25, site page 136 (PDF page 137).
% typescript. Draft compared: none.
% AI-transcribed by Claude Opus 5.5 (claude-opus-5-5, adaptive thinking, effort high), Message Batches API, batch ega25b (msgbatch_01A7n5YuXubUPNK1zsepi8RG), 2026-11-13; not yet checked by eye.
\scanpage{25/136}
\mspage{\hand{39}}

\hand{§20 \quad \unsure{Compléments} \unsure{appl.} \unsure{rationnelles} et \unsure{morphismes} \unsure{birationnels} et}
\hand{\unsure{éclatements} :}

\hand{§21 \quad Diviseurs de Cartier, Weil, Picard,}
\hand{Syst. lin., \unsure{Brauer}-Severi.}

\margin{$\|$} \hand{§22 \quad \unsure{Grassmanniennes}}

\hand{§23 \quad \unsure{Sections} \unsure{hyperplanes}}

\hand{§24 \quad Ext. Inf.}

% NOTE: the following handwritten lines in the lower part of the page are written upside down (the sheet turned by 180 degrees); they are transcribed here in the order in which they appear going down the page as oriented above.
\hand{\unsure{filtration} $\unsure{F}$}
\hand{\unsure{pour tout}} \qquad \hand{\unsure{pour $n$}}
\hand{(two strokes from the summation sign up to ``\unsure{pour tout}'' and one stroke up to ``\unsure{pour $n$}'')}
\hand{$\unsure{p+q=n}$}
\hand{$H_{n}(X) \simeq \sum H^{\unsure{q}}(X, \Omega^{p})$} \quad \hand{\deleted{$H(X$}}
\hand{$\longleftarrow$}
\hand{\deleted{\unsure{Donner} \unsure{une} \illegible{1} \unsure{des} \illegible{1}}}

\mspage{\hand{136}}

\clearpage
% Nouveau EGA V (cote 25): archive cote 25, site page 137 (PDF page 138).
% typescript. Draft compared: none.
% AI-transcribed by Claude Opus 5.5 (claude-opus-5-5, adaptive thinking, effort high), Message Batches API, batch ega25b (msgbatch_01A7n5YuXubUPNK1zsepi8RG), 2026-11-13; not yet checked by eye.
\scanpage{25/137}
\mspage{\hand{40}}

IV 21 Faisceaux inversibles et diviseurs relatifs sur les fibrés projec-%
tifs et multiprojectifs; systêmes linéaires de diviseurs.

1. Faisceaux inversibles sur \xout{lesxf} un fibré projectif; application aux
morphismes d'un fibré projectif dans un autre.

2. Représentabilité de $\underline{\mathrm{Div}}^{L}_{X/S}$ ; diviseurs \inserted[typed]{relatifs} sur les fibrés projectifs et
multiprojectifs.

3. Systêmes linéaires de diviseurs$_{\mathrm{x}}$ et morphismes dans les fibrés projec-%
tifs.

4. Systêmes linéaires de diviseurs$_{\mathrm{x}}$ et Modules inversibles.

Les résultats développés dans le présent paragraphe et les suivants
sont déjà en partie de nature globale, et donnent notamment des complé-%
ments desur les schémas projectifs, utilisant \xout{de façon essentielle} les
constructions globales du Chap II, et même ici certains résultats du
Chap III, en plus des résultats purement locaux du présent Chapitre IV.

Un des buts du présent paragraphe est de développer le langage des
``systêmes linéaires de diviseurs'', lié d'une part à la classification des
morphismes dans un fibré projectif, d'autre part à celle des Modules inver
sibles sur un préschéma relatif donné. Notons d'ailleurs ici que les
``schémas de paramètres'' vraiment naturels pour les systêmes linéaires de
diviseurs sont les schémas de \uline{Brauer-Séveri}, qui généralisent les fibrés
projectifs, et peuvent se définir par exemple comme les fibrés qui devien-%
nent isomorphes à un fibré projectif après extension étale surjective de
la base. Comme leur étude utilise la théorie de la descente, développée
au Chap V, et que d'ailleurs leur classification équivaut à celle des \xout{fibr}
\xout{\unsure{fibrés principaux homogènes}} torseurs de groupe le groupe projectif, nous
renvoyons pour l'étude de ces schémas et ses liens avec la notion de sys-%
tême linéaire de diviseurs au Chap VI$_{\mathrm{x}}$ de notre savant ouvrage.

\mspage{\hand{137}}

\clearpage
% Nouveau EGA V (cote 25): archive cote 25, site page 138 (PDF page 139).
% typescript. Draft compared: none.
% AI-transcribed by Claude Opus 5.5 (claude-opus-5-5, adaptive thinking, effort high), Message Batches API, batch ega25b (msgbatch_01A7n5YuXubUPNK1zsepi8RG), 2026-11-13; not yet checked by eye.
\scanpage{25/138}
\mspage{\hand{41}}

Du point de  vue technique, le résultat principal du présent paragraphe
est le théorème 1.1. \inserted[typed]{déterminant le groupe de Picard d'un} \xout{\unsure{permettant de classifier les faisceaux inversibles}}
\xout{\unsure{sur un}} fibré projectif  en termes de celui de sa base, et ses premiers
coroallaires, développés dans les n$^{\underline{o}}$s 1 et 2.

\margin{138}

\clearpage
% Nouveau EGA V (cote 25): archive cote 25, site page 139 (PDF page 140).
% typescript. Draft compared: none.
% AI-transcribed by Claude Opus 5.5 (claude-opus-5-5, adaptive thinking, effort high), Message Batches API, batch ega25b (msgbatch_01A7n5YuXubUPNK1zsepi8RG), 2026-11-13; not yet checked by eye.
\scanpage{25/139}
% NOTE: the whole typed text of this page is an extremely faint impression (very pale carbon); almost every reading below is conjectural and must be checked on the original
\unsure{Nouveau} \uline{\unsure{EGA V}}

\unsure{IV 16.}

\unsure{Un nouveau n° :} \uline{\unsure{Ensemble singulier et composantes singulières d'une fonction,}}

\uline{\unsure{et critères différentiels}}.

\unsure{Ce n° sera utilisé au par. 20 des sections hyperplanes, mais sa}

\unsure{place naturelle me semble bien au par. 16.}

\uline{\unsure{Définition}} \unsure{1.} \unsure{Soient X un préschéma régulier, f une section de} $\unsure{\mathcal{O}_X}$ \unsure{. Un}

\unsure{point} $\unsure{x\in X}$ \unsure{est dit} \uline{\unsure{zéro}} \illegible{3} \uline{\unsure{singulier}} \unsure{pour f}

\unsure{si on a} $\unsure{f_x \in \mathfrak{m}_x^2}$ \unsure{; zéro} \uline{\unsure{singulier ordinaire}} \unsure{si c'est un zéro singulier}

\unsure{et si de plus l'élément de} $\unsure{\mathfrak{m}_x^2/\mathfrak{m}_x^3 \subset \mathrm{Sym}(\mathfrak{m}_x/\mathfrak{m}_x^2)}$ \unsure{(bien défini),}

\unsure{interprété comme une forme quadratique sur le dual} $\unsure{t_x}$ \unsure{de} $\unsure{\mathfrak{m}_x/\mathfrak{m}_x^2}$ \unsure{sur} $\unsure{k(x)}$,

\unsure{est une forme non dégénérée.} \unsure{(Un zéro singulier qui n'est pas} \illegible{2}

\unsure{lier est parfois appelé} \uline{\unsure{zéro singulier ordinaire}}\unsure{).}

\uline{\unsure{Remarques}} \unsure{2.} \unsure{Si} $\unsure{x\in V(f)}$ \unsure{, alors x est un} \inserted{\unsure{non}} \unsure{zéro singulier de f ssi}

$\unsure{f_x \neq 0}$ \unsure{et x est un point} \illegible{1} \unsure{non singulier (= régulier) de} $\unsure{V(f)}$

\unsure{i.e. que x est un point régulier de} $\unsure{V(f)}$ \unsure{et} $\unsure{V(f)\neq X}$ \unsure{au voisinage de}

% NOTE: the following five lines are enclosed in a drawn rectangle crossed by faint diagonal strokes (possibly a deletion); the frame and strokes are too faint to be certain
\unsure{x. On peut prouver que si x est singulier pour f, il est}

\unsure{ordinaire ssi} $\unsure{f_x\neq 0}$ \unsure{et la multiplicité de} $\unsure{\mathcal{O}_{V(f),x}}$ \unsure{est égale à 2}

\unsure{dans ce cas on dit aussi parfois que x est un} \uline{\unsure{point singulier ordinaire}}

\unsure{du préschéma} $\unsure{V(f)}$ \unsure{(i.e.} \unsure{ssi} $\unsure{x\in V(f)}$ \unsure{est non} \illegible{3}

$\unsure{f_x\neq 0}$ \unsure{et la multiplicité de} $\unsure{\mathcal{O}_{V(f),x}}$ \unsure{est} \unsure{2)}

\uline{\unsure{Définition}} \unsure{3.} \unsure{Soient X un préschéma lisse sur S, f une}

\unsure{section de} $\unsure{\mathcal{O}_X}$ \unsure{,} $\unsure{x\in V(f)}$ \unsure{. On dit que x est un zéro singulier sur}

\uline{\unsure{singulier}} \unsure{sur S} \unsure{(resp.} \uline{\unsure{génériquement singulier ordinaire}}\unsure{) de f sur S si}

\unsure{pour toute extension} \illegible{1} \unsure{de} $\unsure{k(s)}$ \unsure{et tout point} $\unsure{x'}$ \unsure{de} $\unsure{X'}$ \unsure{au-dessus de x}

\unsure{la localisée de} $\unsure{x'}$ \unsure{du point correspondant} \unsure{et de} $\unsure{X_s}$ \unsure{est un zéro}

\unsure{singulier (resp. singulier ordinaire) de} $\unsure{f_{s}}$.

\mspage{\hand{139}}

\clearpage
% Nouveau EGA V (cote 25): archive cote 25, site page 140 (PDF page 141).
% typescript. Draft compared: none.
% AI-transcribed by Claude Opus 5.5 (claude-opus-5-5, adaptive thinking, effort high), Message Batches API, batch pilot (msgbatch_0189Nhf1xB615YWnztubtLfc), 2026-11-12; not yet checked by eye.
\scanpage{25/140}
\margin{\uline{nouveau EGA V}}

IV 16

Un nouveau n$^{\text{o}}$ : \uline{Ensemble singulier et supersingulier d'une fonction},
\uline{et critères différentiels}.

Ce n$^{\text{o}}$ sera utilisé au par.20 des sections hyperplanes, mais sa
place naturelle me semble bien au par. 16.

\uline{Définition} 1 Soient X un préschéma régulier, $\phi$ une section de $\underline{O}_X$ . Un
point $x \in X$ est dit \uline{zéro} \xout{\illegible{3}} \uline{singulier} pour $\phi$
si on a $\phi_x \in \underline{m}_x^2$, zéro \uline{supersingulier} si c'est un zéro singulier
et si de plus l'élément de $\underline{m}_x^2/\underline{m}_x^3 \simeq \mathrm{Sym}(\underline{m}_x/\underline{m}_x^2)$ qu'il définit,
interprété comme une forme quadratique sur le dual $\underline{t}_x$ de $\underline{m}_x/\underline{m}_x^2$ sur k(x),
est une forme \deleted{non} dégénérée.( Un zéro singulier qui n'est pas supersingu-%
lier est parfois appelé \uline{zéro singulier ordinaire}).

\uline{Remarques} 2 Si $x \in V(\phi)$\hand{,} alors x est un zéro \inserted{non} singulier de $\phi$ sss
$\phi_x \neq 0$ et x est un point \xout{\unsure{régxxx}} non singulier i.e. régulier de $V(\phi)$,
i.e. sss x est un point régulier de $V(\phi)$ et $V(\phi) \neq X$ au voisinage de
x . \deleted{On peut prouver que si x est singulier pour $\phi$ , il est singulier
ordinaire sss $\phi_x \neq 0$ et la multiplicité de $\underline{O}_{V(\phi),x}$ est égale à 2
(ou, comme on dit encore parfois, x est un \uline{point singulier ordinaire}
\uline{du préschéma} $V(\phi)$ ); donc un $x \in V(\phi)$ est non supersingulier \inserted{pour $\phi$} sss
$\phi_x \neq 0$ et la multiplicité de $O_{V(\phi),x}$ est $\leq 2$.}

\uline{Définition} 3. Soient X un préschéma lisse sur un corps k, $\phi$ une
section de $\underline{O}_X$ , $x \in V(\phi)$ . On dit que x est un zéro \uline{géométriquement}
\uline{singulier} \xout{\unsure{dexx}} (resp. \uline{géométriquement \unsure{supersingulier}} ) de $\phi$ rel à k,
si pour toute extension k' de k et tout point \hand{$\xi$} de X à valeurs dans
k\hand{,} localisé en x, la \xout{\illegible{1}} point correspondant x' de $X_{k'}$, est un zéro
\xout{\unsure{géom}} singulier (resp. supersingulier) de $\phi_{k'}$ .

\mspage{140}

\clearpage
% Nouveau EGA V (cote 25): archive cote 25, site page 141 (PDF page 142).
% typescript. Draft compared: none.
% AI-transcribed by Claude Opus 5.5 (claude-opus-5-5, adaptive thinking, effort high), Message Batches API, batch ega25b (msgbatch_01A7n5YuXubUPNK1zsepi8RG), 2026-11-13; not yet checked by eye.
\scanpage{25/141}
\mspage{\hand{2}}

\uline{Remarques }4. \inserted[typed]{a)} Du \xout{xx} critère qui sera développé tout de suite, il résul-%
te que dans la définition 3, il suffit de tester avec un seul point
à valeurs dans un k' - on peut\xout{s} par exemple \unsure{prendre|pirendre} $k'=k(x)$ ou $\overline{k(x)}$, et
le point canonique à valeurs dans ce k'. b) Il résulte de la remar
que 2 que x est géométriquement non singulier pour $\phi$ sss $\phi_x \neq 0$ et
$V(\phi)$ est \uline{lisse} sur k en x. c) Supposons que l'on ait \inserted[typed]{un} préschéma X lisse
sur un autre \unsure{S|Y}, une section $\phi$ de $\underline{O}_X$ et un\xout{e} $x \in V(\phi)$, alors on dit que x
est un zéro singulier (resp. supersingulier) \uline{rel à \unsure{S|Y}}, s'il l'est
rel à k(s) sur la fibre $X_s$ (s étant l'image de x dans \unsure{S|Y}). d) Sous
les conditions de déf 1, on voit tout de suite que la singularité resp.
\margin{$\|$} supersingularité d'un $x \in V(\phi)$ pour $\phi$ n'est pas modifiée, quand on
remplace $\phi$ par $\phi'=u\phi$ , où u est une unité en x. Il s'ensuit aussitôt que
la \xout{xxxxxx} définition 1 et par suite aussi la déf. 3 s'étendent de façon
évidente au cas où $\phi$ est une section d'un Module inversible $\underline{L}$ (de façon
à redonner les définitions initiales pour $L = \underline{O}_X$) .

\xout{Considérons la section xxxxx d$^2$} Soit X un préschéma lisse sur un
autre Y, et soit $\phi$ une section de $\underline{O}_X$ , d'où une section $d^2_{X/Y}\phi$ de
$P^2_{X/Y}$ , se réduisant en la section $d^1_{X/Y}\phi$ de $P^1_{X/Y}$ , laquelle elle-même
% NOTE: from "Ceci posé" to the end of the page the text is enclosed by hand in a frame, with small tick marks along its lines; the marginal note on the left refers to it
se réduit en la section $d^0\phi$ \inserted[typed]{$=\phi$} de $P^0_{X/Y} = \underline{O}_X$ . Ceci posé, on a la
\uline{Proposition} 5 L'ensemble des zéros de $d^{\hand{0}}\phi$ ( resp $d^1\phi$ \hand{)} \deleted{resp. $d^2\phi$)}
\margin{garder le texte ci-contre} est respectivement égal à l'ensemble $V(\phi)$ des zéros de $\phi$ (resp. à
l'ensemble $V(\phi)^{\mathrm{sing}}$ des zéros de $\phi$ singuliers rel à S\hand{)} \deleted{resp. à l'ens.}
\deleted{$V(\phi)^{\mathrm{sup\,sing}}$ des zéros de $\phi$ supersinguliers rel à S)}.

La première assertion est triviale, la deuxième n'est autre
que le critère jacobien, ou si on préfère résulte de l'isomorphisme
canonique $\underline{m}_x/\underline{m}_x^2 \simeq \underline{\Omega}^1_{X/k}(x)$ lorsque x est un point rat/k d'un

\mspage{\hand{141}}

\clearpage
% Nouveau EGA V (cote 25): archive cote 25, site page 142 (PDF page 143).
% typescript. Draft compared: none.
% AI-transcribed by Claude Opus 5.5 (claude-opus-5-5, adaptive thinking, effort high), Message Batches API, batch ega25b (msgbatch_01A7n5YuXubUPNK1zsepi8RG), 2026-11-13; not yet checked by eye.
\scanpage{25/142}
\mspage{\hand{3}}

préschéma x sur k, \xout{enfin la dernière assertion résulte de la première}
\xout{précédente et de l'isomorphisme $\underline{m}_x$ analogue $gr^2(P^2_{X/k})(x)$ Sy}

Notons que $gr^1(\underline{P}^1_{X/Y}) \simeq \underline{\Omega}^1_{X/Y}$ , par suite la restriction
$d^1\phi\,|V(\phi)$ peut être interprété comme une section de $\underline{\Omega}^1_{X/Y} \otimes \underline{O}_{V(\phi)}$ ,
qui n'est autre d'ailleurs que la restriction de $d_{X/Y}\,\phi$ à $V(\phi)$ . Nous
pouvons alors considérer le \uline{préschéma des zéros de cette section} , que
nous noterons $V(\phi)^{sing}$ , et dont l'ensemble sous-jacent n'est autre
que l'ensemble des zéros de $\phi$ singuliers rel à Y en vertu de prop.5 .
( NB si \xout{\unsure{$\phi$}} $\mu$ est une section d'un Module localement libre \inserted[typed]{de type fini} $\underline{E}$ sur un
\margin{$\|$}\inserted[typed]{préschéma} \xout{espace annelé} X, on définit de façon évidente le sous-préschéma d es
zéros de $\mu$ , par exemple comme défini par l'Idéal image de $\check{\underline{E}} \longrightarrow \underline{O}_X$
transposé de $\mu$ ; lorsque $\underline{E} = \underline{O}_X{}^n$ et $\mu = \xout{\unsure{(\mu_i)}}\,(\mu_1,..,\mu_n)$, alors cet
idéal n'est autre que $\sum \mu_i\,\underline{O}_X$ ,qui définit $V(\mu_1,..,\mu_n)$ ). Prenant
maintenant la restriction $d^2\phi\,|V(\phi)^{sing}$ , et notant que
$gr^2(P^2_{X/Y}) \simeq Sym\xout{m}^2(\underline{\Omega}^1_{X/Y})$, on trouve une section canonique $M(\phi)$ de
$Sym^2(\underline{\Omega}^{\,1}_{X/\unsure{Y|X}}) \otimes \underline{O}_{V^{sing}}$ . On vérifie immédiatement que (prenant des
\margin{$\|$}point de X à val\xout{eu}rs dans des corps ...) cette section est précisément
celle qui détermine les formes quadratiques envisagées dans déf 1
(dans le cas \xout{\illegible{1}} d'un $X_k$ déduit de X/S par Spec(k)$\rightarrow$ S). On
en déduit une description d e l'ensemble $V(\phi)^{sup\ sing}$ en termes de
cette section, de la façon suivante: interprétant $M(\phi)$ comme définissan
% NOTE: the end of this line is cut by the right edge of the scan (presumably "définissant" / "définissant un")
homomorphisme
\[
M(\phi)\hand{$'$} : \ \hand{$\mathcal{T}$}_{X/Y} \otimes_{\underline{O}_X} \underline{O}_{V^{sing}} \longrightarrow \underline{\Omega}^1_{X/Y} \otimes_{\underline{O}_X} \underline{O}_{V^{sing}} \quad ,
\]
il faut prendre l'ensemble des points en lesquels cet homomorphisme
n'est pas un isomorphisme. Cela montre en particulier que $V(\phi)^{sup\ sing}$
est un ensemble fermé. On peut \deleted{le} préciser \inserted{ce dernier}\deleted{, en} introduisant
\[
D(\phi) = \text{dét}\ M(\phi) \ \in \ \Gamma\,(\underline{\Omega}^d_{X/Y}{}^{\otimes 2} \otimes \underline{O}_{V^{sing}}\ ) \quad ,
\]

\mspage{\hand{142}}

\clearpage
% Nouveau EGA V (cote 25): archive cote 25, site page 143 (PDF page 144).
% typescript. Draft compared: none.
% AI-transcribed by Claude Opus 5.5 (claude-opus-5-5, adaptive thinking, effort high), Message Batches API, batch ega25b (msgbatch_01A7n5YuXubUPNK1zsepi8RG), 2026-11-13; not yet checked by eye.
\scanpage{25/143}
\mspage{\hand{4}}

en supposant X de dimension relative d sur Y en chaque point.  Il y a
lieu de  désigner par \xout{XXX$^{\text{xxp}}$} $V(\phi)^{\text{sup sing}}$ le sous-préschéma fermé de
\xout{Xxxxx} $V(\phi)^{\text{sing}}$ \inserted[typed]{donc de X} défini par l'annulation de cette section (d'un Module
inversible cette fois ci),  dont l'ensemble  sous-jacent est ce qu'il
faut. Il y a lieu de résumer cette construction dans une

\uline{Proposition} 6.

Dans le cas général, on ne peut \xout{dire} dire rien de plus précis sur
$V(\phi)^{\text{sing}}$ et $V(\phi)^{\text{sup sing}}$ .  Nous allons  examiner maintenant un cas
particulier\hand{,} intéressant \xout{\unsure{dans}}pour certaines applications. Nous supposerons
que Y est également lisse sur un préschéma S, avec une dimension relative
constante m pour fixer les idées.  De plus, nous supposons que $V(\phi)^{\text{sing}}$,
que nous notons simplement $V^1$ pour simplifier, défini  par l'annulation
de la section $d^1\phi$  du Module $\underline{P}^1_{X/Y}$ \xout{de rang d} loc libre de \xout{rang} rang d+1,
est  lisse sur S de dimension relative  $(m+d) - (d+1) = m-1$ . (NB  Bien
entendu, les notations $V(\phi)^{\text{sing}}$ et $V(\phi)^{\text{sup sing}}$ sont ambigues en ce sens
qu'elles ne font pas intervenir le préschéma de base \unsure{X|Y} auxquelles elles
se rapportent; dans le cas actuel,  il est sous-entendu que c'est Y,
et on notera d'ailleurs \inserted[typed]{qu'il résulte aussitôt des hypothèses} que tout zéro de $\phi$ \inserted[typed]{singulier} est \uline{non singulier relative-}%
\uline{ment à} S ).  Dans cette situation, nous pouvons écrire  le diagramme
suivant \inserted{de Modules loc. libres sur $V^1$} :

\margin{(D)}
\begin{diagramblock}\begin{tikzcd}
 & & 0 & & \\
 & & \underline{\Omega}^1_{X/Y}\otimes\underline{O}_{V^1} \arrow[u] & & \\
0 \arrow[r] & \underline{P}^{\hand{\vee}1}_{X/Y}\otimes\underline{O}_{V^1} \arrow[r] \arrow[ur, "\hand{\mu}"] & \underline{\Omega}^1_{X/S}\otimes\underline{O}_{V^1} \arrow[r] \arrow[u] & \underline{\Omega}^1_{V^1/S} \arrow[r] & 0 \\
 & \hand{\underline{O}_{V^1}} \arrow[u] \arrow[r] & \underline{\Omega}^1_{Y/S}\otimes\underline{O}_{V^1} \arrow[u] \arrow[ur, "\hand{\nu}"'] & & \\
\hand{L=\omega^{-2}_{X/Y}\otimes\underline{O}_{V^1}} \arrow[ur] & & 0 \arrow[u] & &
\end{tikzcd}\end{diagramblock}

\margin{143}

\clearpage
% Nouveau EGA V (cote 25): archive cote 25, site page 144 (PDF page 145).
% typescript. Draft compared: none.
% AI-transcribed by Claude Opus 5.5 (claude-opus-5-5, adaptive thinking, effort high), Message Batches API, batch ega25b (msgbatch_01A7n5YuXubUPNK1zsepi8RG), 2026-11-13; not yet checked by eye.
\scanpage{25/144}
\mspage{\hand{5}}

La colonne provient de la suite exacte de transtitivité pour les
morphismes lisses $X\to Y$ et $Y\to S$ , en \xout{xx} tensorisant par $\underline{O}_{V1}$ (cela reste
exact, car tous les modules de la suite sont loc libres). La ligne
horizontale est un\xout{e} cas particulier d'une suite exacte obtenue chaque
fois que sur X \deleted{\xout{lisse}} sur S, on a une section \unsure{$\phi$} d'un module loc libre F
et qu'on prend le schéma W des zéros; on trouve une suite exacten
\[ \overset{\hand{\vee}}{F}\otimes \underline{O}_X \longrightarrow \underline{\Omega}^1_{X/S}\otimes \underline{O}_W \longrightarrow \underline{\Omega}^1_{W/S} \longrightarrow 0 \]
et \inserted{(si X/S est lisse)} le premier homomorphisme est injectif exactement en les points où
W est lisse sur X avec la ``bonne '' \unsure{dimension} relative (i.e. partout dans
\inserted[typed]{\xout{pratiquement}}
le cas actuel). \xout{Cet homomorphisme} Cette suite exacte est \xout{écrite en}
\xout{par 16 dans le cas où F = $\underline{O}_X$, la généralisation} \xout{(à ceci près} une
conséquence immédiate de la suite exacte
\[ \underline{J}/\underline{J}^2\longrightarrow \underline{\Omega}^1_{X/S}\otimes \underline{O}_W \longrightarrow \underline{\Omega}^1_{X/S} \longrightarrow 0 \]
\margin{16 \textbar} qui figure au par.16 (on y pourrait mettre en corollaire la version
signalée ici). La caractérisation de l'ensemble des points où on
peut mettre un zéro à gauche est contenue dans le critère \over{j}{J}acobien.

Notons que l'on a un isomorphisme canonique
\[ \underline{P}^1_{X/S} = \underline{\Omega}^1_{X/S} + \underline{O}_X \ ,\ \text{d'où}\ \check{\underline{P}}^1_{X/S} = \hand{\unsure{$\mathcal{G}$|$\mathcal{T}$}}_{X/S} + \underline{O}_X \ , \]
d'autre part on \unsure{vérifi\over{e}{é}} que l'homomorphisme \inserted{composé} \hand{$\mu$} du diagramme \hand{\unsure{(1)}} est nul
sur le facteur $\underline{O}_{V1}$ , et sur le facteur $\hand{\unsure{$\mathcal{G}$|$\mathcal{T}$}}_{X/S}\otimes\underline{O}_{V1}$ \inserted{il} se réduit à
l'homomorphisme \inserted{$M(\phi)'$} déduit de la section $M(\phi)$ de $\underline{\mathrm{Sym}}^2(\hand{$\underline{\Omega}$}^1_{X/S})\otimes \underline{O}_{V1}$ déjà
envisagée. Donc en un point x de X, $M(\phi)$ est non dégénéré i.e. $M(\phi)'$
est surjectif sss \hand{$\mu$} est surjectif \inserted{en x}, et on voit sur le diagramme \hand{(1)},
que cela équivaut aussi à dire que \unsure{\hand{$\nu$}} est surjectif \inserted{en x} (car l'un et l'autre
signifie\inserted{nt} que l'homomorphisme canonique de\xout{xx deux} la somme des deux sous
Modules envisagés de $\underline{\Omega}^1_{X/S}\otimes \underline{O}_{V1}$ dans ce derniers est surjectif en
x). On trouve ainsi :

\margin{144}

\clearpage
% Nouveau EGA V (cote 25): archive cote 25, site page 145 (PDF page 146).
% typescript. Draft compared: none.
% AI-transcribed by Claude Opus 5.5 (claude-opus-5-5, adaptive thinking, effort high), Message Batches API, batch ega25b (msgbatch_01A7n5YuXubUPNK1zsepi8RG), 2026-11-13; not yet checked by eye.
\scanpage{25/145}
\margin{6}

\uline{Proposition} 7. Sous les conditions précédentes (à rappeler), l'ensemble
sous-jacent à $V(\phi)^{\text{sup sing}}$ n'est autre que l'ensemble des points de
$V(\phi)^{\text{sing}}$ \xout{\illegible{1}} où le morphisme $V(\phi)^{\text{sing}}\longrightarrow Y$ (de préschémas lisses
sur S de dimension relative m-1 et m respectivement) est \uline{ramffié}.

En termes de papa (qu'il faudrait signaler en remarque), un point
$x \in V(\phi)$ est donc supersingulier rel à Y sss ``il consiste en au
moins deux points singuliers confondus'' ....

On peut et doit préciser prop.7 du point de vue \inserted{d'une identité de} \xout{des} sous-présché-%
mas et non seulement de sous-ensembles. En effet, $V(\phi)^{\text{sup sing}}$ a été
défini comme un sous-préschéma fermé de X , or on peut également définir
un sous-préschéma fermé naturel de $V^1$ de telle façon que l'ensemble
sous-jacent soit l'ensemble des points de ramification par rapport à Y.
En effet, il faut exprimer l'ensemble des points où un certain homomor-%
phisme de modules loc libres $Q$ ( $=\underline{\Omega}^1_{Y/S}\otimes\underline{O}_{V^1}$ ) $\longrightarrow$ $R$ ( $=\underline{\Omega}^1_{V^1/S}$ )
n'est pas surjectif. Si q et \unsure{m|r} sont les rang respectifs, c'est aussi
l'ensemble d es points où $\bigwedge^{\unsure{m|r}} Q \rightarrow \bigwedge^{\unsure{m|r}} R$ n'est pas surjectif, ce qui est
aussi l'ensemble des zéros de la section évidente de \uline{Hom} ( $\bigwedge^{\unsure{m|r}} Q$ , $\bigwedge^{\unsure{m|r}} R$)
$\simeq (\bigwedge^{q-r} Q)\otimes(\bigwedge^{r} R)$\hand{\illegible{1}}$\otimes(\bigwedge^{q} Q)$\hand{$^{\vee}$} , donc l'ensemble sous-jacent au \inserted{sous-}préschéma
fermé des zéros de cette section, soit \uline{Ram}($V^1/Y$) . Je dis que ce dernier
sous-schéma est identique à $V(\phi)^{\text{sup sing}}$ . C'est un simple exercice
sur le diagramme plus haut, compte tenu que $V(\phi)^{\text{sup sing}}$ est défini
par le même procédé que celui explicité pour $Q\longrightarrow R$, mais en termes de
l'homomorphisme $P$ ( $= \overset{\vee}{\underline{P}}{}^1_{X/Y}\otimes\underline{O}_{V^1}$ ) $\longrightarrow$ $S$ ($=\underline{\Omega}^1_{X/Y}\otimes\underline{O}_{V^1}$) , comme
il résulte de la description de \hand{$\mu$} donnée plus haut. On est alors
ramené à la situation générale suivante:

On a sur un \deleted{préschéma} \inserted{\unsure{espace annelé}} W un module localement libre \inserted{M} de rang m,
et deux sous-modules localement libres P et Q de rangs respectifs

\mspage{145}

\clearpage
% Nouveau EGA V (cote 25): archive cote 25, site page 146 (PDF page 147).
% typescript. Draft compared: none.
% AI-transcribed by Claude Opus 5.5 (claude-opus-5-5, adaptive thinking, effort high), Message Batches API, batch ega25b (msgbatch_01A7n5YuXubUPNK1zsepi8RG), 2026-11-13; not yet checked by eye.
\scanpage{25/146}
\mspage{7}

p et q tels que $p+q=m+1$ , on utilise la construction précédente relati
vement aux morphismes $P\longrightarrow M/Q = S$ et $Q\longrightarrow M/P = R$ , pour trouver
des sections \xout{dx} a de $P \otimes \det S \otimes \det P^{\unsure{\hand{-}}1} = P \otimes \det M \otimes \det P^{-1} \otimes$
$\det Q^{-1}$ et b de $.. = Q \otimes \det M \otimes \det P^{-1} \otimes \det Q^{-1}$ , qu'on peut
aussi envisager comme des homomorphismes de $\underline{L} = \det P \otimes \det Q \otimes \det M^{-1}$
dans P resp. Q . (NB on dénote, pour un Module loc libre F, par dét F
sa puissance extérieure maxima, et on utilise que, pour une suite
exacte $0 \quad F' \quad F \quad F'' \quad 0$ de tels Modules, on a un isomorphisme
canonique $\det F \quad \det F' \otimes \det F''$ ). Ceci posé, \uline{on a commutativité}
\uline{dans le diagramme}

\begin{diagramblock}
\begin{tikzcd}
 & P \arrow[r, hook] & M \\
\det P\otimes\det Q\otimes\det M^{-1}=L \arrow[r, phantom] & L \arrow[u, "a"] \arrow[r, "b"] & Q \arrow[u, hook]
\end{tikzcd}
\end{diagramblock}
,

(évdentuellement au signe près, cela dépend peut-être des conventions
adoptées pour définir certains des isomorphismes canoniques utilisés,
qui comportent le choix d'un signe ...), donc les idéaux $\unsure{\underline{I}}(a)$ et $\unsure{\underline{I}}(b)$
sont identiques (puisque P et Q sont loc facteurs directs dans M). \hand{NB en}
\hand{les points $\notin V(a)=V(b)$ , L est exactement l'intersection de P et Q dans M.}

Il reste, pour \unsure{clariffer|clarifier} complètement la situation particulière
envisagée avec X/Y/S et $\phi$ , décrite dans le diagramme \inserted{(D)} plus haut, \xout{dx}
d'expliciter a et b . On trouve d'abord
\[ \det P = \hand{\omega}_{X/Y}^{\unsure{\hand{-}}1} \otimes \underline{O}_{\unsure{X|V}}1 \quad , \quad \det Q = \hand{\underline{\omega}}_{Y/S} \otimes \underline{O}_{V}1 \]
\[ \det M = \hand{\omega}_{X/S} \otimes \underline{O}_{V}1 \ = \hand{\omega}_{X/Y} \otimes \hand{\omega}_{Y/S} \otimes \underline{O}_{V}1 \]
(le dernier isomorphisme provenant de la suite exacte de transitivité
pour les $\hand{\underline{\Omega}}^1$ pour les morphismes\xout{x} lisses $X-Y$ et $Y-S$), d'où
\[ \underline{L} \hand{\simeq} \hand{\omega}_{X/Y}{}^{-2} \otimes \underline{O}_{V}1 \]
(NB $\hand{\omega}_{X/Y}$ désigne \xout{lxxp} dét $\hand{\underline{\Omega}}^1_{X/Y}$ ) . Ceci dit, \uline{l'\unsure{h}omorphisme} \uline{a} ce \hand{---}
\uline{factorise par la composante $\underline{O}_{V1}$}

\mspage{146}

\clearpage
% Nouveau EGA V (cote 25): archive cote 25, site page 147 (PDF page 148).
% typescript. Draft compared: none.
% AI-transcribed by Claude Opus 5.5 (claude-opus-5-5, adaptive thinking, effort high), Message Batches API, batch ega25b (msgbatch_01A7n5YuXubUPNK1zsepi8RG), 2026-11-13; not yet checked by eye.
\scanpage{25/147}
\mspage{8}

\uline{de} $\overset{\hand{$\vee$}}{\underline{P}}_{X/Y}\otimes\underline{O}_{V^1}$ , \uline{et en tant qu'homomorphisme dans c ette composante,}
\uline{il est déterminé précisément par lasedtion} $D = \text{dét}\,M(\phi)$ \uline{déjà envisagée}
\uline{par ailleurs} (tjrs modulo un signe à déterminer). Bien entendu,
l'homomorphisme b) s'en déduit, compte tenu de la commutativité énoncée
dans le lemme. \deleted{Aux}\inserted{Les} points non supersinguliers$_{\text{x}}$ sont donc exacte\unsure{m}ent
ceux en lesquels $\underline{a}$: $\underline{\omega}^{-2}_{X/Y}\otimes\underline{O}_{V^1}\longrightarrow\underline{O}_{\unsure{V}^1}$ est un isomorphisme, et en
ces points, \xout{\unsure{xxxxxxdéfinitxxxxisomorphismexxdu}} le facteur $\underline{O}_{\unsure{V}^1}$ s'identif
fie donc exactement à l'intersection de $\underline{P}^1_{X/Y}\otimes\underline{O}_{V^1}$ et $\underline{\Omega}^1_{Y/S}\otimes\underline{O}_{V^1}$
dans $\underline{\Omega}^1_{X/S}\otimes\underline{O}_{V^1}$ . Donc sur l'ouvert de\xout{x} $V^1$ formé des points non
supersinguliers, on a une suite exacte
\[ 0\longrightarrow\underline{O}_{V^1}\longrightarrow\underline{\Omega}^1_{Y/S}\otimes\underline{O}_{V^1}\longrightarrow\underline{\Omega}^1_{V^1/S}\longrightarrow 0 \qquad \xout{x}, \]
ainsi que la suite exacte
\[ 0\longrightarrow\underline{O}_{V^1}\longrightarrow\underline{\Omega}^1_{X/S}\otimes\underline{O}_{V^1}\longrightarrow\xout{x}(\underline{\Omega}^1_{X/Y}+\underline{\Omega}^1_{Y/S})\otimes\underline{O}_{V^1}\longrightarrow 0 \]

\margin{|} Il faudrait résumer  toute cette situation dans une proposition
récapitulative, et sous une forme légèrement plus générale, en partant
d'une section $\phi$ d'un \uline{Module inversible} \xout{xxx} K \xout{sur} X - ce qui oblige
à tordre par K un certain\xout{s} nombre\xout{s} des Modules introduits dans les
considérations précéde\unsure{nt|mm}es. C'est en effet  cette situation qu'on
rencontre  dans le par. des sections hyperplanes. Signaler aussi que
\margin{|} $V(\phi)^{\text{sing}}$ \xout{xxxxx} et $V(\phi)^{\text{sup sing}}$ ne changent pas en r emplaçant $\phi$ par u $\phi$,
où u est une unité, \hand{en fait $d\phi|V(\phi)$ et $M(\phi)$ sont multipliés par \unsure{$u$},}
\inserted{et $D(\phi)$ est multiplié par $u^n$ ($n$ = dim. relative de $X$ sur $Y$),}
\uline{Remarque} Sous les conditions envisagées sur X,Y,S,$\phi$ , comme \xout{\unsure{x(x}})
$U = V(\phi)^{\text{sing}}- V(\phi)^{\text{sup sing}}$ est non ramffié sur Y, ses fibres sur Y sont
discrètes, ce qui implique que pour tout $y\in Y$, tout zéro singulier ordi-%
naire de $\phi_y$ dans $X_y$ est isolé dans l'ensemble des zéros singuliers, \deleted{et}
\deleted{à fortiori est}\inserted{i.e. est} isolé dans l'ensemble des points non lisses de $V(\phi_y)/k(y)$.
\xout{Parxsuite} On en conclut \inserted{p.ex.} aussitôt que si $d\geqslant 2$, c'est un point géométri-%

\margin{147}

\clearpage
% Nouveau EGA V (cote 25): archive cote 25, site page 148 (PDF page 149).
% typescript. Draft compared: none.
% AI-transcribed by Claude Opus 5.5 (claude-opus-5-5, adaptive thinking, effort high), Message Batches API, batch ega25b (msgbatch_01A7n5YuXubUPNK1zsepi8RG), 2026-11-13; not yet checked by eye.
\scanpage{25/148}
\mspage{\hand{9}}

ment normal de \xout{X}$V_y$ sur k(y), (puisque de toutes façons \xout{X}  $V_y$ est CM ).
On peut montrer  plus généralement, sous les conditions de déf.1 , que si
x est un zéro singulier ordinaire, alors x est une singularité isolée
de V($\phi$)  en ce sens que toute générisation \xout{\unsure{de}} x' de x dans V($\phi$) est un
point régulier de V($\phi$)  (donc si le lieu singulier de V($\phi$) est fermé, p.
ex. si X est ``ex\xout{c}cellent'' , x est \xout{bien} un point isolé de ce lieu singuli-%
er). \xout{Lorsque} (Pour bien faire, il faudrait inclure ce résultat au Chap
IV , mais  le par. 16 ne semble guère le lieu. Où le mettre ?).  Il
s'ensuit que si X est lisse sur un corps k, et si x est un zéro géométri%
quement singulier ordinaire, alors  x est un point isolé de l'ensemble
des points de non lissité  de V(\over{\unsure{4}}{$\phi$}) , donc est \inserted{(géom.)} normal si  $\dim_x(X) \geqslant 2$ ,
(mais \inserted[typed]{nécessairement} non normal si $\dim_x X = 1$ , comme on le \deleted{vérifie} \inserted{constate} par exemple sur
le cas type $\phi = xy$ , x et y fonctions coordonnées dans le plan affine).

\hand{Voici le diagramme analogue à (D), mais pour}
\hand{une section $\varphi$ d'un Module inversible K :}

\hand{(D')}
\begin{diagramblock}
\begin{tikzcd}
 & & 0 & & \\
 & & \Omega^1_{X/Y}\otimes\underline{O}_{V'} \arrow[u] & & \\
0 \arrow[r] & P^1_{X/Y}(K)^{\vee}\otimes\underline{O}_{V'} \arrow[r] \arrow[ur, "\mu"] & \Omega^1_{X/S}\otimes\underline{O}_{V'} \arrow[r] \arrow[u] & \Omega^1_{V'/S} \arrow[r] & 0 \\
 & K^{-1}\otimes\underline{O}_{V'} \arrow[u] \arrow[r] & \Omega^1_{Y/S}\otimes\underline{O}_{V'} \arrow[u] \arrow[ur, "\nu"'] & & \\
\omega^{\unsure{-2}}_{X/Y}\otimes K^{\unsure{-2n-1}}\otimes\underline{O}_{V'} \arrow[ur] & & 0 \arrow[u] & &
\end{tikzcd}
\end{diagramblock}

\mspage{\hand{148}}

\clearpage
% Nouveau EGA V (cote 25): archive cote 25, site page 149 (PDF page 150).
% typescript. Draft compared: none.
% AI-transcribed by Claude Opus 5.5 (claude-opus-5-5, adaptive thinking, effort high), Message Batches API, batch ega25b (msgbatch_01A7n5YuXubUPNK1zsepi8RG), 2026-11-13; not yet checked by eye.
\scanpage{25/149}
\mspage{10}
% NOTE: a long handwritten diagonal stroke crosses the whole text, from ``Arrivé'' at the top down to the foot of the page. It may cancel the page or mark it as used. The text under it is transcribed as live text.

Arrivé à ce point \unsure{admirable|admitable} de généralité, je m'aperçois malheureuse-%
ment que ces développements sont encore trop particuliers, et qu'au lieu
d'un sous-préschéma défini par une seule section de $\underline{O}_X$ , il y a lieu de
considérer celui défini par m sections, ou ce qui\hand{/}revient au même, par un
homomorphisme $\underline{E} \longrightarrow \underline{O}_X$ , où $\underline{E}$ est un Module loc libre de rang m. Je te
laisse le soin de décider s'il y a lieu dès le début de commencer avec le
cas général, suivant le principe qu'il faut procéder du général au parti-%
culier.

Introduisons d'abord une terminologie pour des anneaux A à
augmentation $A \to A/I$ (qui sera utilisée également dans le contexte
faisceautisé). L'augmentation est dite \uline{régulière} (ou devons nous dire
quasi-régulière, avec les définitions adoptées au par.19 ?) si \inserted{a)} $I/I^2$ est
un A/I -Module projectif de type fini, et \inserted{b)} l'homomorphisme
\[ \theta\colon \mathrm{Sym}\,(I/I^2) \longrightarrow \mathrm{gr}_I(A) \]
est un isomorphisme. Nous dirons que l'augmentation est \inserted{de type} ``\uline{quadratique}
\uline{ordinaire}'' lorsque a) ci-dessus est vérifié, ainsi que b') l'homomor-%
phisme canonique précédent à un noyau engendré par $\theta^2$ , lequel noyau $\underline{K}$
est un \inserted{sous-}Module inversible ``non dégénéré'' de $\mathrm{Sym}^2(I/I^2)$ . Par ``non dégé%
néré'' nous entendons que l'homomorphisme correspondant
\[ \underline{K} \otimes (\underline{I}/\underline{I}^2)^{\hand{\vee}} \longrightarrow \underline{I}/\underline{I}^2 \]
(déduit en utilisant l'homomorphisme canonique
\[ \mathrm{Sym}^2(I/I^2) \longrightarrow \mathrm{Hom}\ (\ (I/I^2)^{\hand{\vee}}\ ,\ I/I^2\ )\quad ) \]
est bijectif, ou ce qui revient au même surjectif (car il s'agit de
Modules projectifs \inserted{sur A/I} de même rang en chaque point), ou encore : prenant
localement une base \inserted{$\phi$} de $\underline{K}$ , \xout{l'homomorphisme correspondant} la forme
quadratique correspondante sur $(I/I^2)^{\hand{\vee}}$ est non dégénérée. Il suffit

\mspage{149}

\clearpage
% Nouveau EGA V (cote 25): archive cote 25, site page 150 (PDF page 151).
% typescript. Draft compared: none.
% AI-transcribed by Claude Opus 5.5 (claude-opus-5-5, adaptive thinking, effort high), Message Batches API, batch ega25b (msgbatch_01A7n5YuXubUPNK1zsepi8RG), 2026-11-13; not yet checked by eye.
\scanpage{25/150}
\mspage{11}
% NOTE: a single handwritten diagonal stroke runs across the whole page, from the top right (above "conditions") to the bottom left (near "k >=1"); it cancels the whole text of the page, which stays legible under it.
d'ailleurs de vérifier ces dernières conditions \xout{xx}\inserted{pour} les fibres réduites
en les différents idéaux premiers (ou même seulement maximaux) de A/I .
\uline{Proposition} 8. Pour que l'augmentation A$\rightarrow$A/I soit \inserted{de \unsure{type}} quadratique
ordinaire, il f et s que les conditions suivantes soient vérifiées:
1º) \xout{xxx2} \xout{I/Ixxxxxx} \inserted[typed]{Les $gr^{n}_{I}(A)$ sont} projectifsde type fini sur A/I .

2º) Si d\xout{x} désigne le rang de $gr^{1}_{I}(A) = I/I^{2}$ en un $\underline{p} \in$Spec(A/I),
alors le rang de $gr^{n}_{I}(A)$ en ce même point est égal (ou simplement
$\geqslant$ ) à $\binom{n+d}{d} - \binom{n-2+d}{d}$ .

3º) (1º et 2º impliquent déjà que $\underline{K}$ = Ker $\underline{\theta}^{\unsure{2}}$ est un Module
inversible localement facteur direct dans \xout{xxxx} Sym$^{2}$(I/I$^{2}$) ). Le sous
Module $\underline{K}$ = Ker $\theta^{2}$ de Sym$^{2}$(I/I$^{2}$) est non dégénéré.

La démonstration est immédiate à partir du lemme suivant; \deleted{que je te
laisse le soin de démontrer:}

\xout{Propxsixxxx} \uline{Lemme} 8.1. Soient B un anneau, M un B-module projectif de
type fini, $\phi \in$ Sym$^{2}$(M) un élément ``non dégénéré'' , alors la multiplica-%
tion par $\phi$
\[ \mathrm{Sym}^{n\hand{\unsure{-2}}}(M) \longrightarrow \mathrm{Sym}^{n\deleted{\unsure{+2}}}(M) \]
est un isomorphisme sur un sous-module facteur direct, donc le \unsure{conoyau}
est un Module projectif de rang $\binom{n+d}{d} - \binom{n-2+d}{d}$ si M est de rang d.

On est immédiatement ramené au cas où B est le spectre d'un corps,
où c'est trivial (en fait, le ``non dégénéré'' nous a servi uniquement
\margin{||} pour assurer que les $\phi(x)$, $x \in$ Spec(B), sont tous $\neq 0$, et on pourrait
formuler le lemme plus raisonnablement en remplaçant 2 par un entier
k $\geqslant$1 quelconque, sans parler de non dégénérescence).

\uline{Proposition} 9. Soit A un anneau augmenté vers A/I, à augmentation
régulière, I/I$^{2}$ étant projectif de rang n sur A/I . Soit E un A-module

\mspage{150}

\clearpage
% Nouveau EGA V (cote 25): archive cote 25, site page 151 (PDF page 152).
% typescript. Draft compared: none.
% AI-transcribed by Claude Opus 5.5 (claude-opus-5-5, adaptive thinking, effort high), Message Batches API, batch ega25b (msgbatch_01A7n5YuXubUPNK1zsepi8RG), 2026-11-13; not yet checked by eye.
\scanpage{25/151}
% NOTE: a long diagonal pencil stroke crosses the whole page, from the lower left (near "(ii) Il existe") to the upper right (near "homomorphisme" in the first line); it runs through the text without hiding it
\mspage{12}

projectif de rang m,  $u{:}E \to I$ un homomorphisme, d'où un homomorphisme
\[ u\otimes id_{A/I} : E\otimes A/I \longrightarrow I/I^2 \ . \]
\xout{Conditions Posons} \inserted[typed]{Soit} $B = A/u(E)A$ , muni de l'augmentation $B\to B/J = A/I$
(où $J = I/u(E)A$ ). Les deux conditions suivantes sont équivalentes:

\margin{$\|$} (i) L'anneau \xout{augmenté} B \xout{\unsure{B/I est}} a une augmentation $B\to B/J$ \inserted{de type} quadra-%
tique ordinaire, de rang $n-(m-1)= n-m+1$ (le \hand{``}\uuline{rang de l'augmentation}\hand{''}
est \inserted{par définition} celui du Module projectif $J/J^2$ ).

\margin{$|$} (ii) $K = \mathrm{Ker}\ u\otimes id_{A/I}$ est un Module inversible, et l'homomor
\unsure{phisme} naturel $K \longrightarrow \mathrm{Sym}^2(J/J^2) = \mathrm{Sym}^2( (I/I^2)/\mathrm{Im}\ E\otimes A/I )$ est
``non dégénéré'' .

Il serait peut-être plus clair d'introduire une notation $V = I/I^2$ ,
et $W = J/J^2 = V / \mathrm{Im}\ E\otimes A/I$ . L'homomorphisme qualifié de naturel
dans (ii) s'obtient en notant que K s'envoie en tous cas dans
\margin{$|$} $I^2/(I^3 +$ \xout{Ixm(E)} $u(IE)) = (I^2/I^3)/\mathrm{Im}\ IE \simeq \mathrm{Sym}^2(I/I^2)/\ \mathrm{Im}\ (I/I^2)\otimes(E\otimes A/$
% NOTE: the end of this line runs off the right edge of the scan
où $I/I^2 \otimes(E \otimes A/I) \to \mathrm{Sym}^2(I/I^2)$ est défini via l'homomorphisme
canonique de multiplication $I/I^2 \otimes I/I^2 \to \mathrm{Sym}^2(I/I^2 )$ . Ces explici-%
tations posées, la démonstration de la prop.9 est à peu près triviale.

\uline{Corollaire} 10. Soit A un anneau local régulier de dimension n, $\underline{m}$ son
idéal maximal, \inserted[typed]{$k=A/\underline{m}$,} E un A-module libre de rang m, $u{:}E \to \underline{M}$ un homomorphisme,
\xout{Conditions équivalentes:} $B = A/u(E)A$, $\underline{n} = \underline{m}/u(E)A$ l'idéal maximal de B,
$V = \underline{m}/\underline{m}^2$ , $W = \underline{n}/\underline{n}^2 =$ \xout{km/} $V/\ \mathrm{Im}\ E\otimes k$ . Conditions équivalentes:

\margin{?\}} (i) B est de dimension n-m (i.e. , en termes d'une base \xout{de} $(\phi_i)$ de
\margin{$\|$} E, les\xout{t} images des $\phi_i$ dans A forment une suite A-régulière et \xout{est} le
\xout{singularité qu} point fermé de Spec(B) est une \uline{singularité quadratique}
ordinaire (par quoi nous entendons que l'augmentation $B\to B/\underline{n}$ est \inserted{de type} quadra-%
tique ordinaire).

(ii) Il existe une base $\phi_1 ,\ldots., \phi_m$ de E telle que \hand{\unsure{$\sqrt{\ }$}}$A/$\hand{$\sum_{i=1}^{m-1}$}$u(\phi i)A$

\mspage{151}

\clearpage
% Nouveau EGA V (cote 25): archive cote 25, site page 152 (PDF page 153).
% typescript. Draft compared: none.
% AI-transcribed by Claude Opus 5.5 (claude-opus-5-5, adaptive thinking, effort high), Message Batches API, batch ega25b (msgbatch_01A7n5YuXubUPNK1zsepi8RG), 2026-11-13; not yet checked by eye.
\scanpage{25/152}
\mspage{13}

% NOTE: a long handwritten diagonal stroke runs from the top edge (above "b) l'image") down to the left margin just above "Notes que", across the whole passage from "soit" to "non dégénérée."; it is transcribed as a deletion of that passage
\deleted{\inserted[typed]{= $A_{m-1}$}}
\deleted{soit \xout{\unsure{singulier}} régulier de dimension $n-(m-1)$, et b) l'image de $u(\phi_m)$}
\deleted{dans $A_{m-1}$ admette le point fermé de Spec($A_{m-1}$) comme \uline{zéro singulier}}
\deleted{\uline{ordinaire} (cf définition 1).}

\deleted{(ii bis) \xout{Pour toute} Le rang de Im($E\otimes k \to V$) est m-1, et pour}
\deleted{toute base $(\phi_i)_{1\leq i\leq m}$ de E telle que les images des $\phi_i$ $(1\leq i\leq m-1)$ \inserted[typed]{dans V}}
\deleted{sont linéairement indépendantes,\xout{\unsure{e}} les conditions a) et b) de (ii) sont}
\deleted{satisfaites.}

\deleted{(iii) Le noyau \xout{de} K de $E\otimes k \to V$ est de dimension 1, et si $\phi\in K$\hand{$-\{0\}$,}}
\deleted{\xout{la forme} l'élément de Symm$^2$(W) défini par $\phi$ est une forme quadratique}
\deleted{non dégénérée.}

% NOTE: a handwritten double vertical line in the left margin runs from "Notes que" down to "= dim O_{X,x}", with a hook-shaped mark at its top
Notes que l'analogue des conditions (ii) et (ii bis) aurait déjà
pu remonter dans la proposition 9. Il y a lieu d'introduire la terminolo-%
gie suivante: si X est un préschéma régulier, $\underline{E}$ un Module localement libre
\xout{de rang m} sur X, u: $\underline{E}\to\underline{O}_X$ un homomorphisme i.e. u une section de
$\underline{E}$\hand{$^\vee$} , on dira que l'élément x de $V(u)=V(u(\underline{E})\underline{O}_X)$ est un \xout{zéro non sing}
``\uline{zéro non singulier}'' pour u\hand{,} si u est $\underline{O}_X$-régulier i.e. dim $\underline{O}_{V(u),x}$
= dim $\underline{O}_{X,x}$ - rang$_x$ $\underline{E}$ , et si de plus x est un point régulier de V(u),
ce qui en termes d'une base $\phi_1,\ldots,\phi_m$ de $\underline{E}$ au voisinage de x s'exprime
aussi en disant que les $u(\phi_i)$ forment une partie d'un systême régulier
de paramètres de $\underline{O}_{X,x}$ ; dans le cas contraire, x est dit un \uline{zéro}
\uline{singulier}. On dit que x est un \uline{zéro singulier ordinaire} s'il satisfait
aux conditions équivalentes du corollaire 10 (ce qui implique que c'est
bien un zéro singulier), où évidemment $A=\underline{O}_{X,x}$ etc. On dit encore que
x est un \uline{zéro supersingulier} s'il est singulier, sans \hand{\^{e}}tre singulier
ordinaire. Donc non supersingulier = non singulier ou singulier ordinaire
Enfin, lorsque X est lisse sur un corps k, on introduit les variantes
géométriques de ces notions de façon évidente. On notera que dans ces

\mspage{152}

\clearpage
% Nouveau EGA V (cote 25): archive cote 25, site page 153 (PDF page 154).
% typescript. Draft compared: none.
% AI-transcribed by Claude Opus 5.5 (claude-opus-5-5, adaptive thinking, effort high), Message Batches API, batch ega25b (msgbatch_01A7n5YuXubUPNK1zsepi8RG), 2026-11-13; not yet checked by eye.
\scanpage{25/153}
\mspage{\hand{14}}

notions, mis à part \xout{xxx}\inserted[typed]{la} condition dimensionelle = la condition de
$\underline{O}_X$-régularité, le caractère de x se voit \xout{xxx} \unsure{en|eh} fait sur la nature du
point x dans V(u), i.e. sur l'anneau local $\underline{O}_{V(u),x}$ (éventuellement consi-%
déré comme algèbre sur k). Lorsque (dans le cas d'un corps de base k)
x est rationnel sur k, la notion relative à k coincide avec la notion
absolue, comme on voit \inserted{\unsure{précédemment|préalablement}} comme dans le cas où $\underline{E}$ est inversible. Enfin,
lorsque X est lisse sur \over{S}{Y}, et u: $\underline{E}\to\underline{O}_X$ est encore un homomorphisme avec
$\underline{E}$ loc libre de type fini, on définit les notions correspondantes relative-%
ment à S, en regardant fibre par fibre. On remarquera que si x est
non supersingulier relativement à \over{S}{Y}, alors V(u) est plat sur S \xout{relativ} en
x (car u est $\underline{O}_X$-régulier relativement à \over{S}{Y} ...).

\margin{$\big|$} Bien entendu (à signaler tout au moins en r emarque après les
définitions en forme) on pouvait tout aussi bien partir d'une section u
d'un Module loc libre $\underline{E}$ , les définitions précédentes s'appliquent en
interprétant u comme un homomorphisme $\check{\underline{E}}\longrightarrow\underline{O}_X$ .

Lorsque X est lisse sur \over{S}{Y}, u:$\underline{E}\longrightarrow\underline{O}_X$ est donné avec $\underline{E}$ loc libre
de rang m, alors l'ensemble des \xout{points} \inserted[typed]{zéros de u} singuliers rel à S est l'ensembl
sous-jacent à un sous-préschéma fermé $V(u)^{\mathrm{sing}} = V^1$ bien déterminé de
$V(u)=V^o$ , obtenu ainsi: on considère $d^1u$ section de $P^1_{X/\over{S}{Y}}(\check{\underline{E}})$ , sa
restriction à $V^o$ peut alors s'interpréter comme un homomorphisme
\[ \underline{E}\otimes\underline{O}_{V^o} \longrightarrow \underline{\Omega}^1_{X/\over{S}{Y}}\otimes\underline{O}_{V^o} \]
et $V(u)^{\mathrm{sing}}=V^1$ désignera le préschéma des zéros de \xout{\unsure{ce dernier. Rem}}
\xout{Considérant de même $d^2$ section de $P^2_{X/S}(\underline{E})$ puis $d^2$ $V^1$ on}
\xout{trouve un homomorphisme canonique} l'homomorphisme correspondant
\[ \bigwedge^m\underline{E}\otimes\underline{O}_{V^o} \to \underline{\Omega}^m_{X/\over{S}{Y}}\otimes\underline{O}_{V^o} \]
i.e. d'une section bien déterminée de $\underline{\Omega}^m_{X/\over{S}{Y}}\otimes\mathrm{dét}(E)^{-1}\otimes\underline{O}_{V^o}$ .

\margin{153}

\clearpage
% Nouveau EGA V (cote 25): archive cote 25, site page 154 (PDF page 155).
% typescript. Draft compared: none.
% AI-transcribed by Claude Opus 5.5 (claude-opus-5-5, adaptive thinking, effort high), Message Batches API, batch ega25b (msgbatch_01A7n5YuXubUPNK1zsepi8RG), 2026-11-13; not yet checked by eye.
\scanpage{25/154}
\mspage{\hand{15}}

Pour caractériser par voie différentielle  les  zéros singuliers ordinaires
relativement à \over{\$}{$S$},  on note qu'il faut simplement exprimer  les conditions
de la proposition 9  pour  l'homomorphisme \xout{\unsure{naturel}} déduit de u
\[ \xout{$d^{\infty}$\illegible{1}} \qquad P^{\infty}_{X/\over{\$}{S}}(\underline{E})_{V^1}\to(P^{\infty}_{X/\over{\$}{S}})_{V^1} \quad , \]
qui se réduit d'ailleurs à une condition  ``d'ordre 2'' i.e. portant seule-%
ment sur l'homomorphisme correspondant
\[ P^{2}_{X/\over{\$}{S}}(\underline{E})_{V^1}\to(P^{2}_{X/\over{\$}{S}})_{V^1} \ . \]
Il faut exprimer en premier lieu  que  le noyau \inserted{$\underline{K}$} de
\[ (\underline{x}\ ) \qquad \underline{E}_{V^1} \longrightarrow (\underline{\Omega}^{1}_{X/\over{\$}{S}})_{V^1} \]
en x est inversible et facteur direct, i.e. que l'homomorphisme envisagé
(qui en vertu de  $x\in V^1$ est  de rang  m-1 en x) est exactement de rang
\margin{(} m-1 en x  (NB \uline{le rang} d'un  homomorphisme de modules loc libres en un
point est par définition le rang de 'lhomomorphisme de vectoriels sur k(x)
correspondant), c e qui s'exprime en considérant
\[ (\bigwedge^{m-1} \underline{E})_{V^1}\to(\underline{\Omega}^{m-1}_{X/\over{\$}{S}})_{V^1} \quad , \]
le  schéma  V'' des zéros de cet homomorphisme  ( $V''\subset V^1$ )  et en écrivant
simplement $x\notin$ V'' .  La définition de V'' assure que  K $|V^1-V''$  est un
Module inversible,  et que  le conoyau  $\underline{W}$  de ($\underline{x}$ )  est tel que
$\underline{W}\,|V^1-V''$  est  localement libre de rang égal à  \unsure{d|n}-(m-1)  (où \unsure{d|n} est la
dimension relative de X sur S, égale au rang de $\underline{\Omega}^{1}_{X/\over{\$}{S}}$).  Appliquant la
construction expliquée avec la prop.9\hand{$^{*}$}, on trouve \xout{\unsure{maintenant}} un homomor-%
phisme
\[ M\ : \qquad \underline{K} \longrightarrow \underline{\mathrm{Symm}}^{2}(\underline{W}) \]
et il reste à exprimer que ce dernier est ``non dégénéré'' en x. Pour ceci,
introduisons le  discriminant
\[ D\ \in\ \Gamma\,(V^1-V'',\ \underline{K}^{\unsure{-n+m-1}}\otimes \text{dét}(\underline{W})^{2}\ ) \]
(où les exposants désignent des puissances tensorielles de Modules

\mspage{\hand{154}}

\clearpage
% Nouveau EGA V (cote 25): archive cote 25, site page 155 (PDF page 156).
% typescript. Draft compared: none.
% AI-transcribed by Claude Opus 5.5 (claude-opus-5-5, adaptive thinking, effort high), Message Batches API, batch ega25b (msgbatch_01A7n5YuXubUPNK1zsepi8RG), 2026-11-13; not yet checked by eye.
\scanpage{25/155}
\mspage{16}

inversibles). Posant $V^2 = V(D)$ , qui est un sous-préschéma fermé de
$V^1-V''$ , on trouve donc que l'ensemble des \xout{\unsure{points}} zéros singuliers
ordinaires rel.\xout{Y}S n'est autre que $V^1 - V'' - V^2$ , i.e. l'ensemble des
zéros supersinguliers \xout{\unsure{sur}} rel à S est $V''\ \ V^2$ .

\margin{$\|$} Au moment de l'introduction de V", en affirmant que $\underline{W}$ et $\underline{K}$ sont
localement libres sur $V^1-V''$, j'ai utilisé en f'ait général, que je dégage
ici en lemme (où $V^1-V''$ est devenu X):

\uline{Lemme} Soient $v\colon \underline{E} \to \underline{F}$ un homomorphisme de Modules localement libres
sur le préschéma X, \xout{\unsure{et soit}} m un entier $>0$, supposons que $\bigwedge^{m} u$ soit
nul. \xout{\unsure{(ce qui implique que le rang de u en chaque point est $\leq$ m) et que}} \inserted[typed]{Alors}
$\bigwedge^{m-1} u$ ne s'annule en aucun point (du point de vue fibres réduites, bien
sur) \xout{\unsure{xxxAlorsx}} \inserted[typed]{sss} u($\underline{E}$) est un sous-Module localement facteur direct de $\underline{F}$,
localement libre de rang m-1, \hand{ou encore ssi u est de rang \unsure{m}-1 en}
\hand{chaque pt de X.}

Il y a lieu de résumer ces constructions dans une proposition
qui généralise prop.6 . Il faut voir comment on peut généraliser de même
prop.7. Pour ceci, on suppose que Y lui-même est lisse sur un préschéma
S , \inserted[typed]{de dimension relative N,} donc X est lisse sur S, de dimension relative N+d\unsure{\hand{$^*$}}; où d est celle de X
sur Y. On supposera de plus avec les notations précédentes $V''=\emptyset$, i.e.
$\underline{W}$ localement libre de rang d-m+1 \inserted[typed]{(c'est tjrs le cas si m=1 !)}. Enfin, on supposera les sous-présché-%
mas fermés \xout{\unsure{Vx}} \inserted[typed]{$V^o$et $V^1$} de X lisses \uline{sur} S , de dimension relative \xout{minima} en chaque
point minima (pour les données d,N,m) , \xout{\unsure{exxqmx}} c'est à dire, comme
on constate aussitôt, de dimensions relatives (N+d) - m et
(N+d) - m - (d-m +1) = N - 1 . (Du point de vue rédaction, le mieux
semble de poser brutalement ces nombres, peut être plutôt en termes
de codimension fibre par fibre, celle de $V^o$ dans X étant m, celle de
$V^1$ dans $V^o$ étant d-m+1 ; et de signaler en remarque la justification
de ce choix par le principe \unsure{du|de} la dimension minima ou dimension \xout{générique}

\margin{155}

\clearpage
% Nouveau EGA V (cote 25): archive cote 25, site page 156 (PDF page 157).
% typescript. Draft compared: none.
% AI-transcribed by Claude Opus 5.5 (claude-opus-5-5, adaptive thinking, effort high), Message Batches API, batch ega25b (msgbatch_01A7n5YuXubUPNK1zsepi8RG), 2026-11-13; not yet checked by eye.
\scanpage{25/156}
\mspage{\hand{17}}

\xout{xg} ``générique'')\hand{\unsure{d}}.   Ceci posé, il faudrait vérifier que l'ensem ble $V^2$
des points supersinguliers de u rel. à Y n'est autre que l'ensemble des
points en lesquels $V^{(1)} \longrightarrow Y$ est ramifié, et préciser ce point
comme une identité de sous-préschémas fermés de X , $V^{\unsure{\over{2}{2}}}$ \over{=}{=} \hand{``}sous-préschéma
de ramification\hand{''} de $V^1 \longrightarrow Y$. Je n'ai pas fait l'exercice en détail,
mais ne doute pas qu'il ne puisse se faire par essentiellement \over{la}{le} même
\xout{méthode} \inserted{genre de dévissage} que cel\over{le}{ui} développé\deleted{e} après la prop.7.

\hand{\unsure{d}} Si $E \longrightarrow F$ est un homomorphisme de Modules localement libres \unsure{sur} \inserted{le préschéma régulier}
X, de rangs m et d ($m \over{<}{\leqslant} d$), et si \over{\unsure{Y}}{Z} est l'ensemble des points de X en
lesquels le rang de u est $\hand{\leqslant}$ m-1, alors \inserted{on peut montrer que} Z est de codimension $\hand{\geqslant}$ d-m+1
en chacun de ses points, \xout{comme on vérifie immédiatement par récurrence sur}
\xout{m, à partir du cas m=1 (}

\mspage{\hand{156}}

\clearpage
% Nouveau EGA V (cote 25): archive cote 25, site page 157 (PDF page 158).
% typescript. Draft compared: none.
% AI-transcribed by Claude Opus 5.5 (claude-opus-5-5, adaptive thinking, effort high), Message Batches API, batch ega25b (msgbatch_01A7n5YuXubUPNK1zsepi8RG), 2026-11-13; not yet checked by eye.
\scanpage{25/157}
% NOTE: the sheet has no legible text of its own; the only written mark is the handwritten page number at the foot, left of centre. Very faint typed text is visible across the whole sheet. It is read as show-through from the next sheet and is not transcribed (see NOTES).
\mspage{\hand{157}}

\clearpage
% Nouveau EGA V (cote 25): archive cote 25, site page 158 (PDF page 159).
% typescript. Draft compared: none.
% AI-transcribed by Claude Opus 5.5 (claude-opus-5-5, adaptive thinking, effort high), Message Batches API, batch ega25b (msgbatch_01A7n5YuXubUPNK1zsepi8RG), 2026-11-13; not yet checked by eye.
\scanpage{25/158}
\margin{nouveau EGA V} \mspage{1}

IV 17.15. \uline{Formes lisses}; \inserted{\unsure{augmentations|argumentations} et singularités élémentaires.} \deleted{\xout{x} formes quadratiques lisses}

NB Je viens de m'apercevoir que la terminologie introduite dans
ma rédaction \inserted[typed]{SUP} de la Supersingularité est déraisonnable, et conflicte
en particulier avec la terminologie recue. (De toutes façons, tu a\unsure{\xout{s}} dû
no\over{t}{t}er que dans la déf.1 de ce texte, il faut lire ``dégénérée'' au lieu de
``non dégénérée''). Le canular provient des variétés de dimension paire
en car.2 , \xout{\unsure{une telle variété}} définie\hand{s} par exemple par une équation $\phi=0$
dans une variété ambiante de dim impaire, car avec la terminologie de mes
notes, une telle variété ne peut avoir de \xout{\unsure{zéro sing}} \hand{``}point singulier
or\over{d}{d}inaire\hand{''} (= singularité quadratique ordinaire), i.e. $\phi$ ne peut avoir
de \hand{``}zéro singulier ordinaire\hand{''}, d\over{u}{û} au fait qu'en car 2, une forme quadratique
à un nombre impair de variables est toujours ``dégénérée''. Or tout le
monde a toujours considéré que même en car. 2, l'origine du cône affine
$x^2 + yz = 0$ est une singularité quadratique ordinaire. Dans les présen-%
tes notes, je dégage la notion de \xout{f}forme quadratique lisse (ou ``ordinaire''),
(sur un Module loc libre de type fini), de telle façon que
non dégénéré \hand{$\Longrightarrow$} lisse, la réciproque étant vraie si E est de rang pair
\uline{ou} les car résiduelles de la base S sont toutes $\neq 2$ . Cette notion étant
posée, je propose de garder la terminologie ``supersingulière'' (qui ne
conflicte avec aucune terminologie reçue) \xout{\unsure{xxxxxx de mon texte}} de SUP,
comme correspondant à la notion de forme quadratique non dégéné\over{f}{r}ée; on
parlera de même de point supersingulier d'un préschéma loc noeth. , ou
géométriquement supersingulier pour un préschéma loc de type fini sur un
corps, dans le même sens (en regardant l'hom $\mathrm{Sym}(\underline{M}/\underline{m}^2) \hand{\longrightarrow} \mathrm{gr}_{\underline{m}}(A)$ et
\unsure{en|ex} disant que le noyau ne peut être engendré par une forme quadratique
non dégénérée); donc ce qui dans SUP s'appelait à tort ``zéro singulier
ordinair\hand{e}'' s'appelera ``zéro singulier non supersingulier'' , ou ``zéro

\margin{158}

\clearpage
% Nouveau EGA V (cote 25): archive cote 25, site page 159 (PDF page 160).
% typescript. Draft compared: none.
% AI-transcribed by Claude Opus 5.5 (claude-opus-5-5, adaptive thinking, effort high), Message Batches API, batch ega25b (msgbatch_01A7n5YuXubUPNK1zsepi8RG), 2026-11-13; not yet checked by eye.
\scanpage{25/159}
\mspage{\hand{2}}

singulier quadratique  non dégénéré" sin on tient à une terminologie
voisine de la terminologie \xout{xé} recue en car.0 ; de même, on pourra
parler de "singularité quadratique non dégénérée" d'un préschéma loc
noeth , resp. "géommétriquement  quadratique non dégénérée" quand on est
loc de type fini sur un corps k.  Par contre, la terminologie
"zéro singulier ordinaire" ou mieux "zéro quadratique ordinaire",
et de même "singularité quadratique ordinaire","singularité géométrique-%
ment quadratique ordinaire",  devra s'entendre (conformément à l'usage)
comme correspondant à la notion de forme quadratique lisse. D'ailleurs,
il peut être indiqué de remplacer le terme "ordinaire" par "élémentaire",
et d'étendre cette terminologie aux singularités pas nécessairement
quadratiques, mais de multiplicité quelconque. Tu me diras ton impression
à cet égard.   -    Il semble bien d'autre part que le t exte SUP soit
formellement correct, en particulier dans\hand{/}prop.7  la notion qui s'introduit
vraiment est celle de zéro supersingulier.  Cela n'empêche qu'il y aurait
lieu, au moins en remarque, d'introduire aussi \xout{$V(\emptyset)$xxxxxxx$^{\text{sinxxél}}$} le sous-%
préschéma  $V(\emptyset)^{\text{ult sing}}$  de $V(\emptyset)^{\text{sup sing}}$, correspondant à la considéra-%
tion des zéros \unsure{ultra-singuliers|ultra-singulaers}, i.e. \xout{xxx} qui sont singuliers sans être
quadratiques élémentaires; il se décrit essentiellement par le même
procédé que $V(\emptyset)^{\text{sup sing}}$  , en prenant le préschéma des zéros du
"discriminant corrigé"  introduit plus bas, au lieu du discriminant ordi-%
naire.  Donc les seules corrections à apporter à SUP semblent d'ordre
terminologique (et elles doivent porter également, évidemment, sur \xout{xxxxx}\inserted[typed]{la}
\inserted[typed]{terminologie}\hand{,}
introduite \hand{$\checkmark$}dans SUP p.10\hand{/}\inserted{ci-dessus}).  Par contre,  IV 23.9.2. est faux tel quel
sauf si k est de car $\neq 2$ \uline{ou} les composantes irréductibles de X  de
dimension paire; dans le cas général, il faut supprimer "$f_n$ satisfait
aux conditions équivalentes de 8.8., en particulier \deleted{\illegible{1}}" ; le reste de la

\mspage{\hand{159}}

\clearpage
% Nouveau EGA V (cote 25): archive cote 25, site page 160 (PDF page 161).
% typescript. Draft compared: none.
% AI-transcribed by Claude Opus 5.5 (claude-opus-5-5, adaptive thinking, effort high), Message Batches API, batch ega25b (msgbatch_01A7n5YuXubUPNK1zsepi8RG), 2026-11-13; not yet checked by eye.
\scanpage{25/160}
\mspage{\hand{3}}

proposition semble correct , et devrait se prouver de façon toute analogue
en montrant que pour $x\in X$\inserted{$^{(k)}$} donné, il y a une hypersurface \inserted{$H$} de degré $n\geqslant 2$
donné qui est tangente à X en x, et telle que x soit un point quadratique
élémentaire de $X$\hand{\unsure{$\cap$}}$H$ , de dimension une de moins .... Bien mettre les
pieds dans le plat, en faisant remarquer que si k \uline{est} de car.2 et X
connexe \inserted{non vide} de dimension impaire, alors les conditions de 8.8. \uline{ne sont pas}
vérifiées i.e. L/K est nécessairement inséparable (je crois, de degré
2 exactement si $X\subset P^r$ , mais sans garantie ...).

15.1. Soient S un préschéma, $\underline{E}$ un Module loc libre de type fini
sur S, $\phi$ une section de $\underline{\mathrm{Sym\xout{m}}}^n(\underline{E})$ , i.e. une ``n-forme sur $\check{\underline{E}}$ '' , où $n$\over{$>$}{\unsure{$\geqslant$}}$0$.
La donnée de $\phi$ équivaut à la donnée d'une section de $\underline{O}_P(n)$ sur $P=P(\underline{E})$
(réf àIII), et définit donc un sous-préschéma $V(\phi)$ de P. Nous dirons que
la forme $\phi$ est \uline{lisse} si $V(\phi)$ est lisse, et si de plus pour tout $s\in S$,
$V(\phi)_s\neq P_s$ (i.e. $\phi(s)\in \mathrm{Sym}^n_{k(s)}E\otimes k(s)$ est \over{$=$}{\unsure{$\neq$}}$0$ ). On voit aussitôt
que $\phi$ est lisse sss pour tout \xout{x} $\in$S, $\phi(s)$ est lisse, \xout{\illegible{3}}
\xout{\unsure{d'un corps}}(ce qui nous ramène au cas où S est le spectre d'un corps) et
que la notion de forme lisse est invariante par changement de corps de
base (ce quin nous ramène en particulier au cas d'un corps de base alg
clos); on peut résumer ces deux propriétés en disant que si $S'\to S$ est
un morphisme surjectif, alors $\phi$ est lisse sss son image inverse $\phi'$ sur S'
l'est. Bien entendu, si E est un module projectif de type fini sur un
anneau A, et $\phi\in\mathrm{Sym}^n_A(E)$, on dit que $\phi$ est lisse si .... Lorsque $E=A^{\unsure{N}+1}$,
la donnée de $\phi$ équivaut à celle d'un\xout{e forme} polynôme homogène de degré n
en des variables $X_0,..,X_r$ , et le critère jacobien implique immédiatement
que \inserted{\unsure{$n$ invers.}} $\phi$ est lisse sss l'idéal engendré par $\phi$ et les $\dfrac{\partial\phi}{\partial X_i}$ contient
une puissance de l'idéal d'augmentation $(X_0,\dots,X_r)$ , i.e. contient une

\margin{160} \margin{\unsure{Car mais on} trouve les sections inversibles de $\underline{O}_S$ !}

\clearpage
% Nouveau EGA V (cote 25): archive cote 25, site page 161 (PDF page 162).
% typescript. Draft compared: none.
% AI-transcribed by Claude Opus 5.5 (claude-opus-5-5, adaptive thinking, effort high), Message Batches API, batch ega25b (msgbatch_01A7n5YuXubUPNK1zsepi8RG), 2026-11-13; not yet checked by eye.
\scanpage{25/161}
\mspage{4}

puissance de chaque v\xout{X}ariable $X_i$ . Plus précis\unsure{\xout{x}}ément, le sous-\xout{pré}schéma
de $P^r$ défini par l'idéal homogène engendré par \xout{le} $\phi$ et les $\frac{\partial\phi}{\partial X_i}$ est
exactement formé des points de V($\phi$) en lesquels V($\phi$) n'est pas lisse sur
Spec(A) avec une dimension relative r-1.
15.2. Soi\unsure{e}nt X un préschéma, $\underline{J}$ un Idéal \inserted{qu coh} \xout{dexxxx} sur X, nous dirons que
l'homomorphisme ``d'augmentation'' $\underline{O}_X \to \underline{O}_X/\underline{J}$ est une ``augmentation
élémentaire de multiplicité n'' si elle satisfait aux conditions suivantes
a) $\underline{J}/\underline{J}^2 = \underline{N}_{Y/X}$ est loc libre de type fini sur Y = V($\underline{J}$).
b) Le noyau de l'homomorphisme canonique
\[ \hand{\mu} : \underline{Sym}_{\underline{O}_Y}(\underline{N}) \longrightarrow gr_{\underline{J}}(\underline{O}_X) \]
est engendré par le noyau \inserted{K} de $\hand{\mu}^n$, lequel est un sous-Module inversible
localement facteur direct de $\underline{Sym}^n_{\underline{O}_Y}(\underline{N})$ (i.e. est engendrable localement
sur Y par une section $\phi$ qui ne\xout{xxx} s'annule en aucun point).
c) Ladite $\phi$ (qui n'est définie, localement, que modulo multiplication
par une unité) est une forme \uline{lisse} . (NB On aurait pu introduire dans
\margin{\textbar}15.1. la notion de sous-Module inversible lisse de $\underline{Sym}^n(E)$, à vrai dire
géométriquement plus importante que celle de \uline{section} lisse, car V($\phi$) est
en fait défini par un tel sous-Module ...\xout{;} cela aurait l'avantage de
permettre de formuler b) et c) en une seule condition).

Lorsque n=2, on parlera de ``augmentation \inserted{de type} quadratique élémentaire'' ,
Si A est un anneau, \xout{lxxxlxxxxxthéxiexx} muni d'une \unsure{augmentation|augmantation} A$\to$A/J,
on conviendra encore de dire que cette augmentation est élémentaire de
multiplicité \unsure{n} , s'il en est ainsi pour .... \xout{E}Lorsque A est un anneau
local, on dira simplement par abus de langage que A est ``élémentaire de
multiplic\uline{i}té n'' si l'aug\xout{m}mentation A $\to$ A/$\underline{r}$(A) est élémentaire de
multiplicité \unsure{n}; noter que cela implique $n \geqslant 2$ , et si A est noethérien,
A est nécessairement non régulier. Si X est un préschéma, n un entier

\mspage{161}

\clearpage
% Nouveau EGA V (cote 25): archive cote 25, site page 162 (PDF page 163).
% typescript. Draft compared: none.
% AI-transcribed by Claude Opus 5.5 (claude-opus-5-5, adaptive thinking, effort high), Message Batches API, batch ega25b (msgbatch_01A7n5YuXubUPNK1zsepi8RG), 2026-11-13; not yet checked by eye.
\scanpage{25/162}
\margin{\unsure{5}}

$\geqslant$ 2, on dira que x \hand{$\in$} X est un point singulier élémentaire de multiplicité
n, si son anneau local $\underline{O}_{X,x}$ est élémentaire de multiplicité n. (Signaler
en remarque que cette \xout{\illegible{1}} terminologie est en accord avec la notion
générale de multiplicité due à Samuel). Pour n=2, on parlera en particu-%
lier de singularité quadratique élémentaire (ou ordinaire, dans la termino
logie classique). On introduit de même, conformément à l'usage général,
les variantes ``géométriques'' pour X localement de type fini sur un corps:
x \inserted{$\in X$} est dit singularité géométriquement élémentaire de multiplicité n, si
pour toute (ou,ce qui revient au même\unsure{,} une)extension K de k, et tout
(ou un) point$^{z}$ de $X_{K}$ au dessus de x, rationnel sur K, z est une
singularité élémentaire de multiplicité n.

15.3. ``Généraliser''SUP prop.8 au cas de la multiplicité n.

15.4. ``Généraliser''SUP prop.9 et cor.10. au cas de la multiplicité n.
(Il ne s'agit pas vraiment de généralisations, mais de variantes..)

15.5. Introduire la notion de zéro singulier élémentaire de multiplicité
n (pour n=2, zéro singulier quadratique élémentaire, ou ordinaire) d'une
section $\emptyset$ de $\underline{O}_{X}$ , ou plus généralement d'un Module localament libre, sur
un X localement noethérien, et la notion géométrique correspondante
(sur un corps de base k).

15.6. Soit X un sous-préschéma d'un préschéma loc noeth régulier X', et
soit x \hand{$\in$}X. Pour que \unsure{\over{X}{x}|x} soit une singularité élémentaire de multiplicité
n de X, il f et s qu'il existe un\xout{e} voisinage ouvert U de x, et une
suite $\underline{O}_{U}$-régulière ($\emptyset_{1},..,\emptyset_{d}$) qui soit telle que $X|U = V(\emptyset_{1},..,\emptyset_{d})$ et
que x soit un zéro singulier élémentaire de multiplicité n de ($\emptyset_{1}$,..).
En particulier, si A est anneau local noethérien qui est élémentaire
\margin{$|$} de multiplicité n, alors A est un anneau ``intersection complète''; d'ail\xout{\illegible{1}}-%
leurs, on voit grâce à 15\unsure{.}4. que l'on peut trouver, au voisinage de x,

\mspage{162}

\clearpage
% Nouveau EGA V (cote 25): archive cote 25, site page 163 (PDF page 164).
% typescript. Draft compared: none.
% AI-transcribed by Claude Opus 5.5 (claude-opus-5-5, adaptive thinking, effort high), Message Batches API, batch ega25b (msgbatch_01A7n5YuXubUPNK1zsepi8RG), 2026-11-13; not yet checked by eye.
\scanpage{25/163}
\mspage{\hand{6}}

un sous-préschéma X" de X, régulier, \xout{\illegible{1}} contenant X, et tel que X
puisse se décrire dans X"\xout{\illegible{1}} par une seule équation $\phi = 0$ , admettant x comme
zéro \inserted{sing} élémentaire de multiplicité n.

On conclut en particulier que si A est local noeth élémentaire de
\xout{15.7 On peut donner des variantes géométriques de 15.\illegible{1} lorsque \illegible{1}}
mult. n, alors A est Coh. Mac. \xout{\unsure{Il est clair}} Nous montrerons au par.20
\xout{\illegible{2} (point de \unsure{x}, pas nécessairement rationnel \unsure{sur} k)}
(par un calcul facile d'éclatement) que le point fermé de Spec(A) est
le seul point singulier de Spec(A); il en résulte que \xout{si dim} A est normal
sss dim A $\geqslant$ 2, réduit sss dimA $\geqslant$1. Si dim A = 0, alors  A est élémentaire
de multiplicité n sss  $m/m^2$ est de \over{dim}{\unsure{rang}}1 et  n est le plus petit entier
tel que $\underline{m}^n = 0$; on notera que ces anneaux sont aussi lesquotients d'anneaux
de valu\xout{\unsure{at}}ion discrète par la\inserted[typed]{\xout{\illegible{1}}} puissance n.ème de \over{l'}{leur} idéal maximal.

15.7.   Il y a lieu de dégager les variantes ``géométriques'' de 15.6. ,
on trouvera en particulier que si X est un préschéma loc de type fini sur
le corps k, $x \in X$, et si \over{n}{d} $= \dim_x X$ , alors  x est une singularité \xout{\unsure{ordinaire}}
\inserted[typed]{élémentaire}
\xout{\illegible{1}} de multiplicité  n de X, sss il existe un voisinage ouvert U de x
tel que U se plonge comme sous-préschéma d'un préschéma lisse X' sur k,
\xout{\illegible{1}} conne xe et de dimension \over{n}{d}+1 , défini par une équation $\phi$=0 , x étant
un zéro \inserted{géom.} singulier ordinaire de multiplicité n de $\phi$ . On aura dit d'ailleurs
\margin{$\surd$} dans 15.5. que cela signifie  que  la \inserted{valeur en x de} partie principale $d^n\phi$ de $\phi$ , \xout{restrei}
\xout{nte à X=V($\phi$)} (qui à priori est un\unsure{e} \xout{\illegible{1}} \inserted{élément} de $\underline{P}^n_{X'/k}\otimes \unsure{\underline{k}(x)}$ ,et plus
précisément de son Idéal d'augmentation $\underline{I}$ )  est en fait un\deleted{\unsure{e section}} \inserted{élément de}
\xout{de} $\underline{I}^n = \underline{Sym}^n_{\underline{O}_U}(\underline{\Omega}^1_{X'/k}\otimes k(\hand{x}))$), et  à ce titre est une \xout{\illegible{1}} \inserted{forme} lisse \xout{en x}
\margin{|} \hand{[}NB Oubli dans 15.1. : l'ensemble des points de lissité d'une forme est
\margin{|} ouvert\hand{]}. On remarquera de même qu'un tel point est isolé dans l'ensemble
des points non lisses, il est géom normal \xout{\illegible{1}} (resp. géom réduit) sss
\over{n}{d} $\geqslant 2$  (resp. \over{n}{d} $\geqslant 1$). Si \over{n}{d} = 0 et k(x) = k (\inserted{i.e. x  isolé et} x rationnel sur k) alors \deleted{x}
point singulier élémentaire  de mult\hand{,} \unsure{n} signifie que $\underline{O}X,x$  est  k-\deleted{\unsure{isomor}}

\mspage{\hand{163}}

\clearpage
% Nouveau EGA V (cote 25): archive cote 25, site page 164 (PDF page 165).
% typescript. Draft compared: none.
% AI-transcribed by Claude Opus 5.5 (claude-opus-5-5, adaptive thinking, effort high), Message Batches API, batch ega25b (msgbatch_01A7n5YuXubUPNK1zsepi8RG), 2026-11-13; not yet checked by eye.
\scanpage{25/164}
\mspage{7}

k-isomorphe à $k[T]/\unsure{(}T^n)$ . Si $d=1$ ,\inserted[typed]{x rat. sur k,} la notion de point singulier
\xout{ordinaire} élémentaire de multiplicité n correspond à la terminologie
classique \xout{\illegible{1}}``point singulier ordinaire à n tangente\unsure{\xout{s}|s} disctinctes'' .
\margin{(}(On aura \inserted{déjà} explicité dans 15.2. que si x est un point de X rat sur k, alors
la notion de singularité élémentaire de mult. n, au sens absolu ou au
sens relatif, \xout{\unsure{elle}} est la même; cela reste d'ailleurs valable si k(x)
est une extension finie séparable de k - sans doute d'ailleurs ``finie''
est inutile - et ce fait mérite d'être inséré en corollaire \xout{\unsure{si on arrive}}.
ou en proposition\hand{)}.

NB J'avais inclus au par.20 un nº sur l'éclatement d'un
préschéma X le long d'un sous-préschéma \inserted{fermé} Y tel que l'augmentation
$\underline{O}_X \longrightarrow \underline{O}_Y$ soit élémentaire de mult.n (auquel j'ai fait allusion dans
15.6.). Je signale qu'il s'agit d'un nº court et bien venu, dont les
seuls ingrédients sont les résultats généraux de lissité du par.17,et
la définition \inserted{du début de} 17.15.\unsure{2}. \inserted[typed]{(et encore pourrait-on éviter cette dernière).} Rien n'empêcherait d'incopporer ces résultats
dès ici, \xout{\unsure{par dans}} par exemple immédiatement après la définition
\xout{17}15.1. ci-dessus; alors \inserted[typed]{(la gamme ``géométrique''de 15.2. et la suite} 15.\unsure{3}. \unsure{à} 15.7. pourraient être séparés
de ces \xout{\unsure{calculs}} résultats, en les groupant dans un Nº 17.16, \xout{\illegible{2}}
\xout{\unsure{lorsque}} De toutes façons, je m'aperçois en cours de rédaction que le
\margin{(}nouveau nº 15 prévu sur les formes lisses éclate en \inserted{au moins} deux numéros assez
distincts et indépendants, l'un contenant des sorites généraux pour
les trucs ``élémentaires de mult n'' , n quelconque, l'autre \deleted{qui} cont\over{ient}{enant} la
caractérisation des formes \uline{quadratiques} lisses, et qui n'emprunte au
précédent que le nº 15.1. , i.e. pratiquement la définition de la lissité
d'une forme.

\mspage{164}

\clearpage
% Nouveau EGA V (cote 25): archive cote 25, site page 165 (PDF page 166).
% typescript. Draft compared: none.
% AI-transcribed by Claude Opus 5.5 (claude-opus-5-5, adaptive thinking, effort high), Message Batches API, batch ega25b (msgbatch_01A7n5YuXubUPNK1zsepi8RG), 2026-11-13; not yet checked by eye.
\scanpage{25/165}
% NOTE: the typed text of the whole page is a very faint carbon and almost unreadable on this scan; most words could not be read with any certainty
\illegible{4}

\illegible{2}

\uline{\unsure{Formalisme différentiel et morphismes lisses}}

\unsure{Pages 14 et 15} des notes \unsure{précédentes}, \illegible{4} \unsure{Page}
\unsure{15 ligne 20} le mot ``\illegible{1}'' est \unsure{absurde}, c'est la \unsure{deuxième assertion} qui est la
plus précise et celle à laquelle on songe.

\unsure{La généralisation de prop.} à laquelle il est fait allusion pages \unsure{14}
et 15 s'obtient à l'aide du diagramme suivant, qui généralise celui de la
page \unsure{4} :
\begin{diagramblock}
\begin{tikzcd}
 & \unsure{\Omega^1_{?}\otimes\mathcal{O}_X} \arrow[d] & & \\
\unsure{0} \arrow[r] & \unsure{\underline{N}_{?}} \arrow[r] & \unsure{\Omega^1_{X/S}\otimes\mathcal{O}_{?}} \arrow[r] & \unsure{\Omega^1_{X/Y}} \arrow[r] & \unsure{0} \\
 & \unsure{\Omega^1_{Y/S}\otimes\mathcal{O}_X} \arrow[u] \arrow[urr] & &
\end{tikzcd}
\end{diagramblock}
\illegible{3} diagramme \illegible{3} \illegible{4} diagramme \illegible{1}
\illegible{4} \unsure{le lemme de la page} \illegible{2}
\illegible{3} dans le diagramme \unsure{précédent}. \illegible{1}
\unsure{les conditions de ce lemme}, on est ramené à prouver que la section
\illegible{1} \unsure{(?,?)} \unsure{déduit de} \illegible{2} \unsure{i.e.} \illegible{1} \unsure{$\to \Omega^1_{X/S}\otimes\mathcal{O}_?$}
\illegible{2} \illegible{1} est celui qui \illegible{4}
\illegible{2} satisfait \illegible{3} \illegible{4} de la
page \unsure{15}. Or les \illegible{3} \illegible{1}
\[ \unsure{V^1 \;\;\; V'} \;\illegible{3} \]
\margin{\unsure{||}}\unsure{donnent une suite exacte} \illegible{2}
\[ \unsure{0 \to \mathrm{Hom}_{?}(?)} \to \illegible{1} \to \illegible{1} \to 0 \; ; \]
\unsure{D'autre part} on a des \unsure{isomorphismes canoniques}
\[ \illegible{2} \simeq \illegible{1} \quad , \quad \illegible{1} \simeq \illegible{1} \quad , \]
\illegible{3} \unsure{(cf.} \illegible{1}\unsure{)} \illegible{4}

\mspage{\hand{165}}

\clearpage
% Nouveau EGA V (cote 25): archive cote 25, site page 166 (PDF page 167).
% typescript. Draft compared: none.
% AI-transcribed by Claude Opus 5.5 (claude-opus-5-5, adaptive thinking, effort high), Message Batches API, batch ega25b (msgbatch_01A7n5YuXubUPNK1zsepi8RG), 2026-11-13; not yet checked by eye.
\scanpage{25/166}
\margin{//} \hand{nouveau EGA V} \mspage{1}

\margin{(}EGA IV par ?

\uline{Ensemble singulier et supersingulier}.

Pages 14 et 15 des notes précédentes, remplacer partout S par Y. Page
15 ligne 10 le \hand{``}i.e.\hand{''} est ab\unsure{us}if, c'est la deuxième assertion qui est la
plus précise et celle à laquelle on en a.

La généralisation de prop.7 à laquelle il est fait allusion pages 16
et 17 s'obtient à l'aide du diagramme suivant, qui généralise celui de la
page 4 :
\begin{diagramblock}\begin{tikzcd}
 & & 0 & & \\
 & & \underline{\Omega}^1_{X/Y}\otimes\underline{O}_{V^1} \arrow[u] & & \\
0 \arrow[r] & \underline{N}_{V^1/X} \arrow[r] \arrow[ur, "\hand{\mu}"] & \underline{\Omega}^1_{X/S}\otimes\underline{O}_{V^1} \arrow[r] \arrow[u] & \underline{\Omega}^1_{V^1/S} \arrow[r] & 0 \\
 & & \underline{\Omega}^1_{Y/S}\otimes\underline{O}_{V^1} \arrow[u] \arrow[ur, "\hand{\wr}"'] & & \\
 & & \hand{0} \arrow[u] & &
\end{tikzcd}\end{diagramblock}

où $\underline{N}$ désigne le faisceau conormal. Tous les faisceaux du diagramme sont
localement libres\inserted{sur $V^1$}, et on applique le lemme de la page 7 à $M=\underline{\Omega}^1_{X/S}\otimes\underline{O}_{V^1}$
et aux deux sous-Modules figurant dans le diagramme précédent. Avec
les notations de ce lemme, on est ramené à prouver que la section a
\xout{a} de \uline{Hom}(L,P) déduit de \inserted{$\mu$:} $P\rightarrow M/Q$ i.e. de \inserted{$\mu$:} $\underline{N}_{V^1/X}\longrightarrow\underline{\Omega}^1_{X/Y}\otimes\underline{O}_{V^1}$ ,
dont le \inserted{pré}schéma des zéros V(a) est celui qui rend \hand{$\mu$} strictement de rang
\hand{$<$} d = rang M/Q , satisfait V(a) = V(D), où D est défini au bas de la
page \over{\unsure{1}}{1}5. Or les inclusions de SS-préschémas lisses
\margin{||}
\[ V^1 \subset V^o \subset X \]
donnent une suite exacte (réf ?)
\[ 0 \rightarrow \underline{N}_{V^o/X}\otimes\underline{O}_{V^1} \rightarrow \underline{N}_{V^1/X} \longrightarrow \underline{N}_{V^1/V^o} \rightarrow 0 \ , \]
d'autre part on a des isomorphismes canoniques
\[ \underline{N}_{V^o/X} \simeq \underline{E}\otimes\underline{O}_{V^1} \ , \quad \underline{N}_{V^1/V^o} \simeq \uline{Hom}(W,K) \ , \]
où W,K sont définis (page 15) comme noyau et conoyau de

\margin{166}

\clearpage
% Nouveau EGA V (cote 25): archive cote 25, site page 167 (PDF page 168).
% typescript. Draft compared: none.
% AI-transcribed by Claude Opus 5.5 (claude-opus-5-5, adaptive thinking, effort high), Message Batches API, batch ega25b (msgbatch_01A7n5YuXubUPNK1zsepi8RG), 2026-11-13; not yet checked by eye.
\scanpage{25/167}
\mspage{\hand{2}}

\[ v^1 : E\otimes\underline{O}_{V^1} \longrightarrow \underline{\Omega}^1_{X/Y}\otimes\underline{O}_{V^1} \]
\xout{\unsure{en termes duquel}} (déduit  lui-même de l'homomorphisme de modules
localement libres sur $V^o$
\[ v : \underline{E\otimes O}_{V^o} \longrightarrow \underline{\Omega}^1_{X/Y}\otimes \underline{O}_{V^o} \quad , \]
en termes duquel  $V^1$ a été défini comme\inserted{le} sous-préschéma de $V^o$ sur lequel
le rang de v est strictement \hand{\unsure{$\leq$}} m-1 = rang $E\otimes\underline{O}_{V^o}$ - 1). Pour la définition
du premier isomorphisme, cf page 5, pour celle du second, df compléments
différentiels (ci joints) à `` variation du rang d'un homomorphisme de
Modules loc libres''. On obtient donc une suite exacte:
\[ 0 \to \underline{E}\otimes\underline{O}_{V^1} \xrightarrow{\ i\ } \underline{N}_{V^1/X} \xrightarrow{\ \ j\ \ } \underline{\mathrm{Hom}}_{\underline{O}_{V^1}}(W,K) \longrightarrow 0 \qquad . \]
Il faut alors vérifier \xout{\unsure{1\textsuperscript{o}) que}}
1\textsuperscript{o}) que
\[ \hand{\mu}i : \underline{E}\otimes\underline{O}_{V^1} \longrightarrow \underline{\Omega}^1_{X/Y}\otimes\underline{O}_{V^1} \]
n'est autre que $v^1$ , ce qui implique alors , grâce à l'hypothèse que
$v^1$ est strictement de rang m-1, que \xout{\unsure{V(u)}} $\hand{\mu}i$ définit par passage au
quotient
\[ \hand{/\mu'} : \quad \underline{\mathrm{Hom}}\,(W,K) \longrightarrow W = \mathrm{Coker}\ v^1 \ , \]
et que V(u) est aussi le sous-préschéma de $V^1$ qui rend $\hand{\mu'}$ strictement
de rang \hand{\unsure{$\leq$}}rang W \deleted{\illegible{1}} ; d'autre part, il' \xout{peut aussi s'interpréter} faut véri-%
fier

2\textsuperscript{o}) que $\hand{\mu'}$ n'est autre que l'homomorphisme déduit de la forme
quadratique \xout{\illegible{1}} M sur W à valeurs dans K, du bas de page 15. Il en résulte
aussitôt que V($\hand{\mu'}$) = V(\over{\illegible{1}}{D}) \xout{x}, où D est le discriminant de M, ce \deleted{\unsure{qui}} qui
achève la démonstration.

La vérification de 1\textsuperscript{o}) et 2\textsuperscript{o}) (où peut-être il y a à introduire un
signe ?) ne peut être que essentiellement méchanique, et est laissée à
Dieudonné.

\margin{167}

\clearpage
% Nouveau EGA V (cote 25): archive cote 25, site page 168 (PDF page 169).
% typescript. Draft compared: none.
% AI-transcribed by Claude Opus 5.5 (claude-opus-5-5, adaptive thinking, effort high), Message Batches API, batch ega25b (msgbatch_01A7n5YuXubUPNK1zsepi8RG), 2026-11-13; not yet checked by eye.
\scanpage{25/168}
\mspage{3}

EGA IV ..... Ensemble singulier et\xout{x} supersingulier.

Il me semble que le changement de point de vue possible signalé dans ma
dernière lettre peut-être récupéré après-coup, et qu'il peut y avoir intérêt
à partir d'abord de la situation où on a  un homomorphisme $\underline{E}\to\underline{O}_X$ , d'où
un $V^o = V(u) = V(u(\underline{E})\underline{O}_X)$, etc. Une fois cette étude faite, on fera les
remarques suivantes: supposons donné un X lisse(de dimension relative \hand{d} ..)
sur Y, un sous-préschéma $V^o = V(\underline{J})$ de X, tel que $\underline{J}$\xout{$/\underline{J}^2$xxxxNxxx} \xout{\unsure{K$\otimes$}xxxxxxxx} soit engendrab-%
le localement \inserted{\unsure{en les pts de $V^o$}} par m éléments (\unsure{$m \geq 0$} étant donné), où ce qui revient au
même, tel que $V^o$ soit localement de présentation finie sur X (ou encore sur
Y) et que \inserted{$\underline{N} =$} $\underline{N}_{V^o/X} = \underline{J}/\underline{J}^2$  soit localement engendrable par m éléments.  L'hy-%
pothèse peut aussi s'exprimer en disant que localement sur X, $\underline{J}$ peut être
représenté comme image d'un homomorphisme  u: $\underline{E}\to\underline{O}_X$, $\underline{E}$ étant localement
libre (ou même libre, tant qu'on \xout{m}reste dans le local !) de rang m.  En d'aut\unsure{x|r}
res termes, localement on est dans la situation de départ envisagée dans
les pages 14 et suivantes. Utilisant une telle représentation, on construit
alors des \inserted{sous-}préschémas \inserted{fermés} $V^1 \subset V^o$ , $V'' \subset V^1$, $V^2 \subset V^{\hand{1}}-V''$ . Il faut alors noter
(en proposition en forme) que ces sous-préschémas (et à fortiori les ensem-%
bles sous-jacents) sont en fait indépendants du choix de u:$\underline{E}\to\underline{O}_X$ définiss-%
ant $V^o$, \deleted{mais} \inserted{et} dépendent seu\over{\unsure{x}}{l}ement de \deleted{\unsure{X/Xxx}} $V^o$ ( et bien sur de X/Y et de
m !).  En effet, considérons d'abord  l'homomorphisme
\[ (\hand{$*^o$}) \qquad \underline{N}\otimes\underline{O}_{V^o}\to\underline{\Omega}^1_{X/Y}\otimes\underline{O}_{V^o} \ , \]
d'où  un homomorphisme
\[ (\hand{$*_m$}) \qquad \bigwedge^m\underline{N}\otimes\underline{O}_{V^o}\to\underline{\Omega}^{\over{\unsure{1}}{m}}_{X/Y}\otimes\underline{O}_{V^o}\ ; \]
on constate aussitôt que $V^1$ est le préschéma des zéros de cet homomorphisme
(qui ``existe'' grâce au fait que $\underline{\Omega}^{\hand{m}}_{X/Y}\otimes\underline{O}_{\unsure{V^o}}$  sur  $V^o$ est localement libre
de rang fini, et se définit comme réprésentant le foncteur qu'on devine
(Sch)$^o_{\hand{\unsure{/V^o}}}$ $\longrightarrow$ (Ens) ; c'est aussi le sous-préschéma de \over{V}{\unsure{$V^o$}}

\margin{168}

\clearpage
% Nouveau EGA V (cote 25): archive cote 25, site page 169 (PDF page 170).
% typescript. Draft compared: none.
% AI-transcribed by Claude Opus 5.5 (claude-opus-5-5, adaptive thinking, effort high), Message Batches API, batch ega25b (msgbatch_01A7n5YuXubUPNK1zsepi8RG), 2026-11-13; not yet checked by eye.
\scanpage{25/169}
\mspage{\hand{\unsure{4}}}

image inverse de la section nulle du fibré vectoriel $V((\underline{\Omega}^m_{X/Y}\otimes\underline{O}_{V^o})^{\hand{\vee}}\otimes\bigwedge^m\underline{N}\otimes\underline{O}_{V^o})$
par la section correspondant à l'homomorphisme $(*)_{m}$ ). Comme le premier memb-%
re de $(*_m)$ est de type fini, il s'ensuit d'ailleurs que $V^1$ est de présantati%
on finie sur $V^o$ . De même, $V''$ se définit comme le préschéma des zéros de
\[ (*_{m-1}) \qquad \bigwedge^{m-1}\underline{N}\otimes\underline{O}_{V^1} \longrightarrow \underline{\Omega}^{m-1}_{X/Y}\otimes\underline{O}_{V^1} \]
déduit de $(*)$ en passant à la puissance extérieure m-1 ème et induisant sur
$V^1$ . Pour l'invariance da sous-préschéma $V^2\subset V^1-V''$ , que je n'ai pas vrai-%
ment vérifiée et que je ne peux donc que conjecturer, je pense qu'on doit
pou\xout{x}voir préciser ce point de la façon suivante. L'homomorphisme
\[ (*\,*^1) \qquad \underline{E}\otimes\underline{O}_{V^1}\longrightarrow\underline{\Omega}^1_{X/Y}\otimes\underline{O}_{V^1} \]
déduit de $(*^o)$ en \xout{\unsure{induisant sur V,}} en composant avec l'homomorphisme
\uline{surjectif} $\underline{E}\otimes\underline{O}_{V^o}\to\underline{N}$ et en \xout{\unsure{composant}} induisant sur $V^1$ , a même conoyau
$\underline{W}$ que
\[ (*^1) \qquad \underline{N}\otimes\underline{O}_{V^1}\to\underline{\Omega}^1_{X/Y}\otimes\underline{O}_{V^1} \quad , \]
lequel \inserted{$\underline{W}$} est donc localement libre sur $V^1-V''$ ,\hand{(}de rang d\deleted{im}+1 en un point
x de $V^1-V''$ en lequel X est de dim rel \hand{d} sur Y\hand{)}. Il s'ensuit que\hand{``}la formation
du noyau de $(*^1)$ \over{c}{C}ommute à tout changement de base $T\to V^1$\inserted{$-V''$}\hand{''}\over{.}{,} \over{s}{S}oit $\underline{k}$ le
noyau de $(*^1)$. Si $\underline{K}$ est le noyau de $(*\,*^1)$ , qui est donc inversible sur
$V^1-V''$ , on voit \xout{donc} que $\underline{k}$ est un quotient de $\underline{K}$ . D'autre part, on a défini
page 15 un homomorphisme
\[ \hand{\underline{M}} : \ \underline{K}\longrightarrow\underline{\mathrm{Sym}}^2(\underline{W}) \]
au dessus de $V^1-V''$ (il peut, être commode de désigner par $\underline{K}$,$\underline{W}$ les Modules
sur $V^1-V''$ exclusivement ...), \deleted{et} \over{l}{L}e point qu'on aurait envie de vérifier
alors est qu'en fait cet homomorphisme se factorise par le quotient $\underline{k}$ de $\underline{K}$:
\[ \hand{m} : \ \underline{k}\longrightarrow\underline{\mathrm{Sym}}^2(\underline{W}) \ \deleted{$*$} ; \]
et que ce dernier homomorphisme (entre deux Modules dont la définition ne
\xout{si ce point est vérifié, alors} dépend plus de u) est en fait indépendant de
u. Si ce point est vérifié (et \uline{s'il est vrai}, un calcul élémentaire doit le

\mspage{\hand{169}}

\clearpage
% Nouveau EGA V (cote 25): archive cote 25, site page 170 (PDF page 171).
% typescript. Draft compared: none.
% AI-transcribed by Claude Opus 5.5 (claude-opus-5-5, adaptive thinking, effort high), Message Batches API, batch ega25b (msgbatch_01A7n5YuXubUPNK1zsepi8RG), 2026-11-13; not yet checked by eye.
\scanpage{25/170}
\mspage{\hand{5}}
donner) alors $V^2$ peut se décrire intrinsèquement (i.e. indépendammant de u)
\xout{xxxxx} en termes de l'homomorphisme précédent, comme étant``le sous-préschéma
qui rend l'homomorphisme précédent non dégénéré'' , i.e. qui représente le
foncteur suivant $\emptyset : (\mathrm{Sch})^{o}_{/V^1 - V''} \longrightarrow (\mathrm{Ens})$ \xout{xxx\unsure{F(T)}xxxxxxxxxxxx}, sous-fonc-%
teur du foncteur final, pour lequel $\emptyset(T)\neq \emptyset$ sss l'homomorphisme $\underline{m}_T$ déduit
de $\underline{m}$ par le changement de base $T \to V^1 - V''$ est tel que $\underline{k}_T$ soit inversible et
$\underline{m}_T$ corresponde\xout{xxx} à une fomme quadratique non dégénérée sur $\underline{W}_T$ à valeurs dans
$\underline{k}_T^{-1}$ .

De plus, il y a lieu \xout{xx} d'expliciter l'ensemble $V^o - V'' - V^2$ \xout{xxxxxx}
(\unsure{points} non supersinguliers de $V^o$ rel à Y, et pour la codimension m dans X)
comme étant formé des $x \in V^o$ tels que a) $V^o$ soit plat sur Y en x , et
b) $V^o_y$ admette x comme point géométriquement régulier (i.e. lisse) \uline{ou}
géométriquement \unsure{singulier|siggulier} quadratique ordinaire, \hand{c)}de codimension m dans $X_y$
au point x . (NB y dédigne le point de Y sous x). Dire également qu'en ces
poihts, $V^o \longrightarrow X$ est une immersion transversalement régulière relativement à
Y, de codimension m , (et on aura déjà dit, comme conséquence triviale de
l'égalité des rangs, que quand il y a un $\underline{E}$,u , alors $\underline{E}\otimes\underline{O}_{V^o} \to \underline{N}$ est un
isomorphisme en les poihts en question).

Bien entendu, une fois prouvé les invariances que je viens de signa-%
ler, la définition de $V^1, V'', V^2$ , \xout{xxxdxxxxxxxxaxixx} est claire, indépendamment
de l'existence globale d'un homomorphisme surjectif u:$\underline{E}$ \quad $\underline{J}$ ; d'autre part,
la proposition esquissée fin de la page 16 pour le cas où Y lui-même est
lisse sur un S, où V''= $\emptyset$ et où \inserted{$V^o$} $V^1$ sont lisses sur S avec la \hand{``}bonne\hand{''} dimension
relative sur S, \xout{xixxxxixxxdødixxxxxixx} savoir la caractérisation de $V^2$ comme le
préschéma de ramification de $V^1 \longrightarrow Y$ , prend également un sens intrinsèque,
indépendammant du choix d'un $\underline{E}$.

Pour terminer, on peut expliciter le cas où $V^o$ est défini comme image
inverse, par un \inserted{Y-}morphisme u: $X \to Z$, d'un sous-préschéma $Z^o$ de Z, en suppo-%
\mspage{\hand{170}}

\clearpage
% Nouveau EGA V (cote 25): archive cote 25, site page 171 (PDF page 172).
% typescript. Draft compared: none.
% AI-transcribed by Claude Opus 5.5 (claude-opus-5-5, adaptive thinking, effort high), Message Batches API, batch ega25b (msgbatch_01A7n5YuXubUPNK1zsepi8RG), 2026-11-13; not yet checked by eye.
\scanpage{25/171}
\mspage{6}

sa\over{\unsure{i}}{n}t qu'en tout point de $u(V^{o})$, $Z$ et $Z^{o}$ sont lisses sur $M$, $Z^{o}$ étant de codi
mension $m$ dans $Z$ en \over{s}{c}es points. Ceci implique automatiquement que $V^{o}$ est
loc de présentation finie sur $Y$, et que $\underline{N}_{V^{o}/X}$ est l'image de
\xout{\unsure{xxx}}
\[ \underline{E}^{o} = u\,(\underline{N}_{Z^{o}/Z}) , \]
qui est localement libre de rang n. Donc dans ce cas, les hypothèses préli-%
minaires faites plus haut sont vérifiées (j'ai oublié de dire qu'on supposait
X encore lis se sur Y), \over{D}{\unsure{d}}e plus on ne dispose pas\inserted{\unsure{,}} à priori d'un homomorphisme
surjectif $u:\underline{E}\rightarrow\underline{J}$ global, avec $\underline{E}$ loc libre de rant m sur X, mais en revanche
on dispose de l'homomorphisme global surjectif
\[ (\unsure{$\nmid$|$\dagger$}) \qquad \underline{E}^{o}\longrightarrow \underline{N}_{\xout{\illegible{1}}} \quad , \]
qui localement peut être considéré comme déduit d'un homomorphisme sur\over{g}{j}ectif
$\underline{E}\rightarrow\underline{J}$ par restriction à $V^{o}$ . Alors $V^{1}$, $V''$, \inserted{\unsure{$V^{2}$}}\unsure{p}euvent se définir directement
en termes de (\unsure{$\nmid$|$\dagger$}) , ainsi que $\underline{K}$ , $\underline{W}$ , de sorte qu'on trouve directement un
\[ \underline{M} : \underline{K}\rightarrow \underline{\mathrm{Sym}}^{2}(\underline{W}) \qquad\qquad \text{sur } V^{1}-V'' \]
intrinsèque à la situation, d'où \xout{\unsure{disc}} un discriminant
\[ D \in \Gamma(V^{1}-V'',\ \underline{K}^{-n+m-1}\hand{$\otimes$}\ \text{dét}(\underline{W})^{\hand{$\otimes$}2}\ ), \]
tel que $V(D) = V^{2}$ .

Je propose de nous contenter d'ajouter toutes ces considérations d'inva-%
riance en petits caractères, après le cas des sections de Modules loc libres;
je te laisse le choix de donner ou non une terminologie de transversalité dans
le dernier cas envisagé. Il reste la question dans quel paragraphe doit être
la place du fourbi supersingulier. Du point de vue résultats, le plus naturel
semblerait le par.19, mais du point de vue technique, dans les pages 18,19
j'ai utilisé un sorite différentiel qui, \xout{w}semble-t-il, viendrait plutôt au par
22 des grassma\unsure{nniennes}. Si tu ne remplaces cette démonstration savante par un
calcul bébête de matrices (certainement possible, mais moins instructif) il
semblerait qu'on devra rejeter la supersingularité au par.22.

\margin{171}

\clearpage
% Nouveau EGA V (cote 25): archive cote 25, site page 172 (PDF page 173).
% typescript. Draft compared: none.
% AI-transcribed by Claude Opus 5.5 (claude-opus-5-5, adaptive thinking, effort high), Message Batches API, batch ega25b (msgbatch_01A7n5YuXubUPNK1zsepi8RG), 2026-11-13; not yet checked by eye.
\scanpage{25/172}
\mspage{7}

EGA IV 19

Remords à ``Ensemble singulier et supersingulier ...'' .

Venant de rédiger le\xout{s} sorite de la transversalité, je m'aperçois que l\unsure{e}
contexte adopté dans mes notes sur les \unsure{ensembhes|ensembles} singuliers et \unsure{supersinguliers}
savoir les sections des Modules loc libres i.e. des bons fibrés vectoriels,
n'est pas encore le ``bon''. Si $u$ est une section d'un fibré vectoriel
$E = W(\underline{E})$ (associé à $\underline{E}$ loc libre de type fini) sur X lisse sur Y, alors $V(u)$
n'est autre que l'image inverse \unsure{du} \deleted{la} \inserted{sous-préschéma $E_{o}$ de $E$} ``section nulle'', et un point $x\in V(u)$
est non singulier \inserted{pour u} rel. à S sss u est \uline{transversal} à $E_{o}$ en x \unsure{rel.} à Y.
Si cette dondition n'est pas vérifiée, i.e. si x est singulier pour u rel à
Y, nous étudions les conditions moyennant lesquelles x n'est pas trop singu-%
lier rel.à Y, i.e. u ne dévie pas trop de la transversalité en x. Je n'ai
aucun doute que tous les développements de mes notes, avec des changements
essentiellement triviaux, s'étendent au cas où on a un X-préschéma E, un
sous-préschéma $E_{o}$ de E, une section \xout{\unsure{dx}} u de E sur X, et qu'on considère l'ima-%
ge inverse $X_{o}$ de $E_{o}$ par u, \inserted[typed]{\unsure{On} suppos\xout{ant} E et $E_{o}$ lisses sur X en les points de $v(X)$.} il est possible (mais \unsure{seul} une rédaction détail-%
lée permettra de le déceler) qu'il faille supposer parfois que la dimension
relative de $E_{o}$ sur X en les points de $v(X_{o})$ soit nulle, i.e. que $E_{o}$ est
même étale sur X en ces points (ce qui sera le cas lorsque $E_{o}$ provient comme
ci-dessus d'une section de E sur X, soit $v_{o}$). Ainsi, x sera un point
``\xout{\illegible{1}} \inserted{de non transversalité à $E_{o}$} ordinaire pour v, relativement à Y'' si et seulement si \deleted{\illegible{1}} \inserted{x} est
géométriquement un point quadratique ordinaire de $X_{o\,y}$ , et si $\mathrm{codim}_{x}(X_{oy},X_{y})$
$= c$ , où c est la dodimension de $E_{o}$ dans E (en $v(x)$) ; cela implique que
$X_{o}\longrightarrow X$ est \inserted{\unsure{un immers.|une immers.}} transversalement régulier\xout{e} en x rel. à Y .... Le rôle du Module
$\underline{E}\otimes\underline{O}_{V(u)}$ sera ici joué par $v_{o}^{*}(\underline{N}_{E_{o}/E})$ , où $v_{o}:X_{o}\longrightarrow E_{o}$ est induit par v. Je
pense que tu n'auras pas de mal à tout écrire dans ce contexte,bien plus
naturel somme toute - quitte à faire à la fin les traductions dans le cas
des fibrés vectoriels. S'il y a des difficultés de démarrage déjà, nous
pourrions les travailler ensemble à ton prochain passage à Paris ... \hand{Je vais t'envoyer des notes succintes ces jours-ci.}
\margin{172}

\clearpage
% Nouveau EGA V (cote 25): archive cote 25, site page 173 (PDF page 174).
% typescript. Draft compared: none.
% AI-transcribed by Claude Opus 5.5 (claude-opus-5-5, adaptive thinking, effort high), Message Batches API, batch ega25b (msgbatch_01A7n5YuXubUPNK1zsepi8RG), 2026-11-13; not yet checked by eye.
\scanpage{25/173}
% NOTE: the sheet itself carries no typed or handwritten text except the page number at the foot. Very faint typed text from the next sheet shows through the paper. It reads the right way round, so it is not a mirrored image of the reverse side. It seems to begin with a heading "16.2. Formes quadratiques libres." and to contain a "Proposition 16.2.?". It is not transcribed as belonging to this page.
\mspage{\hand{173}}

\clearpage
% Nouveau EGA V (cote 25): archive cote 25, site page 174 (PDF page 175).
% typescript. Draft compared: none.
% AI-transcribed by Claude Opus 5.5 (claude-opus-5-5, adaptive thinking, effort high), Message Batches API, batch ega25b (msgbatch_01A7n5YuXubUPNK1zsepi8RG), 2026-11-13; not yet checked by eye.
\scanpage{25/174}
\margin{\unsure{IV|1/U}} \margin{nouveau EGA V} \mspage{1}

17.16. \uline{Formes quadratiques lisses}.

16.1.Soit $Q$ \unsure{une} section de $\underline{Sym}^2(\underline{E})$, ($\underline{E}$ loc libre de rang fini),
de sorte que Q \xout{\unsure{définit}} (qu'on peut aussi interpréter comme une forme
quadratique sur $\overset{\hand{$\vee$}}{\underline{E}}$ ) \xout{\illegible{3} \unsure{section}} définit une forme biliné-%
aire symétrique B sur $\overset{\hand{$\vee$}}{\underline{E}}$, d'où un homomorphisme $\overset{\hand{$\vee$}}{\underline{E}} \longrightarrow \underline{E}$ , et passant aux
Modules déterminants on trouve $\text{dét}\ \overset{\hand{$\vee$}}{\underline{E}} \longrightarrow \det \underline{E}$ , d'où enfin une section
\[ d(Q) \in \Gamma(\ \text{dét}\ (E)^{\otimes 2}) \]
appelée (sauf erreur) \uline{discriminant de la forme quadratique} Q. Lorsque
$\underline{E} = \underline{O}_S^{\ \hand{$n$}}$ , c'est simplement la section de $\underline{O}_S$ déterminant de la matrice
à n lignes et n colonnes, à coeff dans $\Gamma\underline{O}_S$ , qui exprime Q . Notons tout
de suite:

\uline{Proposition} 16.2. Supposons que pour tout $s \in S$, on ait car $k(s) \neq 2$ \uline{ou}
rang. $E\otimes k(s)$ pair. Alors, pour que Q soit lisse, il f et s que Q soit
``non dégénérée'' i.e. que d(Q) soit inversible. \hand{(}D'ailleurs la condition est
suffisante sans restriction sur S ni $\underline{E}$\hand{)}.

La dernière assertion n'est mise que pour la commodité des référen-%
ces éventuelles, car elle est trivialement contenue dans la première,
l'hypothèse que Q est non dégénérée implique en effet que $\underline{E}$ est de rang
pair en tout point dont la car est 2. D'ailleurs, elle se démontre direc-%
tement sans considération de caractéristique, car (se plaçant sur un corps
\inserted{\unsure{alg clos}} et prenant une base) elle exprime que les $\frac{\hand{$\partial$} Q}{\hand{$\partial$} X_i}$ n'ont pas de zéro\xout{s} commun\xout{s}
non trivial, à fortiori \unsure{ils} n'ont pas de zéro commun non trivial avec Q.
Si car $k \neq 2$, la réciproque est vraie, car en vertu de la formule
\[ 2Q = \hand{$\sum$} X_i \frac{\hand{$\partial$} Q}{\hand{$\partial$} X_i} \]
on voit alors que tout zéro commun aux $\frac{\hand{$\partial$} \unsure{Q|E}}{\hand{$\partial$} X_i}$ est aussi un zéro de Q. Enfin,
si car $k = 2$ , \xout{si E est de rang pair} alors la forme bilinéaire B

\margin{174}

\clearpage
% Nouveau EGA V (cote 25): archive cote 25, site page 175 (PDF page 176).
% typescript. Draft compared: none.
% AI-transcribed by Claude Opus 5.5 (claude-opus-5-5, adaptive thinking, effort high), Message Batches API, batch ega25b (msgbatch_01A7n5YuXubUPNK1zsepi8RG), 2026-11-13; not yet checked by eye.
\scanpage{25/175}
\mspage{\hand{2.}}

associée à Q (définie par la matrice des $\frac{\hand{\partial} Q}{\hand{\partial} X_i}$ ) est alternée, donc son
``noyau''  $N\subset\overset{\hand{\vee}}{E}$  est tel que  $\overset{\hand{\vee}}{E}/N$  soit de rang \uline{pair} , donc si E est lui-%
même de rang pair, il en est de même de N. Donc si $N\neq 0$, le rang de N
est au moins 2, d'où s'ensuit qu'il existe au moins un zéro non banal de
Q dans N, i.e. Q est non lisse.  NB Ce recours aux coordonnées est décidé
ment indigne de nous, il faudrait les utiliser uniquement pour prouver
simplement le
\uline{Lemme} \inserted[typed]{16.3.} E vectoriel de rang fini sur un corps k, \inserted[typed]{$Q\in \mathrm{Sym}^2(E)$} x un élément non nul de
$\overset{\hand{\vee}}{E}$ , \xout{\unsure{tel que Q(x}} x' son image dans P(E), supposons  $Q(x,x) = 0$ i.e.
$x'\in V(Q)$. Alors pour que \xout{\unsure{x' soit un point lisse de}} V(Q) soit lisse en x',
et de dim \xout{\illegible{1}} rang E - 1   (i.e. $Q\neq 0$), il f et s  que x n'appar-%
tienne pas au noyau de l'homomorphisme $\overset{\hand{\vee}}{E}\longrightarrow E$ défini par la forme bilinéai-%
re symétrique B associée à Q/
16.4.  L'étude qui suit est destinée essentiellement à \xout{\unsure{traiter le cas}}
donner un critère de lissité pour une forme quadratique dans le cas non
couvert par 16.2. , i.e. essentiellement dans le cas d'un vectoriel de
dimension impaire sur un corps de car 2.\xout{\unsure{Alors}} \inserted[typed]{Dans ce} cas \uline{toute} forme quadra-%
tique  Q  sur $\overset{\hand{\vee}}{E}$ est dégénérée, mais on voit aisément (prenant notamment
% NOTE: a line of typed text overtyped with x stands just below "est dégénérée," between this line and the next
\xout{\illegible{2}}
la``forme standart'' ) qu'elle peut fort bien être lisse.

Introduisons, pour tout entier  $n>0$, la ``forme quadratique standart''
$Q_n$ sur $\underline{Z}^n$ , comme une forme à coéfficients entiers en n variables
$X_1,..,X_n$ ,  en distinguant les deux cas:
a)  $n = 2m$  ,  $Q_{2m}(X_1,..,X_{2m}) = X_1X_2 + X_3X_4 + \ldots$
b)  $n = 2m+1$,  $Q_{2m+1}(X_1,..,X_{2m+1}) =$\xout{\unsure{$(X_1X_2+X$}} $Q_{2m}(X_1,..,X_{2m}) + X_{2m+1}^2$.
\uline{Lemme} 16.5  Soit n un entier impair, et considérons
\[ \xout{\unsure{M}}Q(X_{\xout{\unsure{2}}1},..,X_n) = \sum_i a_iX_i^{\,2} + \hand{\sum_{i<j}} b_{ij}X_iX_j \ , \]
une forme quadratique en n \unsure{indétermihées|indéterminées} $X_i$ , à coéff indéterminés
$a_i$, $b_{ij}$, de sorte que le discriminant  d(Q)  est un polynôme à

\mspage{\hand{175}}

\clearpage
% Nouveau EGA V (cote 25): archive cote 25, site page 176 (PDF page 177).
% typescript. Draft compared: none.
% AI-transcribed by Claude Opus 5.5 (claude-opus-5-5, adaptive thinking, effort high), Message Batches API, batch ega25b (msgbatch_01A7n5YuXubUPNK1zsepi8RG), 2026-11-13; not yet checked by eye.
\scanpage{25/176}
\mspage{\hand{8}}

coéfficients entiers en les $a_i,b_{ij}$ . Alors le ``contenu'' de ce polynôme
est égal à 2 (i.e. le pgcd de ses coéfficients est 2).
\inserted[typed]{multiple de 2.}
\xout{Cela} Prouvons d'abord que ce contenu est \xout{de la forme 2\deleted{$^r$}\illegible{1}}
En effet, cela signifie aussi que \deleted{\xout{\illegible{1}}} lorsqu'on spécialise ce polynôme
en \xout{la}un corps $Z/\over{p}{2}Z$\hand{)} \deleted{avec $p\neq 2$}, il \xout{n'}est pas identiquement nul, ce qui \xout{xxx}
résulte du fait qu'une forme quadratique de degré imp\over{q}{a}ir sur un corps de
car 2 est tjrs dégénérée, i.e. à un discriminant nul. Pour prouver
que le contenu \inserted{c} est exactement 2, il suffit de fai\unsure{r}e le calcul du discri-%
minant \inserted{de la forme} standart de degré n, d'une part elle doit être multiple de c,
d'autre part le calcul donne 2, (car c'est 1 pour le degré pair, et 2 pour
le degré 1 ...).

\xout{\unsure{On en conclut immédiatement}} Il y a lieu alors d'introduire le polynôme
\[ \overset{\hand{\sim}}{d}\xout{*}(Q) = \tfrac{1}{2}d(Q) \]
en les coéfficients de Q, qui est encore à coéfficients entiers, donc
prend une valeur bien déterminée quand on spécialise les $a_i,b_{ij}$ dans un
anneau A quelconque, i.e. quand on prend une forme quadratique \deleted{\xout{Q}}\inserted{q}de degr\unsure{é}
n à coéfficients dans A. On appela $\overset{\hand{\sim}}{d}\xout{*}(Q)$ le``\uline{polynôme} \unsure{\hand{``}}\uline{discrimihant}
\uline{corrigé de la forme quadratique indétérminée Q à n variables}'' , et sa
valeur relativement aux coéfficients d\xout{e}'une forme quadratique q à coeff
dans un anneau A quelconque sera appelé le ``\uline{discriminant corrigé de la}
\uline{forme quadratique} q''. Plus généralement, on déduit de façon essentielle
ment triviale de 16.5. l'énoncé suivant:

\uline{Proposition} 16.6. On peut, \inserted{et} d'une seule manière, associéer à tout\xout{e}
préschéma S, muni d'un Module loc libre \xout{d}e rang fini \uline{E} et d'une forme \inserted{qu.}
$Q \in \unsure{\Gamma}\,\underline{\mathrm{Sym}}^2(E)$, associer une section $\overset{\hand{\sim}}{d}(Q)$ de dét$(E)^{\otimes 2}$ , de façon à satis-%
faire aux conditions suivantes:

(\xout{1}a) Compatibilité avec changement de base\hand{, et fonctorialité pour isomorphismes en \uline{E}.}

(b) Si E est partout de rang pair, on a $\overset{\hand{\sim}}{d}(Q) =$ d(Q) (ou d(Q)
désigne le discriminant de Q).

\mspage{\hand{176}}

\clearpage
% Nouveau EGA V (cote 25): archive cote 25, site page 177 (PDF page 178).
% typescript. Draft compared: none.
% AI-transcribed by Claude Opus 5.5 (claude-opus-5-5, adaptive thinking, effort high), Message Batches API, batch ega25b (msgbatch_01A7n5YuXubUPNK1zsepi8RG), 2026-11-13; not yet checked by eye.
\scanpage{25/177}
\mspage{\hand{4}}

(c) Si E est partout de rang impair, on a $2\hand{$\widetilde{\phantom{d}}$}\!d(Q)^{\unsure{\hand{2}}} = d(Q)$.

\xout{\unsure{De plus, le discriminant corrigé d a ainsi introduit satisfait à}}

\xout{\unsure{les conditions suivantes}} \quad NB Il n'est pas raisonnable d'énoncer la

propr\over{\illegible{1}}{\unsure{ié}}té de compatibilité avec le changement de base, s ans énoncer en

même temps\hand{,} ou même avant\hand{,} la propriété de fonctorialité par isomorphismes,

étant donné que le changement de base \inserted{sur $\underline{E}$} lui-même n'est \over{p}{g}ère défini que

modulo isomorpphisme; il n'en s erait plus de même si on se restreignait

au cas $E = \underline{O}_S{}^n$ ,(ce qui ne serait pas commode pour les références).

\xout{Corollaire 16.7.} \quad \uline{Définition} 16.7. La section $\hand{$\widetilde{\phantom{d}}$}\!d(Q)$ de $\text{dét}(E)^{\otimes 2}$

s'appelera le \uline{discriminant corrigé} de la forme quadratique Q.

\uline{Corollaire} 16.8. \unsure{Soinnt|Soient} Q,Q' des formes relatives à $\underline{E}$, $\underline{E}'$ , supposons

la parité des rangs de $\underline{E}$ (resp. $\underline{E}'$) constante sur S. Alors

a) Si $\underline{E}$ et $\underline{E}'$ ne sont pas tous deux de rang impair, on a
\[ \hand{$\widetilde{\phantom{d}}$}\!d(Q \oplus Q') = \hand{$\widetilde{\phantom{d}}$}\!d(Q)\hand{$\widetilde{\phantom{d}}$}\!d(Q') \]

b) Si $\underline{E}$ et $\underline{E}'$ sont tous d eux de rang impair, on a
\[ \hand{$\widetilde{\phantom{d}}$}\!d(Q \oplus Q') = \hand{2}\,\hand{$\widetilde{\phantom{d}}$}\!d(Q)\hand{$\widetilde{\phantom{d}}$}\!d(Q') \]

La vérification est triviale.

Nous pouvons maintenant énoncer le résultat principal du présent

numéro:

\uline{Théorème} 16.9. Soient $\underline{E}$ un Module loc libre de rang fini sur le préschém-%

ma S, $Q \in \Gamma\,\uline{Sym}^2(\underline{E})$. Les conditions suivantes sont équivalentes:

(i) Q est lisse.

(ii) Le d\unsure{\over{s}{i}}scriminant modifié $\hand{$\widetilde{\phantom{d}}$}\!d(Q) \in \Gamma\,\text{dét}(\underline{E})^{\otimes 2}$ est inversible.

(iii) \inserted[typed]{(Lorsque $\underline{E}$ est d e rang constant\xout{s} n)} Il existe un morphisme surjectif S' \quad S tel que la forme $Q_{S'}$

déduite de Q par changement de base soit isomorphe à la forme quadrati-%

que standart en n variables.

(i\xout{x}i bis). Comme (iii), mais en supposant S'\hand{$\rightarrow$}S fidèlement plat de

présentation finie.

\mspage{\hand{177}}

\clearpage
% Nouveau EGA V (cote 25): archive cote 25, site page 178 (PDF page 179).
% typescript. Draft compared: none.
% AI-transcribed by Claude Opus 5.5 (claude-opus-5-5, adaptive thinking, effort high), Message Batches API, batch ega25b (msgbatch_01A7n5YuXubUPNK1zsepi8RG), 2026-11-13; not yet checked by eye.
\scanpage{25/178}
\uline{Corollaire} 16.10. Soient k un corps, E un vectoriel de dimensionfinie
sur k, $Q \in \mathrm{Sym}^2(E)$. Conditions équivalentes:

(i) Q est lisse.

(ii) $\overset{\hand{$\sim$}}{d}(Q) \neq 0$.

(iii) Il existe une extension \xout{finie} k' de k telle que $Q\otimes k'$ soit
isomorphe à la forme standart.

(iii bis) Comme (iii), avec k' une extension finie de k.

C'est une conséquence triviale de 16.9. On notera que lorsque k est
algébriquement clos, alors (iii) et (iii bis) se remplacent par la
condition plus frappante: Q est isomorphe à la forme standart.

Démonstration de 16.9. On \xout{a évidemment} peut évidemment
supposer E de rang constant n. On aura évidemment
\[ \text{(iii bis)} \hand{$\Longrightarrow$} \text{(iii)} \hand{$\Longrightarrow$} \text{(ii)} \]
compte tenu que pour la forme standart, le discriminant modifié est 1.
Nous prouverons maintenant (ii)\hand{$\Longrightarrow$}(iii bis), puis (i)\hand{$\Longleftrightarrow$}(ii).

Soit $\underline{E}_o = \underline{O}_S^{\ n}$ , $Q_o$ la forme standart sur $\underline{E}_o$ , Considérons le
foncteur
\[ \underline{\mathrm{Isom}}( (\underline{E}_o,Q_o),(\underline{E},Q)) :\ (\mathrm{Sch})_{/S}^{\ o} \longrightarrow (\mathrm{Ens}) , \]
dont la valeur en tout S' sur S est formé de l'ensemble des isomorphismes
$\underline{E}_{o\,S'} \to \underline{E}_{S'}$ compatibles avec les formes $Q_{o\,S'}$ et $Q_{S'}$ . Il est immédiat
(sans condition sur \inserted{$\underline{E}$ \unsure{ni}} Q) que ce foncteur est représentable par un préschéma
P affine et de présentation finie sur S, qui est un \inserted[typed]{sous-}\deleted{pré}schéma de W =
$W(\underline{\mathrm{Hom}}_{\underline{O}_S}(\underline{E}_o,\underline{E}))$ , et un sous-préschéma fermé de \deleted{$\underline{\mathrm{Isom}}(\underline{E}$} l'ouvert
$\underline{\mathrm{Isom}}\ (\underline{E}_o,\underline{E}_1)$ de W \unsure{/} L'implication (ii)\hand{$\Longrightarrow$}(iii bis) sera prouvée si
nous montrons que lorsque $\overset{\hand{$\sim$}}{d}(Q)$ est inversible, P est fidélement plat
sur S: on prendra en effet S'=P.

Posons pour simplifier \xout{\unsure{$\underline{G}(\underline{E}_o$}x}

\mspage{\hand{178}}

\clearpage
% Nouveau EGA V (cote 25): archive cote 25, site page 179 (PDF page 180).
% typescript. Draft compared: none.
% AI-transcribed by Claude Opus 5.5 (claude-opus-5-5, adaptive thinking, effort high), Message Batches API, batch ega25b (msgbatch_01A7n5YuXubUPNK1zsepi8RG), 2026-11-13; not yet checked by eye.
\scanpage{25/179}
\mspage{\hand{6}}

\[ \underline{Q}(\underline{E}) =\xout{W}(\underline{Sym}^2(\underline{E})) \]
et définissons de même $\underline{Q}(\underline{E}_o)$; par l'opération de ``transport de
structure'' on aura donc un morphisme $u \rightsquigarrow \underline{Q}(u)$ :
\[ \underline{Isom}(\underline{E}_o,\underline{E}) \longrightarrow \underline{Isom}(\underline{Q}(\underline{E}_o),\underline{Q}(\underline{E})), \]
et, utilisant la section $q_o$ de $\underline{Q}(\underline{E}_o)$ correspondant à $Q_o$ , on trouve un
morphisme $u \rightsquigarrow \underline{Q}(u)(q_o)$:
\[ (*) \qquad \underline{Isom}(\underline{E}_o,\underline{E}) \longrightarrow \underline{Q}(\underline{E}) \quad . \]
D'autre part, la forme Q \xout{\unsure{sur E}} correspond à une section q de $\underline{Q}(\underline{E})$ sur S,
et P n'est autre que l'image inverse de q(S) par (*), comme il résulte
trivialement des définitions, (ce qui, d'ailleurs, établit en passant la
représentabilité annoncée de P \xout{z}comme préschéma affine et de \xout{type}\inserted{prés} finie
sur S). Introduisons maintenant l'ouvert \xout{$\underline{Q}$x}
\[ \underline{Q}(\underline{E})^* = \underline{Q}(\underline{E})_{\tilde{d}} \]
de $\underline{Q}(\underline{E})$ , correspondant aux formes quadratiques de discriminant corrigé
inversible (représentan\xout{x}t le foncteur $S' \rightsquigarrow$ ensemble des \xout{formes}
sections de $\underline{Sym}^2(\underline{E}_{S'})$ à discriminant corrigé inversible). En fait,
par transport de structure on a $\underline{Isom}(\underline{E}_o,\underline{E}) \rightarrow \underline{Isom}(\underline{Q}(\underline{E}_o)^*,\underline{Q}(\underline{E})^*)$, et
comme $q_o$ est une section de $\underline{Q}(\underline{E}_o)^*$ (puisque le discriminant \unsure{corrigé|modrigé} de
la forme type $Q_o$ vaut 1), le morphisme (*) se factorise en fait en un
morphisme
\[ (*\,*) \qquad \underline{Isom}(\underline{E}_o,\underline{E}) \rightarrow \underline{Q}(\underline{E})^* , \]
qui est évidemment de présentation finie, puisque les deux membres sont de
présentation finie sur S. Il suffit maintenant de prouver la
\xout{Lemme} \uline{Proposition} 16.11. Avec les notations précédentes (à rappeler),
le morphisme (* *) est \uline{fidèlement plat} .

Il s'ensui\deleted{t}\inserted{vra} que si Q satisfait à (ii), i.e. \inserted{si} q est en fait une sectio[n]
de $\underline{Q}(\underline{E})^*$\hand{\unsure{$\hat{}$}} , alors P (déduit de (**) par \inserted{\unsure{le}} changement de base q) est encore

\mspage{\hand{179}}

\clearpage
% Nouveau EGA V (cote 25): archive cote 25, site page 180 (PDF page 181).
% typescript. Draft compared: none.
% AI-transcribed by Claude Opus 5.5 (claude-opus-5-5, adaptive thinking, effort high), Message Batches API, batch ega25b (msgbatch_01A7n5YuXubUPNK1zsepi8RG), 2026-11-13; not yet checked by eye.
\scanpage{25/180}
\mspage{7}

fidèlement plat sur S, ce qui prouvera  (ii)$\Longrightarrow$(iii bis) dans  16.9.

Démonstration de 16.11.  Comme les deux membres sont lisses sur S,
il suffit de prouver la platitude fibre par fibre, ce qui nous ramène au
cas où S est le spectre d'un corps alg clos k.  On peut \xout{\illegible{1}} supposer
évidemment $\underline{E}_o=\underline{\unsure{B|E}}$, donc $(**)$ prend la forme
$(***)$   $Gl(\underline{E})\longrightarrow\underline{Q}(\underline{E})^*$ ,
ce morphisme étant déduit des opérations naturelles du schéma en groupes
$Gl(\underline{E})$ sur $\underline{Q}(\underline{E})^*$ ,par $u\rightsquigarrow u(q_o)$, où $q_o$ est une section de $\underline{Q}(\underline{E})^*$ corres-%
pondant à la forme standart $Q_o$ . Or par le théorème de platitude \unsure{générique|génériaque}\hand{,}
comme $\underline{Q}(\underline{E})^*$ est lisse sur k,donc réduit, il existe un ouvert dense U de
$\underline{Q}(\underline{E})^*$ au dessus duquel le morphisme précédent soit plat.  Pour tout \margin{\unsure{dégager un lemme sur les morphismes d'espaces homogènes}}
\xout{$u\in Gl(E)(k)$}$u\in Gl(\underline{E})(k)$, i.e. pour tout automorphisme u de $\underline{E}$, u(U) satis-%
fera alors à la même condition, et il suffit\hand{,}\inserted{pour établir la platitude} de prouver\hand{),} \xout{\illegible{1}}
de $(***)$\hand{,}que les u(U) recouvrent $\underline{Q}(\underline{E})^*$ . Or pour ceci il suffit de prouver
que $Gl(\underline{E})(k)=$ Aut$(\underline{E})$ opère transitivement dans $\underline{Q}(\underline{E})^*(k)$, ensemble des
formes quadratiques \xout{\illegible{1}} sur \hand{$\checkmark$}E à discriminant non nul, \xout{\illegible{1}} i.e. que
\deleted{(}$***)$ est surjectif, ce \hand{\unsure{|}}qui prouvera alors 16.11.  On est donc ramené à
prouver le

\uline{Lemme} 16.11.1.  Soit Q une forme quadratique en n variables sur un corps k
alg clos, telle que $\overset{\hand{$\sim$}}{d}(Q)\neq 0$ , alors Q est isomorphe à la forme standart
(i.e. se met sous forme standart par un choix convenable de la base).

Lorsque E est de rang pair ou k est de car $\neq 2$ , alors l'hypothèse
signifie que Q est non dégénérée, et la conclusion se trouve dans Bourbaki
Dans le cas contraire (k de car 2, \xout{\unsure{dim}} rang E impair)  on sait que Q est
dégénérée, \xout{\unsure{donc}} soit $E_1$ une droite dans E qui est dans le noyau de la
forme associée, $E_2$ un supplémentaire, de sorte que Q se met sous forme

\margin{180}

\clearpage
% Nouveau EGA V (cote 25): archive cote 25, site page 181 (PDF page 182).
% typescript. Draft compared: none.
% AI-transcribed by Claude Opus 5.5 (claude-opus-5-5, adaptive thinking, effort high), Message Batches API, batch ega25b (msgbatch_01A7n5YuXubUPNK1zsepi8RG), 2026-11-13; not yet checked by eye.
\scanpage{25/181}
\mspage{8}

d'une somme directe de formes $Q_1 \oplus Q_2$ , $Q_1$ de rang 1 et $Q_2$ de rang pair.
En vertu de 16.8. on voit que $Q_1,Q_2$ sont de discriminant corrigé
$\neq 0$, en vertu de ce qui précède il s'ensuit que $Q_2$ est isomorphe à la
forme standart, et d'autre part $Q_1$ est non nulle, donc isomorphe à la
forme standart de degré 1, donc Q est isomorphe à la forme standart, ce
qui prouve 16.11.1.

Il reste à prouver l'équivalence de (i) et (ii) dans 16.9. On peut
évidemment supposer que S est le spectre d'un corps alg clos, \inserted[typed]{et on est} \xout{d'autre part}
ramené à prouver dans ce cas \deleted{\unsure{que}} \inserted[typed]{Lemme 16.11.2.} Q lisse \hand{$\Longleftrightarrow$} Q standart.
En vertu de 16.2. on peut supposer que k est de car 2 et E est de rang
impair. \xout{\unsure{On sait déjà que (ii) implique que}} \inserted[typed]{Lorsque} Q est isomorphe à la forme
standart, \xout{\illegible{3}} on vérifie aussitôt, grâce à 16.3. , qu'elle
est lisse (le noyau de $\overset{\hand{\vee}}{E}$ \hand{$\longrightarrow$}Ev est de dimension 1, et Q n'est pas identi-%
quement nulle sur ce noyau, donc n'y admet que le zéro trivial); \xout{donc}
\inserted{D'autre part,} \xout{\unsure{(ii) $\Rightarrow$ (i). Or}} toujours par 16.3. , \inserted{Q lisse} \deleted{(ii)} implique que \xout{le noyau} la restric-%
tion de Q au noyau N de Q n'admet quenle zéro trivial, ce qui évidemment
\inserted{signifie} \xout{implique} que N est de dim 1 et que $Q|N \neq 0$. Or prenant un supplémentaire
de N, on voit aussitôt que cela implique que Q est isomorphe à la forme
standar\over{t}{d}\hand{/} Cela \inserted[typed]{prouve 16.11.2. et} achève la démonstration de 16.9.

NB On aurait pû donner 16.11.2. dès le début en corollaire à
16.2. ou 16.3. ( qui auraient dû s'interchanger); ainsi la partie du n\textsuperscript{o}
indépendante du ``discriminant corrigé'' serait bloquée au début du n\textsuperscript{o}.
D'ailleurs tu dois mieux savoir que moi dans quelle mesure la notion
de discriminant corrigé et 16.9. étaient connus, pour en créditer qui de
droit. Peut-être y a-t\hand{\unsure{-}}il déjà une autre terminologie reçue ? Je remarque
d'ailleurs que c'est le discriminant \uline{\over{n}{c}orrigé} qui mérité vraiment le nom
de discriminant, c'est lui qui se généralise en le discriminant d'une
forme quelconque (cf N\textsuperscript{o} suivant). Le discriminant d'une forme quadratique
(dans la terminologie adoptée ici, et qui, si je ne me trompe, est la

\margin{181}

\clearpage
% Nouveau EGA V (cote 25): archive cote 25, site page 182 (PDF page 183).
% typescript. Draft compared: none.
% AI-transcribed by Claude Opus 5.5 (claude-opus-5-5, adaptive thinking, effort high), Message Batches API, batch ega25b (msgbatch_01A7n5YuXubUPNK1zsepi8RG), 2026-11-13; not yet checked by eye.
\scanpage{25/182}
\mspage{\hand{9}}

terminologie reçue  - je n'ai pas Bourbaki sous la main pour vérifier ce
point- ) devrait s'appeler \uline{déterminant} de la forme quadratique, et non
discriminant. J'aimerais connaitre ton opinion sur cette question; si
tu es d'accord, \unsure{neus|nous} profiterions de l'occasion pour rectifier sur ce
point la terminologie reçue, qui induit en erreur (puisque jusqu'à ces
derniers jours, j'ai sans m'en rendre compte confondu discriminant et
déterminant).

\uline{Remarques} 1\over{\unsure{5}}{6}.12. Lorsqu'on fait opérer le schéma en groupes absolu
$Gl(n)$ \xout{xxxxZ} $= \underline{Aut}(\underline{Z}^n )$ sur $\underline{Q}(\underline{Z}^n)$ , le stabilisateur de la forme
quadratique standart est désigné par $\underline{O}(n)$, et s'appelle le ``\uline{groupe ortho-}%
\uline{gonal absolu}''.( Plus généralement, pour toute forme quadratique $Q$ \hand{$\in$}
\unsure{$\Gamma$|T} $\underline{Sym}^2(\underline{E})$ , on introduit le sous-préschéma en groupes fermé de
$Gl(\underline{E})$ \xout{xxx} stabilisateur de $Q$ , noté $\underline{O}(Q)$ et appelé \uline{schéma en groupes}
\uline{orthogonal relatif à la forme Q}; lorsque $Q$ est isomorphe à l'image
inverse de la forme standart \inserted[typed]{à n variables}\hand{,} $\underline{O}(Q)$ est isomorphe à $O(n)_S$ ). Ceci posé,
il est immédiat que deux points \inserted{$Q,Q'$} du premier membre de (**), à valeurs dans
un $S'$, ont même image dans le deuxième membre, \unsure{sss|ssi} \xout{ilsxsontxcongruxxxie}
\xout{deuxième} on a $u' = uv$ , avec $v \in \underline{O}(n)(S')$; compte tenu de 16.11. (et
avec la terminologie qui sera développée de façon plus détaillée dans
EGA V ou VI\inserted{voir en attendant SGAD \uline{IV}}) on voit donc que les opérations naturelles à droite de
$\underline{O}(n)_S$ sur le premier membre de (**), et la projection (**), font de
$\underline{Isom}(\underline{E}_0,\underline{E})$ un \uline{fibré principal homogène sur} $\underline{Q}(\underline{E})^*$ , \uline{de groupe} $\underline{O}(n)_{\underline{Q}(\underline{E})^*}$ ;
plus particulièrement, \xout{Gl(E$_0$)xxxxGl} $Gl(n)$ est un fibré principal
homogène sur $\underline{Q}(n)^* = \underline{Q}(\underline{Z}^n)^*$, de groupe structural $\underline{O}(n)_{\underline{Q}(n)^*}$ . Il
s'ensuit \xout{xxx} d'abord que (avec les notations de la démonstration de 16.9)
P est en fait un fibré principal homogène sous le groupe $\underline{O}(n)_S$ , associé
\xout{xxxxxiqxxxxx} d'ailleurs canoniquement à la forme $Q$ (de façon fonctorielle
pour les isomorphismes de formes, et de façon compatible avec les change\_%

\margin{\hand{182}}

\clearpage
% Nouveau EGA V (cote 25): archive cote 25, site page 183 (PDF page 184).
% typescript. Draft compared: none.
% AI-transcribed by Claude Opus 5.5 (claude-opus-5-5, adaptive thinking, effort high), Message Batches API, batch ega25b (msgbatch_01A7n5YuXubUPNK1zsepi8RG), 2026-11-13; not yet checked by eye.
\scanpage{25/183}
\mspage{\hand{10}.}

ments de base. Des arguments purements formels,\xout{\illegible{2}} joints à la
``théorie de la descente plate'' de \inserted{EGA} V \unsure{1|I} , montreront alors que le foncteur
$(\underline{E},Q) \hand{$\rightsquigarrow$} P$ établit une \uline{équivalence} de la c atégorie fibrée des formes
quadratiques lisses sur des Modules localement libres de rang n, de base
un préschéma S quelconque, avec la catégorie fibrée des \uline{fibrés principaux}
\uline{homogènes} sur un préschéma quelconque S, \uline{de groupe} $\underline{O}(n)_S$ .

16.13. Soit Q une forme quadratique lisse sur un $\underline{E}$ loc libre de
rang fini, alors on vérifie facilement que pour $s \hand{$\in$} S$, $\underline{O}(Q)$ est lisse
sur un voisinage ouvert de s , \xout{\illegible{1} k(s) \illegible{1} de \illegible{1}} sauf exactemnt dans
le cas où k(s) est de car 2 et le rang de $\underline{E}$ en s est impair. Supposons
que cette circonstance ne se produi\xout{\unsure{se}} pour aucun $s \hand{$\in$} S$, alors (et alors
seulement) le P envisagé plus haut est lisse sur S (étant un fibré
principal homogène sous $\underline{O}(n)_S$ , donc lisse sss ce dernier l'est). Utili\xout{\illegible{1}}
sant le ``lemme de Hensel'' , il s'ensuit que \xout{\illegible{1}} l'on peut alors aux
conditions équivalentes de 16.9. adjoindre la condition \xout{\unsure{ex}} équivalente

(iii ter) \xout{\illegible{3}} Comme (iii), en supposant $S' \hand{$\to$} S$
\uline{étale} surjectif.

En fait, il y a mieux (en gardant toujours les hypothèses précéden-%
tes), car il résulte de la théorie générale des schémas en groupes réductifs
et fibrés principaux homogènes sur iceux (cf SGAD XXIV) que \xout{\illegible{1}} si Q est
lisse, alors tout point s admet un voisinage ouvert U, et un morphisme
étale surjectif \uline{fini} $S' \hand{$\to$} U$, tel que $Q_{S'}$ soit la forme standart. Si par exem
% NOTE: right edge of the page cut; the end of "exem-" is not visible
ple S est local, on peut dans (iii ter) supposer en plus $S' \hand{$\to$} S$ \uline{fini}.

\inserted{16.14.} \hand{(}\over{L}{C}es résultats \inserted{signalés dans 16.13.} \hand{(}à commencer par la condition (iii ter) pour une Q
lisse)\hand{)} tombent en défaut quand on abandonne l'hypothèse supplémentaire
de rang et de car. , par exemple lorsque S=Spec(k), k corps \inserted{imparfait} de car.2 , et
que E est de rang 1, car la forme quadratique at$^{2}$ , pour $\hand{$a \in$} k - k^{2}$ , ne peu
% NOTE: right edge of the page cut after "peu"

\margin{183}

\clearpage
% Nouveau EGA V (cote 25): archive cote 25, site page 184 (PDF page 185).
% typescript. Draft compared: none.
% AI-transcribed by Claude Opus 5.5 (claude-opus-5-5, adaptive thinking, effort high), Message Batches API, batch ega25b (msgbatch_01A7n5YuXubUPNK1zsepi8RG), 2026-11-13; not yet checked by eye.
\scanpage{25/184}
\mspage{11}

manifestement se mettre sous forme standart après une extension \uline{séparable}
de k (vu que k($a^{1/2}$) n'est pas séparable sur k). Cependant, dans le cas
\xout{gén} général, on peut trouver un morphisme fini \xout{surjectif} localement
libre et surjectif S'$\to$S (en fait, un fibré principal homogène sous
le \inserted{schéma en} groupes \hand{$\mu$}$_{2}$ S \inserted{$\simeq \underline{O}(n)_S/\underline{SO}(n)_S$} \hand{(}des racines carrées\hand{)} de l'unité sur S\xout{,xxxxxxxx},\hand{--} le chan-%
gement de base S'$\to$S ayant pour effet de réduire le groupe structural de
$\underline{O}(n)$ à $\underline{SO}(n)$, qui est lisse), de sorte que pour Q$_{S'}$ , le résultat de
locale isotrivialité en\over{v}{vi}sagé précédemment soit vrai. En particulier, si
S est local, alors on peut dans (iii bis) supposer S'$\to$S \uline{fini} .

\mspage{184}

\end{document}
